% Preamble for the swarm-transcribed Weber volumes.
% Deliberately close to weber_v1_ge.tex so the verified p1-148 region and the swarm
% region can live in one document.
\documentclass[11pt]{article}
\usepackage[ngerman]{babel}
\usepackage[a4paper,margin=0.9in]{geometry}
\usepackage[T1]{fontenc}
\usepackage[utf8]{inputenc}
\usepackage{lmodern}
\usepackage[leqno]{amsmath}
\usepackage{amssymb,mathtools,mathrsfs}
\usepackage{microtype}
\usepackage{fancyhdr}
\usepackage{array,longtable,enumitem,multicol}
\usepackage{tikz}
\setlength{\parindent}{1.5em}
\setlength{\parskip}{0.18em}
\emergencystretch=6em
\setlength{\headheight}{14pt}
\pagestyle{fancy}
\fancyhf{}
\cfoot{\thepage}

% --- Weber conventions -------------------------------------------------------
% Sperrung (letter-spacing) is Weber's emphasis. There is no bold in his body text.
\providecommand{\gesperrt}[1]{\textls[130]{#1}}
% Printed-page marker: a no-op in the PDF. The machine-readable anchor is the
% "% !PRINTED n" comment in the source, which the collector and Codex both read.
\providecommand{\weberpage}[1]{}
% Displayed equation numbered in the LEFT margin, as Weber prints them.
\providecommand{\weq}[2]{\begin{equation}\tag{#1}#2\end{equation}}
% Figure placeholder. A later pass replaces it with TikZ or a cleaned scan crop.
% A marked gap is honest; a fabricated diagram is not.
\providecommand{\weberfigure}[2]{%
  \begin{center}\fbox{\parbox{0.78\linewidth}{\centering\small [Figur, S.~#1: #2]}}\end{center}}

% Text symbols agents reach for that are not always defined. Providing them in the preamble
% fixes every occurrence at once and leaves the generated fragments untouched, which is the
% right trade whenever the fix is "define the macro" rather than "change the text".
\usepackage{textcomp}
% amsmath caps matrix environments at 10 columns. Weber prints permutation matrices far wider
% than that -- Bd III has an 11-column pmatrix of the w_nu -- and the failure is the opaque
% "Extra alignment tab has been changed to \cr".
\setcounter{MaxMatrixCols}{30}
\providecommand{\textquotedblbase}{\glqq}
\providecommand{\textquotedblleft}{``}
\providecommand{\textquotedblright}{''}
\providecommand{\textquotesinglbase}{\glq}

% --- Greek capitals are SLANTED in this fount -------------------------------------------
% RULING 2026-08-07, after the same observation was reported from five separate pages
% (printed 271, 275, 285, 333 and the nu/v cousin on 172).
%
% Verified at 8x zoom on printed 333 (b1 idx 358): every capital Delta on that page --
% "f(x) - f(alpha) = Delta", "Delta_alpha = f(alpha+delta) - f(alpha)" -- is the INCLINED
% form, which LaTeX calls \varDelta. \Delta, LaTeX's default, is upright and is not what
% Vieweg set. The same inclination is visible on capital Phi (printed 171, 282).
%
% This is a property of the TYPEFACE, not of any occurrence -- the same kind of fact as the
% varkappa ruling -- so it holds wherever the letter appears, and the right fix is one line
% here rather than several thousand edits across the fragments. \let leaves every fragment
% untouched, applies to the whole corpus at once, and is reversible by deleting these lines.
%
% Delta and Phi are the two verified by eye. The rest follow from its being a single fount;
% if any capital is later found upright on the page, delete its line and the fragments still
% say exactly what they said before.
\let\Gamma\varGamma
\let\Delta\varDelta
\let\Theta\varTheta
\let\Lambda\varLambda
\let\Xi\varXi
\let\Pi\varPi
\let\Sigma\varSigma
\let\Upsilon\varUpsilon
\let\Phi\varPhi
\let\Psi\varPsi
\let\Omega\varOmega
% Lower-case rho is the OPEN-TAILED form in this fount. Verified at 10x on printed 419
% ("wenn also überhaupt ein \varrho vorhanden ist"): the tail curves down and left rather than
% descending straight, which is \varrho, not \rho. Same kind of fact as the capitals above, and
% the same one-line treatment. (kappa is already handled: the fragments write \varkappa
% explicitly, on the ruling recorded in resolve_class.py.)
\let\rho\varrho

% Common shorthands inherited from the volume I file, so fragments that use them compile.
\providecommand{\eps}{\varepsilon}
\providecommand{\ii}{\mathrm{i}}
\providecommand{\tg}{\operatorname{tg}}
\providecommand{\ctg}{\operatorname{cotg}}
\providecommand{\cotg}{\operatorname{cotg}}
\providecommand{\sn}{\operatorname{sn}}
\providecommand{\cn}{\operatorname{cn}}
\providecommand{\dn}{\operatorname{dn}}
\providecommand{\am}{\operatorname{am}}
\providecommand{\dd}{\,d}
\providecommand{\tht}{\vartheta}
\providecommand{\vark}{\varkappa}
\providecommand{\Norm}{\operatorname{N}}
\providecommand{\ord}{\operatorname{ord}}
\providecommand{\ind}{\operatorname{ind}}
\providecommand{\abs}[1]{\left|#1\right|}
\providecommand{\sect}[2]{\section*{\S.\,#1. #2}}
\providecommand{\frp}{\mathfrak p}
\providecommand{\frP}{\mathfrak P}
\providecommand{\frG}{\mathfrak G}
\providecommand{\frA}{\mathfrak A}
\providecommand{\frB}{\mathfrak B}

\begin{document}
\section*{Gordan--Kerschensteiner, Vorlesungen über Invariantentheorie --- Erster Band: Determinanten}
\noindent\textit{Vollständiger erster deutscher Transkriptionsdurchgang der erfassten Blätter; ausdrückliche Unsicherheiten bleiben; keine kanonische Neukollation.}\par\bigskip

% ===== pdf idx 15, printed 1.0 =====
% !PRINTED 1
\weberpage{1}

\begin{center}
\textbf{Erster Theil.}

\textbf{Theorie der Determinanten.}
\end{center}

\begin{center}
§. 1. \textbf{Permutation und Substitution.}
\end{center}

1. \textit{Permutation.} Wenn man in einer Gruppe von $n$ Elementen — wir wollen sie mit $a, b, c \dots h$ bezeichnen — alle möglichen Vertauschungen dieser Elemente vornimmt, so entsteht eine Reihe neuer Gruppen oder Complexionen, die sich nur durch die Anordnung ihrer Elemente unterscheiden, und die man bekanntlich Permutationen nennt. Die Anzahl dieser Permutationen ist $N = n(n-1)\dots 3 \, . \, 2 \, . \, 1 = n!$, wie sich leicht erkennen lässt, wenn man bedenkt, dass jede Permutation von $n$ Elementen erhalten wird, indem man jedes der $n$ Elemente mit den sämmtlichen Permutationen von den übrigen $(n-1)$ Elementen verbindet. Dabei sieht man auch, dass von diesen $N$ Permutationen $(n-1)!$ mit $a$, ebenso viele mit $b$, $c$ u. s. w. beginnen, und wir wollen dahin übereinkommen, dass wir jede Permutation, die mit dem nämlichen Element beginnt, auch zur nämlichen Klasse rechnen, so dass wir dadurch die sämmtlichen Permutationen in $n$ Klassen getheilt haben.

2. \textit{Derangement.} Anstatt die Elemente einer Permutation mit $a, b, c \dots h$ zu bezeichnen, zieht man es vor, die Ausdrücke $a_1, a_2, a_3 \dots a_n$ einzuführen, oder noch einfacher die dem Buchstaben $a$ angehängten Zahlen $1, 2, 3 \dots n$, welche Indices oder Suffixe genannt werden. Wir nennen dann eine Permutation, in welcher kein höherer Index links vor einem niedern steht, gut geordnet. Wenn dagegen in einer Permutation ein höherer Index links vor einem niedrigeren sich befindet, so sagt man, dieselbe enthält eine Inversion oder ein Derangement, und zwar ist die Zahl dieser Derangements so gross, als die Zahl, welche angiebt, wie oft rechts von einem höhern Index sich niedrigere befinden. So sind z. B. die Permutationen
\[
(pqrs), \quad (1\;2\;3\;4\;5)
\]
gut geordnet; dagegen enthält die Permutation
\[
(7\;4\;1\;2\;3\;6\;5) \text{ Derangements: } 6 + 3 + 1 = 10
\]
%CONT>>


% ===== pdf idx 16, printed 2.0 =====
% !PRINTED 2
\weberpage{2}

%<<CONT
nämlich 6, von dem Index 7 herrührend, 3 in Folge der Stellung von 4, und 1 in Folge der Stellung des Index 6. Ebenso enthält die Permutation $(3\;5\;4\;1\;2)$ die Derangements $3\,1$, $3\,2$, $5\,4$, $5\,1$, $5\,2$, $4\,1$ und $4\,2$, also im Ganzen 7 an der Zahl.

3. \textit{Abzählung der Derangements in einer Permutation durch Theilung.} Um in einer Permutation die Zahl der Derangements rasch zu bestimmen, kann man sich manchmal mit Vortheil folgender Methode bedienen, besonders wenn die Permutation viele Elemente enthält.

Wir theilen die Permutation in 2 beliebige Theile:
\[
\left.
\begin{array}{l}
\text{die erste Gruppe enthalte der Reihe nach die Indices } i_1\,i_2\dots i_\mu\\[2pt]
\text{die zweite Gruppe enthalte der Reihe nach die Indices } k_1\,k_2\dots k_\nu
\end{array}
\right\}\;\nu+\mu=n\,.
\]

Dann bildet sich die Zahl sämmtlicher Derangements

1) aus jenen, welche von der Gruppe der $i$ allein herrühren,

2) aus jenen, welche von der Gruppe der $k$ allein herrühren,

3) aus jenen, welche die Elemente $i$ in Bezug auf die $k$ bilden.

Für die Anzahl der Derangements dritter Gattung lässt sich nun eine einfache Formel angeben; addiert man hiezu sodann die Anzahl der Inversionen erster und zweiter Gattung, die, weil sie die betreffenden Gruppen weniger Elemente enthalten, leichter sich finden lassen, als die der ganzen Permutation, so erhält man die Gesammtzahl. Diese Summenformel liefert aber der Satz:

\begin{quote}
\glqq Die Anzahl der Derangements dritter Gattung ist gleich der Summe aller Indices der Gruppe der $i$ weniger der Summe der natürlichen Zahlen von 1 bis $\mu$, wo $\mu$ die Anzahl der Elemente der ersten Gruppe bedeutet.\grqq{}
\end{quote}

\textit{Beweis:} Die Anzahl dieser Derangements dritter Gattung ändert
%CONT>>


% ===== pdf idx 17, printed 3.0 =====
% !PRINTED 3
\weberpage{3}

%<<CONT
sich nicht, wenn ich zunächst die Gruppe der $i$ und die Gruppe der $k$ unter sich gut ordne. Die erste Gruppe sei dann $(J_1\,J_2\dots J_\mu)$, die zweite $(K_1\,K_2\dots K_\nu)$, also die Permutation gleich:
\[
(J_1\;J_2\dots J_\mu \mid K_1\;K_2\dots K_\nu) \qquad (\nu+\mu=n).
\]
Diese neue Permutation hat einerseits in den $J_\varrho$, andererseits in den $K_\varrho$ keine Derangements; dagegen noch alle Derangements dritter Gattung. Die Zahl derselben zerfällt in solche die von $J_1$, in solche die von $J_2$ u. s. w., bis in solche die von $J_\mu$ herrühren. Der erste Theil ist offenbar gleich $J_1-1$, da $J_1$ mit \textit{allen} kleineren Indices Inversionen bildet; der zweite Theil gleich $J_2-2$, da $J_2$ Inversionen bildet mit allen kleineren Indices, ausgenommen mit dem vorhergehenden $J_1$; ebenso endlich der $\mu^{\text{te}}$ Theil gleich $J_\mu - \mu$; die Gesammtzahl wird also:
\[
\begin{aligned}
Z &= J_1+J_2+J_3\dots+J_\mu-(1+2+\dots \mu)\\
  &= i_1+i_2+i_3\dots+i_\mu-\frac{\mu(\mu+1)}{2}\,.
\end{aligned}
\]

Beispiele: $\alpha$) Sei gegeben die Permutation
\[
(3\;7\;9\;8\;1\;5\;2\;6\;4).
\]
Wir theilen sie beliebig durch einen Verticalstrich in zwei Gruppen:
\[
(3\;7\;9\;8 \mid 1\;5\;2\;6\;4) = (i_1\,i_2\,i_3\,i_4 \mid k_1\,k_2\,k_3\,k_4\,k_5)
\]
und machen jede derselben gut geordnet:
\[
(3\;7\;8\;9 \mid 1\;2\;4\;5\;6) = (J_1\,J_2\,J_3\,J_4 \mid K_1\,K_2\,K_3\,K_4\,K_5).
\]
Nun ist
\[
i_1+i_2+i_3+i_4 = 3+7+8+9 = 27,
\]
\[
\frac{\mu(\mu+1)}{2} = \frac{4 \cdot 5}{2} = 10,
\]
also die Zahl der Inversionen $Z = 27-10=17$, wie man auch durch directes Abzählen findet. Zählt man hierzu die eine Inversion in der Gruppe der $i$, die drei in der Gruppe der $k$, so ist die Gesammtzahl $= 21$.

$\beta$) Sei gegeben die Permutation:
\[
(7\;3\;6\;2\;4\;1\;5);
\]
wir theilen sie in zwei Theile:
\[
(7\;3\;6 \mid 2\;4\;1\;5),
\]
dann ist die Summe $7+3+6=16$, die Zahl der Elemente $=3$. Also die Zahl der Derangements dritter Gattung gleich:
\[
16 - \frac{3 \cdot 4}{2} = 10.
\]

In der That, wenn wir gut ordnen, so erhalten wir:
\[
(3\;6\;7 \mid 1\;2\;4\;5),
\]
also 2 Derangements von 3 herrührend, 4 von 6 herrührend und $7-3=4$ von 7 herrührend, in Summa 10. Hierzu addirt die Derangements der $i=2$, die der $k=2$, liefert die Gesammtsumme 14.

4. \textit{Begriff der Substitution.} Die Ueberführung von einer Permutation in eine andere geschieht durch eine Operation, die man \textit{Substitution} nennt. Soll z. B. die Permutation $(5\;1\;2\;4\;3)$ übergeführt werden in $(4\;5\;1\;3\;2)$, so geschieht dies durch eine Operation, die wir bezeichnen mit:
\[
S = \begin{pmatrix} 5 & 1 & 2 & 4 & 3\\ 4 & 5 & 1 & 3 & 2 \end{pmatrix},
\]
%CONT>>


% ===== pdf idx 18, printed 4.0 =====
% !PRINTED 4
\weberpage{4}

%<<CONT
welche die Wirkung hat, dass das Element 5 durch 4, das Element 1 durch 5, 2 durch 1, 4 durch 3 und 3 durch 2 ersetzt wird. Sie ist von der Permutation, auf welche sie angewendet wird, vollständig unabhängig, und die nämliche Substitution $S$, die wir oben mit
\[
\begin{pmatrix} 5 & 1 & 2 & 4 & 3\\ 4 & 5 & 1 & 3 & 2 \end{pmatrix}'
\]
bezeichneten, wird genau ebenso zum Ausdruck gebracht durch
\[
\begin{pmatrix} 1 & 2 & 3 & 4 & 5\\ 5 & 1 & 2 & 3 & 4 \end{pmatrix} \text{ oder } \begin{pmatrix} 3 & 2 & 1 & 4 & 5\\ 2 & 1 & 5 & 3 & 4 \end{pmatrix} \text{ u. s. w.}
\]
Daraus erkennt man, dass die verschiedenen Ausdrücke ein und derselben Substitution das eine Gemeinsame besitzen, dass die \textit{unter} einander stehenden Zahlenpaare ungeändert bleiben.

Das Zeichen
\[
S = \begin{pmatrix} 5 & 1 & 2 & 4 & 3\\ 4 & 5 & 1 & 3 & 2 \end{pmatrix} = \begin{pmatrix} 1 & 2 & 3 & 4 & 5\\ 5 & 1 & 2 & 3 & 4 \end{pmatrix} = \begin{pmatrix} 3 & 2 & 1 & 4 & 5\\ 2 & 1 & 5 & 3 & 4 \end{pmatrix}
\]
ist also so zu nehmen, dass, welche Permutation wir auch anschreiben, über denselben Zahlen immer wieder die einmal festgesetzten Zahlen stehen.

Man kann die Substitutionen insbesondere auf zwei Arten anschreiben, indem man entweder die obere oder aber auch die untere Permutation in dem Ausdrucke $S$ gut ordnet, also z. B.
\[
\begin{pmatrix} 1 & 2 & 3 & 4 & 5\\ 5 & 1 & 2 & 3 & 4 \end{pmatrix} \text{ oder } \begin{pmatrix} 2 & 3 & 4 & 5 & 1\\ 1 & 2 & 3 & 4 & 5 \end{pmatrix}.
\]

Rechnet man die specielle Substitution
\[
\begin{pmatrix} 1 & 2 & 3 & 4 & 5\\ 1 & 2 & 3 & 4 & 5 \end{pmatrix},
\]
wobei alles ungeändert bleibt, mit, so erkennt man, \textit{dass es bei $n$ Elementen ebenso viele Substitutionen als Permutationen von $n$ Elementen giebt}; denn ordnet man irgend einer Permutation eine andere zu, so kann diese Zuordnung als Ausdruck einer Substitution angesehen werden. Den $n$ oben erwähnten Klassen von Permutationen entsprechen also auch ebenso viele Klassen von Substitutionen; doch werden wir von dieser Eintheilung wenig Gebrauch machen.

Das wichtigste Eintheilungsprincip für die Substitutionen liefert in der Lehre von den Determinanten das Differenzenproduct:
\[
\begin{aligned}
&(a_1-a_2)(a_1-a_3)(a_1-a_4)\cdots\cdots\cdots(a_1-a_n) \times\\
&\qquad(a_2-a_3)(a_2-a_4)\cdots\cdots\cdots(a_2-a_n) \times\\
&\qquad\qquad(a_3-a_4)\cdots\cdots\cdots(a_3-a_n) \times\\
&\qquad\qquad\qquad\cdots\cdots\cdots\cdots\cdots\cdots\cdots\cdots\\
&\qquad\qquad\qquad\qquad(a_{n-2}-a_{n-1})(a_{n-2}-a_n) \times\\
&\qquad\qquad\qquad\qquad\qquad(a_{n-1}-a_n)\,.
\end{aligned}
\]
% ?? uncertain: the subscripts on the (a_{n-2} ...) row are faint at this resolution; transcribed per the pattern of the preceding rows.


% ===== pdf idx 19, printed 5.0 =====
% !PRINTED 5
\weberpage{5}
Da nämlich dasselbe durch Vertauschung irgend zweier oder mehrerer seiner
Elemente entweder sich gar nicht ändert oder nur sein Vorzeichen, so
können wir alle Substitutionen in \textit{zwei Hauptklassen} theilen. In die
erste Klasse kommen jene Substitutionen, welche dasselbe nicht ändern, in
die zweite Klasse jene, welche es ändern.

5. \textit{Arten der Substitutionen}. Die Substitutionen sind entweder
uneigentliche oder eigentliche. Die ersteren werden so genannt, weil sie
abhängig sind von der speciellen Permutation, auf welche sie angewendet
werden, während die eigentlichen, wie schon früher erwähnt, es nicht
sind. Zu den uneigentlichen Substitutionen gehört

$\alpha$) \textit{die einfache Vertauschung}, d.\,h. die Vertauschung
zweier neben einander stehender Elemente einer bestimmten Permutation.
So wird die Permutation (4\,5\,1\,2\,6\,3) in (4\,5\,1\,6\,2\,3)
übergeführt durch die einfache Vertauschung (2\,6), wobei diese
Vertauschung speciell nur für die vorliegende Permutation zu denken ist.
Ebenso gehen die Permutationen (7\,3\,2\,5), (6\,1\,4\,9) durch die
einfachen Vertauschungen (3\,2), (4\,9) in die Permutationen (7\,2\,3\,5),
(6\,1\,9\,4) über.

$\beta$) \textit{der einfache Cyclus}, der entweder vorschreitend oder
rückschreitend sein kann. Vertauscht man nämlich \textit{k neben einander
stehende} Elemente einer bestimmten Permutation in der Weise, dass an
Stelle des ersten Elementes das zweite, an Stelle des zweiten das dritte
u.\,s.\,w., an Stelle des letzten wieder das erste tritt, dann bilden
diese Vertauschungen einen vorschreitenden Cyclus; rückt dagegen das
vorletzte Element an Stelle des letzten, das zweitvorletzte an Stelle des
vorletzten u.\,s.\,w. und das letzte an die erste Stelle, so hat man
einen rückschreitenden Cyclus vor sich. So stellen z.\,B. die Ausdrücke

\[
\begin{pmatrix}2\,4\,3\,5\\4\,3\,5\,2\end{pmatrix},\quad
\begin{pmatrix}1\,7\,6\,4\,8\,3\\7\,6\,4\,8\,3\,1\end{pmatrix}
\]
vorschreitende, die Ausdrücke
\[
\begin{pmatrix}2\,4\,3\,5\\5\,2\,4\,3\end{pmatrix},\quad
\begin{pmatrix}1\,7\,6\,4\,8\,3\\3\,1\,7\,6\,4\,8\end{pmatrix}
\]
rückschreitende Cyclen dar, welche z.\,B. die speciellen Permutationen
(1\,2\,4\,3\,5\,6), (9\,2\,0\,5\,1\,7\,6\,4\,8\,3) überführen im ersten
Falle in (1\,4\,3\,5\,2\,6), (9\,2\,0\,5\,7\,6\,4\,8\,3\,1), im zweiten
Falle in (1\,5\,2\,4\,3\,6), (9\,2\,0\,5\,3\,1\,7\,6\,4\,8).

Man kann sich einen \textit{einfachen} Cyclus am besten in der Weise
klar machen, dass man sich etwa die sechs Elemente 1\,7\,6\,4\,8\,3
symmetrisch auf zwei concentrische Kreisringe in nebenstehender Weise
%CONT>>


% ===== pdf idx 20, printed 6.0 =====
% !PRINTED 6
\weberpage{6}
%<<CONT
angeordnet denkt. Dreht man den inneren Ring (oder auch den äusseren) um
$60^0$ nach links, so charakterisirt die neue Lage beider Ringe einen
vorschreitenden, bei einer gleichen Drehung nach rechts einen
rückschreitenden einfachen Cyclus.

\weberfigure{6}{Zwei concentrische Kreisringe (äusserer und innerer
Ring), auf denen die sechs Elemente 1, 7, 3, 6, 8, 4 doppelt eingetragen
sind (aussen und innen an denselben Winkelpositionen), mit gestrichelten
Diagonalen, die den Mittelpunkt kreuzen; die Drehung des inneren Ringes
gegen den äusseren veranschaulicht den vorschreitenden bzw.
rückschreitenden einfachen Cyclus.}

Zu den eigentlichen Substitutionen gehört:

$\alpha$) \textit{die gewöhnliche allgemeine Vertauschung} beliebig
zweier Elemente einer Permutation von $n$ Elementen (eigentliche
Substitution im engern Sinne). Ihr Ausdruck ist der bereits früher
erwähnte. Solche allgemeine Vertauschungen sind z.\,B. gegeben durch

\[
S_1=\begin{pmatrix}1\,5\,2\,6\,3\,4\\1\,2\,5\,6\,3\,4\end{pmatrix},\qquad
S_2=\begin{pmatrix}3\,4\,6\,7\,5\,1\,2\\3\,1\,6\,7\,5\,4\,2\end{pmatrix}
\]
und man kann die nämlichen Substitutionen $S_1$ und $S_2$ durch eine
ganze Reihe anderer Ausdrücke darstellen, z.\,B.
\[
S_1=\begin{pmatrix}4\,5\,6\,3\,2\,1\\4\,2\,6\,3\,5\,1\end{pmatrix},\qquad
S_2=\begin{pmatrix}3\,4\,1\,6\,7\,5\,2\\3\,1\,4\,6\,7\,5\,2\end{pmatrix},
\]
worauf wir schon früher aufmerksam gemacht haben.

$\beta$) \textit{Der allgemeine Cyclus}, durch welchen $k$ beliebige
Elemente einer Permutation in der Weise wie beim einfachen Cyclus ihre
Plätze ändern. Ist z.\,B. die Permutation (3\,5\,6\,2\,1\,4) von 6
Elementen gegeben, so schreiben wir zunächst, um den Ausdruck eines
allgemeinen Cyclus zu erhalten, einen einfachen an, etwa
\[
\begin{pmatrix}2\,4\,1\,3\\4\,1\,3\,2\end{pmatrix}.
\]
Wenden wir dann diese Substitutionen auf die obige Permutation an, so
erhalten wir den allgemeinen Cyclus
\[
\begin{pmatrix}3\,5\,6\,2\,1\,4\\2\,5\,6\,4\,3\,1\end{pmatrix}=
\begin{pmatrix}6\,3\,4\,5\,1\,2\\6\,2\,1\,5\,3\,4\end{pmatrix}=\text{u.\,s.\,w.}
\]

\textit{Anmerkung} 1. Einfacher Cyclus sowohl als einfache Vertauschung
erhalten den allgemeinen Charakter der eigentlichen Substitutionen,
sobald wir die Beschränkung fallen lassen, dass sie nur für eine
bestimmte Permutation gelten sollen.

So werden wir unter (3\,5) eine einfache Vertauschung, angewendet etwa
auf (4\,2\,3\,5\,1), verstehen, dagegen die allgemeine darstellen durch
\[
\begin{pmatrix}4\,2\,3\,5\,1\\4\,2\,5\,3\,1\end{pmatrix}=
\begin{pmatrix}3\,4\,1\,5\,2\\5\,4\,1\,3\,2\end{pmatrix}=\text{u.\,s.\,w.}
\]
So wird auch
\[
\begin{pmatrix}6\,1\,4\,3\\1\,4\,3\,6\end{pmatrix}
\]
%CONT>>


% ===== pdf idx 21, printed 7.0 =====
% !PRINTED 7
\weberpage{7}
%<<CONT
den einfachen Cyclus bedeuten, der auf die Permutation (5\,2\,6\,1\,4\,3)
anzuwenden ist, während der allgemeine charakterisirt ist durch
\[
\begin{pmatrix}5\,2\,6\,1\,4\,3\\5\,2\,1\,4\,3\,6\end{pmatrix}=
\begin{pmatrix}1\,6\,4\,5\,2\,3\\4\,1\,3\,5\,2\,6\end{pmatrix}=
\begin{pmatrix}2\,1\,6\,5\,3\,4\\2\,4\,1\,5\,6\,3\end{pmatrix}=\text{u.\,s.\,w.}
\]

\textit{Anmerkung} 2. Eine \textit{beliebige} Substitution $S$ verlangt,
abgesehen von der einen
$S=\begin{pmatrix}1\,2\,3\,\dots\,n\\1\,2\,3\,\dots\,n\end{pmatrix}$,
die Alles ungeändert lässt, allgemeine gewöhnliche, oder auch, wie wir
gleich sehen werden, allgemeine cyclische Vertauschungen.

6. \textit{Herstellung der uneigentlichen Substitutionen durch
einander}. Der einfache Cyclus lässt sich herstellen durch einfache
Vertauschungen.

\textit{Beispiel} 1. Sei gegeben der einfache Cyclus
\[
\begin{pmatrix}2\,5\,3\,1\\5\,3\,1\,2\end{pmatrix}.
\]
Lassen wir nun das Element 2 successive nach rechts wandern, dann
entsteht aus (2\,5\,3\,1) durch die einfache Vertauschung (2\,5) zuerst
(5\,2\,3\,1), die nächste einfache Vertauschung (2\,3) führt zu
(5\,3\,2\,1), endlich die letzte (2\,1) zum gesuchten Resultat
(5\,3\,1\,2); also lässt sich ein einfacher Cyclus von 4 Elementen durch
3 einfache Vertauschungen ersetzen.

\textit{Beispiel} 2. Der einfache Cyclus sei
\[
\begin{pmatrix}4\,6\,5\,8\,7\,3\\6\,5\,8\,7\,3\,4\end{pmatrix}.
\]
Wir lassen wieder das erste Element 4 der Permutation wandern.

\begin{center}
\begin{tabular}{cc}
Einfache Vertauschung. & Permutation. \\
(4\,6) & 6\,4\,5\,8\,7\,3 \\
(4\,5) & 6\,5\,4\,8\,7\,3 \\
(4\,8) & 6\,5\,8\,4\,7\,3 \\
(4\,7) & 6\,5\,8\,7\,4\,3 \\
(4\,3) & 6\,5\,8\,7\,3\,4 \\
\end{tabular}
\end{center}
An Stelle des einfachen Cyclus von 6 Elementen treten somit 5 einfache
Vertauschungen.

Man erkennt also leicht die Richtigkeit des Satzes:

\glqq Ist die Differenz der Stellenzahl des ersten und letzten Elementes
eines einfachen Cyclus $\nu$, so kann man denselben durch $\nu$ einfache
Vertauschungen ersetzen; es spielen dann $\nu+1$ Elemente im Cyclus
mit.\grqq{}

7. \textit{Die allgemeine Vertauschung lässt sich durch zwei einfache
Cyclen herstellen.} Sei gegeben 1) die allgemeine Vertauschung
%CONT>>


% ===== pdf idx 22, printed 8.0 =====
% !PRINTED 8
\weberpage{8}
%<<CONT
\[
S=\begin{pmatrix}2\,4\,1\,5\,3\,6\\6\,4\,1\,5\,3\,2\end{pmatrix},
\]
dann liefert der einfache Cyclus (gebildet aus allen 6 Elementen)
\[
\begin{pmatrix}2\,4\,1\,5\,3\,6\\4\,1\,5\,3\,6\,2\end{pmatrix}
\]
zunächst die Permutation (4\,1\,5\,3\,6\,2).

Wendet man auf diese den einfachen Cyclus
\[
\begin{pmatrix}4\,1\,5\,3\,6\\6\,4\,1\,5\,3\end{pmatrix}
\]
an, gebildet aus fünf Elementen, so führt er zum gesuchten Resultat
(6\,4\,1\,5\,3\,2).

\textit{Beispiel} 2. Die Substitution
\[
S=\begin{pmatrix}5\,7\,8\,2\,4\,1\,3\,6\\5\,4\,8\,2\,7\,1\,3\,6\end{pmatrix}
\]
lässt sich spalten in den einfachen vorschreitenden Cyclus von vier
Elementen
\[
\begin{pmatrix}7\,8\,2\,4\\8\,2\,4\,7\end{pmatrix}.
\]
Durch ihn wird in der Permutation (5\,7\,8\,2\,4\,1\,3\,6) zunächst das
Element 7 an die durch $S$ angedeutete Stelle gebracht, so dass man
erhält (5\,8\,2\,4\,7\,1\,3\,6). Wendet man hierauf den
rückschreitenden einfachen Cyclus von drei Elementen
$\begin{pmatrix}8\,2\,4\\4\,8\,2\end{pmatrix}$ an, so erhält man das
gewünschte Resultat (5\,4\,8\,2\,7\,1\,3\,6). Die weiteren Beispiele
werden wohl ohne weitere Erklärung verständlich sein.

3) Allgemeine Vertauschung:
$S=\begin{pmatrix}3\,0\,2\,7\,1\,5\,6\,4\\3\,0\,2\,5\,1\,7\,6\,4\end{pmatrix}.$

Erster Cyclus: $\begin{pmatrix}7\,1\,5\\1\,5\,7\end{pmatrix}$; Resultat:
(3\,0\,2\,1\,5\,7\,6\,4).

Zweiter Cyclus: $\begin{pmatrix}1\,5\\5\,1\end{pmatrix}$; Resultat:
(3\,0\,2\,5\,1\,7\,6\,4).

4) Allgemeine Vertauschung:
$S=\begin{pmatrix}5\,3\,7\,8\,9\,1\,2\,4\,6\\5\,4\,7\,8\,9\,1\,2\,3\,6\end{pmatrix}.$

Erster Cyclus:
$\begin{pmatrix}3\,7\,8\,9\,1\,2\,4\\7\,8\,9\,1\,2\,4\,3\end{pmatrix}$;
Resultat: (5\,7\,8\,9\,1\,2\,4\,3\,6).

Zweiter Cyclus:
$\begin{pmatrix}7\,8\,9\,1\,2\,4\\4\,7\,8\,9\,1\,2\end{pmatrix}$;
Resultat: (5\,4\,7\,8\,9\,1\,2\,3\,6).

Aus diesen vier Beispielen lässt sich leicht folgendes Gesetz
abstrahiren:

\glqq Ist die Differenz der Stellenzahl der beiden zu vertauschenden
Elemente gleich $\nu$, dann spielen beim ersten Cyclus $\nu+1$ Elemente,
beim
%CONT>>


% ===== pdf idx 23, printed 9.0 =====
% !PRINTED 9
\weberpage{9}
%<<CONT
zweiten $v$ Elemente mit. Daher besteht nach dem vorigen Satze der erste Cyclus aus $v$ einfachen Vertauschungen, der zweite Cyclus aus $v-1$. Wenn man also auf \emph{diesem} Wege die allgemeine Vertauschung in einfache zerlegt, so ist Gesammtsumme der letzteren $2v-1$, also eine ungerade Zahl.\glqq{} Es sei hierbei bemerkt, dass diese Zahl $2v-1$ lediglich von dem oben eingeschlagenen Wege abhängig ist. Irgend ein systematisches oder auch willkürliches anderes Verfahren, die allgemeine Vertauschung auf einfache zurückzuführen, wird im Allgemeinen auf eine andere Zahl führen.

8. \emph{Herstellung des allgemeinen Cyclus aus der allgemeinen Vertauschung.} Wir bestimmen den zum allgemeinen Cyclus gehörigen einfachen Cyclus, lösen diesen in seine einfachen Vertauschungen auf, und bilden alsdann die entsprechenden allgemeinen Vertauschungen. Diese ersetzen gerade den allgemeinen Cyclus.

Beispiel 1. Es sei gegeben der allgemeine Cyclus
\[
S = \begin{pmatrix} 4 & 3 & 5 & 1 & 6 & 2 \\ 1 & 2 & 5 & 3 & 6 & 4 \end{pmatrix}.
\]
Der zugehörige einfache ist $\begin{pmatrix} 4 & 1 & 3 & 2 \\ 1 & 3 & 2 & 4 \end{pmatrix}$, welcher selbst in die einfachen Vertauschungen $(4\,1)$, $(4\,3)$, $(4\,2)$ aufgelöst werden kann.

Hierzu bilden wir der Reihe nach die zugehörigen allgemeinen Vertauschungen und erhalten so entsprechend $(4\,1)$
\[
S_1 = \begin{pmatrix} 4 & 3 & 5 & 1 & 6 & 2 \\ 1 & 3 & 5 & 4 & 6 & 2 \end{pmatrix}\text{, daraus entsprechend } (4\,3)
\]
\[
S_2 = \begin{pmatrix} 1 & 3 & 5 & 4 & 6 & 3 \\ 1 & 4 & 5 & 3 & 6 & 2 \end{pmatrix}\text{, daraus entsprechend } (4\,2) % ?? uncertain: top row digits of S_2 read "1 3 5 4 6 3" as printed
\]
\[
S_3 = \begin{pmatrix} 1 & 4 & 5 & 3 & 6 & 2 \\ 1 & 2 & 5 & 3 & 6 & 4 \end{pmatrix},
\]
womit das gewünschte Resultat erreicht ist.

Beispiel 2. Es ist der allgemeine Cyclus
\[
S = \begin{pmatrix} 3 & 5 & 7 & 9 & 8 & 2 & 6 & 4 & 1 \\ 7 & 5 & 8 & 9 & 6 & 2 & 1 & 4 & 3 \end{pmatrix}
\]
in allgemeine Vertauschungen zu zerlegen. Wir bestimmen aus $S$ den einfachen Cyclus
\[
\begin{pmatrix} 3 & 7 & 8 & 6 & 1 \\ 7 & 8 & 6 & 1 & 3 \end{pmatrix}.
\]
Ihm entsprechen die einfachen Vertauschungen:
\[
(3\,7),\ (3\,8),\ (3\,6),\ (3\,1).
\]
Also zerfällt $S$ in die vier allgemeinen Vertauschungen
\[
S_1 = \begin{pmatrix} 3 & 5 & 7 & 9 & 8 & 2 & 6 & 4 & 1 \\ 7 & 5 & 3 & 9 & 8 & 2 & 6 & 4 & 1 \end{pmatrix}, \qquad
S_2 = \begin{pmatrix} 7 & 5 & 3 & 9 & 8 & 2 & 6 & 4 & 1 \\ 7 & 5 & 8 & 9 & 3 & 2 & 6 & 4 & 1 \end{pmatrix},
\]
%CONT>>


% ===== pdf idx 24, printed 10.0 =====
% !PRINTED 10
\weberpage{10}
%<<CONT
\[
S_3 = \begin{pmatrix} 7 & 5 & 8 & 9 & 3 & 2 & 6 & 4 & 1 \\ 7 & 5 & 8 & 9 & 6 & 2 & 3 & 4 & 1 \end{pmatrix}, \qquad
S_4 = \begin{pmatrix} 7 & 5 & 8 & 9 & 6 & 2 & 3 & 4 & 1 \\ 7 & 5 & 8 & 9 & 6 & 2 & 1 & 4 & 3 \end{pmatrix}.
\]
Man erkennt aus diesen Beispielen, dass der allgemeine Cyclus sich genau in derselben Weise in allgemeine Vertauschungen auflösen lässt, wie der einfache Cyclus in einfache Vertauschungen. Spielen daher beim allgemeinen Cyclus $v+1$ Elemente mit, so kann man ihn in $v$ allgemeine Vertauschungen auflösen. (Im Beispiel 1 waren die Elemente des Cyclus an Zahl 4, die allgemeinen Vertauschungen 3; im Beispiel 2 spielten 5 Elemente im Cyclus mit, also ergaben sich 4 allgemeine Vertauschungen.)

9. \emph{Herstellung einer beliebigen Substitution durch allgemeine Cyclen.} Sei z.\,B. die gegebene Substitution
\[
S = \begin{pmatrix} 9 & 0 & 8 & 4 & 3 & 5 & 2 & 1 & 7 & 6 \\ 1 & 6 & 5 & 0 & 9 & 2 & 8 & 3 & 4 & 7 \end{pmatrix}.
\]
Wir können sie uns entstanden denken dadurch, dass zunächst 9 durch 1, 1 durch 3, 3 wieder durch 9 ersetzt, d.\,h. dass die obere Permutation zuerst dem allgemeinen Cyclus:
\[
S_1 = \begin{pmatrix} 9 & 0 & 8 & 4 & 3 & 5 & 2 & 1 & 7 & 6 \\ 1 & 0 & 8 & 4 & 9 & 5 & 2 & 3 & 7 & 6 \end{pmatrix}
\]
unterworfen wurde. In dem Resultat dieser Substitution, d.\,h. in der untern Permutation dieses Ausdruckes $S_1$, trat hierauf an Stelle der 0 der Index 6, an Stelle von 6 der Index 7, an Stelle von 7 der Index 4, endlich an Stelle von 4 wieder der Index Null, d.\,h. sie wurde einer Veränderung unterworfen, bewirkt durch den allgemeinen Cyclus:
\[
S_2 = \begin{pmatrix} 1 & 0 & 8 & 4 & 9 & 5 & 2 & 3 & 7 & 6 \\ 1 & 6 & 8 & 0 & 9 & 5 & 2 & 3 & 4 & 7 \end{pmatrix}.
\]
Aus der untern Permutation der Substitution $S_2$ gelangt man endlich zur gesuchten untern Permutation im Ausdrucke $S$, indem man 8 durch 5, 5 durch 2 und 2 wieder durch 8 ersetzt, also auf sie den allgemeinen Cyclus anwendet:
\[
S_3 = \begin{pmatrix} 1 & 6 & 8 & 0 & 9 & 5 & 2 & 3 & 4 & 7 \\ 1 & 6 & 5 & 0 & 9 & 2 & 8 & 3 & 4 & 7 \end{pmatrix}.
\]
Wir sehen somit, die allgemeine Substitution $S$ lässt sich ersetzen durch die drei allgemeinen Cyclen $S_1$, $S_2$, $S_3$.

Die Anzahl der Cyclen, die für eine allgemeine Substitution eintreten können, lässt sich nicht von vornherein angeben; sie lässt sich aber rasch ermitteln, wenn man, etwa von der ersten Ziffer links ausgehend, nach der im vorstehenden Beispiele dreimal angegebenen Methode verfährt. Jedesmal wenn wir auf diese Weise zur Ausgangsziffer zurückkehren, ist der erste Cyclus geschlossen, und die Ermittlung des zweiten Cyclus kann dann mit dem nächsten Index, der nicht schon im ersten enthalten ist, begonnen werden.
% --- printed page 11 begins here (running head seen at top of this band) ---

Es kann natürlich eintreten, dass überhaupt nur Ein Cyclus auftritt, d.\,h. dass die allgemeine Substitution selbst einen einzigen Cyclus darstellt, wie in dem Beispiele:
\weq{2}{S = \begin{pmatrix} 5 & 9 & 8 & 2 & 3 & 4 & 1 & 7 & 6 \\ 9 & 2 & 5 & 1 & 4 & 6 & 3 & 8 & 7 \end{pmatrix}}
oder aber auch, dass $S$ Cyclen enthält mit einem einzigen Elemente, wie z.\,B. die Substitution
\weq{3}{S = \begin{pmatrix} 5 & 1 & 8 & 2 & 3 & 4 & 7 & 6 \\ 2 & 1 & 5 & 4 & 3 & 6 & 8 & 7 \end{pmatrix}.}
Sie zerfällt in den Cyclus $\begin{pmatrix} 5 & 2 & 4 & 6 & 7 & 8 \\ 2 & 4 & 6 & 7 & 8 & 5 \end{pmatrix}$ mit sechs Elementen und
%CONT>>


% ===== pdf idx 25, printed 11.0 =====
% !PRINTED 11
\weberpage{11}
%<<CONT
Sie zerfällt in den Cyclus $\begin{pmatrix} 5 & 2 & 4 & 6 & 7 & 8 \\ 2 & 4 & 6 & 7 & 8 & 5 \end{pmatrix}$ mit sechs Elementen und in die beiden eingliedrigen Cyclen $(1)$, $(3)$.

Hat man nun auf die angegebene Art die Zahl $v$ der Cyclen ermittelt, welche eine allgemeine Substitution vertreten können, so erhält man die Zahl der allgemeinen Vertauschungen in folgender Weise. Sei die Zahl der Elemente des ersten Cyclus $\alpha_1$, die des zweiten $\alpha_2 \ldots$ u.\,s.\,w., die des $v^{\mathrm{ten}}$ $\alpha_v$, dann ist, weil ja alle Elemente durch alle Cyclen erschöpft sein müssen, $\Sigma\alpha_\varrho = \alpha_1 + \alpha_2 + \cdots \alpha_v = n$. Nun besteht aber nach Nr.\,8 der allgemeine Cyclus von $\alpha_\varrho$ Elementen aus $\alpha_\varrho - 1$ Vertauschungen, daher ist die Summe aller allgemeinen Vertauschungen, welche die allgemeine Substitution ersetzen können, gleich
\[
\Sigma\alpha_\varrho - (1 + 1 + \cdots v\ \text{mal}) = n - v.
\]
Es mag hier nochmals erwähnt werden, dass diese Anzahl lediglich abhängig ist von der hier eingeschlagenen Methode, sie zu ermitteln.

10. \emph{Bestimmung der Hauptklasse} (siehe Nr.\,4), \emph{der eine Substitution angehört.} Diese Aufgabe könnte am Besten am Differenzenproduct direct erledigt werden, indem man an demselben der Reihe nach die einzelnen Vertauschungen vornimmt. Doch möchte sie vielleicht unter Berücksichtigung einzelner Merkmale der Substitutionen, die wir nun besprochen haben, leichter zu lösen sein. Wir sehen nämlich Folgendes:

$\alpha$) Da sich die einfache Vertauschung auf eine einzige bestimmte, wenn auch willkürlich gegebene Permutation bezieht, so können wir der Einfachheit halber annehmen, diese willkürlich gegebene sei die \emph{gut geordnete} Permutation. Dies vorausgesetzt, kann eine einfache Vertauschung zweier Nachbarindices $\alpha_\varrho$ und $\alpha_{\varrho+1}$, die also natürlich nur um 1 verschieden sein können, angewendet auf das Differenzenproduct der Indices, nur das Vorzeichen der einzigen Differenz $(\alpha_\varrho - \alpha_{\varrho+1})$ ändern, da in allen übrigen Differenzen der Abstand der
% --- printed page 12 begins here (running head seen at top of this band) ---
Indices von dem Index $\alpha_\varrho$ bezw. $\alpha_{\varrho+1}$ mindestens zwei Einheiten beträgt, also eine Veränderung des Minuenden oder Subtrahenden um 1 keine Vorzeichenänderung zur Folge haben kann. Demnach ändert mit diesem einzelnen Factor $(\alpha_\varrho - \alpha_{\varrho+1})$ das ganze Product sein Zeichen, d.\,h.: \glqq Die einfache Vertauschung gehört der zweiten Hauptklasse an.\grqq{} Damit ist aber auch die Classification der übrigen Substitutionen gewonnen.

$\beta$) Der einfache Cyclus von $v+1$ Elementen gehört nach (6) der ersten oder zweiten Hauptklasse an, je nachdem $v$ ($=$ der Zahl seiner einfachen Vertauschungen) gerade oder ungerade ist.

$\gamma$) Die allgemeine Vertauschung gehört immer der zweiten Hauptklasse an, da sie nach Nr.\,7 aus $2v-1$, also einer ungeraden Zahl einfacher Vertauschungen besteht.

$\delta$) Der allgemeine Cyclus verhält sich wie der einfache, da ja die
%CONT>>


% ===== pdf idx 26, printed 12.0 =====
% !PRINTED 12
\weberpage{12}
%<<CONT
allgemeinen Vertauschungen, durch welche ersterer ersetzt werden kann, wie die einfachen Vertauschungen zur zweiten Hauptklasse gehören.

$\varepsilon$) Die beliebige Substitution endlich besteht aus $n-v$ allgemeinen Vertauschungen, wo $n$ gleich der Zahl der Elemente und $v$ gleich der Zahl der die Substitution ersetzenden allgemeinen Cyclen ist. Je nachdem diese Differenz also gerade oder ungerade ist, gehört die Substitution zur ersten oder zweiten Hauptklasse.

11. Wir wenden diese Merkmale auf die bereits behandelten Beispiele an.

$\alpha$) In der Permutation $(5\,3\,4\,2\,1)$ ändert die einfache Vertauschung $(2\,1)$ das Vorzeichen des Differenzenproductes
\[
\begin{aligned}
(1-2)(1-3)(1-4)(1-5) &= P \\
(2-3)(2-4)(2-5) & \\
(3-4)(3-5) & \\
(4-5) &
\end{aligned}
\]
denn in ihm ändert nur der erste Factor sein Zeichen, die übrigen gehen ineinander über; man erhält
\[
\begin{aligned}
(2-1)(2-3)(2-4)(2-5) &= -P. \\
(1-3)(1-4)(1-5) & \\
(3-4)(3-5) & \\
(4-5). &
\end{aligned}
\]

$\beta$) Der erste einfache Cyclus in Nr.\,6:
\[
\begin{pmatrix} 2 & 5 & 3 & 1 \\ 5 & 3 & 1 & 2 \end{pmatrix}
\]
setzte sich zusammen aus drei einfachen Vertauschungen $(2\,5)$, $(2\,3)$, $(2\,1)$,
%CONT>>


% ===== pdf idx 27, printed 13.0 =====
% !PRINTED 13
\weberpage{13}
%<<CONT
also gehört er der zweiten Klasse an, während z.\,B. der einfache Cyclus
\[
\begin{pmatrix} 2 & 5 & 3 & 1 & 4 \\ 5 & 3 & 1 & 4 & 2 \end{pmatrix},
\]
da er aus den vier einfachen Vertauschungen $(2\ 5)$, $(2\ 3)$, $(2\ 1)$, $(2\ 4)$ besteht, unter die erste Hauptklasse zu rechnen ist.

$\gamma$) Die allgemeine Vertauschung
\[
S = \begin{pmatrix} 2 & 4 & 1 & 5 & 3 & 6 \\ 6 & 4 & 1 & 5 & 3 & 2 \end{pmatrix} \ \text{Nr. 7, Beispiel 1,}
\]
wurde ersetzt durch die beiden einfachen Cyclen
\[
S_1 = \begin{pmatrix} 2 & 4 & 1 & 5 & 3 & 6 \\ 4 & 1 & 5 & 3 & 6 & 2 \end{pmatrix} \quad \text{und} \quad S_2 = \begin{pmatrix} 4 & 1 & 5 & 3 & 6 \\ 6 & 4 & 1 & 5 & 3 \end{pmatrix};
\]
der erste besteht aus fünf, der zweite aus vier einfachen Vertauschungen. Da die Summe derselben $= 9$ ist, so gehört $S$ zur zweiten Klasse der Substitutionen; die Behauptung läfst sich leicht direct am Differenzenproduct verificiren.

$\delta$) Die Substitution Nr. 7, Beispiel 2,
\[
S = \begin{pmatrix} 5 & 7 & 8 & 2 & 4 & 1 & 3 & 6 \\ 5 & 4 & 8 & 2 & 7 & 1 & 3 & 6 \end{pmatrix}
\]
zerfiel in die Cyclen
\[
S_1 = \begin{pmatrix} 7 & 8 & 2 & 4 \\ 8 & 2 & 4 & 7 \end{pmatrix}, \qquad S_2 = \begin{pmatrix} 8 & 2 & 4 \\ 4 & 8 & 2 \end{pmatrix},
\]
deren erster aus den drei einfachen Vertauschungen $(7\ 8)$, $(7\ 2)$, $(7\ 4)$, deren zweiter aus den zwei Vertauschungen $(8\ 2)$, $(8\ 4)$ besteht, ihre Summe ist $= 5$, also wiederum eine ungerade Zahl.

$\varepsilon$) Der allgemeine Cyclus
\[
S = \begin{pmatrix} 4 & 3 & 5 & 1 & 6 & 2 \\ 1 & 2 & 5 & 3 & 6 & 4 \end{pmatrix}
\]
in Nr. 8 wurde zerlegt in die drei allgemeinen Vertauschungen
\[
S_1 = \begin{pmatrix} 4 & 3 & 5 & 1 & 6 & 2 \\ 1 & 3 & 5 & 4 & 6 & 2 \end{pmatrix}, \quad S_2 = \begin{pmatrix} 1 & 3 & 5 & 4 & 6 & 2 \\ 1 & 4 & 5 & 3 & 6 & 2 \end{pmatrix}, \quad S_3 = \begin{pmatrix} 1 & 4 & 5 & 3 & 6 & 2 \\ 1 & 2 & 5 & 3 & 6 & 4 \end{pmatrix}.
\]
Für $S_1$ ist die Differenz der zu vertauschenden Stellenzahlen $\nu_1 = 3$, daher die zugehörige Zahl der einfachen Vertauschungen $2\nu_1 - 1 = 5$. Ferner ist das zu $S_2$ gehörige $\nu_2 = 2$, die Zahl der einfachen Vertauschungen $2\nu_2 - 1 = 3$, endlich das zu $S_3$ gehörige $\nu_3 = 4$, die Zahl der einfachen Vertauschungen $2\nu_3 - 1 = 7$. Die Summe aller Vertauschungen ist also $5 + 3 + 7 = 15$.

Daher gehört $S$ der zweiten Hauptklasse an.

Analysiren wir auf die gleiche Weise das im nämlichen Abschnitt gegebene zweite Beispiel, so finden wir
\[
2\nu_1 - 1 = 3, \quad 2\nu_2 - 1 = 3, \quad 2\nu_3 - 1 = 3, \quad 2\nu_4 - 1 = 3.
\]


% ===== pdf idx 28, printed 14.0 =====
% !PRINTED 14
\weberpage{14}
Die Summe dieser vier ungeraden Zahlen giebt eine gerade, d.\,h. diese Substitution gehört der ersten Hauptklasse an.

$\zeta$) Betrachten wir endlich die beliebige Substitution
\[
S = \begin{pmatrix} 9 & 0 & 8 & 4 & 3 & 5 & 2 & 1 & 7 & 6 \\ 1 & 6 & 5 & 0 & 9 & 2 & 8 & 3 & 4 & 7 \end{pmatrix}
\]
in Nr. 9. Wir zerlegten sie in die drei allgemeinen Cyclen
\[
S_1 = \begin{pmatrix} 9 & 0 & 8 & 4 & 3 & 5 & 2 & 1 & 7 & 6 \\ 1 & 0 & 8 & 4 & 9 & 5 & 2 & 3 & 7 & 6 \end{pmatrix},
\]
\[
S_2 = \begin{pmatrix} 1 & 0 & 8 & 4 & 9 & 5 & 2 & 3 & 7 & 6 \\ 1 & 6 & 8 & 0 & 8 & 5 & 2 & 3 & 4 & 7 \end{pmatrix},
\]
\[
S_3 = \begin{pmatrix} 1 & 6 & 8 & 0 & 9 & 5 & 2 & 3 & 4 & 7 \\ 1 & 6 & 5 & 0 & 9 & 2 & 8 & 3 & 4 & 7 \end{pmatrix},
\]
deren einfache sind
\[
\begin{pmatrix} 9 & 3 & 1 \\ 1 & 9 & 3 \end{pmatrix},
\]
\[
\begin{pmatrix} 0 & 4 & 7 & 6 \\ 6 & 0 & 4 & 7 \end{pmatrix},
\]
\[
\begin{pmatrix} 8 & 5 & 2 \\ 5 & 2 & 8 \end{pmatrix}.
\]
Zum ersten Cyclus gehören also zwei, zum zweiten drei, zum dritten wieder zwei allgemeine Vertauschungen. Die Anzahl derselben ist somit $= 7$, und weil jede allgemeine Vertauschung nur in eine ungerade Zahl von einfachen aufgelöst werden kann, so ist die Klasse dieser Substitution die zweite.

$\eta$) Dagegen gehört die Substitution
\[
S = \begin{pmatrix} 5 & 6 & 7 & 3 & 4 & 2 \\ 3 & 4 & 6 & 2 & 7 & 5 \end{pmatrix},
\]
welche sich in zwei allgemeine Cyclen spaltet, denen die drei einfachen
\[
\begin{pmatrix} 5 & 3 & 2 \\ 3 & 2 & 5 \end{pmatrix} \quad \text{und} \quad \begin{pmatrix} 6 & 4 & 7 \\ 4 & 7 & 6 \end{pmatrix}
\]
entsprechen, zur ersten Hauptklasse, weil in den zwei allgemeinen Cyclen je drei Elemente mitspielen, so dafs dieselben zusammen durch $2 + 2 = 4$, d.\,i. eine gerade Anzahl allgemeiner Vertauschungen ersetzt werden können.

12. \textit{Classification der Substitution durch Abzählung der Derangements.} Ein weiteres Mittel die Klasse einer Substitution $S$ zu bestimmen, liefert auch die Abzählung der Derangements in den beiden Permutationen von $S$. Es gilt nämlich der Satz:

\glqq Gehört eine Substitution der ersten Klasse an, so ist die Differenz der Derangements in der obern und untern Permutation eine gerade, gehört sie der zweiten Klasse an, eine ungerade Zahl.\grqq{}


% ===== pdf idx 29, printed 15.0 =====
% !PRINTED 15
\weberpage{15}
Wir können diese Behauptung mit den bisherigen Kriterien leicht in Einklang bringen. Jede einfache Vertauschung an der obern Permutation hat nämlich eine Vermehrung oder Verminderung der Anzahl der Derangements um eine und nur eine Einheit zur Folge. Läfst sich nun die Substitution durch $p$ einfache Vertauschungen ersetzen, so wird auch die Zahl der Derangements bei successiver Anwendung dieser $p$ Operationen stetig um eine Einheit immer vermehrt oder vermindert. Die algebraische Summe dieser Aenderungen ist also gleichzeitig mit $p$ gerade oder ungerade; und damit ist die Behauptung erwiesen.

Um die Klasse zu charakterisiren, welcher eine Substitution angehört, führen wir die Gröfse $\varepsilon_S$ ein, welche wir den Modul der Substitution nennen wollen. Wir bestimmen dann denselben so, dafs er, wenn $S$ der ersten Klasse angehört, den Werth $\varepsilon_S = +1$, im andern Falle den Werth $\varepsilon_S = -1$ habe.

Die oben erwähnte Methode der Klassenbestimmung wollen wir noch am vorletzten Beispiel von Nr. 11 erläutern.

Die gegebene Substitution ist:
\[
S = \begin{pmatrix} 9 & 0 & 8 & 4 & 3 & 5 & 2 & 1 & 7 & 6 \\ 1 & 6 & 5 & 0 & 9 & 2 & 8 & 3 & 4 & 7 \end{pmatrix}.
\]
Wir bestimmen zuerst durch die in (3) erwähnte Theilung die Anzahl der Derangements in der obern Permutation.

Trennen wir:
\[
(9\ 0\ 8\ 4\ 3 \mid 5\ 2\ 1\ 7\ 6),
\]
dann ist: (1.) die Zahl der Derangements vor dem Verticalstrich $4 + 2 + 1 = 7$; (2.) hinter demselben gleich $2 + 1 + 1 = 4$; (3.) die Zahl jener, welche die erste Gruppe in Bezug auf die zweite bildet, gleich $9 + 0 + 8 + 4 + 3 - \dfrac{5 \cdot 6}{2} = 9$; also die Summe aller $= 20$. % ?? uncertain: this "20" is printed in an unusual bold/retouched numeral, cannot fully confirm the glyph

Ebenso findet man die Zahl der Derangements in der untern Permutation $= 13$. % ?? uncertain: the digit "3" of "13" shows an odd cursive stroke, possibly ink retouching in the scan
Die Differenz beider Anzahlen ist $7$, was mit der dort gefundenen Zahl allgemeiner Vertauschungen übereinstimmt. Die Substitution gehört also zur zweiten Klasse und hat somit den Modul $\varepsilon_S = -1$.

\begin{center}
\S.\,2.\\
\gesperrt{Begriff der Determinanten.}
\end{center}

13. \textit{Verwendung der Substitutionen zur Herstellung aller Glieder einer Determinante.} Es seien gegeben die $n^2$ Elemente:
\[
a_{11}, a_{12}, \ldots a_{1n}; \ a_{21}, a_{22}, \ldots a_{2n}; \ \ldots a_{i1}, a_{i2} \ldots a_{ik} \ldots a_{in}; \ \ldots a_{n1}, a_{n2} \ldots a_{nn},
\]
die wir also durch zwei angehängte Zeiger, \textit{Doppelindices} genannt, von einander unterscheiden wollen. Jeder Substitution, gebildet aus den
%CONT>>


% ===== pdf idx 30, printed 16.0 =====
% !PRINTED 16
\weberpage{16}
%<<CONT
$n$ einfachen Indices $1, 2, 3, \ldots n$, kann man ein Product $G$, gebildet aus $n$ Elementen $a_{ik}$ zuordnen, wie es die folgenden Beispiele lehren mögen.

$\alpha$) Sei $n = 6$; die 36 Elemente sind alsdann:
\[
a_{11}, a_{12} \ldots a_{16}; \ a_{21}, a_{22} \ldots a_{26}; \ \ldots a_{i1}, a_{i2} \ldots a_{i6}; \ \ldots a_{61}, a_{62} \ldots a_{66};
\]
ferner sei gegeben die Substitution
\[
S = \begin{pmatrix} 5 & 6 & 2 & 1 & 3 & 4 \\ 3 & 6 & 2 & 4 & 1 & 5 \end{pmatrix},
\]
mit dem Modul
\[
\varepsilon_S = -1,
\]
da die Differenz der Derangements $9 - 8 = 1$ ist. Wir lassen der Substitution $S$ das Product entsprechen
\[
G = -a_{53} a_{66} a_{22} a_{14} a_{31} a_{45}.
\]
Da die Doppelindices der ausgewählten Elemente $a_{ik}$ gerade die untereinander stehenden Zahlenpaare der Substitution $S$ sind, so ist klar, dafs allen $N$ Ausdrücken für ein und dieselbe Substitution immer ein und dasselbe Product $G$ auf diese Art zugeordnet ist. Dasselbe hängt also lediglich von der Substitution ab und nicht von der speciellen Permutation, auf welche diese etwa angewendet wurde.

$\beta$) $n = 5$,
\[
S = \begin{pmatrix} 4 & 1 & 2 & 5 & 3 \\ 1 & 2 & 5 & 3 & 4 \end{pmatrix}, \ \text{also} \ \varepsilon_S = +1;
\]
daher $G = +a_{41} a_{12} a_{25} a_{53} a_{34}$.

$\gamma$) $n = 7$,
\[
S = \begin{pmatrix} 0 & 1 & 3 & 6 & 5 & 4 & 2 \\ 1 & 0 & 6 & 4 & 2 & 3 & 5 \end{pmatrix}, \ \text{also} \ \varepsilon_S = -1;
\]
daher $G = -a_{01} a_{10} a_{36} a_{64} a_{52} a_{43} a_{25}$.

Den beiden Hauptklassen der Substitutionen $S$ entsprechen also auch zwei Hauptklassen von Producten $G$, die erste mit positivem Vorzeichen, die zweite mit negativem. Die Aufgabe, alle Producte $G$ zu bilden, können wir lösen, indem wir zunächst alle Substitutionen $S$ bestimmen. Dabei können wir unter anderm entweder die gut geordnete Permutation als die obere wählen — das liefert eine erste Art die Substitutionen aufzustellen; oder aber sie als die untere verwenden — das liefert eine zweite Art.

Haben wir alsdann den Modul $\varepsilon_S$ jeder Substitution ermittelt, so ordnen wir endlich jeder das Product jener Elemente $a_{ik}$ zu, deren Doppelindices gerade die untereinander stehenden Zahlenpaare in $S$ bilden.

Wir wollen in den folgenden Beispielen die Producte $G$ für $n = 2$, $n = 3$, $n = 4$ berechnen, indem wir dabei die Substitutionen nach der ersten Art aufstellen.


% ===== pdf idx 31, printed 17.0 =====
% !PRINTED 17
\weberpage{17}
% [column spec widened tabular 6->10] a row carried more cells than the spec declared,
% which is a fatal alignment error. The spec was widened rather than a cell
% deleted, so no transcribed content is lost; whether Weber's grid really has
% that many columns is a question for a pass with the scan open.

\[
\alpha)\quad n = 2.
\qquad
\left\{
\begin{aligned}
&\textit{Substitutionen:}\quad \begin{pmatrix}1 & 2\\ 1 & 2\end{pmatrix}, \quad \begin{pmatrix}1 & 2\\ 2 & 1\end{pmatrix}.\\[4pt]
&\textit{Modul } \varepsilon_S\textit{:} \qquad\quad {}+1, \qquad -1.\\[4pt]
&\textit{Product } G\textit{:} \qquad {}+a_{11}a_{22}, \qquad -a_{12}a_{21}.
\end{aligned}
\right.
\]

$\beta)$ \quad $n = 3.$

\[
\begin{aligned}
&\textit{Substitutionen:}\quad \begin{pmatrix}1&2&3\\1&2&3\end{pmatrix} \begin{pmatrix}1&2&3\\1&3&2\end{pmatrix} \begin{pmatrix}1&2&3\\2&1&3\end{pmatrix} \begin{pmatrix}1&2&3\\2&3&1\end{pmatrix} \begin{pmatrix}1&2&3\\3&1&2\end{pmatrix} \begin{pmatrix}1&2&3\\3&2&1\end{pmatrix}.\\
&\textit{Modul } \varepsilon_S\textit{:}\qquad {}+1 \qquad -1 \qquad -1 \qquad {}+1 \qquad {}+1 \qquad -1.\\
&\textit{Producte } G\textit{:}\quad {}+a_{11}a_{22}a_{33}, \quad -a_{11}a_{23}a_{32}, \quad -a_{12}a_{21}a_{33}, \quad {}+a_{12}a_{23}a_{31},\\
&\hphantom{\textit{Producte } G\textit{:}}\quad {}+a_{13}a_{21}a_{32}, \quad -a_{13}a_{22}a_{31}.
\end{aligned}
\]

$\gamma)$ \quad $n = 4.$

\begin{center}
\begin{tabular}{c|c|l||c|c|lcccc}
Substit. & Modul & Product $G_i$ & Substit. & Modul & Product $G_i$ \\
\hline
$\begin{pmatrix}1&2&3&4\\1&2&3&4\end{pmatrix}$ & $+1$ & $+a_{11}a_{22}a_{33}a_{44}$ & $\begin{pmatrix}1&2&3&4\\3&1&2&4\end{pmatrix}$ & $+1$ & $+a_{13}a_{21}a_{32}a_{44}$ \\
$\begin{pmatrix}1&2&3&4\\1&2&4&3\end{pmatrix}$ & $-1$ & $-a_{11}a_{22}a_{34}a_{43}$ & $\begin{pmatrix}1&2&3&4\\3&1&4&2\end{pmatrix}$ & $-1$ & $-a_{13}a_{21}a_{34}a_{42}$ \\
$\begin{pmatrix}1&2&3&4\\1&3&2&4\end{pmatrix}$ & $-1$ & $-a_{11}a_{23}a_{32}a_{44}$ & $\begin{pmatrix}1&2&3&4\\3&2&1&4\end{pmatrix}$ & $-1$ & $-a_{13}a_{22}a_{31}a_{44}$ \\
$\begin{pmatrix}1&2&3&4\\1&3&4&2\end{pmatrix}$ & $+1$ & $+a_{11}a_{23}a_{34}a_{42}$ & $\begin{pmatrix}1&2&3&4\\3&2&4&1\end{pmatrix}$ & $+1$ & $+a_{13}a_{22}a_{34}a_{41}$ \\
$\begin{pmatrix}1&2&3&4\\1&4&2&3\end{pmatrix}$ & $+1$ & $+a_{11}a_{24}a_{32}a_{43}$ & $\begin{pmatrix}1&2&3&4\\3&4&1&2\end{pmatrix}$ & $+1$ & $+a_{13}a_{24}a_{31}a_{42}$ \\
$\begin{pmatrix}1&2&3&4\\1&4&3&2\end{pmatrix}$ & $-1$ & $-a_{11}a_{24}a_{33}a_{42}$ & $\begin{pmatrix}1&2&3&4\\3&4&2&1\end{pmatrix}$ & $-1$ & $-a_{13}a_{24}a_{32}a_{41}$ \\
\hline
$\begin{pmatrix}1&2&3&4\\2&1&3&4\end{pmatrix}$ & $-1$ & $-a_{12}a_{21}a_{33}a_{44}$ & $\begin{pmatrix}1&2&3&4\\4&1&2&3\end{pmatrix}$ & $-1$ & $-a_{14}a_{21}a_{32}a_{43}$ \\
$\begin{pmatrix}1&2&3&4\\2&1&4&3\end{pmatrix}$ & $+1$ & $+a_{12}a_{21}a_{34}a_{43}$ & $\begin{pmatrix}1&2&3&4\\4&1&3&2\end{pmatrix}$ & $+1$ & $+a_{14}a_{21}a_{33}a_{42}$ \\
$\begin{pmatrix}1&2&3&4\\2&3&1&4\end{pmatrix}$ & $+1$ & $+a_{12}a_{23}a_{31}a_{44}$ & $\begin{pmatrix}1&2&3&4\\4&2&1&3\end{pmatrix}$ & $+1$ & $+a_{14}a_{22}a_{31}a_{43}$ \\
$\begin{pmatrix}1&2&3&4\\2&3&4&1\end{pmatrix}$ & $-1$ & $-a_{12}a_{23}a_{34}a_{41}$ & $\begin{pmatrix}1&2&3&4\\4&2&3&1\end{pmatrix}$ & $-1$ & $-a_{14}a_{22}a_{33}a_{41}$ \\
$\begin{pmatrix}1&2&3&4\\2&4&1&3\end{pmatrix}$ & $-1$ & $-a_{12}a_{24}a_{31}a_{43}$ & $\begin{pmatrix}1&2&3&4\\4&3&1&2\end{pmatrix}$ & $-1$ & $-a_{14}a_{23}a_{31}a_{42}$ \\
$\begin{pmatrix}1&2&3&4\\2&4&3&1\end{pmatrix}$ & $+1$ & $+a_{12}a_{24}a_{33}a_{41}$ & $\begin{pmatrix}1&2&3&4\\4&3&2&1\end{pmatrix}$ & $+1$ & $+a_{14}a_{23}a_{32}a_{41}$ \\
\end{tabular}
\end{center}

14. \textit{Definition der Determinante.} Man nennt nun die Summe aller Producte $G$, welche sich auf diese Weise durch die Substitutionen von $n$ Indices bilden lassen: die Determinante des Systemes der $n^2$
%CONT>>


% ===== pdf idx 32, printed 18.0 =====
% !PRINTED 18
\weberpage{18}
%<<CONT
Elemente $a_{ik}$. Eine solche Determinante, die wir mit $\Delta$ bezeichnen wollen, ist also dargestellt in dem Ausdrucke

\weq{1}{\Delta = \sum_{i=1}^{i=N} G_i,}

wobei die $G_i$ in der Weise gewonnen werden, dafs man zuerst alle möglichen Substitutionen der $n$ Indices bildet, den zugehörigen Modul bestimmt, und nun jeder Substitution das entsprechende $G$ zuordnet.

Statt nun aber eine Determinante durch die in (1) angegebene Summe auszudrücken, wählt man meist die folgende Schreibweise

\weq{2}{\Delta = \begin{vmatrix} a_{11} & a_{12} & a_{13} & \ldots & a_{1n} \\ a_{21} & a_{22} & \ldots & \ldots & a_{2n} \\ a_{31} & \cdot & \cdot & \cdot & \cdot & \cdot & \cdot & \cdot & \cdot \\ \cdot & \cdot & \cdot & \cdot & \cdot & \cdot & \cdot & \cdot & \cdot \\ a_{n1} & a_{n2} & \ldots & a_{nn} \end{vmatrix}.}

So ist die Determinante der in Beispiel 4 des vorigen Absatzes entwickelten Glieder gegeben durch

\weq{3}{\Delta = \begin{vmatrix} a_{11} & a_{12} & a_{13} & a_{14} \\ a_{21} & a_{22} & a_{23} & a_{24} \\ a_{31} & a_{32} & a_{33} & a_{34} \\ a_{41} & a_{42} & a_{43} & a_{44} \end{vmatrix} = \Sigma G_i.}

Der rechts befindliche Ausdruck ist nur ein Schema, eine Form für die Anordnung der Elemente, und als solche Form belegen wir sie nicht mit dem Namen Determinante, sondern mit dem Namen \textit{Matrix}. Den wirklichen Werth derselben stellt dann der Ausdruck

\weq{1}{\Delta = \Sigma G_i}

dar. Es ist nicht nöthig, dass in einem solchen Schema stets ebenso viele horizontale Zeilen von Elementen auftreten als verticale Reihen (Colonnen). So kann man auch die Matrix bilden

\[
\begin{vmatrix} a_{11} & a_{12} & \ldots & \ldots & a_{1n} \\ a_{21} & a_{22} & \ldots & \ldots & a_{2n} \\ \cdot & \cdot & \cdot & \cdot & \cdot & \cdot & \cdot & \cdot & \cdot & \cdot \\ a_{m2} & a_{m1} & \ldots & a_{mn} \end{vmatrix}
\]
% [sic] printed order "a_{m2} a_{m1}" (not a_{m1} a_{m2})

wo $m \leqq n$ ist; in ihr ordnen die Elemente sich in $n$ Reihen, dagegen in $m$ Zeilen. In diesem Fall aber repräsentirt die Matrix keinen numerischen Werth, während, wenn die Zahl der Zeilen mit denen der Colonnen übereinstimmt, dieselbe einen solchen besitzt, den man mittels der Summe (1) darstellen kann. Die Producte $G_i$ nennt man Glieder,
%CONT>>


% ===== pdf idx 33, printed 19.0 =====
% !PRINTED 19
\weberpage{19}
%<<CONT
ihre Factoren $a_{ik}$ Elemente der Determinante. Der erste Index eines solchen Elementes zeigt an, welcher Zeile, der zweite, welcher Vertical-Reihe oder Colonne es angehört. Man sagt, die Determinante ist $n$-gliedrig oder $n^{\text{ter}}$ Ordnung, so bald sie $n^2$ Elemente besitzt.

15. \textit{Berechnung des Werthes einer Determinante aus ihrer Matrix.} Ist eine Matrix von $n$ Zeilen und $n$ Colonnen gegeben, so fragt sich, wie bestimmt man daraus den Werth derselben, also die zugehörige Determinante, d. h. wie findet man die einzelnen Glieder $G_i$ einer solchen Summe

\weq{1}{\Delta = \Sigma G_i.}

Nun sieht man aber leicht — man vergleiche z.\,B. die 24 Glieder im Beispiel $\gamma)$ Nr. 13 — dass, weil einmal alle ersten Indices der $n$ Elemente eines Gliedes von einander verschieden sind, dann aber, weil auch keine zwei der zweiten Indices in einem Gliede übereinstimmen, jede Zeile und jede Colonne nur durch ein einziges Element in einem Gliede vertreten sein kann. Also hat man das Gesetz: Jedes Glied enthält ein und nur ein Element jeder Zeile und jeder Colonne, und wir erhalten also alle $G_i$, wenn wir sämmtliche Producte bilden, in denen dieses Gesetz beachtet ist. Wir haben aber bereits in Nr. 13 gesehen, wie man systematisch zu diesen Producten gelangen kann, und können nach dieser an den Beispielen $\alpha)$, $\beta)$, $\gamma)$ entwickelten Methode direct die Summe aller Glieder einer Determinante schematisch anschreiben. Es ist also:

\[
\alpha)\quad \Delta_2 = \begin{vmatrix} a_{11} & a_{12} \\ a_{21} & a_{22} \end{vmatrix} = \sum_{i=1}^{i=2!} \pm G_i = a_{11}a_{22} - a_{12}a_{21},
\]

wenn wir die Glieder aus solchen Substitutionen ableiten, bei denen die gut geordnete Permutation sich oben befindet; oder:
%CONT>>


% ===== pdf idx 34, printed 20.0 =====
% !PRINTED 20
\weberpage{20}
%<<CONT
\[
= a_{11}a_{22} - a_{21}a_{12},
\]

wenn bei den entsprechenden Substitutionen die gut geordnete Permutation unten ist.

\[
\beta)\quad \Delta_3 = \begin{vmatrix} a_{11} & a_{12} & a_{13} \\ a_{21} & a_{22} & a_{23} \\ a_{31} & a_{32} & a_{33} \end{vmatrix}
\]

\[
\begin{aligned}
&= \sum_{i=1}^{i=3!} \pm G_i = {}+a_{11}a_{22}a_{33} - a_{11}a_{23}a_{32} - a_{12}a_{21}a_{33}\\
&\hphantom{= \sum_{i=1}^{i=3!} \pm G_i =} {}+a_{12}a_{23}a_{31} + a_{13}a_{21}a_{32} - a_{13}a_{22}a_{31},
\end{aligned}
\]

oder

\[
\begin{aligned}
&= {}+a_{11}a_{22}a_{33} - a_{11}a_{32}a_{23} - a_{21}a_{12}a_{33}\\
&\hphantom{=} {}+a_{21}a_{32}a_{13} + a_{31}a_{12}a_{23} - a_{31}a_{22}a_{13},
\end{aligned}
\]

wo wieder in der ersten Summe Substitutionen mit geordneter Permutation oben, in der zweiten Summe solche mit gutgeordneten Permutanten benutzt wurden. % [sic] "Permutanten"

\[
\gamma)\quad \Delta_4 = \begin{vmatrix} a_{11} & a_{12} & a_{13} & a_{14} \\ a_{21} & a_{22} & a_{23} & a_{24} \\ a_{31} & a_{32} & a_{33} & a_{34} \\ a_{41} & a_{42} & a_{43} & a_{44} \end{vmatrix} = \sum_{i=1}^{i=4!=24} + G_i.
\]

Die Glieder dieser Summe, nach der ersten Art ermittelt, sind bereits in Beispiel $\gamma)$ Nr. 13 angegeben.

Wir wollen sie hier noch in der Gestalt geben, wie sie die zweite Art liefert. Man erhält:

\[
\begin{aligned}
\Delta_4 = \Sigma + G_i &= a_{11}a_{22}a_{33}a_{44} - a_{21}a_{12}a_{33}a_{44} + a_{31}a_{12}a_{23}a_{44} - a_{41}a_{12}a_{23}a_{34}\\
&- a_{11}a_{22}a_{43}a_{34} + a_{21}a_{12}a_{43}a_{34} - a_{31}a_{12}a_{43}a_{24} + a_{41}a_{12}a_{33}a_{24}\\
&- a_{11}a_{32}a_{23}a_{44} + a_{21}a_{32}a_{13}a_{44} - a_{31}a_{22}a_{13}a_{44} + a_{41}a_{22}a_{13}a_{34}\\
&+ a_{11}a_{32}a_{43}a_{24} - a_{21}a_{32}a_{43}a_{14} + a_{31}a_{22}a_{34}a_{14} - a_{41}a_{22}a_{33}a_{14}\\
&+ a_{11}a_{42}a_{23}a_{34} - a_{21}a_{24}a_{31}a_{43} + a_{31}a_{42}a_{13}a_{24} - a_{41}a_{32}a_{13}a_{24}\\
&- a_{11}a_{42}a_{33}a_{24} + a_{21}a_{24}a_{33}a_{41} - a_{31}a_{42}a_{23}a_{14} + a_{41}a_{32}a_{23}a_{14}.
\end{aligned}
\]

Dabei sind die Glieder, welche der nämlichen Permutationsklasse angehören, immer unter einander geschrieben.

Mit Hilfe dieser Formeln $\alpha)$, $\beta)$, $\gamma)$ können wir nun leicht den Werth einer 2-, 3- oder 4gliedrigen Determinante bestimmen.

\textit{Beispiele:}

1.\quad $\Delta = \begin{vmatrix} 3 & 5 \\ 6 & 12 \end{vmatrix} = 3\cdot 12 - 6\cdot 5 = 6.$

2.\quad $\Delta = \begin{vmatrix} 4 & 1 \\ 0 & 2 \end{vmatrix} = 4\cdot 2 - 1\cdot 0 = 8.$

3.\quad $\Delta = \begin{vmatrix} 0 & 3 \\ 6 & 1 \end{vmatrix} = 0 - 18 = -18.$

\[
\begin{aligned}
4.\quad \Delta = \begin{vmatrix} 2 & 4 & 1 \\ 5 & 1 & -3 \\ 4 & 6 & 7 \end{vmatrix} &= 2\cdot 1\cdot 7 - (-2\cdot 3\cdot 6) - 5\cdot 4\cdot 7\\
&+ 5\cdot 6\cdot 1 + (4\cdot 4\cdot{-3}) - 4\cdot 1\cdot 1 = 14\\
&+ 36 - 140 + 30 - 48 - 4 = -112.
\end{aligned}
\]

5.\quad $\Delta = \begin{vmatrix} 1 & 0 & -2 \\ 0 & 3 & 5 \\ 2 & 4 & 4 \end{vmatrix} = 1\cdot 3\cdot 4 - 1\cdot 4\cdot 5 + 2\cdot 3\cdot 2 = {}+4.$

\[
6.\quad \Delta = \begin{vmatrix} 3 & 1 & 1 & 2 \\ 1 & 5 & 0 & 3 \\ 2 & -2 & 1 & 6 \\ 0 & 4 & -5 & 3 \end{vmatrix}
\]
%CONT>>


% ===== pdf idx 35, printed 21.0 =====
% !PRINTED 21
\weberpage{21}
%<<CONT
\[
\begin{array}{l}
= 3\cdot 5\cdot 1\cdot 3 - 1\cdot 1\cdot 1\cdot 3 + 2\cdot 1\cdot 0\cdot 3 - 0\cdot 1\cdot 0\cdot 6\\
- 3\cdot 5\cdot(-5)\cdot 6 + 1\cdot 1\cdot(-5)\cdot 6 - 2\cdot 1\cdot(-5)\cdot 3 + 0\cdot 1\cdot 1\cdot 3\\
- 3\cdot(-2)\cdot 0\cdot 3 + 1\cdot(-2)\cdot 1\cdot 3 - 2\cdot 5\cdot 1\cdot 3 + 0\cdot 5\cdot 1\cdot 6\\
+ 3\cdot(-2)\cdot(-5)\cdot 3 - 1\cdot(-2)\cdot(-5)\cdot 2 + 2\cdot 5\cdot(-5)\cdot 2 - 0\cdot 5\cdot 1\cdot 2\\
+ 3\cdot 4\cdot 0\cdot 6 - 1\cdot 4\cdot 1\cdot 6 + 2\cdot 4\cdot 1\cdot 3 - 0\cdot(-2)\cdot 1\cdot 3\\
- 3\cdot 4\cdot 1\cdot 3 + 1\cdot 4\cdot 1\cdot 2 - 2\cdot 4\cdot 0\cdot 2 + 0\cdot(-2)\cdot 0\cdot 2\\
= 45 + 450 + 90 - 36 - 3 - 30 - 6 - 20 - 24 + 8\\
+ 30 - 30 - 100 + 24 = 647 - 249 = 398.
\end{array}
\]
% ?? uncertain: the exact digit grouping of the multiplication-dot expansion in the block above is hard to read with full confidence; best reading transcribed as printed

\[
7.\quad \Delta = \begin{vmatrix} 1 & 2 & 3 & 4\\ 0 & 1 & 2 & 3\\ 0 & 5 & 1 & 2\\ 0 & 0 & 2 & 1 \end{vmatrix} = 1\cdot 1\cdot 1\cdot 1 - 1\cdot 1\cdot 2\cdot 2 - 1\cdot 5\cdot 2\cdot 1 + 1\cdot 5\cdot 2\cdot 3 = 1 - 4 - 10 + 30 = 17.
\]

16. \textit{Methode zur Aufstellung von Lehrsätzen über Determinanten.} Diese Art der Werthbestimmung einer Determinante ist indess im Allgemeinen noch ziemlich mühevoll. Zahlreiche Determinantenlehrsätze werden den Weg wesentlich abkürzen. Um aber zu diesen Lehrsätzen über Determinanten zu gelangen, werden wir stets zwei Matrices von derselben Elementenzahl
\[
\Delta = \begin{vmatrix}
a_{11} & a_{12} & a_{13} & \ldots & a_{1n}\\
a_{21} & a_{22} & \ldots & \ldots & a_{2n}\\
a_{31} & \cdot & \cdot & \cdot & \cdot\\
\cdot & \cdot & \cdot & \cdot & \cdot\\
a_{n1} & a_{n2} & \ldots & & a_{nn}
\end{vmatrix}
\quad\text{und}\quad
D = \begin{vmatrix}
b_{11} & b_{12} & b_{13} & \ldots & b_{1n}\\
b_{21} & b_{22} & \ldots & \ldots & b_{2n}\\
b_{31} & \cdot & \cdot & \cdot & \cdot\\
\cdot & \cdot & \cdot & \cdot & \cdot\\
b_{n1} & b_{n2} & \ldots & & b_{nn}
\end{vmatrix}
\]
vergleichen. Beziehungen zwischen den einzelnen Elementen $a_{ik}$ und $b_{ik}$, werden sodann auch Veranlassung geben zu einer Relation zwischen den Determinanten $\Delta$ und $D$ selbst.

17. \textit{Die Transposition.} Setzen wir fest, dass die Elemente in $D$ und $\Delta$ durch die Bedingung $a_{ik}=b_{ki}$ mit einander verknüpft seien, also dass z.\,B. für $n=3$ die 9 Relationen bestehen:
\[
\begin{aligned}
&a_{11}=b_{11}, \quad a_{12}=b_{21}, \quad a_{13}=b_{31}; \quad a_{21}=b_{12}, \quad a_{22}=b_{22}, \quad a_{32}=b_{23};\\
&a_{31}=b_{13}, \quad a_{32}=b_{23}, \quad a_{33}=b_{33},
\end{aligned}
\]
% [sic] printed "a_{32}=b_{23}" appears twice, once in each row of the relation list above, exactly as printed
so lässt sich die Frage stellen: In welcher Beziehung stehen dann die Werthe $\Delta$ und $D$?

Um dieselbe zu erkennen, bilde ich einmal alle Glieder $G_i$ von $\Delta$, indem ich die Substitutionen zu Grunde lege mit gut geordneter Permutation oben; also für $n=3$ die Substitutionen:
\[
\begin{pmatrix}1&2&3\\1&2&3\end{pmatrix}, \quad
\begin{pmatrix}1&2&3\\1&3&2\end{pmatrix}, \quad
\begin{pmatrix}1&2&3\\2&1&3\end{pmatrix}, \quad
\begin{pmatrix}1&2&3\\2&3&1\end{pmatrix}, \quad
\begin{pmatrix}1&2&3\\3&1&2\end{pmatrix}, \quad
\begin{pmatrix}1&2&3\\3&2&1\end{pmatrix};
\]
%CONT>>


% ===== pdf idx 36, printed 22.0 =====
% !PRINTED 22
\weberpage{22}
%<<CONT
und folglich die Glieder $G_i$:
\[
+a_{11}a_{22}a_{33}, \quad -a_{11}a_{23}a_{32}, \quad -a_{12}a_{21}a_{33}, \quad +a_{12}a_{23}a_{31},\\
+a_{13}a_{21}a_{32}, \quad -a_{13}a_{22}a_{31}.
\]
Sodann stelle ich die Glieder $\overline{G_i}$ von $\Delta$ auf unter Benützung der Substitutionen, welche die gut geordneten Permutationen unten haben; also die Substitutionen:
\[
\begin{pmatrix}1&2&3\\1&2&3\end{pmatrix}, \quad
\begin{pmatrix}1&3&2\\1&2&3\end{pmatrix}, \quad
\begin{pmatrix}2&1&3\\1&2&3\end{pmatrix}, \quad
\begin{pmatrix}2&3&1\\1&2&3\end{pmatrix}, \quad
\begin{pmatrix}3&1&2\\1&2&3\end{pmatrix}, \quad
\begin{pmatrix}3&2&1\\1&2&3\end{pmatrix},
\]
und folglich die Glieder $\overline{G_i}$:
\[
+b_{11}b_{22}b_{33}, \quad -b_{11}b_{32}b_{23}, \quad -b_{21}b_{12}b_{33}, \quad +b_{21}b_{32}b_{13},\\
+b_{31}b_{12}b_{23}, \quad -b_{31}b_{22}b_{13}.
\]
Setzt man aber in diese Reihe der Glieder $\overline{G_i}$ die Werthe der $b_{ki}=a_{ik}$ ein, so stimmt dieselbe im Modul wie in den Doppelindices der Elemente gerade mit der Reihe der Glieder $G_i$ überein, und es ist daher nothwendig
\weq{1}{\sum_1^6 G_i = \sum_1^6 \overline{G_i}.}
Genau dieselbe Uebereinstimmung beider Reihen der $G_i$ und $\overline{G_i}$ hinsichtlich der Aufeinanderfolge dieser Glieder, ihrer Modul und der Doppelindices ihrer Elemente liefert diese Methode auch im allgemeinen Falle, d.\,h. es ist
\weq{2}{\sum_1^N G_i = \sum_1^N \overline{G_i}.}
Trägt man nun aber auch in die Matrix der Determinante $D$ die aus den $n^2$ Relationen sich bietenden Werthe von $b_{ik}$ ein, so kommt:
\[
D = \begin{vmatrix}
a_{11} & a_{21} & a_{31} & a_{41} & \ldots & a_{n1}\\
a_{12} & a_{22} & a_{32} & \ldots & \ldots & a_{n2}\\
a_{13} & a_{23} & \cdot & \cdot & \cdot & \cdot\\
\cdot & \cdot & \cdot & \cdot & \cdot & \cdot\\
a_{1n} & a_{2n} & \ldots & & \ldots & a_{nn}
\end{vmatrix}.
\]
Vergleicht man diese Matrix mit der Matrix der Determinante
\[
\Delta = \begin{vmatrix}
a_{11} & a_{12} & a_{13} & a_{14} & \ldots & a_{1n}\\
a_{21} & a_{22} & a_{23} & \ldots & \ldots & a_{2n}\\
a_{31} & a_{32} & \cdot & \cdot & \cdot & \cdot\\
\cdot & \cdot & \cdot & \cdot & \cdot & \cdot\\
a_{n1} & a_{n2} & \ldots & & \ldots & a_{nn}
\end{vmatrix}
\]
so erkennt man, dass nur die in der Diagonale dieses Schemas befind-
%CONT>>


% ===== pdf idx 37, printed 23.0 =====
% !PRINTED 23
\weberpage{23}
%<<CONT
lichen Elemente $a_{11}a_{22}\ldots a_{nn}$ ihre alten Plätze behalten haben, während alle übrigen an solchen Stellen in $D$ stehen, die man aus ihrer alten Lage dadurch erhält, dass man die Matrix von $\Delta$ um ihre erste Diagonale $a_{11}a_{22}\ldots a_{nn}$ umklappt. Dabei sind also die Colonnen zu Zeilen, die Zeilen zu Colonnen geworden, und eine solche Vertauschung der Elemente einer Matrix nennt man \textit{Transposition} der Matrix. Also hat man in Verbindung mit Gleichung (1) den Satz:

\begin{center}
\textit{\glqq Der Werth einer Determinante ändert sich nicht, wenn man\\
ihre Matrix transponirt.\grqq{}}
\end{center}

\textit{Beispiele:}

\[
1.\quad \Delta = \begin{vmatrix}3 & 5\\6 & 12\end{vmatrix} = \begin{vmatrix}3 & 6\\5 & 12\end{vmatrix} = 3\cdot 12 - 5\cdot 6 = 6.
\]
2. Transponirt man die 4. Determinante unter den Beispielen zu Nr. 15, so findet man
\[
\Delta = \begin{vmatrix} 2 & 5 & 4\\ 4 & 1 & 6\\ 1 & -3 & 7 \end{vmatrix} = -112.
\]
\[
3.\quad \Delta = \begin{vmatrix} a & b & c & d\\ 0 & e & f & g\\ p & q & r & s\\ 0 & u & v & w \end{vmatrix} = aerw - aevs - aqfw + aqvg + aufs - aurg + pbfw - pbvg - pecw + pevd + pucg - pufd.
\]
Ganz ebenso liefert aber
\[
\Delta = \begin{vmatrix} a & 0 & p & 0\\ b & e & q & u\\ c & b & r & v\\ d & g & s & w \end{vmatrix} = aerw - aesv - afqw + afsu + agqv - agur + bfpw - bpvg - cepw + cgpu + depv - dfpu.
\]
18. \textit{Andere übersichtliche Bezeichnungen für eine Determinante.} Die Elemente der \textit{ersten Diagonale} $a_{11}, a_{22}, a_{33}\ldots a_{nn}$ behalten bei der Transposition, wie schon erwähnt, ihre Stellen bei. Sie bilden ein ausgezeichnetes Glied in der Summe $\Sigma G_i$, das man nach der Stellung seiner Elemente in der Matrix \textit{Diagonalglied} oder auch \textit{Hauptglied} nennt, aus dem alle andern Glieder durch Permutation aller ersten, oder auch der zweiten Indices abgeleitet werden können. Man schreibt daher eine Determinante auch in der Form
\[
\Delta = \Sigma \pm a_{11}a_{22}a_{33}\cdots a_{nn}.
\]
Häufig macht man auch noch von einer andern Schreibweise Gebrauch, die gleichfalls in einer Auszeichnung des Diagonalgliedes beruht und insbesondere dann angewendet wird, wenn die Elemente einer Matrix
%CONT>>


% ===== pdf idx 38, printed 24.0 =====
% !PRINTED 24
\weberpage{24}
%<<CONT
nicht durch Doppelindices, sondern durch einfache auseinandergehalten sind, dafür aber an Stelle eines zweiten Index ein anderer Buchstabe getreten ist.

So bezeichnet man die Determinanten
\[
\Delta_1 = \begin{vmatrix} a_1 & a_2\\ b_1 & b_2 \end{vmatrix}
\]
mit Rücksicht auf ihre Diagonale durch $\Delta = (a_1 b_2)$,
\[
\Delta_2 = \begin{vmatrix} a_1 & a_2 & a_3\\ b_1 & b_2 & b_3\\ c_1 & c_2 & c_3 \end{vmatrix}
\]
mit $\Delta = (a_1 b_2 c_3)$.

Kann kein Missverständniss eintreten, so lässt man selbst die einfachen Indices weg, und schreibt beispielsweise statt
\[
\Delta = \begin{vmatrix} a_1 & a_2\\ b_1 & b_2 \end{vmatrix} = a_1 b_2 - b_1 a_2,
\]
einfach
\[
\Delta = (ab).
\]
19. \textit{Vertauschung zweier Zeilen.} Seien wieder die beiden Determinanten $\Delta$ und $D$ gegeben; zwischen den Elementen sollen alsdann die $n^2$ Relationen bestehen:
\[
1)\ \begin{cases} a_{1i} = b_{2i},\\ a_{2i} = b_{1i}, \end{cases}
\]
\[
2)\quad a_{ik} = b_{ik}, \ i>2,
\]
d.\,h. es soll sein z.\,B. für $n=3$
\[
1)\quad a_{11} = b_{21}, \quad a_{12} = b_{22}, \quad a_{13} = b_{23};
\]
\[
2)\quad a_{21} = b_{11}, \quad a_{22} = b_{12}, \quad a_{23} = b_{13};
\]
\[
3)\quad a_{31} = b_{31}, \quad a_{32} = b_{32}, \quad a_{33} = b_{33}.
\]
Wir bilden wiederum sowohl die Glieder $G_i$ als die $\overline{G_i}$, und zwar die $G_i$ unter Zugrundelegung der Substitutionen, deren obere Permutation die gut geordnete ist, die $\overline{G_i}$ aber durch Substitutionen, deren obere Permutation aus der gut geordneten durch die einfache Vertauschung (1 2) hervorgeht. So ist also in unserm Beispiel für $n=3$ die Substitutionenreihe für die Glieder $G_i$:
\[
\begin{pmatrix}1&2&3\\1&2&3\end{pmatrix}, \quad
\begin{pmatrix}1&2&3\\1&3&2\end{pmatrix}, \quad
\begin{pmatrix}1&2&3\\2&1&3\end{pmatrix}, \quad
\begin{pmatrix}1&2&3\\2&3&1\end{pmatrix}, \quad
\begin{pmatrix}1&2&3\\3&1&2\end{pmatrix}, \quad
\begin{pmatrix}1&2&3\\3&2&1\end{pmatrix},
\]
für die Glieder $\overline{G_i}$:
\[
\begin{pmatrix}2&1&3\\1&2&3\end{pmatrix}, \quad
\begin{pmatrix}2&1&3\\1&3&2\end{pmatrix}, \quad
\begin{pmatrix}2&1&3\\2&1&3\end{pmatrix}, \quad
\begin{pmatrix}2&1&3\\2&3&1\end{pmatrix}, \quad
\begin{pmatrix}2&1&3\\3&1&2\end{pmatrix}, \quad
\begin{pmatrix}2&1&3\\3&2&1\end{pmatrix}.
\]


% ===== pdf idx 39, printed 25.0 =====
% !PRINTED 25
\weberpage{25}
Bestimmt man nun zu jeder Substitution den Modul und bildet hierauf die entsprechenden Glieder $G_i$ und $\overline{G}_i$, so bemerkt man, dass immer nur die beiden ersten Factoren $a_{ik}$, bezw. $b_{ik}$ in ihren Doppelindices von einander abweichen bei je zwei Gliedern $G_i$ und $\overline{G}_i$, welche dieselbe Permutation der zweiten Indices haben, und dass, weil die Anfangssubstitutionen schon verschiedene Moduln haben und alle sp\"ateren durch die n\"amliche Anzahl einfacher Vertauschungen entstehen, die Vorzeichen je zweier solcher Glieder entgegengesetzt sein m\"ussen. So sind in unserm Beispiel die Glieder
\[
\begin{aligned}
G_i &= + a_{11} a_{22} a_{33}, - a_{11} a_{23} a_{32}, - a_{12} a_{21} a_{33}, + a_{12} a_{21} a_{31}, \\
&\quad - a_{13} a_{21} a_{32}, + a_{13} a_{22} a_{31}, \\
\overline{G}_i &= - b_{21} b_{12} b_{33}, + b_{21} b_{13} b_{32}, + b_{22} b_{11} b_{33}, - b_{22} b_{13} b_{31}, \\
&\quad + b_{23} b_{11} b_{32}, - b_{33} b_{12} b_{31}.
\end{aligned}
\]
Tr\"agt man aber in die Reihe der $\overline{G}_i$ die Werthe der $b_{ik}$ gem\"ass den Bedingungen ein, so erhalten gerade die beiden ersten Factoren jedes Gliedes $\overline{G}_i$ gleichfalls dieselben Doppelindices, wie die entsprechenden Glieder $G_i$, nur die entgegengesetzten Vorzeichen bleiben; nach dieser Substitution besteht also allgemein die Relation
\[
G_i = - \overline{G}_i,
\]
und folglich auch
\[
\sum G_i = - \sum \overline{G}_i,
\]
oder
\[
\varDelta = - D.
\]
Ersetzt man nun aber auch in der Matrix $D$ die Elemente $b_{ik}$ durch ihre entsprechenden Elemente $a_{ik}$, so entsteht eine neue Matrix, die sich von der Matrix zu $\varDelta$ nur dadurch unterscheidet, dass die erste Zeile mit der zweiten vertauscht ist. W\"ahrend also in unserm Beispiele
\[
\varDelta = \begin{vmatrix} a_{11} & a_{12} & a_{13} \\ a_{21} & a_{22} & a_{23} \\ a_{31} & a_{32} & a_{33} \end{vmatrix}
\]
ist, wird
\[
D = \begin{vmatrix} a_{21} & a_{22} & a_{23} \\ a_{11} & a_{12} & a_{13} \\ a_{31} & a_{32} & a_{33} \end{vmatrix}.
\]
Wir sehen also: \glqq Vertauscht man in einer Matrix die erste mit der zweiten Zeile, so \"andert die Determinante ihr Vorzeichen.\grqq{} Ebenso kann man aber zeigen, dass die Vertauschung zweier beliebiger Zeilen, oder mit R\"ucksicht auf den ersten Lehrsatz (16) zweier beliebiger Colonnen die n\"amliche Wirkung haben muss. Daher der Lehrsatz:
%CONT>>


% ===== pdf idx 40, printed 26.0 =====
% !PRINTED 26
\weberpage{26}
%<<CONT
\glqq\textit{Vertauscht man in einer Matrix zwei beliebige Zeilen, oder zwei beliebige Colonnen, so \"andert der Werth der zugeh\"origen Determinante nur das Vorzeichen.}\grqq{}

Beispiel 1:
\[
\varDelta = \begin{vmatrix} p & q & r \\ x & y & z \\ a & b & c \end{vmatrix} = pyc - pbz - xqc + xbr + aqz - ayr,
\]
dagegen:
\[
D = \begin{vmatrix} p & q & r \\ a & b & c \\ x & y & z \end{vmatrix} = pbz - pyc - aqz + ayz + xqc - xbr = - \varDelta.
\]
Beispiel 2:
\[
\varDelta = \begin{vmatrix} A & 0 & 0 & D \\ 0 & F & G & 0 \\ 0 & H & J & C \\ M & 0 & B & L \end{vmatrix} = AFJL - AFBC - AHGL - MFJD + MHGD.
\]
Dagegen:
\[
D = \begin{vmatrix} A & 0 & 0 & D \\ 0 & G & F & 0 \\ 0 & J & H & C \\ M & B & 0 & L \end{vmatrix} = AGHL - AFJL + AFBC + MFJD - MHGD = - \varDelta.
\]
Zusatz: \glqq Besitzt die Matrix zwei gleiche Reihen oder Zeilen, so ist der Werth der zugeh\"origen Determinante Null.\grqq{} Denn wenn man die beiden gleichen Zeilen oder Reihen vertauscht, so muss die Determinante ihr Zeichen \"andern; die Matrix von $\varDelta$ \"andert sich aber nicht; folglich ist einmal $\varDelta = + D$, nach der Vertauschung aber auch $\varDelta = - D$; also $2 \varDelta = 0$, oder $\varDelta = 0$.

20. \textit{Iteration der Vertauschungen zweier Zeilen oder Colonnen.} Wenn man nun eine Reihe von Vertauschungen an einer Matrix $\varDelta$ vorgenommen hat, bald von Zeilen mit Zeilen, bald von Colonnen mit Colonnen, ohne auf die Anzahl der Vertauschungen zu achten, so kann man gleichwohl aus dem Diagonalglied der zuletzt gewonnenen Matrix $D$ sofort ermitteln, ob das Vorzeichen der urspr\"unglichen Determinante sich dabei ge\"andert hat oder nicht. Wir haben eben gezeigt, dass durch eine einfache Vertauschung von Zeilen mit Zeilen, oder Colonnen mit Colonnen, beide Determinanten Glied um Glied noch \"ubereinstimmen m\"ussen, d.\,h. dass keine Glieder nach der Vertauschung entstehen, deren Zusammensetzung aus den Elementen $a_{ik}$ eine andere
%CONT>>


% ===== pdf idx 41, printed 27.0 =====
% !PRINTED 27
\weberpage{27}
%<<CONT
sein k\"onnte, als die vor der Vertauschung. Was sich ge\"andert hat, ist nur das Vorzeichen der Glieder und etwa die Aufeinanderfolge in der Reihe der Glieder. Diese Erscheinung wiederholt sich aber mit jeder Vertauschung, und so sehen wir, dass auch die Glieder der letzten Determinante, die wir erhalten, nicht anders zusammengesetzt sein k\"onnen, als die der ersten. Beide Determinanten m\"ussen bei geeigneter Anordnung Glied um Glied --- vom Vorzeichen abgesehen --- \"ubereinstimmen. Und wenn daher irgend zwei entsprechende Glieder entgegengesetztes Vorzeichen haben, so besitzen es \textit{alle} entsprechenden Glieder. Insbesondere wird das Diagonalglied $\overline{G}_1$ der Matrix $D$ einem Glied $G_i$ der Matrix $\varDelta$ entsprechen. W\"ahrend aber $\overline{G}_1$ in $D$ als Diagonalglied das Vorzeichen $+$ hat, wird $G_i$ in $\varDelta$ im Allgemeinen den noch zu bestimmenden Modul $\varepsilon_S$ besitzen. Ist dieser Modul $\varepsilon_S = + 1$, dann ist auch $D = + \varDelta$; dagegen ist $D = - \varDelta$, wenn $\varepsilon_S = - 1$.

Um aber die Substitution, welche demjenigen Gliede $G_i$ in $\varDelta$ zugeh\"ort, das dem Hauptgliede $\overline{G}_1$ in $D$ entspricht, zu bestimmen, und dadurch auch den betreffenden Modul, brauchen wir nur in $\overline{G}_1$ die Elemente $b$ durch jene $a_{ik}$ zu ersetzen, die ihnen verm\"oge ihrer urspr\"unglichen Pl\"atze entsprechen; aus den Doppelindices des so erhaltenen Gliedes folgt die Substitution und daraus der Modul.

Beispiel. Es sei gegeben:
\[
\varDelta = \begin{vmatrix} \alpha & \beta & \gamma & \delta \\ \varepsilon & \mu & \nu & \vartheta \\ \lambda & \varrho & \varkappa & \iota \\ \varphi & \psi & \chi & \tau \end{vmatrix}.
\]
Durch eine Reihe von Vertauschungen habe man erhalten:
\[
D = \begin{vmatrix} \vartheta & \varepsilon & \nu & \mu \\ \delta & \alpha & \gamma & \beta \\ \tau & \varphi & \chi & \psi \\ \iota & \lambda & \varkappa & \varrho \end{vmatrix}.
\]
Das Diagonalglied in $D$ ist:
\[
\overline{G}_1 = + \vartheta \alpha \chi \varrho.
\]
Dasselbe wird in der urspr\"unglichen Determinante repr\"asentirt durch:
\[
G_i = \varepsilon_S \cdot a_{24} a_{11} a_{43} a_{32}.
\]
Die zugeh\"orige Substitution ist:
\[
S = \begin{pmatrix} 2 & 1 & 4 & 3 \\ 4 & 1 & 3 & 2 \end{pmatrix};
\]
%CONT>>


% ===== pdf idx 42, printed 28.0 =====
% !PRINTED 28
\weberpage{28}
%<<CONT
sie zerf\"allt in den dreigliedrigen Cyclus der zum einfachen $\begin{pmatrix} 2 & 4 & 3 \\ 4 & 3 & 2 \end{pmatrix}$ geh\"ort und den eingliedrigen $(1)$. Die Zahl der allgemeinen Vertauschungen ist also $n - \nu = 4 - 2 = 2$, d.\,h. $\varepsilon_S = + 1$, und also $D = \varDelta$.

21. \textit{Anwendung: Ein beliebiges Element $a_{ik}$ soll durch Vertauschungen an die Stelle von $a_{11}$ bef\"ordert werden.} Die Aufgabe, eine Reihe von Vertauschungen zu machen, tritt insbesondere dann an uns heran, wenn wir gewisse Elemente an besondere Pl\"atze der Matrix, etwa an die durch die Diagonale ausgezeichneten Stellen bringen wollen, ohne dabei --- so weit es m\"oglich ist --- die andern Elemente aus ihrer gegenseitigen Lage und Anordnung zu bringen. Wir benutzen zu dem Zwecke die cyklische Vertauschung.

Haben wir zun\"achst nur ein Element $a_{ik}$ an die erste bisher durch $a_{11}$ ausgef\"ullte Stelle der Diagonalreihe zu bringen, dann vertauschen wir cyklisch die $i^{\text{te}}$ Zeile mit der ersten, ebenso die $k^{\text{te}}$ Colonne mit der ersten Colonne. Dabei ist in der That das Element $a_{ik}$ zum ersten Element in der Diagonale geworden, w\"ahrend sowohl alle Reihen und Colonnen --- je die erste ausgeschlossen ---, als auch insbesondere die Elemente der Diagonale, abgesehen von $a_{ik}$, noch immer gut geordnet sind. Dem Diagonalglied entspricht also die Substitution
\[
S = \begin{pmatrix} i & 1 & 2 & 3 & \ldots & h \\ k & 1 & 2 & 3 & \ldots & h \end{pmatrix},
\]
deren erste Permutation $i - 1$, deren zweite $k - 1$ Derangements aufweist. Ihre Differenz ist $i - k$, daher der Modul $\varepsilon_S = (-1)^{i-k}$, oder besser, da $i - k$ gleichzeitig mit $i + k = \nu$ gerade oder ungerade ist,
\[
\varepsilon_S = (-1)^{i+k} = (-1)^{\nu},
\]
also:
\[
\nu = i + k.
\]
Es ist daher:
\[
\begin{vmatrix}
a_{11} & a_{12} & a_{13} & \ldots & a_{1h} \\
a_{21} & a_{22} & a_{23} & \ldots & a_{2h} \\
a_{31} & a_{32} & & & \\
\vdots & & & & \\
a_{h1} & a_{h2} & \ldots & & a_{hh}
\end{vmatrix}
= (-1)^{\nu}
\begin{vmatrix}
a_{ik} & a_{i1} & a_{i2} & a_{i3} \ldots a_{i,k-1} & a_{i,k+1} \ldots a_{ih} \\
a_{1k} & a_{11} & a_{12} & a_{13} \ldots a_{1,k-1} & a_{1,k+1} \ldots a_{1h} \\
a_{2k} & a_{21} & a_{22} & a_{23} \ldots a_{2,k-1} & a_{2,k+1} \ldots a_{2h} \\
\vdots & \vdots & \vdots & \vdots & \vdots \\
a_{i-1,k} & a_{i-1,1} & & & \\
a_{i+1,k} & a_{i+1,1} & & & \\
\vdots & \vdots & & & \\
a_{hk} & a_{h1} & a_{h2} & \ldots a_{h,k-1} & a_{h,k+1} \ldots a_{hh}
\end{vmatrix}
\]
\[
= (-1)^{i+k} (a_{11} a_{22} \ldots a_{i-1,i-1} a_{i+1,i} a_{i+2,i+1} \ldots a_{k,k-1} a_{k+1,k+1} \ldots a_{hh}),
\]
wenn $i < k$, oder
\[
= (-1)^{i+k} (a_{11} a_{22} \ldots a_{i-1,i-1} a_{i,i+1} a_{i+1,i+2} \ldots a_{k-1,k} a_{k+1,k+1} \ldots a_{hh}),
\]
wenn $k < i$.


% ===== pdf idx 43, printed 29.0 =====
% !PRINTED 29
\weberpage{29}
22. \textit{Wollen wir nun zwei Elemente $a_{ik}$ und $a_{rs}$ in die Diagonalreihe
bringen,} und zwar $a_{ik}$ an die Stelle von $a_{11}$, $a_{rs}$ an die Stelle von
$a_{22}$, so ist n\"othig, dass $i \geqq r$ und $k \geqq s$ ist, d.\,h. dass die beiden
Elemente weder derselben Zeile noch Colonne angeh\"oren. Ist dies der Fall, so haben
wir zun\"achst die $i^{te}$ Zeile mit der ersten Zeile und die $k^{te}$ Colonne mit
der ersten Colonne, dann die $r^{te}$ Zeile mit der $2^{ten}$ Zeile und die $s^{te}$
Colonne mit der $2^{ten}$ cyclisch zu vertauschen; die Elemente des \"ubrigen Theiles
der Matrix behalten dabei ihre gegenseitige Lage m\"oglichst unge\"andert bei;
insbesondere sind im Diagonalglied, in welchem nun $a_{ik}$ die erste und $a_{rs}$
die zweite Stelle behauptet, die \"ubrigen Elemente von $a_{11}$ bis $a_{hh}$ gut
geordnet. Ihm entspricht also die Substitution
\[
S = \begin{pmatrix} i\ r\ 1\ 2\ \dots\ h \\ k\ s\ 1\ 2\ \dots\ h \end{pmatrix}.
\]

Die Derangements der ersten Permutation setzen sich zusammen (nach Nr.\,3)
\begin{enumerate}
\item[1)] aus den von der Permutation $(ir)$ herr\"uhrenden,
\item[2)] aus jenen, welche $i$ und $r$ mit $1\,2\,\dots\,h$ bilden, das sind
\[
i - 1 + r - 2.
\]
\end{enumerate}
Ganz ebenso bestehen die Derangements der unteren Permutation
\begin{enumerate}
\item[1)] aus den von der Permutation $(ks)$ herr\"uhrenden,
\item[2)] aus jenen, welche $k$ und $s$ mit $1\,2\,\dots\,h$ bilden, d.\,s.
$k - 1 + s - 2$.
\end{enumerate}
Die Differenz derselben ist also, wenn wir mit $(ir)$ und $(ks)$ gleich die Zahl
0 oder 1 der zugeh\"origen Derangements bezeichnen:
\[
(ir) - (ks) + i + r - k - s;
\]
daher ist der Modul der Substitution dargestellt durch eine Potenz von $(-1)$,
welche diese Zahl zum Exponenten hat. Da aber dieselbe gleichzeitig mit
\[
(ir) + (ks) + i + r + k + s = \nu
\]
gerade oder ungerade ist, so k\"onnen wir denselben in der Form $(-1)^{\nu}$
geben, wobei wir also unter $\nu$ stets diese Summe verstehen wollen. Es ist also:
\[
\begin{vmatrix}
a_{11} & a_{12} & a_{13} & \dots & a_{1h} \\
a_{21} & a_{22} & a_{23} & \dots & a_{2h} \\
a_{31} & a_{32} & a_{33} & \dots & a_{3h} \\
\vdots & & & & \vdots \\
a_{h1} & a_{h2} & a_{h3} & \dots & a_{hh}
\end{vmatrix}
=
\]
%CONT>>


% ===== pdf idx 44, printed 30.0 =====
% !PRINTED 30
\weberpage{30}
%<<CONT
\[
= (-1)^{\nu}
\left|
\begin{array}{cccccccccccc}
a_{ik} & a_{is} & a_{i1} & a_{i2} & \dots & a_{i,k-1} & a_{i,k+1} & \dots & a_{i,s-1} & a_{i,s+1} & \dots & a_{ih} \\
a_{rk} & a_{rs} & a_{r1} & a_{r2} & \dots & a_{r,k-1} & a_{r,k+1} & \dots & a_{r,s-1} & a_{r,s+1} & \dots & a_{rh} \\
a_{1k} & a_{1s} & a_{11} & a_{12} & \dots & a_{1,k-1} & a_{1,k+1} & \dots & a_{1,s-1} & a_{1,s+1} & \dots & a_{1h} \\
a_{2k} & a_{2s} & a_{21} & a_{22} & \dots & a_{2,k-1} & \cdot & \cdot & \cdot & \cdot & \cdot & a_{2h} \\
\multicolumn{12}{c}{\cdot\ \cdot\ \cdot\ \cdot\ \cdot\ \cdot\ \cdot\ \cdot\ \cdot\ \cdot\ \cdot\ \cdot} \\
a_{i-1,k} & a_{i-1,s} & a_{i-1,1} & a_{i-1,2} & \cdot & \cdot & \cdot & \cdot & \cdot & \cdot & \cdot & \cdot \\
a_{i+1,k} & a_{i+1,s} & a_{i+1,1} & \cdot & \cdot & \cdot & \cdot & \cdot & \cdot & \cdot & \cdot & \cdot \\
\vdots & \vdots & & & & & & & & & & \\
a_{r-1,k} & a_{r-1,s} & \cdot & \cdot & \cdot & \cdot & \cdot & \cdot & \cdot & \cdot & \cdot & \cdot \\
a_{r+1,k} & a_{r+1,s} & \cdot & \cdot & \cdot & \cdot & \cdot & \cdot & \cdot & \cdot & \cdot & \cdot \\
\vdots & \vdots & \vdots & \vdots & & & & & & & & \\
a_{hk} & a_{hs} & a_{h1} & a_{h2} & \dots & a_{h,k-1} & a_{h,k+1} & \dots & a_{h,s-1} & a_{h,s+1} & \dots & a_{hh}
\end{array}
\right|
\]
wobei $\nu = (ir) + (ks) + i + r + k + s.$

23. \textit{Wir sehen nun, wie sich analog der Zahlenwerth $\nu$ ergeben wird,
wenn wir drei Elemente $a_{ik}$, $a_{rs}$, $a_{lm}$ an die Pl\"atze von resp.
$a_{11}$, $a_{22}$, $a_{33}$ durch cyclische Vertauschung bringen.} Stellen wir
wieder die von den drei ersten Indices herr\"uhrende Zahl von Derangements durch
$(irl)$, die von den zweiten herr\"uhrende durch $(ksm)$ dar, eine Zahl, die
offenbar nur die Werthe $0, 1, 2, 3$ annehmen kann, dann ist der Modul der
Substitution, welche dem neuen Diagonalglied zugeh\"ort, durch eine Potenz
$(-1)^{\nu}$ dargestellt, wo
\[
\nu = (irl) + (ksm) + (i + r + l + k + s + m).
\]
Daraus l\"asst sich der Satz abstrahiren: \glqq Bringt man $p$ Elemente, von
denen keine zwei der n\"amlichen Colonne oder Zeile angeh\"oren d\"urfen, der
Reihe nach durch cyclische Vertauschung an die erste, zweite u.\,s.\,w.
$p^{te}$ Stelle der Diagonale einer Matrix, dann setzt sich der Exponent $\nu$
des zum Diagonalglied geh\"origen Substitutionsmoduls $\varepsilon_S = (-1)^{\nu}$
zusammen:
\begin{enumerate}
\item[1)] aus der Zahl der Derangements, welche die $p$ ersten Factoren
$a_{\varrho\sigma}$ durch die ersten Indices,
\item[2)] aus der Zahl, welche sie durch ihre zweiten Indices bilden,
\item[3)] aus der Summe ihrer Indices \"uberhaupt.\grqq{}
\end{enumerate}
\textit{Beispiel.} Sei gegeben die Determinante:
\[
\Delta = \begin{vmatrix}
\varepsilon & \iota & \varphi & \omega & \sigma \\
\iota & \varphi & \omega & \sigma & \omega \\
\varphi & \omega & \sigma & \omega & \varphi \\
\omega & \sigma & \omega & \varphi & \iota \\
\sigma & \omega & \varphi & \iota & \varepsilon
\end{vmatrix}.
\]


% ===== pdf idx 45, printed 31.0 =====
% !PRINTED 31
\weberpage{31}
Wir stellen uns die Aufgabe, die f\"unf Elemente $\sigma$ der zweiten Diagonale
durch cyclische Vertauschung in die Hauptdiagonale zu bringen, so dass jedes
Element $\sigma$ in einer Zeile bleibt; wir erhalten der Reihe nach
% ?? uncertain: exact glyph sequence in the four 5x5 sigma/epsilon/iota/phi/omega
% determinants below (dense symbol grid, image resolution limits confident reading)
\[
\Delta = (-1)^{\nu_1}
\begin{vmatrix}
\sigma & \varepsilon & \iota & \varphi & \omega \\
\omega & \iota & \varphi & \omega & \sigma \\
\varphi & \varphi & \omega & \sigma & \omega \\
\iota & \omega & \sigma & \omega & \varphi \\
\varepsilon & \sigma & \omega & \varphi & \iota
\end{vmatrix}
=
(-1)^{\nu_2}
\begin{vmatrix}
\sigma & \omega & \varepsilon & \iota & \varphi \\
\omega & \sigma & \iota & \varphi & \omega \\
\varphi & \omega & \varphi & \omega & \sigma \\
\iota & \varphi & \omega & \sigma & \omega \\
\varepsilon & \iota & \sigma & \omega & \varphi
\end{vmatrix}
\]
\[
= (-1)^{\nu_3}
\begin{vmatrix}
\sigma & \omega & \varphi & \varepsilon & \iota \\
\omega & \sigma & \omega & \iota & \varphi \\
\varphi & \omega & \sigma & \varphi & \omega \\
\iota & \varphi & \omega & \omega & \sigma \\
\varepsilon & \iota & \varphi & \sigma & \omega
\end{vmatrix}
=
(-1)^{\nu_4}
\begin{vmatrix}
\sigma & \omega & \varphi & \iota & \varepsilon \\
\omega & \sigma & \omega & \varphi & \iota \\
\varphi & \omega & \sigma & \omega & \varphi \\
\iota & \varphi & \omega & \sigma & \omega \\
\varepsilon & \iota & \varphi & \omega & \sigma
\end{vmatrix}.
\]
Da die Elemente $\sigma$ urspr\"unglich die Pl\"atze
\[
a_{51}, a_{42}, a_{33}, a_{24}, a_{15}
\]
inne hatten, so ist demgem\"ass
\[
\nu_4 = (5\,4\,3\,2\,1) + (1\,2\,3\,4\,5) + 2\sum_1^5 i
\]
\[
= 10 + 0 + 2\cdot\frac{5(5+1)}{2} = 40,
\]
also
\[
(-1)^{\nu_4} = +1.
\]

\begin{center}
\S.\,3.\\
\gesperrt{Unterdeterminanten erster Ordnung.}
\end{center}

24. \textit{Begriff der Unterdeterminante.} Aus den $N$ Gliedern der
Determinante $\Delta$ kann man insbesondere diejenigen herausgreifen, welche
den Factor $a_{ik}$ besitzen. Sie seien $G_{i_1}, G_{i_2}, \dots G_{i_\varrho}$,
wo $i_1, i_2, \dots i_\varrho$ irgendwelche Zahlen der Reihe $1, 2, 3 \dots$
bis $N = n!$ sind. Bezeichnet man das Product der andern Elemente in
$G_{i_\varrho}$ mit $H_{i_\varrho}$, dann ist:
\[
G_{i_1} = a_{ik} \cdot H_{i_1},
\]
\[
G_{i_2} = a_{ik} \cdot H_{i_2},
\]
\[
\dots\dots\dots
\]
\[
G_{i_\varrho} = a_{ik} \cdot H_{i_\varrho}.
\]
Bildet man die Summe aller dieser Glieder, so erh\"alt man
\[
G_{i_1} + G_{i_2} + \dots G_{i_\varrho} = a_{ik} H_{i_1} + a_{ik} H_{i_2} +
\dots a_{ik} H_{i_\varrho}
\]
oder
\[
\sum G_{i_\varrho} = a_{ik} \sum H_{i_\varrho}.
\]


% ===== pdf idx 46, printed 32.0 =====
% !PRINTED 32
\weberpage{32}
Bezeichnet man $\sum H_{i_\varrho}$ mit $A_{ik}$*), so erh\"alt man
\[
\sum G_{i_\varrho} = a_{ik} \cdot A_{ik}.
\]
Die Gr\"osse $A_{ik}$ nennt man \textit{Unterdeterminante (Minor)} des
Elementes $a_{ik}$.

25. \textit{Berechnung von $A_{11}$.} Die Aufgabe, die Unterdeterminante
$A_{ik}$ zu berechnen, wollen wir zun\"achst an dem speciellen Falle von
$i = k = 1$ durchf\"uhren, d.\,h. wir wollen $A_{11}$ bestimmen. Jedem Glied,
das den Factor $a_{11}$ besitzt, muss eine Substitution zugeordnet sein, in
der die zwei Indices 1 stets \"uber einander stehen, d.\,h. ein einziges Paar
bilden. Alle \"ubrigen Indices, unter denen sich also 1 nicht mehr befindet,
liefern durch ihre Permutationen, die wir der Reihe nach mit der gut
geordneten $(2\,3\,\dots\,h)$ zu einer Substitution verbinden, diejenigen
Substitutionen, die an erster Stelle das Zahlenpaar $\binom{1}{1}$ haben. Den
Permutationen der $n-1$ Elemente $(2\,3\,\dots\,h)$ entsprechen aber die
$(n-1)!$ Substitutionen, welche zu den Gliedern $H_{i_\varrho}$ Veranlassung
geben. Die Summe aller dieser $H_{i_\varrho}$ ist also selbst eine
Determinante:
\[
\Delta' = \sum_{i_\varrho = 1}^{i_\varrho = (n-1)!} H_{i_\varrho},
\]
und zu ihr geh\"ort die Matrix:
\[
\Delta' =
\begin{vmatrix}
a_{22} & a_{23} & a_{24} & \dots & a_{2h} \\
a_{32} & a_{33} & a_{34} & \dots & a_{3h} \\
a_{42} & a_{43} & a_{44} & \dots & a_{4h} \\
\vdots & \vdots & & \dots & \\
a_{h2} & a_{h3} & a_{h4} & \dots & a_{hh}
\end{vmatrix}.
\]
Vergleicht man aber diese Matrix $\Delta'$ mit der Matrix der gegebenen
Determinante $\Delta$, so erkennt man: $\Delta'$ entsteht aus $\Delta$, indem
man die erste Zeile und erste Colonne streicht, oder: \textit{Die
Unterdeterminante $A_{11}$ des Elementes $a_{11}$ entsteht, indem man jene
Zeile und Colonne streicht, die sich in $a_{11}$ schneiden.}

26. \textit{Berechnung von $A_{ik}$.} Nun k\"onnen wir die allgemeine Aufgabe,
$A_{ik}$ zu bestimmen, leicht auf diese zur\"uckf\"uhren. Ich bringe zu dem
Zwecke durch zwei Cyclen das Element $a_{ik}$ an die Stelle von $a_{11}$. Zu
der dadurch entstehenden Matrix geh\"ort eine Determinante vom Werthe
\[
D = (-1)^{\nu}\Delta, \quad \text{wo } \nu = i + k,
\]
%CONT>>

*) Es ist $\dfrac{\partial \Delta}{\partial a_{ik}} = A_{ik}$.


% ===== pdf idx 47, printed 33.0 =====
% !PRINTED 33
\weberpage{33}
% running head: \S\ 3. Unterdeterminanten erster Ordnung.

und wir sahen in Nr. 21, dass bei einer solchen Reihe von Vertauschungen die
Glieder von $D$ in passender Ordnung genommen dieselben Werthe haben wie die
von $\Delta$, nur noch multiplicirt mit $(-1)^\nu$. Insbesondere entsprechen
den Gliedern mit dem Factor $a_{ik}$ in $\Delta$ die Glieder mit dem Factor
$b_{11}$ in $D$; ihre Summe
\[
G_{i_1}+G_{i_2}+\cdots G_{i_\varrho}=a_{ik}\sum H_{i_\varrho}=a_{ik}\cdot A_{ik}
\]
ist bis auf einen Factor $(-1)^\nu$ gleich der entsprechenden
\[
\overline{G}_{i_1}+\overline{G}_{i_2}+\cdots \overline{G}_{i_\varrho}=b_{11}\sum \overline{H}_{i_\varrho}=b_{11}\cdot B_{11},
\]
wenn die $B_{ik}$ die Unterdeterminanten der $b_{ik}$ bedeuten. Die Matrix zu
$B_{11}$ entsteht aber aus der Matrix zu $D$, wenn man die erste Zeile und
erste Colonne streicht; die nämliche Matrix erhält man aber aus $\Delta$,
wenn man in ihr jene Zeile und Colonne streicht, die sich in $a_{ik}$
schneiden. Daher der Satz: \glqq Lässt man in der Matrix zu $\Delta$ die
$i^{\text{te}}$ Zeile und die $k^{\text{te}}$ Colonne weg, so erhält man eine
neue Matrix $\Delta'$. Multiplicirt man den Werth der entsprechenden
Determinante mit $(-1)^\nu$, wo $\nu=i+k$, so erhält man die zum Element
$a_{ik}$ gehörige Unterdeterminante $A_{ik}$.\grqq{}

27. \textit{Beispiel:} 1) Sei gegeben die Determinante
\[
\begin{vmatrix} p & q & r \\ u & v & w \\ x & y & z \end{vmatrix};
\]
dann ist die Unterdeterminante
\[
\text{zu } q=(-1)^{1+2}\begin{vmatrix} u & w \\ x & z \end{vmatrix}=-(uz-wx),
\]
\[
\text{zu } v=(-1)^{2+2}\begin{vmatrix} p & r \\ x & z \end{vmatrix}=+(pz-rx),
\]
\[
\text{zu } y=(-1)^{3+2}\begin{vmatrix} p & r \\ u & w \end{vmatrix}=-(pw-ur).
\]
In der That, berechnen wir die Determinante nach dem in Nr. 13 oder 15
gegebenen Schema, so sieht man, dass $uz$, $wx$ die einzigen Producte sind,
die mit $q$ multiplicirt auftreten, ebenso $pz$ und $rx$ die einzigen mit $v$
multiplicirten u.\,s.\,w.

2) In der Determinante
\[
\begin{vmatrix}
a_{11} & a_{12} & a_{13} & a_{14} \\
a_{21} & a_{22} & a_{23} & a_{24} \\
a_{31} & a_{32} & a_{33} & a_{34} \\
a_{41} & a_{42} & a_{43} & a_{44}
\end{vmatrix}
\]
% running head foot: Gordan, Determinanten.   3
%CONT>>


% ===== pdf idx 48, printed 34.0 =====
% !PRINTED 34
\weberpage{34}
% running head: Erster Theil. Theorie der Determinanten.
%<<CONT
ist die Unterdeterminante zu $a_{32}$ dargestellt durch die Matrix
\[
(-1)^{3+2}\begin{vmatrix} a_{11} & a_{13} & a_{14} \\ a_{21} & a_{23} & a_{24} \\ a_{41} & a_{43} & a_{44} \end{vmatrix}
\]
\[
=(-1)^5\{a_{11}a_{23}a_{44}-a_{11}a_{43}a_{24}-a_{21}a_{13}a_{44}+a_{21}a_{43}a_{14}+a_{41}a_{13}a_{24}-a_{41}a_{23}a_{14}\}.
\]
In der That, bilden wir die Substitutionen, die diesen Gliedern zugehören:
\[
\begin{pmatrix} 1 & 2 & 4 \\ 1 & 3 & 4 \end{pmatrix}
\begin{pmatrix} 1 & 4 & 2 \\ 1 & 3 & 4 \end{pmatrix}
\begin{pmatrix} 2 & 1 & 4 \\ 1 & 3 & 4 \end{pmatrix}
\begin{pmatrix} 2 & 4 & 1 \\ 1 & 3 & 4 \end{pmatrix}
\begin{pmatrix} 4 & 1 & 2 \\ 1 & 3 & 4 \end{pmatrix}
\begin{pmatrix} 4 & 2 & 1 \\ 1 & 3 & 4 \end{pmatrix},
\]
so erkennt man, dass \textit{sämmtliche} Permutationen der drei Indices
$1\,2\,4$ mit der gut geordneten $(1\,3\,4)$ verbunden sind, wie wir es auch
in unserer obigen Darstellung verlangten.

28. \textit{Berechnung des Werthes einer Determinante mittelst der
Unterdeterminanten, welche den Elementen der ersten Zeile entsprechen.}
Nach den Entwicklungen in Nr. 15 hat jedes Glied der Determinante stets ein
und nur ein Element der ersten Zeile unter seinen Factoren. Ordnen wir also
alle $N$ Glieder in Klassen, indem wir zuerst alle Glieder herausgreifen,
welche den Factor $a_{11}$, sodann alle, welche den Factor $a_{12}$
u.\,s.\,w., endlich jene, welche den Factor $a_{1n}$ besitzen, so ist die
Gesammtheit aller Glieder dadurch erschöpft. Nun bilden aber nach Nr. 24
alle Glieder mit dem Factor $a_{11}$ zusammen die Unterdeterminante $A_{11}$,
alle mit dem Factor $a_{12}$ die Unterdeterminante $A_{12}$ u.\,s.\,w.,
endlich alle mit dem Factor $a_{1n}$ die Unterdeterminante $A_{1n}$; folglich
ist der Werth der Determinante $\Delta$ dargestellt durch die Summe
\weq{1}{\Delta = a_{11}A_{11}+a_{12}A_{12}+a_{13}A_{13}+\cdots+a_{1n}A_{1n}.}
Die hierbei auftretenden Unterdeterminanten $A_{1i}$ sind entnommen aus der
Matrix von $n$ Colonnen und $(n-1)$ Zeilen:
\weq{A}{
\begin{vmatrix}
a_{21} & a_{22} & a_{23} & \cdots & a_{2,n} \\
a_{31} & a_{32} & a_{33} & \cdots & a_{3,n} \\
a_{41} & & & & \\
\vdots & & & & \\
a_{n1} & a_{n2} & a_{n3} & \cdots & a_{nn}
\end{vmatrix}.
}
Die Matrix der Unterdeterminante $A_{11}$ entsteht dann hieraus, indem man
die erste Colonne weglässt; analog ist $A_{12}$ dargestellt durch die nach
dem Streichen der zweiten Colonne übrig bleibende Matrix multiplicirt mit
$(-1)^\nu$, wo $\nu=1+2=3$ ist u.\,s.\,w., endlich die Unterdeterminante
$A_{1n}$ durch die Matrix, welche die letzte Colonne nicht enthält,
multiplicirt mit $(-1)^\nu$, wo $\nu=1+n$. Man
%CONT>>


% ===== pdf idx 49, printed 35.0 =====
% !PRINTED 35
\weberpage{35}
% running head: \S\ 3. Unterdeterminanten erster Ordnung.
%<<CONT
sieht hiebei, indem man der Zahl $n$ alle Werthe von 1 bis $n$ beilegt, dass
der Factor $(-1)^\nu$ abwechselnd positiv und negativ ist, d.\,h. die
Glieder der Summe (1) haben abwechselnd das Zeichen plus und minus.

29. \textit{Entwicklung der Determinanten nach Unterdeterminanten der
Elemente einer beliebigen Zeile oder Colonne.} Ganz ebenso, wie wir hier den
Werth der Determinante $\Delta$ nach den Unterdeterminanten der ersten
Zeile entwickelt haben, können wir ihn auch berechnen nach den
Unterdeterminanten einer beliebigen Zeile oder Colonne, die wir ja nach den
vorausgehenden Sätzen jederzeit zur ersten Zeile oder Colonne machen können.

Wir haben also ganz allgemein für den Werth der Determinante $\Delta$ die
beiden Formeln
\weq{B}{
\left\{
\begin{array}{l}
\Delta = a_{i1}A_{i1}+a_{i2}A_{i2}+a_{i3}A_{i3}\cdots a_{in}A_{in}, \\
\Delta = a_{1i}A_{1i}+a_{2i}A_{2i}+a_{3i}A_{3i}\cdots a_{ni}A_{ni}.
\end{array}
\right.
}
30. \textit{Beispiele.} 1) Berechnen wir nach dieser Methode noch einmal die
viergliederige Determinante Beispiel 6 in Nr. 15 nach den
Unterdeterminanten, die zu den Elementen der ersten Zeile gehören. Wir
erhalten:
\[
\Delta=\begin{vmatrix} 3 & 1 & 1 & 2 \\ 1 & 5 & 0 & 3 \\ 2 & -2 & 1 & 6 \\ 0 & 4 & -5 & 3 \end{vmatrix}
=3\begin{vmatrix} 5 & 0 & 3 \\ -2 & 1 & 6 \\ 4 & -5 & 3 \end{vmatrix}
-1\begin{vmatrix} 1 & 0 & 3 \\ 2 & 1 & 6 \\ 0 & -5 & 3 \end{vmatrix}
+1\begin{vmatrix} 1 & 5 & 3 \\ 2 & -2 & 6 \\ 0 & 4 & 3 \end{vmatrix}
-2\begin{vmatrix} 1 & 5 & 0 \\ 2 & -2 & 1 \\ 0 & 4 & -5 \end{vmatrix}
\]
\[
=3\{15+150+30-12\}-\{3+30-30\}+\{-6-24-30+24\}
\]
\[
-2\{10-4+50\}=549-3-36-112=398,
\]
was mit dem früheren Resultat übereinstimmt.

2) Fassen wir umgekehrt die einzelnen Glieder der ausgerechneten
viergliederigen Determinante Nr. 17 Beispiel 3 in der Weise zusammen, dass
wir sie nach den Elementen der ersten Colonne ordnen, so ergiebt sich:
\[
a[erw-evs-qfw+qvg+ufs-urg] % ?? uncertain: last letter of bracket smudged/annotated in scan, read as "g" by cofactor-expansion sense (matches erw-esv-fqw+fsu+gqv-gru)
+p[bfw-bvg-ecw+evd+ucg-ufd].
\]
Die beiden Klammerfactoren dieser Summe sind aber offenbar nur die
Entwicklungen der beiden dreigliederigen Determinanten (vergleiche das
Schema in Nr. 15, $\beta$)
\[
A_{11}=\begin{vmatrix} e & f & g \\ q & r & s \\ u & v & w \end{vmatrix}
\quad \text{und} \quad
A_{31}=\begin{vmatrix} b & c & d \\ e & f & y \\ u & v & w \end{vmatrix}.
\]
% running head foot: 3*


% ===== pdf idx 50, printed 36.0 =====
% !PRINTED 36
\weberpage{36}
% running head: Erster Theil. Theorie der Determinanten.

\textit{Beispiel} 3. Sind in einer Zeile oder Reihe alle Elemente bis auf
eines gleich null, so wird man am Besten bei einer Berechnung der
Determinante zuerst nach dieser entwickeln, da sich die Summe (B) dann auf
ein einziges Glied reducirt.

Es sei:
\[
\Delta=\begin{vmatrix} 3 & 2 & 5 & 1 \\ 0 & 0 & 4 & 0 \\ 4 & 6 & 2 & 3 \\ 1 & 2 & 1 & 0 \end{vmatrix}.
\]
Wir entwickeln nach der zweiten Zeile; es wird also:
\[
\Delta=(-1)^\nu 4\begin{vmatrix} 3 & 2 & 1 \\ 4 & 6 & 3 \\ 1 & 2 & 0 \end{vmatrix}, \quad \nu=2+3=5.
\]

\textit{Beispiel} 4. Sind alle Elemente auf einer Seite der Diagonale einer
Matrix null, so reducirt sich der Werth der Determinante auf das
Diagonalglied; denn man hat:
\[
\Delta=\begin{vmatrix}
a & b & c & d & e \\
0 & f & g & h & i \\
0 & 0 & k & l & m \\
0 & 0 & 0 & n & 0 \\
0 & 0 & 0 & 0 & p
\end{vmatrix}
=a\begin{vmatrix} f & g & h & i \\ 0 & k & l & m \\ 0 & 0 & n & 0 \\ 0 & 0 & 0 & p \end{vmatrix}
=af\begin{vmatrix} k & l & m \\ 0 & n & 0 \\ 0 & 0 & p \end{vmatrix}
=afk\begin{vmatrix} n & 0 \\ 0 & p \end{vmatrix}
=afknp.
% ?? uncertain: several zero entries in this display carry a small circular ink mark in the scan (over the "0" in row4/col5, and inside the two smaller minors); read as plain 0, mark may be a printing/scan blemish rather than a diacritic
\]

\textit{Beispiel} 5. Man kann umgekehrt jede Determinante als
Unterdeterminante betrachten, deren Element $a_{11}=1$, deren Elemente
$a_{i1}=0$, wenn $i>1$, und deren Elemente $a_{1i}$, wenn $i>1$, beliebig
sind; so ist offenbar:
\[
\Delta=\begin{vmatrix} p & q & r \\ u & v & w \\ x & y & z \end{vmatrix}
=\begin{vmatrix}
1 & \varkappa & \lambda & \mu \\ % ?? uncertain: first entry of row 1 after "1" -- read as \varkappa, glyph partly unclear in scan
0 & p & r & q \\ % [sic] printed order "p, r, q" (not "p, q, r" as in the 3x3 above)
0 & u & v & w \\
0 & x & y & z
\end{vmatrix}.
\]

31. \textit{Beziehungen zwischen den Unterdeterminanten der Elemente der
ersten Zeile.} Bezeichnen wir wie früher mit $D$ eine Determinante von
$n^2$ Elementen $b_{ik}$, und setzen wir aber nun fest, dass zwischen den
Elementen von $D$ und den Elementen $a_{ik}$ der Determinante $\Delta$ die
Relationen bestehen:
\[
\alpha)\ b_{1i}=a_{2i}, \qquad \beta)\ b_{ki}=a_{ki},\ k>1.
\]
Dann ist, weil die Matrix zu $D$ zwei gleiche — nämlich die erste und
zweite — Zeilen hat, der Werth der Determinante $D$ gleich null.


% ===== pdf idx 51, printed 37.0 =====
% !PRINTED 37
\weberpage{37}
(Nr. 19, Zusatz.) Entwickeln wir daher die Determinante $D$ nach den
Unterdeterminanten der ersten Zeile, so erhalten wir die Relation:

\weq{1}{b_{11}B_{11}+b_{12}B_{12}+\cdots+b_{1n}B_{1n}=0.}

Nun ist aber wegen der Bedingung ($\alpha$)

\[
b_{11}=a_{21},\ b_{12}=a_{22},\cdots b_{1n}=a_{2n}.
\]

Ebenso sind, wegen Bedingung ($\beta$), die Unterdeterminanten $B_{1i}$ der
gleichen Matrix

\[
(A)\quad
\begin{vmatrix}
b_{21} & b_{22} & b_{23} & \ldots & b_{2n} \\
b_{31} & b_{32} & b_{33} & \ldots & b_{3n} \\
\cdot & \cdot & \cdot & \cdot & \cdot \\
b_{n1} & b_{n2} & b_{n3} & \ldots & b_{nn}
\end{vmatrix}
=
\begin{vmatrix}
a_{21} & a_{22} & a_{23} & \ldots & a_{2n} \\
a_{31} & a_{32} & a_{33} & \ldots & a_{3n} \\
\cdot & \cdot & \cdot & \cdot & \cdot \\
a_{n1} & a_{n2} & a_{n3} & \ldots & a_{nn}
\end{vmatrix}
\]

entnommen wie die Unterdeterminante $A_{1i}$, d. h. es ist auch

\[
B_{11}=A_{11},\ B_{12}=A_{12}\cdots B_{1n}=A_{1n}\, .
\]

Tragen wir also diese Werthe von $b_{1i}$ und $B_{1i}$ in die Relation (1)
ein, so erhalten wir:

\weq{2}{a_{21}A_{11}+a_{22}A_{12}+a_{23}A_{13}+\cdots a_{2n}A_{1n}=0.}

Diese Gleichung ist eine Identität, da ja zwischen den Elementen $a_{ik}$
der Determinante $\Delta$ keine Beziehungen bestehen sollen; sie sagt:

\glqq Multiplicirt man die $n$ Determinanten der Matrix $(A)$ Nr. 28 der
Reihe nach mit den Elementen der ersten Zeile dieser Matrix und summirt
diese Producte, so ist diese Summe identisch null.\grqq{}

32. \textit{Verallgemeinerung der Relation (2) Nr. 31.} Nun ist aber jede
Zeile einer Determinante mit $n^2$ willkürlichen Elementen
gleichberechtigt, und man kann schon daraus schliessen, dass die Elemente
$a_{2i}$ in (2) durch die Elemente jeder folgenden Zeile ersetzt werden
können, ohne dass unser Satz seine Gültigkeit verliert. Strenger folgt
dieses, wenn wir statt der Bedingung $\alpha$) $b_{1i}=a_{2i}$ die
Bedingung einführen

\[
b_{1i}=a_{\varrho i}\, .
\]

Dann bleiben immer noch die Matrix $(A)$ und mit ihr die
Unterdeterminanten $A_{1i}$ die nämlichen, und nur an Stelle der Elemente
$a_{2i}$ in (2) treten die Elemente $a_{\varrho i}$, $\varrho=3,4,5\cdots
n$. Wir haben sonach eine erste Erweiterung des Satzes:

\glqq Multiplicirt man die Determinanten der Matrix $(A)$ der Reihe nach
mit den Elementen einer beliebigen Zeile derselben, so ist die Summe
aller dieser Producte identisch null.\grqq{}


% ===== pdf idx 52, printed 38.0 =====
% !PRINTED 38
\weberpage{38}
Wir haben bisher immer die Determinanten der Matrix $(A)$, d. i. die
Unterdeterminanten der ersten Zeile der Determinante $\Delta$
ausgezeichnet. Da wir aber jede Zeile oder Colonne zur ersten Zeile
machen können, ohne dass sich der absolute Werth von $\Delta$ ändert, so
sind wir berechtigt, ganz allgemein den Satz aufzustellen:

\glqq Multiplicirt man die Elemente einer Zeile bezw. Colonne mit der je
entsprechenden Unterdeterminante einer andern Zeile bezw. Colonne, so ist
die Summe aller dieser Producte identisch null.\grqq{}

Wir erhalten also, wenn wir gleichzeitig das in Nr. 28 enthaltene
Resultat berücksichtigen, folgendes System von Relationen:

\[
(C)\quad
\begin{cases}
a_{i1}A_{k1}+a_{i2}A_{k2}+\cdots+a_{in}A_{kn}=\Delta, & \text{wenn } i=k,\\
a_{i1}A_{k1}+a_{i2}A_{k2}+\cdots+a_{in}A_{kn}=0, & \text{wenn } i\gtrless k,\\
a_{1i}A_{1k}+a_{2i}A_{2k}+\cdots+a_{ni}A_{nk}=\Delta, & \text{wenn } i=k,\\
a_{1i}A_{1k}+a_{2i}A_{2k}+\cdots+a_{ni}A_{nk}=0, & \text{wenn } i\gtrless k.
\end{cases}
\]

33. \textit{Beispiel.} Sei gegeben die Determinante

\[
\Delta=
\begin{vmatrix}
a_{11} & a_{12} & a_{13} \\
a_{21} & a_{22} & a_{23} \\
a_{31} & a_{32} & a_{33}
\end{vmatrix};
\]

die zur letzten Colonne gehörige Matrix ist:

\[
\begin{vmatrix}
a_{11} & a_{12} \\
a_{21} & a_{22} \\
a_{31} & a_{32}
\end{vmatrix};
\]

sie enthält die Unterdeterminanten:

\[
A_{33}=\begin{vmatrix} a_{11} & a_{12} \\ a_{21} & a_{22} \end{vmatrix},\quad
A_{23}=-\begin{vmatrix} a_{11} & a_{12} \\ a_{31} & a_{32} \end{vmatrix},\quad
A_{13}=\begin{vmatrix} a_{21} & a_{22} \\ a_{31} & a_{32} \end{vmatrix}.
\]

Multipliciren wir sie mit den Elementen der ersten Colonne, so erhalten
wir:

\[
\begin{aligned}
a_{11}A_{13}+a_{21}A_{23}+a_{31}A_{33} &= a_{11}(a_{21}a_{32}-a_{31}a_{22})-a_{21}(a_{11}a_{32}-a_{12}a_{31})\\
&\quad{}+a_{31}(a_{11}a_{22}-a_{21}a_{12}).
\end{aligned}
\]

Die rechte Seite dieser Gleichung verschwindet aber identisch.

34. \textit{Anwendungen der Identitäten in $(C)$.} Aus der zweiten
specielleren Fassung des soeben erhaltenen allgemeinen Satzes ergeben
sich für die $n$ Determinanten $A_{1k}$, welche in der Matrix $(A)$
enthalten sind, folgende $(n-1)$ Relationen. Bezeichnen wir nämlich mit
Ein-
%CONT>>


% ===== pdf idx 53, printed 39.0 =====
% !PRINTED 39
\weberpage{39}
%<<CONT
schluss des Vorzeichens diese $n$ Grössen der Reihe nach mit $p_1p_2\ldots
p_n$, so haben wir die Beziehungen:

\[
(D)\quad
\begin{cases}
a_{21}p_1+a_{22}p_2+\cdots+a_{2n}p_n=0,\\
a_{31}p_1+a_{32}p_2+\cdots+a_{3n}p_n=0,\\
a_{41}p_1+a_{42}p_2+\cdots+a_{4n}p_n=0,\\
\cdot\ \cdot\ \cdot\ \cdot\ \cdot\ \cdot\ \cdot\ \cdot\ \cdot\ \cdot\ \cdot\\
a_{n1}p_1+a_{n2}p_2+\cdots+a_{nn}p_n=0.
\end{cases}
\]

Mit Hilfe dieser Relationen wollen wir noch zwei Identitäten ableiten,
die später von Wichtigkeit sein werden.

1) Sei gegeben die Matrix:

\[
\begin{vmatrix}
a_1 & b_1 & c_1 \\
a_2 & b_2 & c_2
\end{vmatrix};
\]

schreiben wir die mit $p_1p_2p_3$ bezeichneten Minoren in der Form
$(bc)$, $(ca)$, $(ab)$ --- vergleiche Nr. 18 --- so werden die Relationen
$(D)$:

\[
\begin{gathered}
a_1(bc)+b_1(ca)+c_1(ab)=0,\\
a_2(bc)+b_2(ca)+c_2(ab)=0.
\end{gathered}
\]

Multipliciren wir die erste Gleichung mit $x_1$, die zweite mit $x_2$ und
addiren, führen sodann für $a_1x_1+a_2x_2$, $b_1x_1+b_2x_2$,
$c_1x_1+c_2x_2$ bezw. die Symbole $a_x$, $b_x$, $c_x$ ein, so erhalten
wir die Identität

\weq{I}{a_x(bc)+b_x(ca)+c_x(ab)=0.}

2) Ebenso liefert die Matrix:

\[
\begin{vmatrix}
a_1 & b_1 & c_1 & d_1 \\
a_2 & b_2 & c_2 & d_2 \\
a_3 & b_3 & c_3 & d_3
\end{vmatrix},
\]

wenn wir wieder die Unterdeterminanten $p_1p_2p_3p_4$ mit $(bcd)$,
$-(cda)$, $(dab)$, $-(abc)$ bezeichnen, die drei Relationen:

\[
\begin{gathered}
a_1(bcd)-b_1(cda)+c_1(dab)-d_1(abc)=0,\\
a_2(bcd)-b_2(cda)+c_2(dab)-d_2(abc)=0,\\
a_3(bcd)-b_3(cda)+c_3(dab)-d_3(abc)=0.
\end{gathered}
\]

Multiplicirt man auch hier die drei Gleichungen der Reihe nach mit
$x_1x_2x_3$ und setzt $a_1x_1+a_2x_2+a_3x_3=a_x$,
$b_1x_1+b_2x_2+b_3x_3=b_x$, $c_1x_1+c_2x_2+c_3x_3=c_x$,
$d_1x_1+d_2x_2+d_3x_3=d_x$, nachdem man die drei Gleichungen addirt hat,
so kommt die Identität:

\weq{II}{a_x(bcd)-b_x(cda)+c_x(dab)-d_x(abc)=0.}


% ===== pdf idx 54, printed 40.0 =====
% !PRINTED 40
\weberpage{40}
35. \textit{Addition von zwei Determinanten.} Ergänzt man die Matrix
$(A)$ auf drei verschiedene Arten zur Matrix einer Determinante

1) indem man zu den $(n-1)$ Zeilen die Zeile $a_{11}\,a_{12}\ldots
a_{1n}$,

2) indem man die Zeile $b_{11}\,b_{12}\ldots b_{1n}$, endlich

3) indem man die Zeile $a_{11}+b_{11},\,a_{12}+b_{12}\ldots
a_{1n}+b_{1n}$ hinzufügt,

so erhält man dadurch die Matrix von drei verschiedenen Determinanten
$\Delta_0$, $\Delta_1$, $\Delta_2$. Ihre Werthe sind nach 29 $(B)$

\[
\left.
\begin{aligned}
\Delta_0 &= a_{11}p_1+a_{12}p_2+\cdots a_{1n}p_n\\
\Delta_1 &= b_{11}p_1+b_{12}p_2+\cdots b_{1n}p_n\\
\Delta_2 &= (a_{11}+b_{11})p_1+(a_{12}+b_{12})p_2+\cdots (a_{1n}+b_{1n})p_n
\end{aligned}
\right\}
\quad (E).
\]

Man erkennt, dass der Werth der dritten Determinante gleich ist der
Summe der Werthe der beiden andern; daher der Satz:

\glqq Die Summe zweier Determinanten, deren Matrices in $(n-1)$ Zeilen
gleiche, in der $n^{\text{ten}}$ Zeile aber verschiedene Elemente haben,
ist dargestellt durch Eine Determinante, deren Matrix die $(n-1)$
gemeinschaftlichen Zeilen besitzt und deren $n^{\text{te}}$ Zeile die
Summe je zweier entsprechender Elemente der beiden nicht
gemeinschaftlichen Zeilen enthält.\grqq{}

36. \textit{Addition von $m$ Determinanten mit $(n-1)$ gemeinschaftlichen
Zeilen.} Wir hätten ebenso die Matrix $(A)$ auf $m$ Arten zu den
Matrices von $m$ verschiedenen Determinanten ergänzen können. Die Summe
aller dieser Determinanten ist dann dargestellt durch eine neue
Determinante, in deren einer Zeile jedes Element aus der Summe von $m$
entsprechenden Elementen der einzelnen Determinanten besteht, während
die übrigen Zeilen nur die allen gemeinschaftlichen Elemente enthalten.
Aus der Entwicklung des Satzes geht unmittelbar auch seine
Umkehrbarkeit hervor, d. h. tritt in einer Determinante eine Zeile auf,
deren Elemente algebraische Summen sind, so lässt sich dieselbe immer
in eine Summe von Determinanten zerlegen.

37. \textit{Beispiele.} 1) Es sei gegeben die Matrix

\[
\begin{vmatrix}
a_1 & b_1 & c_1 & d_1 \\
a_2 & b_2 & c_2 & d_2 \\
a_3 & b_3 & c_3 & d_3
\end{vmatrix}.
\]
% ?? uncertain: scan shows a stray mark over the d_2 subscript in row 2 (ink flaw/correction, not legible as separate text)

Ergänzen wir sie durch 1) Anfügen der ersten Zeile

\[
\alpha_1\ \beta_1\ \gamma_1\ \delta_1
\]

zur Determinante $\Delta_1$, 2) durch Anfügen der ersten Zeile

\[
-\alpha_2,\ -\beta_2,\ -\gamma_2,\ -\delta_2
\]
%CONT>>


% ===== pdf idx 55, printed 41.0 =====
% !PRINTED 41
\weberpage{41}
%<<CONT
zur Determinante $\varDelta_2$, 3) endlich durch Anfügen der Zeile
\[
+\alpha_3, - \beta_3, +\gamma_3, -\delta_3
\]
zur Determinante $\varDelta_3$, dann ist
\[
\varDelta = \varDelta_1+\varDelta_2+\varDelta_2 =
\begin{vmatrix}
\alpha_1-\alpha_2+\alpha_3, & \beta_1-\beta_2-\beta_3, & \gamma_1-\gamma_1+\gamma_3, & \delta_1-\delta_2-\delta_3\\
a_1 & b_1 & c_1 & d_1\\
a_2 & b_2 & c_2 & d_2\\
a_3 & b_3 & c_3 & d_3
\end{vmatrix}.
\]
2) Die Determinante:
\[
\varDelta = \begin{vmatrix} a+\lambda, & b, & c\\ p & q & r\\ u & v & w\end{vmatrix}
= \begin{vmatrix} a & b & c\\ p & q & r\\ u & v & w\end{vmatrix}
+ \begin{vmatrix} \lambda & 0 & 0\\ p & q & r\\ u & v & w\end{vmatrix}.
\]
3) Die Determinante:
\[
\varDelta = \begin{vmatrix} a_{11}+\lambda, & a_{12}, & a_{13}\\ a_{21}, & a_{22}+\lambda, & a_{23}\\ a_{31} & a_{32} & a_{33}+\lambda\end{vmatrix}
= \begin{vmatrix} a_{11} & a_{12} & a_{13}\\ a_{21} & a_{22}+\lambda & a_{23}\\ a_{31} & a_{32} & a_{33}+\lambda\end{vmatrix}
+ \begin{vmatrix} \lambda & a_{12} & a_{13}\\ 0 & a_{22}+\lambda & a_{23}\\ 0 & a_{32} & a_{33}+\lambda\end{vmatrix}
= \varDelta_1+\varDelta_2.
\]
Jede dieser Determinanten $\varDelta_1$ und $\varDelta_2$ lässt sich wieder zerlegen:
\[
\varDelta_1 = \begin{vmatrix} a_{11} & a_{12} & a_{13}\\ a_{21} & a_{22} & a_{23}\\ a_{31} & a_{32} & a_{33}+\lambda\end{vmatrix}
+ \begin{vmatrix} a_{11} & 0 & a_{13}\\ a_{21} & \lambda & a_{23}\\ a_{31} & 0 & a_{33}+\lambda\end{vmatrix}
= \varDelta_1'+\varDelta_1'',
\]
und
\[
\varDelta_2 = \begin{vmatrix} \lambda & a_{12} & a_{13}\\ 0 & a_{22} & a_{23}\\ 0 & a_{32} & a_{33}+\lambda\end{vmatrix}
+ \begin{vmatrix} \lambda & 0 & a_{13}\\ 0 & \lambda & a_{23}\\ 0 & 0 & a_{33}+\lambda\end{vmatrix}
= \varDelta_2'+\varDelta_2''.
\]
Zerlegt man endlich jede dieser vier Determinanten, so erhält man $\varDelta$ dargestellt durch die Summe:
\[
\begin{aligned}
\varDelta ={}& \begin{vmatrix} a_{11} & a_{12} & a_{13}\\ a_{21} & a_{22} & a_{23}\\ a_{31} & a_{32} & a_{33}\end{vmatrix}
+ \begin{vmatrix} a_{11} & a_{12} & 0\\ a_{21} & a_{22} & 0\\ a_{31} & a_{32} & \lambda\end{vmatrix}
+ \begin{vmatrix} a_{11} & 0 & a_{13}\\ a_{21} & \lambda & a_{23}\\ a_{31} & 0 & a_{33}\end{vmatrix}
+ \begin{vmatrix} a_{11} & 0 & 0\\ a_{21} & \lambda & 0\\ a_{31} & 0 & \lambda\end{vmatrix}\\
&+ \begin{vmatrix} \lambda & a_{12} & a_{13}\\ 0 & a_{22} & a_{23}\\ 0 & a_{32} & a_{33}\end{vmatrix}
+ \begin{vmatrix} \lambda & a_{12} & 0\\ 0 & a_{22} & 0\\ 0 & a_{32} & \lambda\end{vmatrix}
+ \begin{vmatrix} \lambda & 0 & a_{13}\\ 0 & \lambda & a_{23}\\ 0 & 0 & a_{33}\end{vmatrix}
+ \begin{vmatrix} \lambda & 0 & 0\\ 0 & \lambda & 0\\ 0 & 0 & \lambda\end{vmatrix}.
\end{aligned}
\]


% ===== pdf idx 56, printed 42.0 =====
% !PRINTED 42
\weberpage{42}
38. \textit{Multiplication einer Determinante mit einer beliebigen Grösse $\lambda$.} Gehen wir wiederum von der Matrix $(A)$ aus und ergänzen sie zur Matrix einer Determinante $\varDelta_1$, indem wir die Elemente $\lambda a_{11}$, $\lambda a_{12}$, $\lambda a_{13} \ldots \lambda a_{1n}$ als erste Zeile anfügen. Dann erhalten wir nach Nr. 28, (1)
\[
\varDelta_1 = \lambda a_{11}\,p_1 + \lambda a_{12}\,p_2 + \lambda a_{13}\,p_3 + \cdots \lambda a_{1n}\,p_n
\]
\[
= \lambda(a_{11}\,p_1 + a_{12}\,p_2 + \cdots + a_{1n}\,p_n) = \lambda \cdot \varDelta.
\]
Andererseits ist $\varDelta_1$ dargestellt durch die Matrix:
\[
\varDelta_1 = \begin{vmatrix}
\lambda a_{11}, & \lambda a_{12}, & \lambda a_{13} \ldots \lambda a_{1n}\\
a_{21} & a_{22} & a_{23} \ldots . a_{2n}\\
a_{31} & \ldots\ldots\ldots\ldots\ldots & a_{3n}\\
\ldots\ldots\ldots\ldots\ldots\ldots\ldots\ldots\\
a_{n1} & a_{n2} & a_{n3} \ldots a_{nn}
\end{vmatrix} = \lambda \cdot \varDelta.
\]
Man sieht daraus: wenn alle Elemente der ersten Zeile einen gemeinschaftlichen Factor $\lambda$ besitzen, so ist die Determinante selbst ein Product, dessen einer Factor $\lambda$ ist. Dies gilt offenbar auch, wenn $\lambda$ Factor aller Elemente einer \textit{beliebigen} Zeile oder Colonne ist. Ebenso erkennt man umgekehrt den Satz:
\begin{quote}
\glqq Eine Determinante wird mit einem Factor $\lambda$ multiplicirt, indem man alle Elemente einer Zeile oder Colonne mit demselben multiplicirt.\grqq{}
\end{quote}
39. \textit{Verbindung des Additions- und Multiplicationssatzes.} Wir hatten in Nr. 35 die Gleichungen $(E)$ erhalten:
\[
\varDelta_0 = a_{11}\,p_1 + a_{12}\,p_2 + \cdots a_{1n}\,p_n,
\]
\[
\varDelta_1 = b_{11}\,p_1 + b_{12}\,p_2 + \cdots b_{1n}\,p_n,
\]
\[
\varDelta_2 = \varDelta_0 + \varDelta_1.
\]
Setzen wir nun fest, dass die Elemente $b_{ii}$ den Bedingungen genügen:
\[
\alpha)\qquad b_{11} = \lambda a_{21}, \quad b_{12} = \lambda a_{22}, \quad b_{13} = \lambda a_{23} \cdots b_{1n} = \lambda a_{2n},
\]
dann wird wegen der ersten Relation in Nr. 34 $(D)$
\[
\varDelta_1 = \lambda(a_{21}\,p_1 + a_{22}\,p_2 + \cdots a_{2n}\,p_n) = 0,
\]
folglich
\[
\varDelta_2 = \varDelta_0;
\]
das heisst:
\begin{quote}
\glqq Addirt man die mit der willkürlichen Grösse $\lambda$ multiplicirt gedachten Elemente der zweiten Zeile zu den entsprechenden der ersten Zeile, so ändert sich der Werth der Determinante nicht.\grqq{}
\end{quote}


% ===== pdf idx 57, printed 43.0 =====
% !PRINTED 43
\weberpage{43}
Auch hier ist die Verallgemeinerung sofort einzusehen. Wir ersetzen die Bedingungen $\alpha$) durch die allgemeineren
\[
\beta)\qquad b_{11} = \lambda a_{\varrho 1}, \quad b_{12} = \lambda a_{\varrho 2} \cdots b_{1n} = \lambda a_{\varrho n},
\]
und erinnern uns an die Vertauschbarkeit von Zeilen und Colonnen. Dann ist der Satz evident:
\begin{quote}
\glqq Der Werth einer Determinante bleibt ungeändert, wenn man die mit dem nämlichen Factor multiplicirt gedachten Elemente einer Zeile oder Colonne zu den entsprechenden Elementen einer andern Zeile bezw. Colonne addirt.\grqq{}
\end{quote}
Dass dieser Process ohne Einfluss auf den Werth der Determinante beliebig wiederholt werden kann, bedarf wohl nur der Erwähnung.

40. \textit{Anwendungen.} \textit{Beispiel} 1. Die Determinante
\[
\begin{vmatrix} x^2+y^2 & x & y & 1\\ a_3 & -a_2 & -a_1 & 1\\ b_3 & -b_2 & -b_1 & 1\\ c_3 & -c_2 & -c_1 & 1\end{vmatrix}
= \begin{vmatrix} 0 & x & y & 1\\ a_3+a_2 x+a_1 y & -a_2 & -a_1 & 1\\ b_3+b_2 x+b_1 y & -b_2 & -b_1 & 1\\ c_3+c_2 x+c_1 y & -c_2 & -c_1 & 1\end{vmatrix}
\]
\[
= \begin{vmatrix} 0 & x & y & 1\\ a_3+a_2 x+a_1 y & -(x+a_2) & -(y+a_1) & 0\\ b_3+b_2 x+b_1 y & -(x+b_2) & -(y+b_1) & 0\\ c_3+c_2 x+c_1 y & -(x+c_2) & -(y+c_1) & 0\end{vmatrix}
\]
\[
= -\begin{vmatrix} a_3+a_2 x+a_1 y & x+a_2 & y+a_1\\ b_3+b_2 x+b_1 y & x+b_2 & y+b_1\\ c_3+c_2 x+c_1 y & x+c_2 & y+c_1\end{vmatrix}.
\]
Hier wurde die zweite Determinante aus der ersten erhalten, indem man die mit $x$ multiplicirt gedachte zweite Colonne und die mit $y$ multiplicirt gedachte dritte Colonne von der ersten subtrahirte. Daraus ging die dritte hervor durch Addition der mit $-1$ multiplicirten ersten Zeile zu den drei übrigen. Die letzte Determinante liefert die Entwicklung der eben erhaltenen nach den Elementen der letzten Colonne.

\textit{Beispiel} 2.
\[
\varDelta = \begin{vmatrix} 1 & 1 & 1\\ \alpha & \beta & \gamma\\ \alpha^2 & \beta^2 & \gamma^2\end{vmatrix}
= \begin{vmatrix} 1 & 0 & 0\\ \alpha & \beta-\alpha & \gamma-\alpha\\ \alpha^2 & \beta^2-\alpha^2 & \gamma^2-\alpha^2\end{vmatrix}
=
\]
%CONT>>


% ===== pdf idx 58, printed 44.0 =====
% !PRINTED 44
\weberpage{44}
%<<CONT
\[
= (\beta-\alpha)(\gamma-\alpha) \begin{vmatrix} 1 & 0 & 0\\ \alpha & 1 & 1\\ \alpha^2 & \beta+\alpha & \gamma+\alpha\end{vmatrix}
\]
\[
= (\beta-\alpha)(\gamma-\alpha) \begin{vmatrix} 1 & 0 & 0\\ \alpha & 1 & 0\\ \alpha^2 & \beta+\alpha & \gamma-\beta\end{vmatrix}
\]
\[
= (\beta-\alpha)(\gamma-\alpha)(\gamma-\beta) \begin{vmatrix} 1 & 0 & 0\\ \alpha & 1 & 0\\ \alpha^2 & \beta+\alpha & 1\end{vmatrix} = (\beta-\alpha)(\gamma-\alpha)(\gamma-\beta).
\]
Ganz in derselben Weise erhält man der Reihe nach:

\textit{Beispiel} 3.
\[
\varDelta = \begin{vmatrix} 1 & 1 & 1 & 1 & 1\\ \alpha_1 & \alpha_2 & \alpha_3 & \alpha_4 & \alpha_5\\ \alpha_1^2 & \alpha_2^2 & \alpha_3^2 & \alpha_4^2 & \alpha_5^2\\ \alpha_1^3 & \alpha_2^3 & \alpha_3^3 & \alpha_4^3 & \alpha_5^3\\ \alpha_1^4 & \alpha_2^4 & \alpha_3^4 & \alpha_4^4 & \alpha_5^4\end{vmatrix}
= \begin{vmatrix}
1 & 0 & 0 & 0 & 0\\
\alpha_1, & \alpha_2-\alpha_1, & \alpha_3-\alpha_1, & \alpha_4-\alpha_1, & \alpha_5-\alpha_1\\
\alpha_1^2, & \alpha_2^2-\alpha_1^2, & \alpha_3^2-\alpha_1^2, & \alpha_4^2-\alpha_1^2, & \alpha_5^2-\alpha_1^2\\
\alpha_1^3, & \alpha_2^3-\alpha_1^3, & \alpha_3^3-\alpha_1^3, & \alpha_4^3-\alpha_1^3, & \alpha_5^3-\alpha_1^3\\
\alpha_1^4, & \alpha_2^4-\alpha_1^4, & \alpha_3^4-\alpha_1^4, & \alpha_4^4-\alpha_1^4, & \alpha_5^4-\alpha_1^4
\end{vmatrix}
\]
\[
= (\alpha_2-\alpha_1)(\alpha_3-\alpha_1)(\alpha_4-\alpha_1)(\alpha_5-\alpha_1) \times
\]
% ?? uncertain: the row separators in the following two 4-column blocks mix commas and semicolon-like gaps in the scan; reproduced as printed, spacing approximate
\[
\times \begin{vmatrix}
1 & 0 & 0 & 0\\
\alpha_1 & 1 & 1 & 1\\
\alpha_1^2, & \alpha_1+\alpha_2 & {}; \alpha_1+\alpha_3 & , \alpha_1+\alpha_4, \alpha_1+\alpha_5\\
\alpha_1^3, & \alpha_1^2+\alpha_1\alpha_2+\alpha_2^2 & , \alpha_1^2+\alpha_1\alpha_3+\alpha_3^2 & , \cdots, \cdots\\
\alpha_1^4, & \alpha_1^3+\alpha_1^2\alpha_2+\alpha_1\alpha_2^2+\alpha_2^3 & , \alpha_1^3+\alpha_1^2\alpha_3+\alpha_1\alpha_3^2+\alpha_3^3 & , \cdots, \cdots
\end{vmatrix}
\]
\[
= (\alpha_2-\alpha_1)(\alpha_3-\alpha_1)(\alpha_4-\alpha_1)(\alpha_5-\alpha_1) \times
\]
\[
\times \begin{vmatrix}
1 & 0 & 0 & 0\\
\alpha_1 & 1 & 0 & 0\\
\alpha_1^2, & \alpha_1+\alpha_2, & \alpha_3-\alpha_2, & \alpha_4-\alpha_2, \quad \alpha_5-\alpha_2\\
\alpha_1^3, & \dfrac{\alpha_2^2-\alpha_1^2}{\alpha_2-\alpha_1}, & (\alpha_3-\alpha_2)\displaystyle\sum_{1}^{3}\alpha_i, & \alpha_4^2-\alpha_2^2+\alpha_1(\alpha_4-\alpha_2), \cdots\\
\alpha_1^4, & \dfrac{\alpha_2^4-\alpha_1^4}{\alpha_2-\alpha_1}, & \cdots, & \cdots, \cdots
\end{vmatrix}
\]
\[
= (\alpha_2-\alpha_1)(\alpha_3-\alpha_1)(\alpha_4-\alpha_1)(\alpha_5-\alpha_1)(\alpha_3-\alpha_2)(\alpha_4-\alpha_2)(\alpha_5-\alpha_2) \times
\]
\[
\times \begin{vmatrix}
1 & 0 & 0 & 0 & 0\\
\alpha_1 & 1 & 0 & 0 & 0\\
\alpha_1^2, & \alpha_1+\alpha_2, & 1 & 1 & 1\\
\alpha_1^3, & \dfrac{\alpha_2^2-\alpha_1^2}{\alpha_2-\alpha_1}, & \displaystyle\sum_{1}^{3}\alpha_i, & \alpha_1+\alpha_2+\alpha_4, & \alpha_1+\alpha_2+\alpha_5\\
\alpha_1^4, & \dfrac{\alpha_2^3-\alpha_1^3}{\alpha_2-\alpha_1}, & \cdots, & \cdots, & \cdots
\end{vmatrix}
\]
%CONT>>


% ===== pdf idx 59, printed 45.0 =====
% !PRINTED 45
\weberpage{45}
%<<CONT
{\footnotesize
\[
\begin{aligned}
={}& (\alpha_2-\alpha_1)(\alpha_3-\alpha_1)(\alpha_4-\alpha_1)(\alpha_5-\alpha_1)(\alpha_3-\alpha_2)(\alpha_4-\alpha_2)(\alpha_5-\alpha_2) \times\\
&
\left|
\begin{array}{lllll}
1 & 0 & 0 & 0 & 0 \\
\alpha_1 & 1 & 0 & 0 & 0 \\
\alpha_1^2, & \alpha_1+\alpha_2, & 1 & 0 & 0 \\
\alpha_1^3, & \dfrac{\alpha_2^2-\alpha_1^2}{\alpha_2-\alpha_1}, & \displaystyle\sum_1^3\alpha_i, & \alpha_4-\alpha_3, & \alpha_5-\alpha_3 \\
\alpha_1^4, & \dfrac{\alpha_2^4-\alpha_1^4}{\alpha_2-\alpha_1}, & \cdots, & (\alpha_4-\alpha_3)\displaystyle\sum_1^4\alpha_i, & (\alpha_5-\alpha_3)\left(\displaystyle\sum_1^5\alpha_i-\alpha_4\right)
\end{array}
\right|
\end{aligned}
\]
}
\[
= \prod_1^{5\,(i>k)} \frac{\alpha_i-\alpha_k}{\alpha_5-\alpha_4} \times
\left|
\begin{array}{lllll}
1 & 0 & 0 & 0 & 0 \\
\alpha_1 & 1 & 0 & 0 & 0 \\
\alpha_1^2, & \alpha_1+\alpha_2, & 1 & 0 & 0 \\
\alpha_1^3, & \dfrac{\alpha_2^2-\alpha_1^2}{\alpha_2-\alpha_1}, & \displaystyle\sum_1^3\alpha_i, & 1, & 1 \\
\alpha_1^4, & \dfrac{\alpha_2^4-\alpha_1^4}{\alpha_2-\alpha_1}, & \cdots, & \displaystyle\sum_1^4\alpha_i, & \displaystyle\sum_1^5\alpha_i-\alpha_4
\end{array}
\right|
\]
\[
= \prod_1^{5\,(i>k)} (\alpha_i-\alpha_k) \times
\left|
\begin{array}{lllll}
1 & 0 & 0 & 0 & 0 \\
\alpha_1 & 1 & 0 & 0 & 0 \\
\alpha_1^2, & \alpha_1+\alpha_2, & 1 & 0 & 0 \\
\alpha_1^3, & \dfrac{\alpha_2^2-\alpha_1^2}{\alpha_2-\alpha_1}, & \displaystyle\sum_1^3\alpha_i, & 1, & 0 \\
\alpha_1^4, & \dfrac{\alpha_2^4-\alpha_1^4}{\alpha_2-\alpha_1}, & \cdots, & \displaystyle\sum_1^4\alpha_i, & 1
\end{array}
\right|.
\]

Die Determinante im letzten Ausdrucke besitzt aber (vergl. Nr.\,30, Beispiel 4) den Werth 1, also hat man

\[
\Delta = (\alpha_2-\alpha_1)(\alpha_3-\alpha_1)(\alpha_4-\alpha_1)(\alpha_5-\alpha_1)(\alpha_3-\alpha_2)(\alpha_4-\alpha_2)
(\alpha_5-\alpha_2)(\alpha_4-\alpha_3)(\alpha_5-\alpha_3)(\alpha_5-\alpha_4).
\]

\textit{Anmerkung.} Dergleichen Determinanten wie (2) und (3) werden übrigens rascher auf folgende Weise berechnet. Man denkt sich $\alpha_1$ bis $\alpha_5$ als Veränderliche, und frägt sich, unter welcher Bedingung $\Delta=0$ wird. Dies geschieht aber, so oft zwei Colonnen in $\Delta$ einander gleich werden, d. i. also, wenn $\alpha_2=\alpha_1$, $\alpha_3=\alpha_1$ $\cdot$ $\cdot$ u. s. w., d. h. $\alpha_2-\alpha_1$, $\alpha_3-\alpha_1$ u. s. w. müssen Factoren von $\Delta$ sein. Auf diese Weise erhält man sofort
%CONT>>


% ===== pdf idx 60, printed 46.0 =====
% !PRINTED 46
\weberpage{46}
%<<CONT
\[
\Delta = \varrho \prod_1^5 (\alpha_i-\alpha_k), \quad i>k.
\]

Es handelt sich noch $\varrho$ zu berechnen. In unserm Falle kann $\varrho$ nur noch eine Constante sein und zwar gleich $+1$, wie man durch Vergleichung des Diagonalgliedes von $\Delta$, mit dem Product aller Minuenden in $\Pi$ erkennt.

\textit{Beispiel} 4. Die Sätze dieses Paragraphen, insbesondere der letzte, leisten bei Berechnung von Determinanten, namentlich wenn die Elemente grosse Zahlen sind, werthvolle Dienste. Wir wollen dies in dem folgenden Beispiele zeigen.

Es ist:
\[
\Delta =
\left|
\begin{array}{cccc}
18 & 36 & 58 & 50 \\
26 & 39 & 80 & 78 \\
17 & 39 & 55 & 45 \\
9 & 16 & 27 & 23
\end{array}
\right|
\]
zu berechnen.

Subtrahirt man die zweifache letzte Zeile von der ersten, ihr dreifaches von der zweiten und ihr zweifaches von der dritten, so kommt:
\[
\Delta =
\left|
\begin{array}{rrrr}
0 & 4 & 4 & 4 \\
-1 & -9 & -1 & 9 \\
-1 & 7 & 1 & -1 \\
9 & 16 & 27 & 23
\end{array}
\right|.
\]

Subtrahirt man von der letzten Zeile die vierfache erste, addirt sodann die dritte zur zweiten und subtrahirt endlich die letzte Colonne von der zweiten und dritten, dann erhält man der Reihe nach:
\[
\Delta =
\left|
\begin{array}{rrrr}
0 & 4 & 4 & 4 \\
-1 & -9 & -1 & 9 \\
-1 & 7 & 1 & -1 \\
9 & 0 & 11 & 7
\end{array}
\right|
=
\left|
\begin{array}{rrrr}
0 & 4 & 4 & 4 \\
-2 & -2 & 0 & 8 \\
-1 & 7 & 1 & -1 \\
9 & 0 & 11 & 7
\end{array}
\right|
\]
\[
=
\left|
\begin{array}{rrrr}
0 & 0 & 0 & 4 \\
-2 & -10 & -8 & 8 \\
-1 & 8 & 2 & -1 \\
9 & -7 & 4 & 7
\end{array}
\right|
= +4
\left|
\begin{array}{rrr}
2 & 10 & 8 \\
-1 & 8 & 2 \\
9 & -7 & 4
\end{array}
\right|.
\]

Dividirt man in der letzten Determinante einmal die letzte Colonne, dann die erste Zeile durch den Factor 2, addirt alsdann die erste Zeile zur zweiten und subtrahirt das neunfache der ersten Zeile von der dritten, so erhält man endlich:
%CONT>>


% ===== pdf idx 61, printed 47.0 =====
% !PRINTED 47
\weberpage{47}
%<<CONT
\[
\Delta = 16
\left|
\begin{array}{rrr}
1 & 5 & 2 \\
-1 & 8 & 1 \\
9 & -7 & 2
\end{array}
\right|
= 16
\left|
\begin{array}{rrr}
1 & 5 & 2 \\
0 & 13 & 3 \\
0 & -52 & -16
\end{array}
\right|
= -16
\left|
\begin{array}{rr}
13 & 3 \\
52 & 16
\end{array}
\right|
\]
\[
= -64
\left|
\begin{array}{rr}
13 & 3 \\
13 & 4
\end{array}
\right|
= -64\cdot 13
\left|
\begin{array}{rr}
1 & 3 \\
1 & 4
\end{array}
\right|
= -64\cdot 13(4-3) = -64\cdot 13 = -832.
\]

\textit{Beispiel} 5. Man berechne die Determinante:
\[
\Delta =
\left|
\begin{array}{cccc}
0 & \varrho & \sigma & \tau \\
\varrho & 0 & \tau & \sigma \\
\sigma & \tau & 0 & \varrho \\
\tau & \sigma & \varrho & 0
\end{array}
\right|.
\]

Addirt man die zweite, dritte und vierte Zeile zur ersten, so erhält man:
\[
\Delta =
\left|
\begin{array}{cccc}
\varrho+\sigma+\tau & \varrho+\sigma+\tau & \varrho+\sigma+\tau & \varrho+\sigma+\tau \\
\varrho & 0 & \tau & \sigma \\
\sigma & \tau & 0 & \varrho \\
\tau & \sigma & \varrho & 0
\end{array}
\right|
= (\varrho+\sigma+\tau)
\left|
\begin{array}{cccc}
1 & 1 & 1 & 1 \\
\varrho & 0 & \tau & \sigma \\
\sigma & \tau & 0 & \varrho \\
\tau & \sigma & \varrho & 0
\end{array}
\right|.
\]

Subtrahirt man ferner die zweite Colonne von der ersten, die dritte von der zweiten, die vierte von der dritten, so kommt:
\[
\Delta = (\varrho+\sigma+\tau)
\left|
\begin{array}{cccc}
0 & 0 & 0 & 1 \\
\varrho & -\tau & \tau-\sigma & \sigma \\
\sigma-\tau & \tau & -\varrho & \varrho \\
\tau-\sigma & \sigma-\varrho & \varrho & 0
\end{array}
\right|
= (\varrho+\sigma+\tau)
\left|
\begin{array}{ccc}
\varrho & -\tau & \tau-\sigma \\
\sigma-\tau & \tau & -\varrho \\
\tau-\sigma & \sigma-\varrho & \varrho
\end{array}
\right|.
\]

Addirt man hier die dritte zur zweiten Zeile, dann entsteht:
\[
\Delta = (\varrho+\sigma+\tau)
\left|
\begin{array}{ccc}
\varrho & -\tau & \tau-\sigma \\
0 & \tau+\sigma-\varrho & 0 \\
\tau-\sigma & \sigma-\varrho & \varrho
\end{array}
\right|
= (\varrho+\sigma+\tau)(\tau+\sigma-\varrho)
\left|
\begin{array}{cc}
\varrho & \tau-\sigma \\
\tau-\sigma & \varrho
\end{array}
\right|
\]
\[
= (\varrho+\sigma+\tau)(\tau+\sigma-\varrho)(\varrho+\tau-\sigma)(\varrho-\tau+\sigma).
\]

\begin{center}
\textbf{§ 4. Partialdeterminanten höherer Ordnung.}
\end{center}

41. \textit{Begriff der Unterdeterminante zweiter Ordnung einer Determinante $\Delta$.} Es möge noch einmal daran erinnert werden, dass in keinem Gliede $G_\varrho$ einer Determinante zwei Elemente $a_{ik}$ und $a_{rs}$ als Factoren auftreten können, für welche $i=r$, oder $k=s$ wäre, da jedes Glied aus jeder Zeile und jeder Colonne ein und nur ein Element besitzt. Wählen wir nun aus allen $N$ Gliedern $G_\varrho$ diejenigen heraus, welche
%CONT>>


% ===== pdf idx 62, printed 48.0 =====
% !PRINTED 48
\weberpage{48}
%<<CONT
das Product $a_{ik}\cdot a_{rs}$ ($i\gtrless r$, $k\gtrless s$) als Factor haben und bezeichnen sie der Reihe nach mit

\weq{1}{G_{k_1},\ \ G_{k_2},\ \ G_{k_3},\ \cdots G_{k_\sigma}.}

Das Product aller übrigen Factoren eines dieser ausgewählten Glieder sei \textit{inclusive Modul} $\varepsilon_S$ mit $K_{k_\varrho}$ bezeichnet, so dass also die Beziehungen gelten:

\weq{2}{G_{k_1}=a_{ik}a_{rs}\cdot K_{k_1},\quad G_{k_2}=a_{ik}a_{rs}K_{k_2},\cdots \text{u.\,s.\,w.},\quad G_{k_\sigma}=a_{ik}a_{rs}K_{k_\sigma}.}

Die Summe aller dieser Glieder ist daher dargestellt durch

\weq{3}{\sum_{\varrho=1}^{\varrho=\sigma} G_{k_\varrho} = a_{ik}\,a_{rs}\sum_{\varrho=1}^{\varrho=\sigma} K_{k_\varrho},}

oder wenn wir

\weq{4}{\sum_{\varrho=1}^{\varrho=\sigma} K_{k_\varrho} = A_{\substack{ir\\ks}}}

setzen, so erhalten wir:

\weq{5}{\sum_{\varrho=1}^{\varrho=\sigma} G_{k_\varrho} = a_{ik}\,a_{rs}\cdot A_{\substack{ir\\ks}}.}

Wir nennen nun $A_{\substack{ir\\ks}}$ die zum Producte $a_{ik}a_{rs}$ gehörige \textit{Partialdeterminante zweiter Ordnung} von $\Delta$, oder Minor zweiter Ordnung.

42. \textit{Beziehung des Minors zweiter Ordnung $A_{\substack{ir\\ks}}$ zum Minor erster Ordnung $A_{ik}$.} Die oben ausgewählten Glieder $G_{k_\varrho}$ müssen, da sie ja den Factor $a_{ik}$ enthalten, sämmtlich unter den bereits früher (Nr.\,24) herausgegriffenen Gliedern $G_{i_\varrho}$ auftreten, welche die Form besassen:

\weq{6}{G_{i_\varrho} = a_{ik}\cdot H_{i_\varrho}.}

Jedes Glied $G_{i_\varrho}$, dessen Factor $H_{i_\varrho}$ das Element $a_{rs}$ enthält, ist auch ein Glied $G_{k_\varrho}$. Bezeichnen wir also jeden solchen Factor $H_{i_\varrho}$, in welchem $a_{rs}$ auftritt, mit $H_{k_\varrho}$, dann besitzt dieser selbst wieder die Form

\weq{7}{H_{k_\varrho} = a_{rs}\cdot K_{k_\varrho},}

wenn wir unter $K_{k_\varrho}$ die Gesammtheit aller anderen Elemente in $H_{k_\varrho}$ verstehen, wie wir bereits in Nr.\,41 festgesetzt haben.

Wir haben also auch wegen Gleichung (4) in Nr.\,41 für die Summe aller dieser Factoren $H_{k_\varrho}$ die Beziehung

\[
\sum_{\varrho=1}^{\varrho=\sigma} H_{k_\varrho} = \sum_{\varrho=1}^{\varrho=\sigma} a_{rs} K_{k_\varrho} = a_{rs}\sum_{\varrho=1}^{\varrho=\sigma} K_{k_\varrho} = a_{rs}\,A_{\substack{ir\\ks}}.
\]
%CONT>>


% ===== pdf idx 63, printed 49.0 =====
% !PRINTED 49
\weberpage{49}
\S.\,4. Partialdeterminanten höherer Ordnung.

Nun ist aber anderntheils die Summe aller Grössen $H_{i_\varrho}$, aus denen die Grössen $H_{k_\varrho}$ gewählt sind, gerade die Determinante $A_{ik}$, und zwar der Minor zum Elemente $a_{ik}$ in der Determinante $\varDelta$. Ganz ebenso ist die Summe aller Grössen $K_{k_\varrho}$ die Determinante $A_{\substack{ir\\ks}}$, und zwar der Minor zum Elemente $a_{rs}$ in der Determinante $A_{ik}$. Wir erkennen sonach: \textit{Die Partialdeterminante $A_{\substack{ir\\ks}}$ zweiter Ordnung von $\varDelta$ ist Minor erster Ordnung von $A_{ik}$, oder allgemeiner ausgesprochen: Unterdeterminanten von Unterdeterminanten sind selbst wieder Minoren der ursprünglichen Determinante.}

Es mag hierbei noch erwähnt sein, dass der Minor $A_{\substack{ir\\ks}}$ nicht nur als Unterdeterminante zum Elemente $a_{rs}$ in der Determinante $A_{ik}$ dargestellt werden kann, sondern es kann genau ebenso gezeigt werden, dass $A_{ir}$ auch Unterdeterminante zu $a_{ik}$ in der Unterdeterminante $A_{rs}$ ist\footnote{Es ist also offenbar: $A_{\substack{ir\\ks}} = \dfrac{\partial A_{ik}}{\partial a_{rs}} = \dfrac{\partial A_{rs}}{\partial a_{ik}} = \dfrac{\partial^2 \varDelta}{\partial a_{ik}\,\partial a_{rs}}$.}.

43. \textit{Beispiel.} Greifen wir aus der Determinante
\[
\varDelta = \begin{vmatrix} p & q & r & s\\ a & b & c & d\\ t & u & v & w\\ k & x & y & z \end{vmatrix}
\]
jene Glieder $G_{i_\varrho}$ heraus, welche das Element $r$ enthalten, so ist die Summe aller dieser Glieder nach Nr. 26 dargestellt durch
\[
\sum G_{i_\varrho} = r \cdot A_{13} = (-1)^{1+3} r \begin{vmatrix} a & b & d\\ t & u & w\\ k & x & z \end{vmatrix}
= r[auz - axw - tbz + txd + kbw - kud].
\]
% ?? uncertain: the printed summation sign carries a small ornamental hook/prime mark above-left throughout this page and the next; rendered here uniformly as \sum.

Wählen wir nun aus diesen alle Glieder $G_{k_\varrho}$, die überdies das Element $d$ besitzen, so erhalten wir
\[
\sum G_{k_\varrho} = r[txd - kud] = rd[xt - ku] = rd \begin{vmatrix} x & u\\ k & t \end{vmatrix} = rd\, A_{\substack{12\\34}}.
\]
Ganz ebenso gelangen wir zu dieser Summe, wenn wir aus $\varDelta$ zunächst alle Glieder $\overline{G}_{i_\varrho}$ mit dem Factor $d$ entnehmen. Es ist dann:
%CONT>>


% ===== pdf idx 64, printed 50.0 =====
% !PRINTED 50
\weberpage{50}
Erster Theil. Theorie der Determinanten.

%<<CONT
\[
\sum \overline{G}_{i_\varrho} = d \cdot A_{24} = (-1)^{2+4} d \cdot \begin{vmatrix} p & q & r\\ t & u & v\\ k & x & y \end{vmatrix}
= d[puy - pxv - tqy + txr + kqv - kur].
\]
Sondert man in der letzten Form die Glieder mit dem Factor $r$ ab, so kommt:
\[
\sum G_{k_\varrho} = \sum \overline{G}_{k_\varrho} = d[txr - kur] = dr(tx - ku) = dr \begin{vmatrix} x & u\\ k & t \end{vmatrix}.
\]

44. \textit{Begriff des Minors dritter Ordnung.} Die Erweiterung der vorhergehenden Entwicklung hat nun keine Schwierigkeiten mehr. Wählt man aus den $N$ Gliedern $G_\varrho$ jene aus, welche den Factor
\[
a_{ik}\, a_{rs}\, a_{lm}
\]
besitzen, wobei wieder $i, r, l$ einerseits und $k, s, m$ andrerseits von einander verschieden sein müssen, so bilden ganz wie vorhin die übrigen Factoren aller dieser Glieder selbst wieder eine Determinante, die wir analog incl. Vorzeichen mit $A_{\substack{irl\\ksm}}$ bezeichnen. Wir nennen sie Partialdeterminante oder Minor dritter Ordnung, und die Summe aller dieser Glieder mit dem Factor $a_{ik} \cdot a_{rs} \cdot a_{lm}$ ist dargestellt durch
\[
a_{ik} . a_{rs} . a_{lm} . A_{\substack{irl\\ksm}}.
\]
Wir können zu diesem Minor auf verschiedene Arten gelangen; man kann ihn auffassen als \textit{Minor erster Ordnung} zum Elemente $a_{lm}$ in der Partialdeterminante $A_{\substack{ir\\ks}}$, oder zum Elemente $a_{rs}$ im Minor $A_{\substack{il\\km}}$ oder zum Elemente $a_{ik}$ in der Unterdeterminante $A_{\substack{rl\\sm}}$. Man kann ihn aber auch deuten als \textit{Minor zweiter Ordnung} von $A_{ik}$, oder von $A_{rs}$, oder von $A_{lm}$. Wir sehen also wieder den Satz bestätigt: Unterdeterminanten von Unterdeterminanten sind selbst Minoren der ursprünglichen Determinante.

45. \textit{Begriff des Minors $p^{\text{ter}}$ Ordnung.} Greift man auf diese Weise ganz allgemein aus den $N$ Gliedern $G_\varrho$ diejenigen heraus, welche $p$ gemeinschaftliche Elemente als Factoren besitzen, etwa die Elemente:
\[
a_{i_1k_1}\ a_{i_2k_2}\ a_{i_3k_3} \ldots a_{i_pk_p},
\]
so umfasst die zugehörige Partialdeterminante $p^{\text{ter}}$ Ordnung alle jene Elemente, welche die nicht gemeinschaftlichen Factoren der ausgewählten Glieder bilden. Diese Partialdeterminante $p^{\text{ter}}$ Ordnung von $\varDelta$, welche wir inclusive Vorzeichen mit
\[
A_{\substack{i_1 i_2 \ldots i_p\\ k_1 k_2 \ldots k_p}}
\]
%CONT>>


% ===== pdf idx 65, printed 51.0 =====
% !PRINTED 51
\weberpage{51}
\S.\,4. Partialdeterminanten höherer Ordnung.

%<<CONT
bezeichnen, ist Partialdeterminante $(p-1)^{\text{ter}}$ Ordnung von $A_{i_1k_1}$, $(p-2)^{\text{ter}}$ Ordnung von $A_{\substack{i_1 i_2\\k_1 k_2}}$ u.\,s.\,w., kurz sie ist Minor jeder vorhergehenden Unterdeterminante, d.\,h. jedes Minors niedrigerer Ordnung.

46. \textit{Berechnung der Partialdeterminante zweiter Ordnung $A_{\substack{12\\12}}$.} Der Minor zweiter Ordnung $A_{\substack{12\\12}}$, also jener Minor, der zum Producte $a_{11}a_{22}$ gehört, kann nach den Entwicklungen in Nr. 42 aufgefasst werden als Unterdeterminante erster Ordnung des Elementes $a_{22}$ im Minor erster Ordnung $A_{11}$, dessen Matrix entsteht durch Unterdrückung der ersten Zeile und Colonne in $\varDelta$. Es muss daher $A_{\substack{12\\12}}$ genau ebenso aus $A_{11}$ erhalten werden, wie $A_{11}$ aus $\varDelta$; d.\,h. da die Unterdeterminante $A_{11}$ dargestellt ist durch
\[
A_{11} = \begin{vmatrix} a_{22} & a_{23} & a_{24} \ldots a_{2n}\\ a_{32} & a_{33} & a_{34} \ldots a_{3n}\\ a_{42} & . & . & . & . & . & .\\ . & . & . & . & . & . & .\\ a_{n2} & a_{n3} & a_{n4} & & & & a_{nn} \end{vmatrix},
\]
so erhält man daraus die Matrix von $A_{\substack{12\\12}}$ durch Unterdrückung der ersten Zeile und Colonne, d.\,i. jener beiden Reihen, die sich im Elemente $a_{22}$ schneiden. Dieselbe ist also:
\[
A_{\substack{12\\12}} = \begin{vmatrix} a_{33} & a_{34} & a_{35} \ldots a_{3n}\\ a_{43} & a_{44} & a_{45} \ldots a_{4n}\\ a_{53} & . & . & . & . & . & .\\ . & . & . & . & . & . & .\\ a_{n3} & a_{n4} & a_{n5} \ldots a_{nn} \end{vmatrix},
\]
eine Matrix, die offenbar auch aus $\varDelta$ erhalten werden kann durch Weglassung jener Zeilen und Colonnen, in deren Schnittpunkten die Elemente $a_{11}$ $a_{22}$ sich befinden.

47. \textit{Berechnung des Minors $A_{\substack{ir\\ks}}$.} Ist nun im allgemeinen Falle der Minor zweiter Ordnung $A_{\substack{ir\\ks}}$ zu berechnen, so bringen wir zunächst durch cyclische Vertauschung die Elemente $a_{ik}$ und $a_{rs}$ in der Determinante $\varDelta$ an die Stellen von $a_{11}$, bezw. $a_{22}$. Denken wir uns hierauf die erhaltene neue Determinante $D$ in den Elementen $b_{ik}$ geschrieben, so dass also
%CONT>>


% ===== pdf idx 66, printed 52.0 =====
% !PRINTED 52
\weberpage{52}
Erster Theil. Theorie der Determinanten.

%<<CONT
\[
D = \begin{vmatrix} b_{11} & b_{12} & b_{13} \ldots b_{1n}\\ b_{21} & b_{22} & . & . & b_{2n}\\ b_{31} & . & . & . & . & . & .\\ . & . & . & . & . & . & .\\ b_{n1} & . & . & . & . & b_{nn} \end{vmatrix},
\]
wobei
\[
\text{(I)}\quad
\begin{cases}
b_{11} = a_{ik},\quad b_{12} = a_{is},\quad b_{13} = a_{i1}\ldots\\
b_{21} = a_{rk},\quad b_{22} = a_{rs},\quad b_{23} = a_{r1}\ldots\\
b_{31} = a_{1k},\ldots\ \text{u.\,s.\,w.,}
\end{cases}
\]
dann stimmen noch immer, wie wir früher (Nr. 20) gesehen, beide Determinanten $D$ und $\varDelta$ in allen einzelnen Gliedern überein, abgesehen im Allgemeinen vom Vorzeichen. Letzteres hat sich nach Nr. 22 um den Factor $(-1)^\nu$, wo
\[
\nu = (ir) + (ks) + i + k + r + s,
\]
geändert, indem $D$ durch Substitution der Werthe von $b_{ik}$ aus (I) bis auf diesen Factor $(-1)^\nu$ in $\varDelta$ übergeht.

Während nun aber die Summe der Glieder, welche $b_{11}\,b_{22}$ in $D$ als Factor haben, dargestellt ist durch
\weq{1}{\sum \overline{G} = b_{11}\ b_{22} \cdot B_{\substack{12\\12}},}
besitzen die nämlichen Glieder, entnommen aus der Determinante $\varDelta$, die Factoren $a_{ik}\ a_{rs}$, und ihre Summe ist
\weq{2}{\sum G = a_{ik}\ a_{rs}\ A_{\substack{ir\\ks}}.}
Nun ist aber anderntheils
\[
\sum G = (-1)^\nu \sum \overline{G},
\]
daher die Beziehung:
\[
a_{ik}\ a_{rs}\ A_{\substack{ir\\ks}} = (-1)^\nu\ b_{11}\ b_{22}\ B_{\substack{12\\12}},
\]
oder weil ja $a_{ik} = b_{11}$ und $a_{rs} = b_{22}$, so wird endlich:
\weq{3}{A_{\substack{ir\\ks}} = (-1)^\nu\ B_{\substack{12\\12}}, \quad \nu = (ir) + (ks) + i + k + r + s.}
Die Matrix des Minors $B_{\substack{12\\12}}$ entsteht aber nach dem Vorigen durch Unterdrückung der ersten und zweiten Zeile und Colonne; das sind aber gerade jene Zeilen, die sich in $a_{ik}$ und $a_{rs}$ schneiden. Daher der Satz:
%CONT>>


% ===== pdf idx 67, printed 53.0 =====
% !PRINTED 53
\weberpage{53}
\S.\,4. Partialdeterminanten h\"oherer Ordnung.

\glqq Die Unterdeterminante zweiter Ordnung $A_{\substack{ir\\ks}}$ zum Producte $a_{ik}\,a_{rs}$
in der Determinante $\varDelta$ wird erhalten, wenn man in $\varDelta$ die $i^{\text{te}}$ und
$r^{\text{te}}$ Zeile, die $k^{\text{te}}$ und $s^{\text{te}}$ Colonne unterdr\"uckt und die
Determinante der so erhaltenen Matrix mit $(-1)^\nu$ multiplicirt, wo $\nu=(ir)+(ks)+i+k+r+s$.\grqq{}

48. \textit{Berechnung des Minors $p^{\text{ter}}$ Ordnung.} Ganz auf dieselbe Weise zeigt man,
dass die Matrix des Minors $p^{\text{ter}}$ Ordnung von $\varDelta$, welcher zum Producte

\[
a_{i_1 k_1}\, a_{i_2 k_2}\, a_{i_3 k_3} \dots a_{i_p k_p}
\]

geh\"ort, erhalten wird, indem man die $p$ Zeilen und Colonnen unterdr\"uckt, die sich in diesen
Elementen schneiden. Die Determinante, deren Matrix sich aus den restirenden Colonnen und Zeilen
zusammensetzt, stellt, wenn man sie noch mit $(-1)^\nu$ multiplicirt, wo nach Nr. 23

\[
\nu = (i_1\,i_2\,i_3 \cdots i_p) + (k_1\,k_2 \cdots k_p) + \sum_1^p i_\varrho + \sum_1^p k_\varrho \,,
\]

den Minor $p^{\text{ter}}$ Ordnung dar.

\textit{Beispiel.} Der Minor vierter Ordnung zu dem Producte $a_{62}\,a_{53}\,a_{14}\,a_{31}$ in

\[
\varDelta = \left|
\begin{array}{cccccc}
a_{11} & a_{12} & a_{13} & a_{14} & a_{15} & a_{16} \\
a_{21} & a_{22} & a_{23} & a_{24} & a_{25} & a_{26} \\
a_{31} & a_{32} & a_{33} & a_{34} & a_{35} & a_{36} \\
a_{41} & a_{42} & a_{43} & a_{44} & a_{45} & a_{46} \\
a_{51} & a_{52} & a_{53} & a_{54} & a_{55} & a_{56} \\
a_{61} & a_{62} & a_{63} & a_{64} & a_{65} & a_{66}
\end{array}
\right|
\]

ist dargestellt durch:

\[
A_{\substack{6513\\2341}} = (-1)^\nu \left| \begin{array}{cc} a_{25} & a_{26} \\ a_{45} & a_{46} \end{array} \right|\,,
\]

wo

\[
\begin{aligned}
\nu &= (6513)+(2341)+(6+5+1+3)+(2+3+4+1) \\
    &= 5+3+15+10 = 33\,,
\end{aligned}
\]

also:

\[
A_{\substack{6513\\2341}} = -\left(a_{25}\,a_{46} - a_{45}\,a_{26}\right)\cdot
\]

49. \textit{Die correspondirende Determinante des Minors zweiter Ordnung.} Wir hatten die Matrix
des Minors $A_{\substack{ir\\ks}}$ erhalten durch Unterdr\"uckung jenes Colonnen- und
Zeilenpaares, das sich in den Elementen $a_{ik}$ und $a_{rs}$ schneidet. Die n\"amlichen Zeilen
und Colonnen schneiden sich aber auch noch in zwei andern Elementen, n\"amlich in $a_{is}$ und
$a_{rk}$. Zum
%CONT>>


% ===== pdf idx 68, printed 54.0 =====
% !PRINTED 54
\weberpage{54}
Erster Theil. Theorie der Determinanten.

%<<CONT
Producte dieser beiden geh\"ort also ein Minor $A_{\substack{ir\\sk}}$, der die n\"amliche Matrix
besitzt. W\"ahrend aber die Unterdeterminante $A_{\substack{ir\\ks}}$ den Factor $(-1)^\nu$, wo

\[
\nu = (ir)+(ks)+i+r+k+s\,,
\]

enth\"alt, ist der Minor $A_{\substack{ir\\sk}}$ multiplicirt mit $(-1)^{\bar\nu}$, wo

\[
\bar\nu = (ir)+(sk)+i+r+s+k\,.
\]

Die beiden Gr\"ossen $\nu$ und $\bar\nu$ sind nothwendig um eine und nur eine Einheit
verschieden, d.\,h. es besteht die Relation

\[
\bar\nu = \nu \pm 1\,,
\]

oder

\[
(-1)^{\bar\nu} = -(-1)^\nu\cdot
\]

Daraus folgt, dass die beiden so erhaltenen Minoren nur durch das Vorzeichen sich unterscheiden,
d.\,h. dass

\[
A_{\substack{ir\\ks}} = -A_{\substack{ir\\sk}}\,.
\]

S\"ammtliche Glieder der Determinante $\varDelta$, welche die Glieder dieser beiden durch die
n\"amliche Matrix dargestellten Minoren zu Factoren haben, sind daher dargestellt durch den
Ausdruck:

\[
(A)\quad a_{ik}\,a_{rs}\,A_{\substack{ir\\ks}} + a_{is}\,a_{rk}\,A_{\substack{ir\\sk}}
= \left(a_{ik}\,a_{rs} - a_{is}\,a_{rk}\right) A_{\substack{ir\\ks}} \cdot
\left|
\begin{array}{ccccccccc}
a_{11} & a_{12} & \dots & a_{1,k-1} & a_{1,k+1} & \dots & a_{1,s-1} & a_{1,s+1} & \dots\, a_{1n} \\
a_{21} & a_{32} & \dots & a_{2,k-1} & a_{2,k+1} & \dots & a_{2,s-1} & a_{2,s+1} & \dots\, a_{2n} \\
a_{31} & \cdot & \cdot & \cdot & \cdot & \cdot & \cdot & \cdot & \cdot \\
\cdot & \cdot & \cdot & \cdot & \cdot & \cdot & \cdot & \cdot & \cdot \\
a_{i-1,1} & a_{i-1,2} & \cdot & \cdot & \cdot & \cdot & \cdot & \cdot & a_{i-1,n} \\
a_{i+1,1} & a_{i+1,1} & \cdot & \cdot & \cdot & \cdot & \cdot & \cdot & a_{i+1,n} \\
\cdot & \cdot & \cdot & \cdot & \cdot & \cdot & \cdot & \cdot & \cdot \\
a_{r-1,1} & a_{r-1,2} & \cdot & \cdot & \cdot & \cdot & \cdot & \cdot & a_{r-1,n} \\
a_{r+1,1} & a_{r+1,2} & \cdot & \cdot & \cdot & \cdot & \cdot & \cdot & a_{r+1,n} \\
\cdot & \cdot & \cdot & \cdot & \cdot & \cdot & \cdot & \cdot & \cdot \\
a_{n1} & a_{n2} & a_{n3} & \cdot & \cdot & \cdot & \cdot & \cdot & a_{nn}
\end{array}
\right|
\]
% [sic] second entry of the a_{i+1,\cdot} row prints "a_{i+1,1}" a second time where the
% pattern of the surrounding rows would give a_{i+1,2}; reproduced as printed.

\[
= (-1)^\nu \left|\begin{array}{cc} a_{ik} & a_{is} \\ a_{rk} & a_{rs} \end{array}\right| \cdot
\left|
\begin{array}{ccccccccc}
a_{i-1,1} & a_{i-1,2} & \cdot & \cdot & \cdot & \cdot & \cdot & \cdot & a_{i-1,n} \\
a_{i+1,1} & a_{i+1,1} & \cdot & \cdot & \cdot & \cdot & \cdot & \cdot & a_{i+1,n} \\
\cdot & \cdot & \cdot & \cdot & \cdot & \cdot & \cdot & \cdot & \cdot \\
a_{r-1,1} & a_{r-1,2} & \cdot & \cdot & \cdot & \cdot & \cdot & \cdot & a_{r-1,n} \\
a_{r+1,1} & a_{r+1,2} & \cdot & \cdot & \cdot & \cdot & \cdot & \cdot & a_{r+1,n} \\
\cdot & \cdot & \cdot & \cdot & \cdot & \cdot & \cdot & \cdot & \cdot \\
a_{n1} & a_{n2} & a_{n3} & \cdot & \cdot & \cdot & \cdot & \cdot & a_{nn}
\end{array}
\right|\,.
\]
% [sic] second entry of the a_{i+1,\cdot} row prints "a_{i+1,1}" a second time here as well.

Ist hierbei $i<r$, $k<s$ --- eine Voraussetzung, die immer erlaubt ist und von der ich auch, wie
man erkennt, in der zweiten Matrix der Relation $(A)$ stillschweigend Gebrauch gemacht habe ---
so sind in beiden Matrices Reihen und Colonnen gut geordnet, in der ersten nach Voraussetzung, in
der zweiten a priori, da sie aus der gut ge-
%CONT>>


% ===== pdf idx 69, printed 55.0 =====
% !PRINTED 55
\weberpage{55}
\S.\,4. Partialdeterminanten h\"oherer Ordnung.

%<<CONT
ordneten Matrix von $\varDelta$ nur durch Unterdr\"uckung von zwei Zeilen und Colonnen entsteht.
In diesem Falle reducirt sich $\nu$ auf die Summe der Indices allein, d.\,h. es ist:

\[
\nu = i+r+k+s\,,
\]

und man nennt nun die Determinante:

\[
(-1)^\nu \left|\begin{array}{cc} a_{ik} & a_{is} \\ a_{rk} & a_{rs} \end{array}\right|
\]

die \textit{correspondirende Determinante} zu der durch die Matrix von $A_{\substack{ir\\ks}}$
dargestellten Determinante.

50. \textit{Die correspondirende Determinante des Minors dritter Ordnung.} Die Matrix des Minors
dritter Ordnung $A_{\substack{irl\\ksm}}$ ergab sich aus der Matrix von $\varDelta$ durch
Unterdr\"uckung jener drei Zeilen und jener drei Colonnen, die sich in den Elementen $a_{ik}\,
a_{rs}\,a_{lm}$ schneiden. Genau auf dieselbe Matrix werden wir aber gef\"uhrt, wenn wir die
Minoren zu irgend einem der Producte bestimmen, das aus $a_{ik}\cdot a_{rs}\cdot a_{lm}$ durch
Vertauschung der ersten oder auch der zweiten Indices entsteht, also zu den f\"unf weiteren
Producten:

\[
a_{ik}\,a_{ls}\,a_{rm}\,,\quad a_{rk}\,a_{is}\,a_{lm}\,,\quad a_{rk}\,a_{ls}\,a_{im}\,,\quad
a_{lk}\,a_{is}\,a_{rm}\,,\quad a_{lk}\,a_{rs}\,a_{im}\,;
\]

denn dieselben enthalten keine andern Indices als die bereits in $a_{ik}\,a_{rs}\,a_{lm}$
auftretenden $i,\,r,\,l,\,k,\,s,\,m$, so dass bei Bildung der Matrix des jeweiligen Minors immer
dieselben Zeilen und Colonnen aus $\varDelta$ zu streichen sind. Die Matrix zu
$A_{\substack{irl\\ksm}}$ ist also auch Matrix f\"ur die Minoren, die aus $A_{\substack{irl\\ksm}}$
durch Permutation etwa der ersten Indices entstehen, also zu:

\[
A_{\substack{ilr\\ksm}}\,,\quad A_{\substack{ril\\ksm}}\,,\quad A_{\substack{rli\\ksm}}\,,\quad
A_{\substack{lir\\ksm}}\,,\quad A_{\substack{lri\\ksm}}\,.
\]

W\"ahrend nun aber der Minor $A_{\substack{irl\\ksm}}$ den Faktor $(-1)^\nu$ besitzt, wo

\[
\nu = (irl)+(ksm)+(i+r+l)+(k+s+m)\,,
\]

muss das Vorzeichen der \"ubrigen Minoren eben wegen der Permutation der Indices $(irl)$
wechseln, d.\,h. es m\"ussen die Relationen bestehen:

\[
A_{\substack{irl\\ksm}} = -A_{\substack{ilr\\ksm}} = -A_{\substack{ril\\ksm}} =
+A_{\substack{rli\\ksm}} = +A_{\substack{lir\\ksm}} = -A_{\substack{lri\\ksm}}\,.
\]

Es ist daher die Summe aller jener Glieder, welche ein Glied irgend eines dieser sechs Minoren
als Factor haben, dargestellt durch den Ausdruck:
%CONT>>


% ===== pdf idx 70, printed 56.0 =====
% !PRINTED 56
\weberpage{56}
Erster Theil. Theorie der Determinanten.

%<<CONT
\[
A_{\substack{irl\\ksm}} \cdot \bigl(
a_{ik}\,a_{rs}\,a_{lm} - a_{ik}\,a_{ls}\,a_{rm} - a_{rk}\,a_{is}\,a_{lm} + a_{rk}\,a_{ls}\,a_{im}
+ a_{lk}\,a_{is}\,a_{rm} - a_{lk}\,a_{rs}\,a_{im}
\bigr)
\]

\[
\left|
\begin{array}{ccccccccccc}
a_{11} & a_{22} & \dots & a_{1,k-1} & a_{1,k+1} & \dots & a_{1,s-1} & a_{1,s+1} & \dots & a_{1,m-1} & a_{1,m+1} \dots\, a_{1n} \\
\cdot & \cdot & \cdot & \cdot & \cdot & \cdot & \cdot & \cdot & \cdot & \cdot & \cdot \\
a_{i-1,1} & a_{i-1,2} & \cdot & \cdot & \cdot & \cdot & \cdot & \cdot & \cdot & \cdot & a_{i-1,n} \\
a_{i+1,1} & a_{i+1,2} & \cdot & \cdot & \cdot & \cdot & \cdot & \cdot & \cdot & \cdot & a_{i+1,n} \\
\cdot & \cdot & \cdot & \cdot & \cdot & \cdot & \cdot & \cdot & \cdot & \cdot & \cdot \\
a_{r-1,1} & a_{r-1,2} & \cdot & \cdot & \cdot & \cdot & \cdot & \cdot & \cdot & \cdot & a_{r-1,n} \\
a_{r+1,1} & a_{r+1,2} & \cdot & \cdot & \cdot & \cdot & \cdot & \cdot & \cdot & \cdot & a_{r+1,n} \\
\cdot & \cdot & \cdot & \cdot & \cdot & \cdot & \cdot & \cdot & \cdot & \cdot & \cdot \\
a_{l-1,1} & a_{l-1,2} & \cdot & \cdot & \cdot & \cdot & \cdot & \cdot & \cdot & \cdot & a_{l-1,n} \\
a_{l+1,1} & a_{l+1,2} & \cdot & \cdot & \cdot & \cdot & \cdot & \cdot & \cdot & \cdot & a_{l+1,n} \\
\cdot & \cdot & \cdot & \cdot & \cdot & \cdot & \cdot & \cdot & \cdot & \cdot & \cdot \\
a_{n1} & a_{n2} & \cdot & \cdot & \cdot & a_{n,s-1} & a_{n,s+1} & \cdot & \cdot & \cdot & a_{nn}
\end{array}
\right|
\]

\[
= (-1)^\nu \left|\begin{array}{ccc} a_{ik} & a_{is} & a_{im} \\ a_{rk} & a_{rs} & a_{rm} \\
a_{lk} & a_{ls} & a_{lm} \end{array}\right| \cdot
\]

Ist hier wiederum $i<r<l$, und $k<s<m$, dann reducirt sich der Exponent $\nu$ auch hier auf die
Summe der Indices, so dass

\[
\nu = i+r+l+k+s+m\,,
\]

und die Determinante

\[
(-1)^\nu \left|\begin{array}{ccc} a_{ik} & a_{is} & a_{im} \\ a_{rk} & a_{rs} & a_{rm} \\
a_{lk} & a_{ls} & a_{lm} \end{array}\right|
\]

wird dann die \textit{correspondirende Determinante} zu der durch die Matrix von
$A_{\substack{irl\\ksm}}$ dargestellten Determinante genannt.

51. \textit{Die correspondirende Determinante des Minors $p^{\text{ter}}$ Ordnung.} Die beiden
vorhergehenden speciellen F\"alle von correspondirenden Determinanten zeigen, dass deren Matrix
gerade durch jene Elemente gebildet ist, in welchen sich die bei Bildung des zugeh\"origen Minors
zweiter oder dritter Ordnung unterdr\"uckten Zeilen und Colonnen schneiden. Die Erweiterung der
Definition einer correspondirenden Determinante ergiebt sich daher in folgender Weise. Greift man
aus der Matrix einer Determinante $\varDelta$ irgend welche $\varrho$ Zeilen und $\varrho$
Colonnen heraus, so erhalten wir zwei kleinere Determinanten, und zwar 1) diejenige, deren Matrix
aus allen Elementen besteht, in welchen sich die $\varrho$ Colonnen und $\varrho$ Zeilen
schneiden; 2) diejenige, welche sich aus den in $\varDelta$ \"ubrig bleibenden Elementen
zusammensetzt. Die erste hat $\varrho$ Zeilen und $\varrho$ Colonnen, die zweite $\sigma$ Zeilen
und $\sigma$ Colonnen, so
%CONT>>


% ===== pdf idx 71, printed 57.0 =====
% !PRINTED 57
\weberpage{57}
%<<CONT
dass $\varrho + \sigma = n$. Bildet man das Product beider Determinanten, indem man die Summe der Glieder der einen mit der Summe der Glieder der andern multiplicirt, so bietet dasselbe --- vom Vorzeichen abgesehen --- alle Glieder der ursprünglichen Determinante, welche sowohl die Glieder der ersten als die der zweiten Determinante als Factor haben. Sind nun in den beiden Determinanten Zeilen und Colonnen gut geordnet und multiplicirt man die eine von ihnen mit der entsprechenden Potenz $(-1)^{\nu}$, so dass das Product derselben die betreffenden Glieder der ursprünglichen Determinanten \textit{inclusive} Vorzeichen liefert, dann nennt man die beiden Determinanten von $\varrho^2$, bezw. $\sigma^2$ Elementen \textit{correspondirende Determinanten}. Als Exponent $\nu$ kann hierbei die Summe aller Indices, welche entweder die Diagonalelemente der einen oder die der anderen Determinante besitzen, genommen werden.

Charakterisiren wir die Determinante $\Delta^{(\sigma)}$, deren Matrix nach Unterdrückung der $\varrho$ Colonnen und Zeilen sich aus den übrig bleibenden Elementen $a_{ik}$ zusammensetzt, durch die Elemente ihrer Diagonale setzen wir also
\[
\Delta^{(\sigma)} = + (a_{i_1 k_1}\, a_{i_2 k_2} \cdots a_{i_\sigma k_\sigma}),
\]
dann ist die correspondirende Determinante $\Delta^{(\varrho)}$ inclusive Vorzeichen dargestellt durch
\[
A_{\substack{i_1 i_2 \cdots i_\sigma \\ k_1 k_2 \cdots k_\sigma}},
\]
so dass also:
\[
\Delta^{(\varrho)} = A_{\substack{i_1 i_2 \cdots i_\sigma \\ k_1 k_2 \cdots k_\sigma}} = (-1)^{\nu} (a_{i'_1 k'_1}\, a_{i'_2 k'_2} \cdots a_{i'_\varrho k'_\varrho}),
\]
\[
\nu = \sum_1^\sigma i_\mu + \sum_1^\sigma k_\mu
\]
und
\[
\varrho + \sigma = n.
\]

52. \textit{Beispiel 1.} Sei gegeben die Determinante
\[
\Delta = \begin{vmatrix} a_{11} & a_{12} & a_{13} & x_1 \\ a_{21} & a_{22} & a_{23} & x_2 \\ a_{31} & a_{32} & a_{33} & x_3 \\ x_1 & x_2 & x_3 & 0 \end{vmatrix}.
\]
Entwickelt man die Determinante in die Summe ihrer 24 Glieder, so werden Glieder mit $x_1^2$, $x_2^2$, $x_3^2$ auftreten. Interessirt es uns nur, die Co\"efficienten von $x_1^2$, $x_2^2$, $x_3^2$ kennen zu lernen, so werden sie uns rascher durch die in Nr. 49 gegebene Methode geliefert. Darnach sind die Glieder mit $x_1^2$ enthalten in
%CONT>>


% ===== pdf idx 72, printed 58.0 =====
% !PRINTED 58
\weberpage{58}
%<<CONT
\[
(-1)^{\nu_1} b_{14}\, b_{41} \cdot B_{\substack{14 \\ 41}} = (-1)^{\nu_1} x_1^2 \begin{vmatrix} a_{22} & a_{23} \\ a_{32} & a_{33} \end{vmatrix},
\]
wo
\[
\nu_1 = (14) + (41) + 1 + 4 + 4 + 1 = 11,
\]
es ist also der Co\"efficient von $x_1^2$
\[
= - (a_{22} a_{33} - a_{32} a_{23});
\]
die Glieder, welche $x_2^2$ und $x_3^2$ enthalten, werden geliefert durch
\[
(-1)^{\nu_2} b_{24}\, b_{42} \cdot B_{\substack{24 \\ 42}} = (-1)^{\nu_2} x_2^2 \begin{vmatrix} a_{11} & a_{13} \\ a_{31} & a_{33} \end{vmatrix},
\]
\[
\nu_2 = (24) + (42) + 2 + 4 + 4 + 2 = 13;
\]
\[
(-1)^{\nu_3} b_{34}\, b_{43} \cdot B_{\substack{34 \\ 43}} = (-1)^{\nu_3} x_3^2 \begin{vmatrix} a_{11} & a_{12} \\ a_{21} & a_{22} \end{vmatrix},
\]
\[
\nu_3 = (34) + (43) + 3 + 4 + 4 + 3 = 15.
\]
Daher ist der Co\"efficient von $x_2^2$:
\[
- (a_{11} a_{33} - a_{31} a_{13}),
\]
und der Co\"efficient von $x_3^2$:
\[
- (a_{11} a_{22} - a_{21} a_{12}).
\]
\textit{Beispiel 2.} Sei gegeben die Determinante:
\[
\Delta = \begin{vmatrix}
a_1 & a_2 & a_3 & +\lambda^3 & 0 \\
b_1 & b_2 & b_3 & -3\lambda^2 & +\lambda^3 \\
c_1 & c_2 & c_3 & +3\lambda & -3\lambda^2 \\
d_1 & d_2 & d_3 & -1 & +3\lambda \\
e_1 & e_2 & e_3 & 0 & -1
\end{vmatrix}.
\]
% ?? uncertain: exact column split of the staggered $\lambda$-terms in this 5x5 determinant is reconstructed from the printed row-by-row layout (each row's third-named entry followed by two signed terms); the arithmetic (binomial-type staircase $\lambda^3,-3\lambda^2,3\lambda,-1$ shifted one row per column) is consistent with what is printed.
Wir bestimmen die correspondirenden Determinanten zu den Minoren zweiter Ordnung
\weq{1}{
\Delta_1^{(2)} = \begin{vmatrix} b_1 & b_2 & b_3 \\ c_1 & c_2 & c_3 \\ d_1 & d_2 & d_3 \end{vmatrix} = + (bcd)
}
und
\[
\Delta_2^{(2)} = \begin{vmatrix} c_1 & c_2 & c_3 \\ d_1 & d_2 & d_3 \\ e_1 & e_2 & e_3 \end{vmatrix} = + (cde).
\]
Die correspondirende Determinante zu $\Delta_1^{(2)}$ ist
\[
A_{\substack{234 \\ 123}} = \Delta_1^{(3)} = (-1)^{\nu_1} \begin{vmatrix} \lambda^3 & 0 \\ 0 & -1 \end{vmatrix},
\]
\[
= + \lambda^3
\]
\[
\nu_1 = 2 + 3 + 4 + 1 + 2 + 3 = 15,
\]
%CONT>>


% ===== pdf idx 73, printed 59.0 =====
% !PRINTED 59
\weberpage{59}
%<<CONT
oder auch
\[
\nu_1 = 1 + 5 + 4 + 5 = 15.
\]
Die correspondirende Determinante zu $\Delta_2^{(2)}$ ist
\[
A_{\substack{345 \\ 123}} = \Delta_2^{(3)} = (-1)^{\nu_2} \begin{vmatrix} \lambda^3 & 0 \\ -3\lambda^2 & +\lambda^3 \end{vmatrix},
\]
\[
= \lambda^6
\]
\[
\nu_2 = 3 + 4 + 5 + 1 + 2 + 3 = 18,
\]
oder auch
\[
\nu_2 = 1 + 2 + 4 + 5 = 12.
\]

\S.\,5.

\gesperrt{Entwicklung der Determinante in eine Summe von Producten correspondirender Determinanten.}

53. \textit{Summe aus Producten von Minoren $(n-2)^{\text{ter}}$ Ordnung und ihren correspondirenden $A_{\substack{ir \\ ks}}$.} Wir hatten in Nr. 28 und 29 bereits Entwicklungen der Determinante nach Minoren erster Ordnung gegeben und erhielten als eine derselben:
\weq{1}{
\Delta = a_{11} A_{11} + a_{12} A_{12} + \cdots a_{1n} A_{1n}.
}
Die Grössen $A_{1,k}$ sind nun aber selbst Determinanten und wir können sie wie die Determinante $\Delta$ nach ihren ersten Minoren, d. h. nach den zweiten Minoren der ursprünglichen Determinante entwickeln. Es ist dann:
\weq{2}{
\begin{aligned}
A_{11} &= a_{22} A_{\substack{12\\12}} + a_{23} A_{\substack{12\\13}} + \cdots\cdots a_{2n} A_{\substack{12\\1n}} \\
A_{12} &= a_{21} A_{\substack{12\\21}} + a_{23} A_{\substack{12\\23}} + \cdots\cdots a_{2n} A_{\substack{12\\2n}} \\
A_{13} &= a_{21} A_{\substack{12\\31}} + a_{22} A_{\substack{12\\32}} + a_{24} A_{\substack{12\\34}} \cdots a_{2n} A_{\substack{12\\3n}} \\
&\cdots\cdots\cdots\cdots\cdots\cdots\cdots\cdots\cdots\cdots\cdots\cdots \\
A_{1k} &= a_{21} A_{\substack{12\\k1}} + \cdots a_{2,k-1} A_{\substack{12\\k,k-1}} + a_{2,k+1} A_{\substack{12\\k,k+1}} + \cdots a_{2n} A_{\substack{12\\kn}} \ \text{u.\,s.\,w.}
\end{aligned}
}
Tragen wir die so erhaltenen Entwicklungen für die Minoren $A_{1k}$ in Gleichung (1) ein, so bekommen wir:
\weq{3}{
\begin{aligned}
\Delta = a_{11} a_{22} A_{\substack{12\\12}} &+ a_{11} a_{23} A_{\substack{12\\13}} + \cdots + a_{11} a_{2n} A_{\substack{12\\1n}} \\
{}+ a_{12} a_{21} A_{\substack{12\\21}} &+ a_{12} a_{23} A_{\substack{12\\23}} + \cdots \\
\cdots + a_{1k} a_{21} A_{\substack{12\\k1}} \cdots\cdots &+ a_{1n} a_{2n-1} A_{\substack{12\\n,n-1}}.
\end{aligned}
}


% ===== pdf idx 74, printed 60.0 =====
% !PRINTED 60
\weberpage{60}
Fragen wir uns, welche Elemente treten als Factoren der Minoren zweiter Ordnung in dieser Entwicklung auf, so erkennen wir, sie sind dadurch charakterisirt, dass ihre ersten Indices 1 und 2 sind, während in ihren zweiten Indices alle Amben der $n$ Zahlen $1\,2\cdots$ bis $n$ und zwar doppelt vertreten sind, einmal als Combination $(ik)$, dann als Combination $(ki)$. Fassen wir also je zwei Glieder $a_{2k}$ in der Summe
\weq{4}{
\Delta = \sum_{\substack{i=1\\k=1}}^{\substack{i=n\\k=n}} a_{1i}\, a_{2k}\, A_{\substack{12\\ik}},
}
wo
\[
i \geqq k,
\]
geeignet zusammen, nämlich solche zwei Glieder, deren eines die Elemente $a_{1i}\, a_{2k}$, und deren anderes die Elemente $a_{1k}\, a_{2i}$ als Factoren besitzt, so lässt sich die Gleichung (4) wegen
\[
A_{\substack{12\\ik}} = - A_{\substack{12\\ki}} \quad \text{(vgl. Nr. 49)}
\]
schreiben:
\weq{5}{
\Delta = \sum (a_{1i} a_{2k} - a_{1k} a_{2i}) A_{\substack{12\\ik}} = \sum_{\substack{i=1\\k=2}}^{\substack{k=n\\i=n-1}} \begin{vmatrix} a_{1i} & a_{2i} \\ a_{1k} & a_{2k} \end{vmatrix} \cdot A_{\substack{12\\ik}}, \quad k > i.
}
Diese Darstellung der Determinante $\Delta$ ist aber nichts Anderes, als eine Entwicklung derselben in eine Summe von Producten je zweier Minoren. Die erste Reihe der Minoren geht aus der Matrix hervor
\weq{A}{
\begin{vmatrix} a_{11} & a_{12} & \ldots & a_{1n} \\ a_{21} & a_{22} & \ldots & a_{2n} \end{vmatrix},
}
die zweite Reihe wird von den mit diesen Determinanten correspondirenden gebildet, deren Werth \textit{inclusive} Vorzeichen $(-1)^{\nu}$ dargestellt ist durch $A_{\substack{12\\ik}}$, so dass
\[
A_{\substack{12\\ik}} = (-1)^{\nu} (a_{3\varrho_1}\, a_{4\varrho_2} \cdots a_{n\varrho_{n-2}}),
\]
wo
\[
\varrho_1 < \varrho_2 \cdots < \varrho_{n-2}
\]
und von $i$ und $k$ verschieden,
und
\[
\nu = 1 + 2 + i + k \ \text{ist}.
\]
Die Summe aller Producte je zweier solcher correspondirender Determinanten liefert $\Delta$.


% ===== pdf idx 75, printed 61.0 =====
% !PRINTED 61
\weberpage{61}
54. \textit{Entwicklung in eine Summe von Producten aus Minoren $(n-3)^{\text{ter}}$
Ordnung und deren correspondirenden Determinanten.} Ganz ebenso lässt
sich $\Delta$ in eine Summe von Producten entwickeln, deren einer Factor
eine Determinante der dreizeiligen Matrix
\[
\begin{vmatrix}
a_{11}\ a_{12}\ \ldots\ a_{1n} \\
a_{21}\ a_{22}\ \ldots\ a_{2n} \\
a_{31}\ a_{32}\ \ldots\ a_{3n}
\end{vmatrix},
\]
und deren anderer Factor die correspondirende Determinante
\[
A_{\substack{1\ 2\ 3\\ i\ k\ r}} = (-1)^{v}\left(a_{4\varrho_1}\, a_{5\varrho_2}\cdots a_{n\varrho_{n-3}}\right), \quad \varrho_1 < \varrho_2 \cdots < \varrho_3
\]
und von $i$, $k$, $r$ verschieden,

und
\[
v = 1+2+3+i+k+r
\]
ist. Wir ersetzen zu dem Zwecke in der Gleichung (3) die Minoren
zweiter Ordnung durch ihre Entwicklung nach Minoren dritter Ord-
nung:
\[
A_{\substack{1\ 2\\ i\ k}} = a_{31}A_{\substack{1\ 2\ 3\\ i\ k\ 1}} + a_{32}A_{\substack{1\ 2\ 3\\ i\ k\ 2}} + \cdots,
\]
dann treten in der Summe nur solche drei Elemente als Factoren der
Minoren dritter Ordnung auf, welche 1) als erste Indices 1, 2, 3 be-
sitzen, 2) als zweite Indices irgend welche Ternen der Zahlen $(1\,2\cdots n)$,
und zwar ist jede solche Combination zu dreien sechs Mal vertreten,
entsprechend der Anzahl der Permutationen, welche drei Indices bilden
können. Fasst man wieder je sechs solche zusammengehörige Glieder
in der Entwicklung von $\Delta$ zusammen, so kommt, unter Berück-
sichtigung der Resultate in Nr. 50:
\[
\Delta = \sum_{\substack{i=1\\ k=1\\ r=1}}^{\substack{i=n\\ k=n\\ r=n}} \bigl(a_{1i}a_{2k}a_{3r} - a_{1i}a_{2r}a_{3k} - a_{1k}a_{2i}a_{3r} + a_{1k}a_{2r}a_{3i} % ?? uncertain: printed exponent-like mark over the k in this term (a_{3k}) not fully legible, read as plain a_{3k}
+ a_{1r}a_{2i}a_{3k} - a_{1r}a_{2k}a_{3i}\bigr)A_{\substack{1\ 2\ 3\\ i\ k\ r}}
\]
\[
= \sum \begin{vmatrix} a_{1i}\ a_{2i}\ a_{3i} \\ a_{1k}\ a_{2k}\ a_{3k} \\ a_{1r}\ a_{2r}\ a_{3r} \end{vmatrix} \cdot A_{\substack{1\ 2\ 3\\ i\ k\ r}},
\]
wie oben behauptet wurde.

55. \textit{Allgemeiner Fall.} So kann endlich jede Determinante $\Delta$ in
eine Summe von Producten entwickelt werden, deren Factoren Unter-
determinanten beliebigen Grades sind. Man nimmt eine Matrix von
%CONT>>


% ===== pdf idx 76, printed 62.0 =====
% !PRINTED 62
\weberpage{62}
%<<CONT
$p$ Zeilen — oder etwa auch von $p$ Colonnen —, bildet daraus die
$\binom{n}{p}$ Determinanten $p^{\text{ten}}$ Grades
\[
\Delta^{(p)} = \left(a_{i_1k_1}\, a_{i_2k_2}\, a_{i_3k_3}\cdots a_{i_pk_p}\right);
\]
zu jeder derselben bestimmt man die correspondirende Determinante
$(n-p)^{\text{ten}}$ Grades
\[
\Delta^{(n-p)} = A_{\substack{i_1\,i_2\cdots i_p\\ k_1\,k_2\cdots k_p}} = (-1)^{v}\left(a_{i_{p+1},\,k_{p+1}},\, a_{i_{p+2},\,k_{p+2}}\cdots a_{i_n,\,k_n}\right),
\]
wobei
\[
v = \sum_1^p i_\varrho + \sum_1^p k_\varrho \, .
\]
Die Summe aller Producte je zweier correspondirender Determinanten
liefert die ursprüngliche Determinante $\Delta$. Man hat also:
\[
\Delta = \sum \left(a_{i_1k_1}\, a_{i_2k_2}\cdots a_{i_pk_p}\right) \cdot A_{\substack{i_1\,i2\cdots i_p\\ k_1\,k_2\cdots k_p}} \cdot
\]

Diese Entwicklung von $\Delta$ nach Producten ihrer Minoren hat zuerst
Laplace gegeben. (Recherches sur le calcul intégral et sur le' système % [sic] stray mark before "système" reproduced as printed
du monde, 1772, Seite 304.)

56. \textit{Anwendungen des Laplace'schen Satzes bei Bildung von Relationen
zwischen Unterdeterminanten.} Wir hatten in Nr 32 Relationen auf-
gestellt, welche zwischen den Elementen einer Zeile oder Colonne und
Unterdeterminanten $1^{\text{ter}}$ Ordnung bestehen, oder, wie wir uns jetzt auch
ausdrücken können, zwischen Minoren $(n-1)^{\text{ter}}$ Ordnung und ihren
correspondirenden Determinanten. Die Methode, welche wir dort be-
nutzten, bestand darin, dass wir die Matrix $(A)$ in Nr. 28, oder eine
ihr analoge durch eine Zeile erweiterten, die bereits in der Matrix
auftrat, und die so erhaltene Determinante nach den Elementen dieser
Zeile entwickelten. Der Satz von Laplace setzt uns nun in den Stand,
auf dieselbe Weise Relationen zwischen Minoren beliebigen Grades
herzustellen. Wir nehmen eine beliebige Matrix $(B)$ von $p$ Zeilen und
$n$ Colonnen, erweitern sie zu einer Matrix von $n$ Zeilen und $n$ Colonnen,
indem wir $n-p$ Zeilen hinzufügen, unter denen wenigstens eine be-
reits in der Matrix $(B)$ selbst auftritt. Die durch diese neue Matrix
dargestellte Determinante verschwindet wegen der gleichen Zeilen.
Entwickelt man also dieselbe in eine Summe von Producten, deren
Factoren etwa die $\binom{n}{n-p}$ Determinanten der hinzugefügten Matrix
und die correspondirenden sind, so erhält man eine Relation zwischen
%CONT>>


% ===== pdf idx 77, printed 63.0 =====
% !PRINTED 63
\weberpage{63}
%<<CONT
Minoren $p^{\text{ter}}$ und $n-p^{\text{ter}}$ Ordnung irgend einer Determinante $\Delta$, von
deren Matrix die oben erwähnte Matrix $(B)$ ein Theil ist.

57. \textit{Beispiel 1.} Sei gegeben die fünfgliedrige Determinante:
\[
\Delta = \begin{vmatrix}
\alpha_1\ \beta_1\ \gamma_1\ \delta_1\ \varepsilon_1 \\
\alpha_2\ \beta_2\ \gamma_2\ \delta_2\ \varepsilon_2 \\
\alpha_3\ \beta_3\ \gamma_3\ \delta_3\ \varepsilon_3 \\
\alpha_4\ \beta_4\ \gamma_4\ \delta_4\ \varepsilon_4 \\
\alpha_5\ \beta_5\ \gamma_5\ \delta_5\ \varepsilon_5
\end{vmatrix} .
\]
Wir erhalten eine Relation zwischen Minoren zweiter und dritter Ord-
nung derselben, wenn wir etwa die Matrix der zweiten, dritten und
vierten Zeile durch die Matrix der ersten und zweiten Zeile zu einer
neuen Matrix
\[
\Delta' = \begin{vmatrix}
\alpha_1\ \beta_1\ \gamma_1\ \delta_1\ \varepsilon_1 \\
\alpha_2\ \beta_2\ \gamma_2\ \delta_2\ \varepsilon_2 \\
\alpha_2\ \beta_2\ \gamma_2\ \delta_2\ \varepsilon_2 \\
\alpha_3\ \beta_3\ \gamma_3\ \delta_3\ \varepsilon_3 \\
\alpha_4\ \beta_4\ \gamma_4\ \delta_4\ \varepsilon_4
\end{vmatrix}
\]
erweitern; die durch sie dargestellte Determinante $\Delta'$ ist Null wegen
der Gleichheit der zweiten und dritten Zeile. Man kann sie nun wieder
mit der Determinante
\[
D = \left(b_{11}\ b_{22}\ b_{33}\ b_{44}\ b_{55}\right)
\]
vergleichen, letztere nach den Determinanten der beiden ersten Zeilen,
also nach den Determinanten der Matrix
\[
(C) \qquad \begin{vmatrix} b_{11}\ b_{12}\ b_{13}\ b_{14}\ b_{15} \\ b_{21}\ b_{22}\ b_{23}\ b_{24}\ b_{25} \end{vmatrix}
\]
entwickeln, und in der Entwicklung hernach die Elemente $b_{ik}$ durch
die entsprechenden Elemente in $\Delta'$ ersetzen. Bezeichnen wir die zehn
Determinanten der Matrix $(C)$ mit
\[
(b_1b_2)_{12},\ (b_1b_2)_{13},\ (b_1b_2)_{14},\ (b_1b_2)_{15},\ (b_1b_2)_{23},\ (b_1b_2)_{24}\cdots (b_1b_2)_{45},
\]
dann ist die Entwicklung dargestellt durch
\[
D = \sum_{\substack{i=1\\ k=2}}^{\substack{k=5\\ i=4}} (b_1b_2)_{ik}\cdot B_{\substack{1\ 2\\ i\ k}};
\]
hier ist
\[
B_{\substack{1\ 2\\ i\ k}} = (-1)^{v_{ik}}\left(b_{3\varrho_1}\, b_{4\varrho_2}\, b_{5\varrho_3}\right),
\]
%CONT>>


% ===== pdf idx 78, printed 64.0 =====
% !PRINTED 64
\weberpage{64}
%<<CONT
wobei
\[
1)\quad v_{ik} = 1+2+i+k = 3+i+k,
\]
\[
2)\quad \varrho_1,\ \varrho_2,\ \varrho_3 \gtrless i \text{ oder } k,
\]
\[
3)\quad \varrho_1 < \varrho_2 < \varrho_3 \, .
\]
Wir erhalten also:
\[
\begin{aligned}
D = {}& (-1)^{3+3}(b_1b_2)_{12}(b_{33}b_{44}b_{55}) + (-1)^{3+4}(b_1b_2)_{13}(b_{32}b_{44}b_{55}) \\
& + (-1)^{3+5}(b_1b_2)_{14}(b_{32}b_{43}b_{55}) + (-1)^{3+6}(b_1b_2)_{15}(b_{32}b_{43}b_{54}) \\
& + (-1)^{3+5}(b_1b_2)_{23}(b_{31}b_{44}b_{55}) + (-1)^{3+6}(b_1b_2)_{24}(b_{31}b_{43}b_{55}) \\
& + (-1)^{3+7}(b_1b_2)_{25}(b_{31}b_{43}b_{54}) + (-1)^{3+7}(b_1b_2)_{34}(b_{31}b_{42}b_{55}) \\
& + (-1)^{3+8}(b_1b_2)_{35}(b_{31}b_{42}b_{54}) + (-1)^{3+9}(b_1b_2)_{45}(b_{31}b_{42}b_{53}) \, .
\end{aligned}
\]
Ersetzt man hier die Werthe von $b_{ik}$ durch die entsprechenden Elemente
% ?? uncertain: leading letter of this word is cropped at the page-image edge; read as "in" giving "in \Delta'"
in $\Delta'$, so kommt:
\[
\begin{aligned}
0 = {}& (\alpha_1\beta_2)(\gamma_2\delta_3\varepsilon_4) - (\alpha_1\gamma_2)(\beta_2\delta_3\varepsilon_4) + (\alpha_1\delta_2)(\beta_2\gamma_3\varepsilon_4) - (\alpha_1\varepsilon_2)(\beta_2\gamma_3\delta_4) \\
& + (\beta_1\gamma_2)(\alpha_2\delta_3\varepsilon_4) - (\beta_1\delta_2)(\alpha_2\gamma_3\varepsilon_4) + (\beta_1\varepsilon_2)(\alpha_2\gamma_3\delta_4) \\
& + (\gamma_1\delta_2)(\alpha_2\beta_3\varepsilon_4) - (\gamma_1\varepsilon_2)(\alpha_2\beta_3\delta_4) \\
& + (\delta_1\varepsilon_2)(\alpha_2\beta_3\gamma_4) \, .
\end{aligned}
\]
58. \textit{Beispiel 2.} Die viergliedrige Determinante
\[
D = \begin{vmatrix}
b_{11}\ b_{12}\ b_{13}\ b_{14} \\
b_{21}\ b_{22}\ b_{23}\ b_{24} \\
b_{31}\ b_{32}\ b_{33}\ b_{34} \\
b_{41}\ b_{42}\ b_{43}\ b_{44}
\end{vmatrix}
\]
giebt nach zweigliedrigen Unterdeterminanten entwickelt:
\[
\begin{aligned}
D = {}& \sum (b_1b_2)_{ik}B_{\substack{1\,2\\ i\,k}} \\
= {}& (-1)^{3+3}(b_1b_2)_{12}(b_3b_4)_{34} + (-1)^{3+4}(b_1b_2)_{13}(b_3b_4)_{24} \\
& + (-1)^{3+5}(b_1b_2)_{14}(b_3b_4)_{23} + (-1)^{3+5}(b_1b_2)_{23}(b_3b_4)_{14} \\
& + (-1)^{3+6}(b_1b_2)_{24}(b_3b_4)_{13} + (-1)^{3+7}(b_1b_2)_{34}(b_3b_4)_{12} \, .
\end{aligned}
\]
Setzen wir nun
\[
b_{1i} = b_{3i} = a_i, \quad b_{2i} = b_{4i} = b_i,
\]
dann erhält $D$ zwei Paar gleiche Zeilen und geht über in
\[
\Delta = \begin{vmatrix}
a_1\ a_2\ a_3\ a_4 \\
b_1\ b_2\ b_3\ b_4 \\
a_1\ a_2\ a_3\ a_4 \\
b_1\ b_2\ b_3\ b_4
\end{vmatrix}
= (ab)_{12}(ab)_{34} - (ab)_{13}(ab)_{24} + (ab)_{14}(ab)_{23}
+ (ab)_{23}(ab)_{14} - (ab)_{24}(ab)_{13} + (ab)_{34}(ab)_{12}
\]
\[
= 2\left\{(ab)_{12}(ab)_{34} - (ab)_{13}(ab)_{24} + (ab)_{14}(ab)_{23}\right\} = 0 \, .
\]


% ===== pdf idx 79, printed 65.0 =====
% !PRINTED 65
\weberpage{65}
Bezeichnen wir noch k\"urzer $(-1)^\nu (ab)_{ik}$ mit $p_{ik}$, so liefert die Matrix $\Delta$
die Relation:

\[
p_{12}p_{34} + p_{13}p_{42} + p_{14}p_{23} = 0.
\]

Sie entsteht aus ihrem ersten Gliede $p_{12}p_{34}$, indem man die 3 Indices 2 3 4 cyclisch
vertauscht.

59. \textit{Berechnung des Werthes einer Determinante mit Hilfe der Laplace'schen
Entwicklung.} \textit{Beispiel} 1. Wir berechnen zun\"achst unter Anwendung dieser
Entwicklung von $\Delta$ noch einmal den Werth der Determinante (6) in Nr. 15:

\[
\begin{aligned}
\Delta &= \begin{vmatrix} 3 & 1 & 1 & 2 \\ 1 & 5 & 0 & 3 \\ 2 & -2 & 1 & 6 \\ 0 & 4 & -5 & 3 \end{vmatrix}
= \begin{vmatrix} 3 & 1 \\ 1 & 5 \end{vmatrix} \cdot \begin{vmatrix} 1 & 6 \\ -5 & 3 \end{vmatrix}
- \begin{vmatrix} 3 & 1 \\ 1 & 0 \end{vmatrix} \cdot \begin{vmatrix} -2 & 6 \\ 4 & 3 \end{vmatrix}\\
&\quad + \begin{vmatrix} 3 & 2 \\ 1 & 3 \end{vmatrix} \cdot \begin{vmatrix} -2 & 1 \\ 4 & -5 \end{vmatrix}
+ \begin{vmatrix} 1 & 1 \\ 5 & 0 \end{vmatrix} \cdot \begin{vmatrix} 2 & 6 \\ 0 & 3 \end{vmatrix}
- \begin{vmatrix} 1 & 2 \\ 5 & 3 \end{vmatrix} \cdot \begin{vmatrix} 2 & 1 \\ 0 & -5 \end{vmatrix}
+ \begin{vmatrix} 1 & 2 \\ 0 & 3 \end{vmatrix} \cdot \begin{vmatrix} 2 & -2 \\ 0 & 4 \end{vmatrix}\\
&= (15-1)(3+30) - (0-1)(-6-24) + (9-2)(10-4)\\
&\quad + (0-5)(6-0) - (3-10)(-10-0) + (3-0)(8-0)\ .\\
&= 14\cdot 33 - 30 + 42 - 30 - 70 + 24 = 398\ .
\end{aligned}
\]

60. \textit{Beispiel} 2. Wesentliche Vortheile f\"ur die Berechnung des Werthes einer
Determinante gew\"ahrt die Methode von Laplace dann, wenn die Matrix der Determinante eine
gr\"ossere Zahl von Elementen enth\"alt, welche Null sind. Hierher geh\"ort in erster Linie
der Fall, wobei jene Elemente verschwinden, in welchen sich $p$ Verticalreihen mit $q$
Zeilen schneiden. Wir k\"onnen dann durch Vertauschung von Zeilen mit Zeilen und Colonnen mit
Colonnen die Elemente in der Matrix der gegebenen Determinante so ordnen, dass alle
verschwindenden Elemente selbst eine zusammenh\"angende Matrix von $p$ Colonnen und $q$
Zeilen bilden. Nehmen wir zun\"achst an $q = n - p$; die gegebene Determinante ist dann:

\[
\Delta = \begin{vmatrix}
a_{11} & \cdot\cdot\cdot & a_{1p} & a_{1p} & a_{1,p+1} & \cdot\cdot\cdot & a_{1n} \\
\cdot & \cdot & \cdot & \cdot & \cdot & \cdot & \cdot \\
a_{p1} & \cdot\cdot\cdot & a_{pp} & a_{p,p+1} & \cdot\cdot\cdot\cdot\cdot & a_{pn} \\
0 & \cdot\cdot\cdot\cdot & 0 & a_{p+1,p+1}, & \cdot\cdot\cdot & a_{p+1,n} \\
0 & \cdot\cdot\cdot\cdot & 0 & \cdot & \cdot\cdot\cdot & \cdot \\
0 & \cdot\cdot\cdot\cdot & 0 & a_{n,p+1} & \cdot\cdot\cdot & a_{nn}
\end{vmatrix}\ .
\]


% ===== pdf idx 80, printed 66.0 =====
% !PRINTED 66
\weberpage{66}
Entwickeln wir sie nach der Methode von Laplace, indem wir die Determinanten aus der Matrix

\[
(A)\quad
\begin{vmatrix}
a_{11} & a_{12} & \cdot\cdot\cdot & a_{1p} \\
a_{21} & \cdot\cdot\cdot\cdot\cdot & & a_{2p} \\
& \cdot & \cdot & \cdot \\
a_{p1} & \cdot & \cdot & a_{pp} \\
0 & 0 & \cdot\cdot & 0 \\
0 & \cdot & \cdot & 0 \\
& \cdot & \cdot & \cdot \\
0 & \cdot & \cdot & 0
\end{vmatrix}
\]

bilden, jede mit der correspondirenden Determinante multipliciren und die Summe \"uber alle
Producte nehmen, so erkennt man, dass diese Summe sich auf das einzige Product

\[
\begin{vmatrix}
a_{11} & \cdot\cdot\cdot & a_{1p} \\
a_{21} & \cdot\cdot\cdot & a_{2p} \\
\cdot & \cdot & \cdot \\
a_{p1} & \cdot\cdot\cdot & a_{pp}
\end{vmatrix}
\cdot
\begin{vmatrix}
a_{p+1,p+1} & \cdot\cdot\cdot & a_{p+1,n} \\
a_{p+2,p+1} & \cdot\cdot\cdot & a_{p+2,n} \\
\cdot & \cdot & \cdot \\
a_{n,p+1} & \cdot & \cdot & a_{nn}
\end{vmatrix}
= \Delta^{(p)} \cdot \Delta^{(n-p)}
\]

reducirt, weil ja alle Determinanten der Matrix $(A)$ --- die in diesem Producte auftretende
ausgenommen --- verschwinden.

61. \textit{Beispiel} 3. Nehmen wir zweitens an $q = n - p + 1$, dann enth\"alt die Matrix
$(A)$ nicht mehr wie vorhin $p$ Colonnen und $p$ Zeilen nicht verschwindender Elemente,
sondern sie enth\"alt nun eine Zeile mehr als sie Colonnen hat. Wir k\"onnen dann aus ihr

\[
\binom{p+1}{p} = p+1
\]

nicht verschwindender Determinanten bilden und die zugeh\"origen correspondirenden
ermitteln; die Summe der $(p+1)$ Producte liefert dann die Determinante $\Delta$.

Sei z. B. gegeben:

\[
\Delta = \begin{vmatrix}
a_1 & b_1 & c_1 & d_1 & e_1 \\
a_2 & b_2 & c_2 & d_2 & e_2 \\
a_3 & b_3 & c_3 & d_3 & e_3 \\
0 & 0 & c_4 & d_4 & e_4 \\
0 & 0 & c_5 & d_5 & e_5
\end{vmatrix},
\]
%CONT>>


% ===== pdf idx 81, printed 67.0 =====
% !PRINTED 67
\weberpage{67}
%<<CONT
dann liefert die Entwicklung nach den Determinanten der Matrix der beiden ersten Colonnen:

\[
\begin{aligned}
\Delta &= \begin{vmatrix} a_1 & b_1 \\ a_2 & b_2 \end{vmatrix} \cdot
\begin{vmatrix} c_3 & d_3 & e_3 \\ c_4 & d_4 & e_4 \\ c_5 & d_5 & e_5 \end{vmatrix}
+ (-1)^{\nu_1} \begin{vmatrix} a_1 & b_1 \\ a_3 & b_3 \end{vmatrix} \cdot
\begin{vmatrix} c_2 & d_2 & e_2 \\ c_4 & d_4 & e_4 \\ c_5 & d_5 & e_5 \end{vmatrix}\\
&\quad + (-1)^{\nu_2} \begin{vmatrix} a_2 & b_2 \\ a_3 & b_3 \end{vmatrix} \cdot
\begin{vmatrix} c_1 & d_1 & e_1 \\ c_4 & d_4 & e_4 \\ c_5 & d_5 & e_5 \end{vmatrix}\ ,
\end{aligned}
\]

wobei

\[
\nu_1 = 1 + 2 + 1 + 3 = 7
\]
\[
\nu_2 = 1 + 2 + 2 + 3 = 8,
\]

also:

\[
\Delta = (ab)_{12}(cde)_{345} - (ab)_{13}(cde)_{245} + (ab)_{23}(cde)_{145}\ .
\]

62. \textit{Entwicklung einer Determinante nach Producten von Elementen zweier sich
schneidender Reihen.} Manche Aufgaben lassen es w\"unschenswerth erscheinen, eine
Determinante nach den Elementen einer Zeile und Colonne gleichzeitig zu entwickeln. Wir
wollen dabei der Einfachheit halber annehmen, diese Elemente \underline{w\"urden} % [sic] underlined in print
die erste Zeile und Colonne bilden, die sich also in dem Elemente $a_{11}$ schneiden. Um zur
Entwicklungsformel zu gelangen, gehen wir wieder aus von der Gleichung in Nr. 28:

\weq{1}{\Delta = a_{11}A_{11} + a_{12}A_{12} + \cdots a_{1n}A_{1n}\ .}

Von den darin auftretenden Minoren $A_{1i}$ entwickeln wir den ersten $A_{11}$ nicht weiter,
die \"ubrigen aber nach den Unterdeterminanten der Elemente in der ersten \textit{Colonne}
(w\"ahrend wir sie in Nr. 53 (2) nach den Unterdeterminanten der Elemente in der ersten
\textit{Zeile} geordnet hatten).

Wir erhalten also:

\[
(2)\quad
\left\{
\begin{aligned}
A_{12} &= a_{21}A_{\substack{12\\21}} + a_{31}A_{\substack{13\\21}} + a_{41}A_{\substack{14\\21}} + \cdots a_{n1}A_{\substack{1n\\21}}\\
A_{13} &= a_{21}A_{\substack{12\\31}} + a_{31}A_{\substack{13\\31}} + a_{41}A_{\substack{14\\31}} + \cdots a_{n1}A_{\substack{1n\\31}}\\
&\cdots\cdots\cdots\cdots\cdots\cdots\cdots\cdots\cdots\\
A_{1k} &= a_{21}A_{\substack{12\\k1}} + a_{31}A_{\substack{13\\k1}} + \cdots + a_{k1}A_{\substack{1k\\k1}} + \cdots + a_{n1}A_{\substack{1n\\k1}}\\
&\cdots\cdots\cdots\cdots\cdots\cdots\cdots\cdots\cdots\\
A_{1n} &= a_{21}A_{\substack{12\\n1}} + a_{31}A_{\substack{13\\n1}} + a_{41}A_{\substack{14\\n1}} \cdots + a_{n1}A_{\substack{1n\\n1}}
\end{aligned}
\right.
\]


% ===== pdf idx 82, printed 68.0 =====
% !PRINTED 68
\weberpage{68}
Tragen wir diese Werthe in $\Delta$ ein, so kommt:

\[
(3)\quad
\begin{aligned}
\Delta &= a_{11}A_{11} + a_{21}a_{12}A_{\substack{12\\21}} + a_{31}a_{12}A_{\substack{13\\21}} + \cdots + a_{n1}a_{12}A_{\substack{1n\\21}}\\
&\quad + a_{21}a_{13}A_{\substack{12\\31}} + a_{31}a_{13}A_{\substack{13\\31}} + \cdots\\
&\quad \cdots\cdots\cdots\cdots\cdots\cdots\cdots\cdots\cdots\\
&\quad + a_{21}a_{1,k-1}A_{\substack{12\\k1}} + \cdots + a_{n1}a_{1n}A_{\substack{1n\\n1}}\ .
\end{aligned}
\]

Man erkennt, dass in dieser Summe durchgehends die Factoren der Minoren $A_{\substack{1i\\k1}}$
Producte von Elementen sind, welche der ersten Zeile und ersten Colonne angeh\"oren, und zwar
tritt jedes solche Product $a_{i1}a_{1k}$ nur einmal auf, indem jedes der $(n-1)$ Elemente der
ersten Zeile mit jedem der $(n-1)$ Elemente der ersten Colonne --- $a_{11}$ ist ja
ausgeschlossen --- multiplicirt wird. Wir k\"onnen also, die $(n-1)^2$ Glieder der obigen
Entwicklung zusammenfassend, die Gleichung (3) schreiben:

\weq{4}{\Delta = a_{11}A_{11} + \sum a_{i1}a_{1k}A_{\substack{1i\\k1}}\ ,}

oder, weil

\[
A_{\substack{1i\\k1}} = -A_{\substack{1i\\1k}}\quad \text{(vergleiche Nr. 49)},
\]

so wird endlich die gesuchte Entwicklung dargestellt durch

\weq{5}{\Delta = a_{11}A_{11} - \sum a_{i1}a_{1k}A_{\substack{1i\\1k}}\ ,}

\[
k\ \text{und}\ i > 1\ ,
\]

wobei erinnert sei, dass $A_{\substack{1k\\1i}}$ der Minor erster Ordnung des Elements
$a_{ik}$ in der Determinante $A_{11}$ ist. Wir erhalten sonach die Regel:

\glqq Um eine Determinante nach den Producten der Elemente zweier sich schneidender Reihen zu
entwickeln, bilde man die Determinante $A_{11}$, so wie ihre $(n-1)^2$ ersten Minoren
$A_{\substack{1k\\1i}}$.

Jeden derselben multiplicire man mit dem entsprechenden Producte $a_{i1}a_{1k}$, addire alle
Producte und subtrahire die Summe von $a_{11}A_{11}$\grqq{}.

63. \textit{Entwicklung einer Determinante nach den Producten von Elementen zweier Paare
sich schneidender Reihen.} Bisweilen tritt auch die Frage auf, wie l\"asst sich eine
Determinante in der Weise entwickeln, dass sie nach Producten geordnet ist, deren beide erste
Factoren zweigliedrige Determinanten sind, entnommen aus der Matrix der zwei ersten Zeilen
und der Matrix der zwei ersten Colonnen. Mit Hilfe der
%CONT>>


% ===== pdf idx 83, printed 69.0 =====
% !PRINTED 69
\weberpage{69}

Zerlegung einer Determinante nach der Methode von Laplace können wir auch diese Frage erledigen.

Wir können zunächst nach Laplace die Determinante entwickeln nach den zweigliedrigen Determinanten der Matrix
\[
\left|\begin{array}{ccccc}
a_{11} & a_{12} & a_{13} & \cdots & a_{1n}\\
a_{21} & a_{22} & a_{23} & \cdots & a_{2n}
\end{array}\right|.
\]
Sondern wir hier von vornherein in dieser Summe die Producte
\[
\left|\begin{array}{cc} a_{11} & a_{12}\\ a_{21} & a_{22}\end{array}\right|\cdot A_{\substack{12\\12}},\quad
\sum\left|\begin{array}{cc} a_{11} & a_{1i}\\ a_{21} & a_{2i}\end{array}\right|\cdot A_{\substack{12\\1i}},\quad
\sum\left|\begin{array}{cc} a_{12} & a_{1i}\\ a_{22} & a_{2i}\end{array}\right|\cdot A_{\substack{12\\2i}}
\]
ab, dann sind die übrigen nur Producte von Determinanten der Matrix
\[
(B)\qquad \left|\begin{array}{cccc} a_{13} & a_{14} & \cdots & a_{1n}\\ a_{23} & a_{24} & \cdots & a_{2n}\end{array}\right|
\]
in die entsprechenden correspondirenden Determinanten
\[
A_{\substack{12\\ik}} = (-1)^{\nu}\,(a_{3\varrho_1} a_{4\varrho_2} a_{5\varrho_3} \cdots a_{n\varrho_{n-2}}),
\]
wobei
\[
\begin{aligned}
&\text{1)}\quad \nu = 1+2+i+k,\\
&\text{2)}\quad k>i>2,\\
&\text{3)}\quad \varrho_1<\varrho_2<\varrho_3\cdots<\varrho_{n-2},
\end{aligned}
\]
und alle $\varrho_n$ von $k$ und $i$ verschieden sind.

Wir haben somit als erstes Resultat:
\weq{1}{
\begin{aligned}
\Delta = {}& \left|\begin{array}{cc} a_{11} & a_{12}\\ a_{21} & a_{22}\end{array}\right|\cdot A_{\substack{12\\12}}
+ \sum_{i=3}^{i=n}\left|\begin{array}{cc} a_{11} & a_{1i}\\ a_{21} & a_{2i}\end{array}\right|\cdot A_{\substack{12\\1i}}\\
&+ \sum_{i=3}^{i=n}\left|\begin{array}{cc} a_{12} & a_{1i}\\ a_{22} & a_{2i}\end{array}\right|\cdot A_{\substack{12\\2i}}
+ \sum_{\substack{i=3\\k=4}}^{\substack{k=n\\i=n-1}}\left|\begin{array}{cc} a_{1i} & a_{1k}\\ a_{2i} & a_{2k}\end{array}\right|\cdot A_{\substack{12\\ik}},\quad k>i>2.
\end{aligned}
}

Jeder der in der letzten Form dieser Entwicklung auftretenden Minoren zweiter Ordnung lässt sich nun aber wieder in eine Summe von Producten entwickeln, deren einer Factor eine Determinante der Matrix
\[
(C)\qquad \left|\begin{array}{cc}
a_{31} & a_{32}\\
a_{41} & a_{42}\\
a_{51} & a_{52}\\
\hdotsfor{2}\\
a_{n1} & a_{n2}
\end{array}\right|
\]
ist, während der andere von der entsprechenden correspondirenden
%CONT>>


% ===== pdf idx 84, printed 70.0 =====
% !PRINTED 70
\weberpage{70}
%<<CONT
Determinante, einem Minor vierter Ordnung von $\Delta$, gebildet wird. Diese Minoren vierter Ordnung sind also der Determinante entnommen
\[
A_{\substack{12\\12}} = \left|\begin{array}{cccc}
a_{33} & a_{34} & \cdots & a_{3n}\\
a_{43} & a_{44} & \cdots & a_{4n}\\
a_{53} & a_{54} & \cdots & a_{5n}\\
\hdotsfor{4}\\
a_{n3} & a_{n4} & \cdots & a_{nn}
\end{array}\right|;
\]
sie sind die sämmtlichen Unterdeterminanten zweiter Ordnung von $A_{\substack{12\\12}}$. Trägt man nun diese Entwicklungen der Minoren $A_{\substack{12\\ik}}$ in den vierten Term der Gleichung (1) ein, so erhalten wir:
\[
\begin{aligned}
\Delta = {}& \left|\begin{array}{cc} a_{11} & a_{12}\\ a_{21} & a_{22}\end{array}\right|\cdot A_{\substack{12\\12}}
+ \sum_{i=3}^{i=n}\left|\begin{array}{cc} a_{11} & a_{1i}\\ a_{21} & a_{2i}\end{array}\right|\cdot A_{\substack{12\\1i}}
+ \sum_{i=3}^{i=n}\left|\begin{array}{cc} a_{12} & a_{1i}\\ a_{22} & a_{2i}\end{array}\right|\cdot A_{\substack{12\\2i}}\\
&+ \sum \left|\begin{array}{cc} a_{1i} & a_{1k}\\ a_{2i} & a_{2k}\end{array}\right|\cdot \left|\begin{array}{cc} a_{\varrho1} & a_{\varrho2}\\ a_{\sigma1} & a_{\sigma2}\end{array}\right|\cdot A_{\substack{12\varrho\sigma\\ik12}},
\end{aligned}
\]
wo
\[
\begin{aligned}
&\sigma>\varrho,\\
&k>i,
\end{aligned}
\]
und $i\varrho$ einestheils, $k\sigma$ anderntheils alle Werthe von 3, resp. 4 bis $n-1$, resp. $n$ annimmt.

Der letzte Term enthält die gesuchte Entwicklung. Er entsteht, indem man einmal alle Determinanten der Matrix (B), und dann alle Determinanten der Matrix (C) bildet, jedesmal also
\[
\binom{n-2}{2} = \frac{(n-2)(n-3)}{2}
\]
an der Zahl. Jede Determinante des einen Systems multiplicire man mit jeder des andern Systems und füge zum Product noch als dritten Factor jene Unterdeterminante vierter Ordnung von $\Delta$, deren Matrix entsteht, wenn man im Minor $A_{\substack{12\\12}}$ die beiden Zeilen und Colonnen streicht, denen die beiden vorhergehenden Factoren entnommen sind. Die Summe aller Producte liefert das gewünschte Resultat.

64. \textit{Specielle F\"alle.} Man wird sich natürlich im Allgemeinen nicht der in 62 und 63 klargelegten Methoden bedienen, um eine Determinante zu berechnen, sondern die bereits früher gegebenen, bequemeren vorziehen. Am interessantesten sind diese Entwicklungen, sobald jene Elemente verschwinden, in welchen sich die betreffenden Zeilen und Colonnen, nach denen zu entwickeln ist, schneiden; also im
%CONT>>


% ===== pdf idx 85, printed 71.0 =====
% !PRINTED 71
\weberpage{71}
%<<CONT
ersten Falle, wenn $a_{11}=0$, im zweiten, wenn $a_{11}=a_{12}=a_{21}=a_{22}=0$. Diese Fälle treten in der That bei manchen Problemen der analytischen Geometrie auf. Ist $a_{11}=0$, so reducirt sich die Entwicklung in Nr. 62 (5) auf
\[
\Delta = -\sum a_{1i}\,a_{k1}\cdot A_{\substack{1k\\1i}},
\]
und im zweiten Falle, sind $a_{11}=a_{12}=a_{21}=a_{22}=0$, so wird
\[
\Delta = \sum \left|\begin{array}{cc} a_{1i} & a_{1k}\\ a_{2i} & a_{2k}\end{array}\right|\cdot \left|\begin{array}{cc} a_{\varrho1} & a_{\varrho2}\\ a_{\sigma1} & a_{\sigma2}\end{array}\right|\cdot A_{\substack{12\varrho\sigma\\ik12}},
\]
da ja die zweigliederigen Determinanten in den vorausgehenden Formen stets eine Colonne von verschwindenden Elementen besitzen.

65. \textit{Beispiel} 1. Die bereits in Nr. 52 als Beispiel verwendete Determinante
\[
\Delta = \left|\begin{array}{cccc}
a_{11} & a_{12} & a_{13} & x_1\\
a_{21} & a_{22} & a_{23} & x_2\\
a_{31} & a_{32} & a_{33} & x_3\\
x_1 & x_2 & x_3 & 0
\end{array}\right|
\]
wollen wir direct nach den Variabeln $x_1\,x_2\,x_3$, also nach den Elementen der letzten Zeile und Colonne, gleichzeitig entwickeln. Die Entwicklungsformel ist gegeben durch:
\[
\Delta = a_{44}\,A_{\substack{44\\44}} - \sum a_{i4}\,a_{4k}\,A_{\substack{4i\\4k}},
\]
oder weil hier $a_{44}=0$
\[
\Delta = -\sum a_{i4}\,a_{4k}\,A_{\substack{4i\\4k}},\quad i \text{ und } k<4.
\]
Wir erhalten sonach, wenn wir die Summe entsprechend $x_i^2$ und $x_i x_k$ in zwei Theile trennen,
\[
\Delta = -\sum_{i=3}^{i=1} a_{i4}\,a_{4i}\,A_{\substack{4i\\4i}} - \sum_{i\gtrless k}^{\substack{i=1\,2\,3\\k=1\,2\,3}} a_{i4}\,a_{4k}\,A_{\substack{4k\\4i}}.
\]
Hierin ist:
\[
A_{\substack{41\\41}} = (-1)^{\nu_1}(a_{22}a_{33}),\quad \nu_1=(41)+(41)+2(4+1)=12,
\]
\[
A_{\substack{42\\42}} = (-1)^{\nu_2}(a_{11}a_{33}),\quad \nu_2=(42)+(42)+2(4+2)=14,
\]
\[
A_{\substack{43\\43}} = (-1)^{\nu_3}(a_{11}a_{22}),\quad \nu_3=(43)+(43)+2(4+3)=16,
\]
\[
A_{\substack{41\\42}} = (-1)^{\nu_4}(a_{21}a_{33}-a_{31}a_{23}),\quad \nu_4=13,\qquad
A_{\substack{42\\43}} = (-1)^{\nu_7}(a_{11}a_{32}-a_{31}a_{12}),\quad \nu_7=15,
\]
\[
A_{\substack{41\\43}} = (1-)^{\nu_5}(a_{21}a_{32}-a_{31}a_{22}),\quad \nu_5=14,
\] % [sic] printed "(1-)" where "(-1)" is meant
\[
A_{\substack{43\\41}} = (-1)^{\nu_8}(a_{12}a_{23}-a_{13}a_{22}),\quad \nu_8=14,
\]
\[
A_{\substack{42\\41}} = (-1)^{\nu_6}(a_{12}a_{33}-a_{13}a_{32}),\quad \nu_6=13,\qquad
A_{\substack{43\\42}} = (-1)^{\nu_9}(a_{11}a_{23}-a_{13}a_{21}),\quad \nu_9=15,
\]
also
%CONT>>


% ===== pdf idx 86, printed 72.0 =====
% !PRINTED 72
\weberpage{72}
%<<CONT
\[
\begin{aligned}
\Delta ={}& -\{x_1^{2}(a_{22}a_{33}-a_{23}a_{32})+x_2^{2}(a_{11}a_{33}-a_{31}a_{13})+x_3^{2}(a_{11}a_{22}-a_{21}a_{12})\}\\
&+\{x_1x_2(a_{21}a_{33}-a_{31}a_{23}+a_{12}a_{33}-a_{32}a_{13})-x_1x_3(a_{21}a_{32}-a_{31}a_{22}\\
&\quad +a_{12}a_{23}-a_{22}a_{13})+x_2x_3(a_{11}a_{32}-a_{31}a_{12}+a_{11}a_{23}-a_{21}a_{13})\}.
\end{aligned}
\]
66. \textit{Beispiel} 2. Für die Determinante
\[
\Delta = \left|\begin{array}{ccccc}
a_1 & a_2 & a_3 & x_1 & y_1\\
b_1 & b_2 & b_3 & x_2 & y_2\\
c_1 & c_2 & c_3 & x_3 & y_3\\
x_1 & x_2 & x_3 & 0 & 0\\
y_1 & y_2 & y_3 & 0 & 0
\end{array}\right|
\]
liefert die Entwicklung
\[
\Delta = \sum \left|\begin{array}{cc} a_{4i} & a_{4k}\\ a_{5i} & a_{5k}\end{array}\right|\cdot \left|\begin{array}{cc} a_{\varrho4} & a_{\varrho5}\\ a_{\sigma4} & a_{\sigma5}\end{array}\right|\cdot A_{\substack{45\varrho\sigma\\ik45}},\quad \left\{\begin{array}{l}\varrho<\sigma<4\\ i<k<4\end{array}\right.
\]
folgenden Werth:
\[
\begin{aligned}
\Delta = {}& (-1)^{\nu_1}\left|\begin{array}{cc} x_1 & y_1\\ x_2 & y_2\end{array}\right|^2 c_3
+ (-1)^{\nu_2}\left|\begin{array}{cc} x_1 & y_1\\ x_3 & y_3\end{array}\right|^2 b_2
+ (-1)^{\nu_3}\left|\begin{array}{cc} x_2 & y_2\\ x_3 & y_3\end{array}\right|^2 a_1\\
&+ \left|\begin{array}{cc} x_1 & y_2\\ x_2 & y_2\end{array}\right|\left\{(-1)^{\nu_4}\left|\begin{array}{cc} x_1 & y_2\\ x_3 & y_3\end{array}\right| b_3 + (-1)^{\nu_5}\left|\begin{array}{cc} x_2 & y_2\\ x_3 & y_3\end{array}\right| a_3\right\}\\
&+ \left|\begin{array}{cc} x_1 & y_1\\ x_3 & y_3\end{array}\right|\left\{(-1)^{\nu_6}\left|\begin{array}{cc} x_1 & y_1\\ x_2 & y_2\end{array}\right| c_2 + (-1)^{\nu_7}\left|\begin{array}{cc} x_2 & y_2\\ x_3 & y_3\end{array}\right| a_2\right\}\\
&+ \left|\begin{array}{cc} x_2 & y_2\\ x_3 & y_3\end{array}\right|\left\{(-1)^{\nu_8}\left|\begin{array}{cc} x_1 & y_1\\ x_2 & y_2\end{array}\right| c_1 + (-1)^{\nu_9}\left|\begin{array}{cc} x_1 & y_1\\ x_3 & y_3\end{array}\right| b_1\right\};
\end{aligned}
\]
\[
\nu_1=(4512)+(1245)+24=4+0+24=28,\quad \nu_2=30,\quad \nu_3=32
\]
\[
\nu_4=(4513)+(1245)+25=4+0+25=29,\quad \nu_5=30
\]
\[
\nu_6=(4512)+(1345)+25=4+0+25=29,\quad \nu_7=31
\]
\[
\nu_8=(4512)+(2345)+26=4+0+26=30,\quad \nu_9=31,
\]
\[
\Delta = c_3(x_1y_2)^2 + b_2(x_1y_3)^2 + a_1(x_2y_3)^2 - (b_3+c_2)(x_1y_2)(x_1y_3)
+ (a_3+c_1)(x_1y_2)(x_2y_3) - (a_2+b_1)(x_2y_3)(x_1y_3).
\]

67. \textit{Begriff der R\"anderung einer Determinante.} Die in Beispiel 1 und 2 aufgestellten Determinanten bezeichnet man als \textit{geränderte} Determinanten. Die Ränderung besteht darin, dass man eine Matrix von $n^2$ Elementen durch eine Zeile und Colonne von $n$ Elementen erweitert und in den Schnittpunkt der beiden Reihen Null setzt. Dabei ist es nicht nothwendig, dass, wie in den beiden Beispielen, die zugefügte Zeile und Colonne dieselben Elemente enthalten. Die Ränderung kann wie in Beispiel (2) mehr als einmal vollzogen werden; man spricht dann von einem zwei- und dreifachen Saum der Determinante, deren
%CONT>>


% ===== pdf idx 87, printed 73.0 =====
% !PRINTED 73
\weberpage{73}
%<<CONT
sämmtliche Elemente in den Schnittpunkten der Saumreihen Null sind.
Die Determinante in Beispiel 1 ist also entstanden durch Ränderung
der Determinante
\[
D = \begin{vmatrix} a_{11} & a_{12} & a_{13} \\ a_{21} & a_{22} & a_{23} \\ a_{31} & a_{32} & a_{33} \end{vmatrix}
\]
mit $x_1\ x_2\ x_3$. Ebenso ist die Determinante in Beispiel 2 durch
Ränderung der nämlichen Determinante $D$ mit $x_1\ x_2\ x_3$, $y_1\ y_2\ y_3$
entstanden.

\begin{center}
\textbf{§ 6. Multiplication von Determinanten.}
\end{center}

68. \textit{Darstellung des Gedankenganges.} Bereits in Nr. 60 wurde ein
Fall besprochen, in welchem eine $n$-gliederige Determinante $\Delta$ in ein
einziges Product zweier Determinanten von $p^2$, bezw. $(n-p)^2$ Elementen
zerfiel. Die Bedingung war, dass sämmtliche Elemente, in denen sich $p$
Colonnen und $(n-p)$ Zeilen schneiden, verschwinden mussten. Man sieht
leicht, dass umgekehrt das Product zweier Determinanten vom $r^{\text{ten}}$
und $s^{\text{ten}}$ Grade immer durch eine einzige Determinante
$(r+s)^{\text{ten}}$ Grades sich darstellen lässt, indem man etwa in der
Matrix der ersten Determinante noch $s$ Zeilen mit verschwindenden
Elementen an die $r^{\text{te}}$ Zeile anfügt, dagegen die Matrix der
zweiten Determinante durch $r$ an die erste angesetzte Zeile mit
Elementen, welche beliebig gewählt werden können, erweitert, und die
beiden so erhaltenen Matrices von $r+s$ Zeilen und $r$ bezw. $s$ Colonnen
zu einer einzigen Matrix mit $r+s$ Colonnen und $r+s$ Zeilen verbindet.
Dabei kann $r\;\begin{smallmatrix}\geqq\\<\end{smallmatrix}\;s$ sein; % ?? uncertain: stacked relation symbol (greater-or-equal over less-than), reconstructed as a smallmatrix stack
ist speciell $r=1$, also $\Delta = a_{11}$, so haben wir den bereits in Nr. 30
Beispiel 5 erwähnten Fall vor uns: wir erweitern die Matrix der
Determinante mit den $s^2$ Elementen um eine Colonne von $s$ Elementen
Null und um eine Zeile von $s$ beliebigen Elementen. An die Stelle, wo
sich Colonne und Zeile schneiden, tritt das eine Element $a_{11}$ der
ersten Determinante $\Delta$. Ist $r=s=n$, dann sind die Erweiterungen
einer wie der andern Matrix selbst Determinantenmatrices, und zwar die
eine mit $n^2$ Elementen Null, die andere mit $n^2$ beliebigen Elementen.

Diese Art Determinanten zu multipliciren hat aber das eine Missliche an
sich, dass das Product eine Determinante höheren Grades ist, als die
beiden Determinanten der Factoren, und es entsteht die Frage, ob es
nicht gelingt, sie auf einen niedrigeren Grad zu reduciren. In der That
zeigt sich, dass es ganz allgemein möglich ist, die Determinante des
Productes bis auf den Grad jener Determinante zu redu-
%CONT>>


% ===== pdf idx 88, printed 74.0 =====
% !PRINTED 74
\weberpage{74}
%<<CONT
ciren, welche die höhere unter den beiden Factoren ist. Zum Beweise
setzen wir zunächst voraus $r=s=n$, d. h. dass beide Determinanten von
gleichem Grade sind, was immer erlaubt ist, da ja nach Beispiel 5 Nr. 30
die Matrix jeder Determinante zur Matrix einer Determinante beliebig
hohen Grades sich erweitern lässt, ohne dass der Werth der Determinante
sich ändert.

69. \textit{Aufstellung der Matrix der Productdeterminante.} Sei also
gegeben:
\[
\Delta = \begin{vmatrix}
a_{11} & a_{12} & a_{13} & \dots & a_{1n} \\
a_{21} & a_{22} & \dots & \dots & a_{2n} \\
a_{31} & \cdot & \cdot & \cdot & \cdot \\
\cdot & \cdot & \cdot & \cdot & \cdot \\
a_{n1} & a_{n2} & a_{n3} & \dots & a_{nn}
\end{vmatrix},
\quad
D = \begin{vmatrix}
b_{11} & b_{12} & b_{13} & \dots & b_{1n} \\
b_{21} & b_{22} & \dots & \dots & b_{2n} \\
b_{31} & \cdot & \cdot & \cdot & \cdot \\
\cdot & \cdot & \cdot & \cdot & \cdot \\
b_{n1} & b_{n2} & \dots & \dots & b_{nn}
\end{vmatrix}.
\]
Wir transponiren zunächst die Determinante $\Delta$ und erweitern die
Matrix von $\Delta$ um $n$ Horizontalreihen von $n^2$ Elementen Null,
die wir an die letzte Zeile ansetzen. Sodann wählen wir die
Determinante der $n^2$ willkürlichen Elemente, wodurch die Matrix von
$D$ erweitert werden soll, so dass die Diagonalglieder durchgehends
$-1$, alle übrigen null sein sollen. Die beiden so erhaltenen Matrices
von $2n$ Zeilen und je $n$ Colonnen vereinigen wir zu einer einzigen
\[
J = \begin{vmatrix}
a_{11} & a_{21} & a_{31} & a_{41} & \dots & a_{n1} & -1 & 0 & 0 & 0 & \dots & 0 \\
a_{12} & a_{22} & a_{32} & \dots & \dots & a_{n2} & 0 & -1 & 0 & \dots & \dots & 0 \\
a_{13} & a_{23} & \dots & \dots & \dots & \dots & 0 & 0 & -1 & \dots & \dots & 0 \\
\vdots & & & & & & & & \vdots & & & \\
a_{1n} & a_{2n} & \dots & \dots & \dots & a_{nn} & 0 & 0 & 0 & \dots & \dots & -1 \\
0 & 0 & 0 & 0 & \dots & 0 & b_{11} & b_{12} & b_{13} & \dots & \dots & b_{1n} \\
0 & 0 & 0 & \dots & \dots & 0 & b_{21} & b_{22} & b_{23} & \dots & \dots & b_{2n} \\
0 & 0 & \dots & \dots & \dots & 0 & b_{31} & b_{32} & \dots & \dots & \dots & b_{3n} \\
\vdots & & & & & & \vdots & & & & & \\
0 & 0 & 0 & \dots & \dots & 0 & b_{n1} & b_{n2} & \dots & \dots & \dots & b_{nn}
\end{vmatrix}.
\]
Die Determinante $J$ ist also nach dem Satze von Laplace gleich dem
Producte $\Delta \cdot D$.

70. \textit{Reduction der Productdeterminante $J$.} Um sie nun auf eine
$n$-gliederige zu reduciren, addiren wir:
\begin{enumerate}
\item[1)] die mit $b_{11}, b_{12}, b_{13} \dots b_{1n}$ bezw.
multiplicirte erste, zweite, dritte $\dots$ $n^{\text{te}}$ Zeile von
$J$ zur $(n+1)^{\text{ten}}$ Zeile,
\item[2)] die mit $b_{21}, b_{22}, b_{23} \dots b_{2n}$ bezw.
multiplicirte erste, zweite, dritte $\dots$ $n^{\text{te}}$ Zeile zur
$(n+2)^{\text{ten}}$ Zeile,
\end{enumerate}
%CONT>>


% ===== pdf idx 89, printed 75.0 =====
% !PRINTED 75
\weberpage{75}
%<<CONT
\begin{enumerate}
\item[3)] die mit $b_{31}, b_{32}, b_{33} \dots b_{3n}$ bezw.
multiplicirte erste, zweite, dritte $\dots$ $n^{\text{te}}$ Zeile zur
$(n+3)^{\text{ten}}$ Zeile u. s. w.,
\item[] endlich: die mit $b_{n,1}, b_{n,2} \dots b_{nn}$ bezw.
multiplicirte $n^{\text{te}}$ Zeile zur $2n^{\text{ten}}$ Zeile.
\end{enumerate}
Dadurch erhalten wir die Matrix: $J=$
\[
\begin{vmatrix}
a_{11}, & a_{21}, & \dots, & a_{n,1}, & -1 & 0 & \dots & 0 \\
a_{12}, & a_{22}, & \dots, & a_{n,2}, & 0 & -1 & \dots & 0 \\
a_{13}, & a_{23}, & \dots, & a_{n,3}, & 0 & 0 & \dots & 0 \\
\vdots & \vdots & & \vdots & & & & \\
a_{1n}, & a_{2n}, & \dots, & a_{n,n}, & 0 & 0 & \dots & -1 \\
a_{11}b_{11}+a_{12}b_{12}\cdots a_{1n}b_{1n}, & a_{21}b_{11}+a_{22}b_{12}\cdots a_{2n}b_{1n}, & \dots, & a_{n1}b_{11}+a_{n2}b_{12}\cdots a_{nn}b_{1n}, & 0 & 0 & \dots & 0 \\
a_{11}b_{21}+a_{12}b_{22}\cdots a_{1n}b_{2n}, & a_{21}b_{21}+a_{22}b_{22}\cdots a_{2n}b_{2n}, & \dots, & a_{n1}b_{21}+a_{n2}b_{22}\cdots a_{nn}b_{2n}, & 0 & 0 & \dots & 0 \\
a_{11}b_{31}+a_{12}b_{32}\cdots a_{1n}b_{3n}, & a_{21}b_{31}+a_{22}b_{32}\cdots a_{2n}b_{3n}, & \dots, & a_{n1}b_{31}+a_{n2}b_{32}\cdots a_{nn}b_{3n}, & 0 & 0 & \dots & 0 \\
a_{11}b_{41}+a_{12}b_{42}\cdots a_{1n}b_{4n}, & \cdot & \cdot & \cdot & & & & \cdot \\
\vdots & & & & & & & \\
a_{11}b_{n1}+a_{12}b_{n2}\cdots a_{1n}b_{nn}, & a_{21}b_{n1}+\cdots a_{2n}b_{nn}, & \dots, & a_{n1}b_{n1}+\cdots a_{nn}b_{nn}, & 0 & 0 & \dots & 0
\end{vmatrix}
\] % ?? uncertain: in the printed row "a_{n1}b_{31}+a_{n2}b_{32}\cdots a_{nn}b_{3n}," the final subscript is smudged in the scan and reconstructed by the row's evident pattern
Die erste Determinante reducirt sich, wie wir in Beispiel 4 Nr. 30
gezeigt haben, auf das Diagonalglied $(-1)^n$; der Exponent $\nu$ aber
hat den Werth
\[
\begin{aligned}
\nu &= (1+2+3\dots+n)+[(n+1)+(n+2)+\dots+2n] \\
&= n(n+1)+n^2 = 2n^2+n.
\end{aligned}
\]


% ===== pdf idx 90, printed 76.0 =====
% !PRINTED 76
\weberpage{76}
Es wird somit
\[
(-1)^{\nu}
\begin{vmatrix}
-1 & 0 & 0 & \dots & 0 \\
0 & -1 & 0 & \dots & 0 \\
0 & 0 & -1 & \dots & 0 \\
\vdots & & & & \vdots \\
0 & 0 & 0 & \dots & -1
\end{vmatrix}
= (-1)^{2n^2+n} \cdot (-1)^n = (-1)^{2(n^2+n)} = +1.
\]
Daher ist das Product der beiden Determinanten:
\[
J = \Delta \cdot D =
\begin{vmatrix}
a_{11}b_{11}+a_{12}b_{12}+\dots a_{1n}b_{1n}, & a_{21}b_{11}\dots+a_{2n}b_{1n}, & \dots, & a_{n1}b_{11}+\dots+a_{nn}b_{1n} \\
a_{11}b_{21}+a_{12}b_{22}+\dots a_{1n}b_{2n}, & a_{21}b_{21}\dots+a_{2n}b_{2n}, & \dots, & a_{n1}b_{21}+\dots a_{nn}b_{2n} \\
a_{11}b_{31}+a_{12}b_{32}+\dots a_{1n}b_{3n}, & \cdot & \cdot & \cdot \\
\vdots & & & \\
a_{11}b_{n1}+a_{12}b_{n2}+\dots a_{1n}b_{nn}, & a_{21}b_{n1}\dots+a_{2n}b_{nn}, & \dots, & a_{n1}b_{n1}+\dots a_{nn}b_{nn}
\end{vmatrix}
\]

71. \textit{Beispiel 1.} Es seien gegeben die beiden Determinanten
\[
\Delta = \begin{vmatrix} a_1 & b_1 & c_1 \\ a_2 & b_2 & c_2 \\ a_3 & b_3 & c_3 \end{vmatrix}, \quad
D = \begin{vmatrix} \alpha_1 & \beta_1 & \gamma_1 \\ \alpha_2 & \beta_2 & \gamma_2 \\ \alpha_3 & \beta_3 & \gamma_3 \end{vmatrix},
\]
alsdann ist das Product der beiden Determinanten
\[
J = \Delta \cdot D =
\begin{vmatrix}
a_1\alpha_1+b_1\beta_1+c_1\gamma_1, & a_2\alpha_1+b_2\beta_1+c_2\gamma_1, & a_3\alpha_1+b_3\beta_1+c_3\gamma_1 \\
a_1\alpha_2+b_1\beta_2+c_1\gamma_2, & a_2\alpha_2+b_2\beta_2+c_2\gamma_2, & a_3\alpha_2+b_3\beta_2+c_3\gamma_2 \\
a_1\alpha_3+b_1\beta_3+c_1\gamma_3, & a_2\alpha_3+b_2\beta_3+c_2\gamma_3, & a_3\alpha_3+b_3\beta_3+c_3\gamma_3
\end{vmatrix}.
\]

\textit{Beispiel 2.} Es seien gegeben die beiden Determinanten
\[
\Delta = \begin{vmatrix} x_1 & x_2 & x_3 & x_4 \\ y_1 & y_2 & y_3 & y_4 \\ z_1 & z_2 & z_3 & z_4 \\ u_1 & u_2 & u_3 & u_4 \end{vmatrix}
\quad \text{und} \quad
D = \begin{vmatrix} r_1 & r_2 \\ s_1 & s_2 \end{vmatrix}.
\]
Um das Product zu bilden, erweitern wir zuerst die Matrix von $D$, so
dass
\[
D = D' = \begin{vmatrix} r_1 & r_2 & 0 & 0 \\ s_1 & s_2 & 0 & 0 \\ 0 & 0 & 1 & 0 \\ 0 & 0 & 0 & 1 \end{vmatrix},
\]
sodann erhalten wir
\[
J = D \cdot \Delta =
\begin{vmatrix}
x_1 r_1 + x_2 r_2, & y_1 r_1 + y_2 r_2, & z_1 r_1 + z_2 r_2, & u_1 r_1 + u_2 r_2 \\
x_1 s_1 + x_2 s_2, & y_1 s_1 + y_2 s_2, & z_1 s_1 + z_2 s_2, & u_1 s_1 + u_2 s_2 \\
x_3, & y_3, & z_3, & u_3 \\
x_4, & y_4, & z_4, & u_4
\end{vmatrix}.
\]


% ===== pdf idx 91, printed 77.0 =====
% !PRINTED 77
\weberpage{77}

72. \textit{Zweite Art, die Productdeterminante zu bilden.} Wir haben in diesen beiden Beispielen nach dem Muster der reducirten Determinante $J$ die Determinante des Productes auf eine erste Art in der Weise hergestellt, dass wir die Elemente der ersten \textit{Zeile} von $\varDelta$ der Reihe nach multiplicirten mit den entsprechenden Elementen aller $n$ \textit{Zeilen} von $D$; das lieferte die $n$ Elemente der ersten \textit{Colonne} des Productes $J$; die $n$ Elemente der zweiten Colonne in $J$ entstanden durch Multiplication der zweiten Zeile in $\varDelta$, mit den entsprechenden aller $n$ Zeilen von $D$, u.\,s.\,f., endlich die Elemente der $k^{\mathrm{ten}}$ Colonne, durch Multiplication der Elemente in der $k^{\mathrm{ten}}$ Zeile von $\varDelta$, mit den entsprechenden aller $n$ Zeilen von $D$. Wir hätten aber auch die Determinante $J$ auf eine \textit{zweite} Art erhalten, wenn wir die Elemente der ersten Colonne von $\varDelta$ mit den entsprechenden aller Colonnen von $D$ multipliciren; dadurch entstehen die Elemente der ersten \textit{Zeile} von $J$; ihre zweite Zeile entsteht dann durch Multiplication der Elemente der zweiten Colonne von $\varDelta$ mit den entsprechenden aller Colonnen von $D$ u.\,s.\,w.

Die Berechtigung dazu erhalten wir unmittelbar durch das Gesetz, dass die Transposition die beiden Determinanten $D$ und $\varDelta$ nicht ändert; durch Transposition ordnen sich aber die Elemente einer bestimmten Colonne in eine Zeile, so dass die zweite Methode nur die auf die transponirten Determinanten angewendete erste ist.

73. \textit{Anmerkung 1. Symmetrische Determinanten.} Bilden wir ähnlich das Quadrat einer Determinante vierten Grades, etwa indem wir Colonnen mit Colonnen multipliciren, so erhalten wir:

\[
\varDelta^2 =
\begin{vmatrix}
a_{11} & a_{12} & a_{13} & a_{14}\\
a_{21} & a_{22} & a_{23} & a_{24}\\
a_{31} & a_{32} & a_{33} & a_{34}\\
a_{41} & a_{42} & a_{43} & a_{44}
\end{vmatrix}^{2}
=
\begin{vmatrix}
\sum a_{i1}^{2}, & \sum a_{i1}a_{i2}, & \sum a_{i1}a_{i3}, & \sum a_{i1}a_{i4}\\
\sum a_{i2}a_{i1}, & \sum a_{i2}^{2}, & \sum a_{i2}a_{i3}, & \sum a_{i2}a_{i4}\\
\sum a_{i3}a_{i1}, & \sum a_{i3}a_{i2}, & \sum a_{i3}^{2}, & \sum a_{i3}a_{i4}\\
\sum a_{i4}a_{i1}, & \sum a_{i4}a_{i2}, & \sum a_{i4}a_{i3}, & \sum a_{i4}^{2}
\end{vmatrix}.
\]

Die so erhaltene Matrix hat die Eigenschaft, dass je zwei Elemente, welche zur Hauptdiagonale symmetrisch liegen, einander gleich sind.


% ===== pdf idx 92, printed 78.0 =====
% !PRINTED 78
\weberpage{78}

Jede Determinante, in deren Matrix also

\[
b_{ik} = b_{ki},
\]

heisst symmetrische Determinante, und das Beispiel lehrt: \glqq{}\textit{Das Quadrat einer Determinante ist eine symmetrische Determinante.}\grqq{} Ist $b_{ik} = -b_{ki}$, also $b_{ii}=0$, so nennt man die Determinante schiefsymmetrisch oder hemisymmetrisch. Es existiren über diese und ähnliche specielle Fälle von Determinanten zahlreiche Sätze, die sich hauptsächlich auf die Auswerthung derselben beziehen; vgl. \gesperrt{Salmon}: Vorlesungen über die Algebra der linearen Transformationen, fünfte Vorlesung; \gesperrt{Mansion}, Elemente der Determinantentheorie, Seite 10, 23, 26, 31, 34, 38; \gesperrt{Günther}, Determinanten Capitel III. Hier sei nur einer Eigenschaft symmetrischer Determinanten Erwähnung gethan, die in den erwähnten Werken nicht besprochen ist.

74. \textit{Anmerkung 2. Lineare Relationen zwischen den Minoren von symmetrischen Determinanten.} Wir haben in Nr. 56 gesehen, wie es immer gelingt, beliebige Relationen zwischen Minoren einer Determinante herzustellen, und dies auch an drei Beispielen dargethan. So ist die in Beispiel 2 erhaltene Relation:

\[
(ab)_{12}(ab)_{14} - (ab)_{13}(ab)_{24} + (ab)_{14}(ab)_{23} = 0
\]

eine Beziehung der vier Minoren zweiter Ordnung, gebildet aus den zwei Zeilen einer Determinante:

\[
\varDelta =
\begin{vmatrix}
a_1 & a_2 & a_3 & a_4\\
b_1 & b_2 & b_3 & b_4\\
u & v & w & m\\
x & y & z & n
\end{vmatrix}.
\]

Während aber die Relationen, zu welchen die Minoren beliebiger Determinanten Veranlassung geben, immer von höherem als vom ersten Grade sind in den Unterdeterminanten, giebt es, wie sich leicht finden lässt (vgl. \gesperrt{Kronecker}, Monatsber. der Berl. Akad. 1882), bei \textit{symmetrischen} Determinanten auch \textit{lineare} Relationen zwischen Minoren. Wir wollen dies nur an einem einfachen Falle zeigen. Sei die symmetrische Determinante:

\[
\varDelta =
\begin{vmatrix}
x & a & b & c\\
a & y & o & l\\
b & o & z & m\\
c & l & m & u
\end{vmatrix};
\]
%CONT>>


% ===== pdf idx 93, printed 79.0 =====
% !PRINTED 79
\weberpage{79}

%<<CONT
greifen wir irgend einen Minor zweiter Ordnung aus derselben heraus, der kein Diagonalglied enthält, etwa

\[
\begin{vmatrix} b & c\\ o & l\end{vmatrix},
\]

so existiren in doppelter Weise noch zwei andere Minoren, deren jeder ein und nur ein Elementenpaar desselben enthält, nämlich

\[
\begin{vmatrix} a & b\\ l & m\end{vmatrix} \text{ und } \begin{vmatrix} o & a\\ m & c\end{vmatrix}.
\]

Die Summe aller drei Determinanten giebt eine der erwähnten Relationen, nämlich:

\[
\begin{vmatrix} b & c\\ o & l\end{vmatrix} + \begin{vmatrix} a & b\\ l & m\end{vmatrix} + \begin{vmatrix} o & a\\ m & c\end{vmatrix} = 0.
\]

75. \textit{Multiplication von Matrices, welche keine Determinanten darstellen; Gedankengang.} Wenn man in der Weise, wie es der Productsatz (Nr. 71 erste Art) für zwei Matrices von je $n^2$ Elementen vorschreibt, aus den beiden Matrices

\[
\begin{vmatrix}
a_{11} & a_{12} & a_{13} & \ldots & a_{1n}\\
a_{21} & a_{22} & . & . & a_{2n}\\
a_{31} & . & . & . & a_{3n}\\
. & . & . & . & .\\
a_{\nu 1} & a_{\nu 2} & . & . & a_{\nu n}
\end{vmatrix}
\text{ und }
\begin{vmatrix}
b_{11} & b_{12} & b_{13} & \ldots & b_{1n}\\
b_{21} & b_{22} & . & . & b_{2n}\\
b_{31} & . & . & . & .\\
. & . & . & . & .\\
b_{\nu 1} & . & . & . & b_{\nu n}
\end{vmatrix},\ \nu \gtrless n,
\]

eine neue bildet, so unterscheiden wir die beiden Fälle, dass die dadurch repräsentirte Determinante von $\nu^2$ Elementen entweder identisch Null oder von Null verschieden ist. Wir werden zunächst den Fall besprechen, in welchem dieselbe \textit{nicht} verschwindet. Zu dem Zwecke bilden wir zuerst
%CONT>>


% ===== pdf idx 94, printed 80.0 =====
% !PRINTED 80
\weberpage{80}

%<<CONT
nach der in Nr. 68 angegebenen Weise aus beiden Matrices, indem wir die erste vorher transponiren, eine Determinante $R$ von $(n+\nu)^2$ Elementen und reduciren sie nach der Methode in Nr. 70. In zweiter Linie entwickeln wir die noch nicht reducirte Determinante $R$ nach der Methode von \gesperrt{Laplace}, wodurch wir ihren Werth ausgedrückt erhalten durch eine Summe von Producten je zweier entsprechender Determinanten der beiden gegebenen Matrices. Da im Allgemeinen keine dieser Determinante $\nu^{\mathrm{ten}}$ Grades verschwindet, so ist der Werth von $R$ in der That von Null verschieden. Hierbei ist jedoch immer vorausgesetzt, dass $\nu < n$, und dass die Elemente der reducirten Productdeterminante $R$, wie wir sehen werden, gebildet sind durch Multiplication von \textit{Zeilen mit Zeilen} der beiden Matrices.

76. \textit{Aufstellung der Productdeterminante $R$.} Wir setzen also voraus $\nu < n$; erweitern sodann die erste Matrix, die wir vorher in die Matrix transponiren:

\[
\begin{vmatrix}
a_{11} & a_{21} & \ldots & a_{\nu 1}\\
a_{12} & a_{22} & . & .\\
a_{13} & . & . & .\\
a_{14} & . & . & .\\
. & . & . & .\\
. & . & . & .\\
a_{1n} & a_{2n} & \ldots & a_{\nu n}
\end{vmatrix},
\]

durch $\nu$ Zeilen mit je $\nu$ verschwindenden Elementen. Die zweite Matrix der Elemente $b_{ik}$ erweitern wir durch $n$ Zeilen von je $n$ willkürlichen Elementen, die wir aber wieder so wählen, dass die Diagonalelemente den Werth $-1$ haben, alle übrigen null sind. Aus den beiden so erhaltenen Matrices bilden wir nun die Determinante

\[
R =
\begin{vmatrix}
a_{11} & a_{21} & a_{31} & \ldots & a_{\nu 1} & -1 & 0 & 0 & 0 & \ldots & 0\\
a_{12} & a_{22} & \ldots & \ldots & a_{\nu 2} & 0 & -1 & 0 & \ldots & \ldots & 0\\
a_{13} & \ldots & \ldots & \ldots & a_{\nu 3} & 0 & 0 & -1 & \ldots & \ldots & 0\\
. & . & . & . & . & . & . & . & . & . & .\\
a_{1n} & a_{2n} & a_{3n} & \ldots & a_{\nu n} & 0 & 0 & 0 & \ldots & \ldots & -1\\
0 & 0 & 0 & \ldots & 0 & b_{11} & b_{12} & b_{13} & \ldots & \ldots & b_{1n}\\
0 & 0 & \ldots & \ldots & 0 & b_{21} & b_{22} & b_{23} & \ldots & \ldots & b_{2n}\\
0 & \ldots & \ldots & \ldots & 0 & b_{31} & b_{32} & \ldots & \ldots & \ldots & b_{3n}\\
. & . & . & . & . & . & . & . & . & . & .\\
. & . & . & . & . & . & . & . & . & . & .\\
0 & 0 & 0 & \ldots & 0 & b_{\nu 1} & b_{\nu 2} & \ldots & \ldots & \ldots & b_{\nu n}
\end{vmatrix}.
\]

77. \textit{Reduction der Determinante $R$.} Wir können nun genau wie früher die Matrix der Determinante $J$ so jetzt die der Determinante $R$ auf weniger Zeilen und Colonnen reduciren. Man addirt zu dem Zwecke wieder die mit $b_{11}, b_{12} \ldots b_{1n}$ bezw. multiplicirt gedachte erste, zweite $\ldots n^{\mathrm{te}}$ Zeile insgesammt zur $n+1^{\mathrm{ten}}$; ebenso dieselben $n$ ersten Zeilen der Reihe nach mit $b_{21}, b_{22} \ldots b_{2n}$ multiplicirt gedacht, zur $n+2^{\mathrm{ten}}$ und so weiter, endlich die mit $b_{\nu 1}, b_{\nu 2} \ldots b_{\nu n}$ multiplicirte erste bis $n^{\mathrm{te}}$ Zeile zur $n+\nu^{\mathrm{ten}}$ Zeile.

Das Resultat dieser Operationen ist dann die Determinante:
%CONT>>


% ===== pdf idx 95, printed 81.0 =====
% !PRINTED 81
\weberpage{81}
% [column spec widened array 5->6; array 4->5] a row carried more cells than the spec declared,
% which is a fatal alignment error. The spec was widened rather than a cell
% deleted, so no transcribed content is lost; whether Weber's grid really has
% that many columns is a question for a pass with the scan open.

{\scriptsize\setlength{\arraycolsep}{3pt}
\[
R=\left|\begin{array}{ccccccccc}
a_{11} & , & a_{21} & \cdot\cdot, & a_{\nu 1} & , & -1 & 0 & 0\cdot\cdot\ 0\\
a_{12} & , & a_{22} & \cdot\cdot, & a_{\nu 2} & , & 0 & -1 & 0\cdot\cdot\ 0\\
a_{13} & , & a_{23} & \cdot\cdot, & a_{\nu 3} & , & 0 & 0 & -1\cdot\cdot\\
\vdots & & \vdots & & \vdots & & \cdot & \cdot & \cdot\ \cdot\ \cdot\ \cdot\ \cdot\ \cdot\\
a_{1,n} & , & a_{2,n} & \cdot\cdot, & a_{\nu n} & , & 0 & 0 & 0\cdot\cdot(-1)\\
a_{11}b_{11}+a_{12}b_{12}+\cdot\cdot a_{1n}b_{1n} & , & a_{21}b_{11}+a_{22}b_{12}\cdot\cdot+a_{2n}b_{1n} & \cdot\cdot, & a_{\nu 1}b_{11}+\cdot\cdot a_{\nu n}b_{1n} & , & 0 & 0 & 0\cdot\cdot\ 0\\
a_{11}b_{21}+a_{12}b_{22}+\cdot\cdot a_{1n}b_{2n} & , & a_{21}b_{21}+a_{22}b_{22}\cdot\cdot\ a_{2n}b_{2n} & \cdot\cdot, & a_{\nu 1}b_{21}+\cdot\cdot a_{\nu n}b_{2n} & , & 0 & 0 & 0\cdot\cdot\ 0\\
a_{11}b_{31}+a_{12}b_{32}\ \cdot\cdot a_{1n}b_{3n} & , & a_{21}b_{31}+\cdot\cdot\cdot\cdot+a_{2n}b_{3n} & \cdot\cdot, & a_{\nu 1}b_{31}+\cdot\cdot a_{\nu n}b_{3n} & , & 0 & 0 & \cdot\ \cdot\ \cdot\ \cdot\\
a_{11}b_{41}\ \cdot\cdot\cdot\cdot+a_{1n}b_{4n}, & \cdot & \cdot & \cdot\ \cdot\ \cdot\ \cdot\ \cdot\ \cdot & \cdot & \cdot & \cdot & 0 & \cdot\ \cdot\ \cdot\ 0\\
\cdot & \cdot & \cdot & \cdot\ \cdot\ \cdot\ \cdot\ \cdot & \cdot & \cdot & \cdot & \cdot & \cdot\\
a_{11}b_{\nu 1}+a_{12}b_{\nu 2}\cdot\cdot\ a_{1n}b_{\nu n} & , & a_{21}b_{11}+\cdot\cdot\cdot\cdot\ a_{2n}b_{\nu n} & \cdot\cdot, & a_{\nu 1}b_{11}+\cdot\cdot a_{\nu n}b_{\nu n} & , & 0 & 0 & 0\cdot\cdot\ 0
\end{array}\right|.
\]
}

Entwickelt man diese Determinante nach den letzten $n$ Colonnen, so
zerf\"allt sie nach dem Satze von Laplace wieder in das Product zweier
Determinanten, und zwar erh\"alt man:
{\scriptsize\setlength{\arraycolsep}{3pt}
\[
R=(-1)^{\bar\nu}\cdot
\left|\begin{array}{cccccc}
-1 & 0 & 0\ldots & & 0\\
& 0 & -1 & 0\ldots & 0\\
& 0 & 0 & -1\ldots & 0\\
& \cdot & \cdot & \cdot & \cdot & \cdot\\
& & & & & 0\\
& 0 & 0 & 0\ldots & -1
\end{array}\right|
\cdot
\left|\begin{array}{ccccc}
a_{11}b_{11}\cdot\cdot+a_{1n}b_{1n} & , & a_{21}b_{11}+\cdot\cdot+a_{2n}b_{1n} & \cdot\cdot, & a_{\nu 1}b_{11}+\cdot\cdot a_{\nu n}b_{1n}\\
a_{11}b_{21}\cdot\cdot+a_{1n}b_{2n} & , & a_{21}b_{21}+\cdot\cdot+a_{2n}b_{2n} & \cdot\cdot, & a_{\nu 1}b_{21}+\cdot\cdot a_{\nu n}b_{2n}\\
a_{11}b_{31}\ \cdot\cdot a_{1n}b_{3n} & , & a_{21}b_{31}+\cdot\cdot\cdot\cdot+a_{2n}b_{3n} & \cdot\cdot, & a_{\nu 1}b_{31}+\cdot\cdot a_{\nu n}b_{3n}\\
a_{11}b_{41}\ \cdot\cdot\cdot\cdot+a_{1n}b_{1n}, & & \cdot\ \cdot\ \cdot\ \cdot\ \cdot\ \cdot\ \cdot\ \cdot\ \cdot\ \cdot\ \cdot\ \cdot & &\\
& \cdot & \cdot & \cdot & \cdot\\
a_{11}b_{\nu 1}\cdot\cdot\ a_{1n}b_{\nu n} & , & a_{21}b_{\nu 1}\ \cdot\cdot\cdot\cdot & \cdot\cdot, & a_{\nu 1}b_{\nu 1}\ \cdot\cdot\cdot a_{\nu n}b_{\nu n}
\end{array}\right|.
\]
}
Dabei ist der Werth der Determinante mit den Diagonalgliedern $(-1)$
gegeben durch $(-1)^n$; und anderntheils ist der Exponent $\bar\nu$ gleich der
Summe aller ersten plus der Summe aller zweiten Indices des Diagonalgliedes
dieser Determinante, also (vgl. Nr. 55)
\[
\bar\nu=[1+2+3\cdots+n]+[(\nu+1)(\nu+2)+\cdots(\nu+n)]
\]
\[
=\frac{n(n+1)}{2}+n\cdot\nu+\frac{n(n+1)}{2}=n(n+1)+n\nu.
\]
Nun ist aber $n(n+1)$ jedenfalls eine gerade Zahl; daher k\"onnen wir setzen:
\[
(-1)^{\bar\nu}=(-1)^{n(n+1)+n\nu}=(-1)^{n\cdot\nu}.
\]


% ===== pdf idx 96, printed 82.0 =====
% !PRINTED 82
\weberpage{82}
% [column spec widened array 4->7; array 4->12; array 4->6] a row carried more cells than the spec declared,
% which is a fatal alignment error. The spec was widened rather than a cell
% deleted, so no transcribed content is lost; whether Weber's grid really has
% that many columns is a question for a pass with the scan open.

Es wird somit
\[
(-1)^{\bar\nu}\cdot
\left|\begin{array}{ccccccc}
(-1) & 0 & 0\ \ldots & 0\\
& 0 & (-1) & 0\ \ldots & 0\\
& 0 & 0 & (-1)\ldots & 0\\
& \cdot & \cdot & \cdot & \cdot & \cdot & \cdot\\
& 0 & 0 & 0\ldots & (-1)
\end{array}\right|
=(-1)^{n\cdot\nu}\cdot(-1)^n=(-1)^{n(\nu+1)}.
\]
Wir erhalten also: $R=$
{\scriptsize\setlength{\arraycolsep}{2pt}
\[
(-1)^{n(\nu+1)}
\left|\begin{array}{cccccccccccc}
a_{11}b_{11}+a_{12}b_{12}\cdot\cdot+a_{1n}b_{1n} & , & a_{21}b_{11}+\cdot\cdot+a_{2n}b_{1n} & \cdot\cdot, & a_{\nu 1}b_{11}+\cdot\cdot a_{\nu n}b_{1n}\\
a_{11}b_{21}+a_{12}b_{22}\cdot\cdot+a_{1n}b_{2n} & , & a_{21}b_{21}+\cdot\cdot+a_{2n}b_{2n} & \cdot\cdot, & a_{\nu 1}b_{21}+\cdot\cdot a_{\nu n}b_{2n}\\
a_{11}b_{31}+a_{12}b_{32}\cdot\cdot\ a_{1n}b_{3n} & , & \cdot & \cdot & \cdot & \cdot & \cdot & \cdot & \cdot & \cdot\\
a_{11}b_{41}+a_{12}b_{42}\cdot\cdot+a_{1n}b_{4n}, & & \cdot & \cdot & \cdot & \cdot & \cdot & \cdot & \cdot & \cdot & \cdot & \cdot\\
& \cdot & \cdot & \cdot & \cdot & \cdot\\
a_{11}b_{\nu 1}+a_{12}b_{\nu 2}\cdot\cdot+a_{1n}b_{\nu n} & , & a_{21}b_{\nu 1}+\cdot\cdot & a_{2n}b_{\nu n} & \cdot\cdot, & a_{\nu 1}b_{\nu 1}+\cdot\cdot a_{\nu n}b_{\nu n}
\end{array}\right|
\]
}

78. \textit{Entwicklung der Determinante $R$ nach Laplace.} Wir k\"onnen
nun aber auch die Determinante $R$, ohne sie vorher auf niedrigeren Grad zu
reduciren, unter Anwendung der Methode von Laplace direct auf eine zweite
Art darstellen. Entwickeln wir sie n\"amlich nach den Determinanten der $\nu$
ersten Colonnen, so erhalten wir eine Summe von Producten je zweier
Determinanten, deren eine aus der Matrix der $\nu$ ersten Colonnen, deren
andere aus der Matrix der $n$ letzten Colonnen entnommen ist. Die Zahl
dieser Producte ist gleich der Anzahl nicht verschwindender Determinanten,
welche die Matrix der $\nu$ ersten Colonnen liefert, also gleich
$\binom{n}{\nu}$. Wir erhalten, wenn wir die Determinanten jeder der
beiden Matrices mit $\varDelta_\varrho^{(\nu)}$ bezw. $D_\varrho^{(n)}$ bezeichnen,
\[
R=\sum_{\varrho=1}^{\varrho=\binom{n}{\nu}}(-1)^\varrho\cdot\varDelta_\varrho^{(\nu)}\cdot D_\varrho^{(n)}, \tag{A}
\]
wobei der Factor $(-1)^\varrho$ der jeweiligen correspondirenden Determinante
$D_\varrho^{(n)}$ entspricht. Enth\"alt $\varDelta_\varrho^{(\nu)}$ die Zeilen
\[
i_1,\ i_2,\ i_3\cdots i_\nu,
\]
ist also die Matrix dieser aus den $\nu$ ersten Colonnen entnommenen
Determinante gegeben durch:
\[
\varDelta_\varrho^{(\nu)}=
\left|\begin{array}{cccccc}
a_{1i_1} & a_{2i_1} & a_{3i_1} & \cdots & a_{\nu i_1}\\
a_{1i_2} & a_{2i_2} & a_{3i_2} & \cdots & a_{\nu i_2}\\
a_{1i_2} & a_{2i_3} & a_{3i_3} & \cdots & a_{\nu i_3}\\
\cdot & \cdot & \cdot & \cdot & \cdot & \cdot\\
a_{1i_\nu} & a_{2i_\nu} & a_{3i_\nu} & \cdots & a_{\nu i_\nu}
\end{array}\right| \tag{B}
\]
% [sic] dritte Zeile druckt $a_{1i_2}$ (statt $a_{1i_3}$) wie im Original


% ===== pdf idx 97, printed 83.0 =====
% !PRINTED 83
\weberpage{83}
% [column spec widened array 4->6] a row carried more cells than the spec declared,
% which is a fatal alignment error. The spec was widened rather than a cell
% deleted, so no transcribed content is lost; whether Weber's grid really has
% that many columns is a question for a pass with the scan open.

so ist in dem Factor $(-1)^\varrho$ der Exponent
\[
\varrho=(1+2\cdots\nu)+(i_1+i_2\cdots+i_\nu)
\]
\[
=\frac{\nu(\nu+1)}{2}+\sum_{\varrho=1}^{\varrho=\nu}i_\varrho \tag{C}
\]
(vergleiche Nr. 55).

79. \textit{Berechnung von $D_\varrho^{(n)}$.} Die Determinanten $D_\varrho^{(n)}$
haben nun aber durchgehends die Eigenschaft, dass ihre $(n-\nu)$ ersten
Zeilen:
\[
i_{\nu+1},\ i_{\nu+2}\cdots i_n
\]
nur verschwindende Elemente besitzen, abgesehen von dem in jeder Zeile
einmal auftretenden Elemente $(-1)$, das in der Zeile $i_\varrho$ von
einer Colonne ausgeschnitten wird, die in der urspr\"unglichen Determinante
den Index $\nu+i_\varrho$ hat. Denkt man sich also in einer solchen
Determinante $D_\varrho^{(n)}$ die Colonnen in der Weise vertauscht, dass
diejenigen, in welchen sich das Element $(-1)$ befindet, die $(n-\nu)$
ersten bilden, wobei indes die Indices der Colonnen, so weit es geht, gut
geordnet bleiben m\"ogen, dann ist die Matrix einer solchen Determinante:
\[
D_\varrho^{(n)}=
\left|\begin{array}{ccccccccc}
-1 & 0 & 0 & \ldots & 0 & 0 & 0 & 0 & \ldots 0\\
0 & -1 & 0 & \ldots & 0 & 0 & 0 & 0 & \ldots 0\\
0 & 0 & -1 & \ldots & 0 & 0 & 0 & 0 & \ldots 0\\
\vdots & \vdots & \vdots & & & & & &\\
0 & 0 & 0 & \ldots & -1 & 0 & 0 & 0 & \ldots 0\\
b_{1i_{\nu+1}} & b_{1i_{\nu+2}} & b_{1i_{\nu+3}} & \ldots & b_{1i_n} & b_{1i_1} & b_{1i_2} & b_{1i_2} & \ldots b_{1i_\nu}\\
% ?? uncertain: $b_{1i_2}$ wiederholt an f\"unfter und siebter Stelle der Zeile, so gedruckt
b_{2i_{\nu+1}} & b_{2i_{\nu+3}} & b_{2i_{\nu+3}} & \ldots & b_{2i_n} & b_{2i_1} & b_{2i_2} & \ldots & \ldots b_{2i_\nu}\\
% ?? uncertain: $b_{2i_{\nu+3}}$ wiederholt an zweiter und dritter Stelle, so gedruckt
\vdots & \vdots & & & & & &\\
b_{\nu i_{\nu+1}} & b_{\nu i_{\nu+2}} & b_{\nu i_{\nu+3}} & \ldots & b_{\nu i_n} & b_{\nu i_1} & b_{\nu i_2} & \ldots & \ldots b_{\nu i_\nu}
\end{array}\right|.
\]
Sie zerf\"allt nach der Entwicklung von Laplace selbst wieder in ein
Product von zwei Determinanten, n\"amlich in das Product
\[
D_\varrho^{(n)}=
\left|\begin{array}{cccccc}
-1 & 0 & 0\ldots & 0\\
& 0 & -1 & 0\ldots & 0\\
& 0 & 0 & -1\ldots & 0\\
& \cdot & \cdot & \cdot & \cdot & \cdot\\
& 0 & 0 & 0\ldots & -1
\end{array}\right|
\cdot B_{\substack{1\ 2\ 3\ \ldots\ n-\nu\\i_{\nu+1},\ i_{\nu+2}\ \cdots\ i_n}}, \tag{D}
\]


% ===== pdf idx 98, printed 84.0 =====
% !PRINTED 84
\weberpage{84}
% [column spec widened array 4->5; array 4->5; array 4->5; array 4->5] a row carried more cells than the spec declared,
% which is a fatal alignment error. The spec was widened rather than a cell
% deleted, so no transcribed content is lost; whether Weber's grid really has
% that many columns is a question for a pass with the scan open.

wobei die Matrix der correspondirenden Determinante
\[
B_{\substack{1\ 2\ 3\ \ldots\ n-\nu\\i_{\nu+1},\ i_{\nu+2}\ \cdots\ i_n}}
=(-1)^{\bar\nu}
\left|\begin{array}{cccc}
b_{1i_1} & b_{1i_2} & \ldots & b_{1i_\nu}\\
b_{2i_2} & b_{2i_2} & \ldots & b_{2i_\nu}\\
% ?? uncertain: zweite Zeile druckt $b_{2i_2}$ zweimal statt $b_{2i_1},b_{2i_2}$, so gedruckt
\cdot & \cdot & \cdot & \cdot\\
b_{\nu i_1} & b_{\nu i_2} & \ldots & b_{\nu i_\nu}
\end{array}\right|.
\]
Die Determinante in Gleichung (D), deren Diagonalreihe durch die $(n-\nu)$
Elemente $(-1)$ ausgef\"ullt ist, hat den Werth
\[
\left|\begin{array}{ccccc}
-1 & 0 & 0 & 0\ldots & 0\\
0 & -1 & 0 & \ldots\ldots & 0\\
0 & 0 & -1 & \ldots\ldots & 0\\
\cdot & \cdot & \cdot & \cdot & \cdot\\
0 & 0 & 0 & 0\ldots & -1
\end{array}\right|
=(-1)^{n-\nu}.
\]
Es handelt sich also nur noch, den Exponenten $\bar\nu$ des Factors
$(-1)^{\bar\nu}$ zu ermitteln. Nach Nr. 48 \S.\,4 ist derselbe aber dargestellt
durch:
\[
\bar\nu=(1\ \ 2\ \ 3\cdots\ (n-\nu))+(i_{\nu+1}\ \ i_{\nu+2}\ \ \cdots\ i_n)
\]
\[
+[1+2+3\cdots+(n-\nu)]+[i_{\nu+1}+i_{\nu+2}+\cdots i_n). % [sic] Klammer ungleich, so gedruckt
\]
Die beiden ersten Terme in dieser Summe haben den Werth Null, da in ihnen
die Indices gut geordnet sind; also wird
\[
\bar\nu=[1+2+3\cdots+(n-\nu)]+[i_{\nu+1}+i_{\nu+2}\cdots+i_n]
\]
\[
=\frac{(n-\nu)(n-\nu+1)}{2}+\sum_{\varrho=1}^{\varrho=n-\nu}i_{\nu+\varrho}.
\]
Demnach wird die Determinante $D_\varrho^{(n)}$ dargestellt durch
\[
D_\varrho^{(n)}=(-1)^k
\left|\begin{array}{ccccc}
b_{1i_1} & b_{1i_2} & \ldots & b_{1i_\nu}\\
b_{2i_1} & b_{2i_2} & \ldots & b_{2i_\nu}\\
\cdot & \cdot & \cdot & \cdot & \cdot\\
b_{\nu i_1} & b_{\nu i_2} & \ldots & b_{\nu i_\nu}
\end{array}\right|,
\]
wobei
\[
k=(n-\nu)+\frac{(n-\nu)(n-\nu+1)}{2}+\sum_{\varrho=1}^{\varrho=n-\nu}i_{\nu+\varrho}. \tag{E}
\]

80. \textit{Substitution der Werthe von $D_\varrho^{(n)}$ und $\varDelta_\varrho^{(\nu)}$ in
die Laplace'sche Entwicklung von $R$.} Tragen wir nun diese Werthe von
$D_\varrho^{(n)}$ und $\varDelta_\varrho^{(\nu)}$ in die Gleichung (A) ein, so
erhalten wir:
\[
R=\sum{}'(-1)^\varrho\cdot(-1)^k
\left|\begin{array}{ccccc}
a_{1i_1} & a_{2i_1} & \ldots & a_{\nu i_1}\\
a_{1i_2} & a_{2i_2} & \ldots & a_{\nu i_2}\\
\vdots & \cdot & \cdot & \cdot & \cdot\\
a_{1i_\nu} & a_{2i_\nu} & \ldots & a_{\nu i_\nu}
\end{array}\right|
\cdot
\left|\begin{array}{ccccc}
b_{1i_1} & b_{1i_2} & \ldots & b_{1i_\nu}\\
b_{2i_2} & b_{2i_2} & \ldots & b_{2i_\nu}\\
% ?? uncertain: $b_{2i_2}$ wiederholt statt $b_{2i_1},b_{2i_2}$, so gedruckt
\vdots & \cdot & \cdot & \cdot & \cdot\\
b_{\nu i_1} & b_{\nu i_2} & \ldots & b_{\nu i_\nu}
\end{array}\right|.
\]


% ===== pdf idx 99, printed 85.0 =====
% !PRINTED 85
\weberpage{85}

Es erübrigt nur noch, den Exponenten von $(-1)^{\varrho+k}$ zu vereinfachen. Wegen der Relationen $(C)$ und $(E)$ ist aber

\[
\varrho+k=\frac{\nu(\nu+1)}{2}+\sum_{\varrho=1}^{\varrho=\nu}i_\varrho+\sum_{\varrho=1}^{\varrho=n-\nu}i_{r+\varrho}+\frac{(n-\nu)(n-\nu+1)}{2}+n-\nu.
\]

Da der zweite und dritte Term rechts zusammen nichts anderes ist, als die Summe aller Indices von 1 bis $n$, so erhalten wir

\[
\begin{aligned}
\varrho+k&=\frac{\nu(\nu+1)}{2}+\frac{n(n+1)}{2}+\frac{(n-\nu)(n-\nu+1)}{2}+n-\nu\\
&=(\nu^2+\nu+n^2+n+n^2-2\nu n+\nu^2+n-\tfrac{1}{2})+n-\nu\\
&=(\nu^2+n^2+n-n\cdot\nu)+n-\nu\\
&=(\nu^2-\nu+n^2+n)+n(1-\nu).
\end{aligned}
\]

Das erste Glied dieser Summe: $(\nu^2-\nu+n^2+n)$ ist jedenfalls eine gerade Zahl, kann also unberücksichtigt bleiben, das zweite Glied $n(1-\nu)$ gerade oder ungerade, ebenso wie $n(1+\nu)$; wir können daher $\varrho+k$ durch $n(1+\nu)$ ersetzen, so dass wir nun als Endresultat erhalten:

\[
R=(-1)^{n(\nu+1)}\sum
\begin{vmatrix}
a_{1i_1} & a_{2i_1} & \dots & a_{\nu i_1}\\
a_{1i_2} & a_{2i_2} & \dots & a_{\nu i_2}\\
\cdot & \cdot & \cdot & \cdot\\
a_{1i_\nu} & a_{2i_\nu} & \dots & a_{\nu i_\nu}
\end{vmatrix}
\cdot
\begin{vmatrix}
b_{1i_1} & b_{1i_2} & \dots & b_{1i_\nu}\\
b_{2i_1} & b_{2i_2} & \dots & b_{2i_\nu}\\
\cdot & \cdot & \cdot & \cdot\\
b_{\nu i_1} & b_{\nu i_2} & \dots & b_{\nu i_\nu}
\end{vmatrix},
\]

oder, wenn wir noch die Matrix mit den Elementen $a_{k i_\varrho}$ transponiren

\weq{F}{
R=(-1)^{n(\nu+1)}\sum
\begin{vmatrix}
a_{1i_1} & a_{1i_2} & \dots & a_{1i_\nu}\\
a_{2i_1} & a_{2i_2} & \dots & a_{2i_\nu}\\
\cdot & \cdot & \cdot & \cdot\\
a_{\nu i_1} & a_{\nu i_2} & \dots & a_{\nu i_\nu}
\end{vmatrix}
\cdot
\begin{vmatrix}
b_{1i_1} & b_{1i_2} & \dots & b_{1i_\nu}\\
b_{2i_1} & b_{2i_2} & \dots & b_{2i_\nu}\\
\cdot & \cdot & \cdot & \cdot\\
b_{\nu i_1} & b_{\nu i_2} & \dots & b_{\nu i_\nu}
\end{vmatrix}
}

81. \textit{Zusammenfassung der Resultate der letzten Untersuchungen.} Wir haben somit für die Determinante $R$ zwei der Form nach verschiedene Entwicklungen erhalten, die beide mit den zwei Matrices

\[
\begin{vmatrix}
a_{11} & a_{12} & a_{13} & \dots & a_{1n}\\
a_{21} & a_{22} & a_{23} & \dots & a_{2n}\\
\cdot & \cdot & \cdot & \cdot & \cdot\\
a_{\nu 1} & a_{\nu 2} & a_{\nu 3} & \dots & a_{\nu n}
\end{vmatrix}
\quad\text{und}\quad
\begin{vmatrix}
b_{11} & b_{12} & b_{13} & \dots & b_{1n}\\
b_{21} & b_{22} & b_{23} & \dots & b_{2n}\\
\cdot & \cdot & \cdot & \cdot & \cdot\\
b_{\nu 1} & b_{\nu 2} & b_{\nu 3} & \dots & b_{\nu n}
\end{vmatrix},
\]

wo $n>\nu$, in einer gewissen Beziehung stehen. Verbindet man nämlich die beiden Matrices einmal in der Weise, wie sie das Gesetz über
%CONT>>


% ===== pdf idx 100, printed 86.0 =====
% !PRINTED 86
\weberpage{86}

%<<CONT
Multiplication von Determinanten gebietet, so erhält man die zuerst gegebene (Nr. 77) Darstellung von $R$; multiplicirt man andrerseits je eine Determinante der ersten Matrix mit der entsprechenden, d. h. jener, die aus Colonnen mit gleichem Index gebildet ist, der andern Matrix, und summirt die so erhaltenen $\binom{n}{\nu}$ Producte, so erhält man die zweite Darstellung von $R$ (Nr. 80, $F$). In beiden Entwicklungen tritt der Factor $(-1)^{n(\nu+1)}$ auf, so dass wir also nach Division mit diesem Factor die Beziehung erhalten:

\[
\begin{vmatrix}
a_{11}b_{11}+a_{12}b_{12}+\cdots a_{1n}b_{1n}, & a_{21}b_{11}+\cdots a_{2n}b_{1n}, & \cdots, & a_{\nu 1}b_{11}+\cdots a_{\nu n}b_{1n}\\
a_{11}b_{21}+a_{12}b_{22}+\cdots a_{1n}b_{2n}, & a_{21}b_{21}+\cdots a_{2n}b_{2n}, & \cdots, & a_{\nu 1}b_{21}+\cdots a_{\nu n}b_{2n}\\
a_{11}b_{31}\cdot\cdot\cdot\cdot\cdot\cdot, & a_{21}b_{31}+\cdot\cdot\cdot, & \cdots, & a_{\nu 1}b_{31}\cdot\cdot\cdot\\
\cdot & \cdot & \cdot & \cdot\\
a_{11}b_{\nu 1}+a_{12}b_{\nu 2}+\cdots a_{1n}b_{\nu n}, & a_{21}b_{\nu 1}\cdots a_{2n}b_{\nu n}, & \cdots, & a_{\nu 1}b_{\nu 1}+\cdots a_{\nu n}b_{\nu n}
\end{vmatrix}
\]
\[
=\sum
\begin{vmatrix}
a_{1i_1} & a_{2i_2} & \dots & a_{1i_\nu}\\
a_{2i_1} & a_{2i_2} & \dots & a_{2i_\nu}\\
\cdot & \cdot & \cdot & \cdot\\
a_{\nu i_1} & a_{\nu i_2} & \dots & a_{\nu i_\nu}
\end{vmatrix}
\cdot
\begin{vmatrix}
b_{1i_1} & b_{1i_2} & \dots & b_{1i_\nu}\\
b_{2i_1} & b_{2i_2} & \dots & b_{2i_\nu}\\
\cdot & \cdot & \cdot & \cdot\\
b_{\nu i_1} & b_{\nu i_2} & \dots & b_{\nu i_\nu}
\end{vmatrix}
\]
\[
=
\begin{vmatrix}
a_{11} & a_{12} & a_{13} & \dots & a_{1n}\\
a_{21} & a_{22} & a_{23} & \dots & a_{2n}\\
\cdot & \cdot & \cdot & \cdot & \cdot\\
a_{\nu 1} & a_{\nu 2} & . & . & a_{\nu n}
\end{vmatrix}
\cdot
\begin{vmatrix}
b_{11} & b_{12} & b_{13} & \dots & b_{1n}\\
b_{21} & b_{22} & b_{23} & \dots & b_{2n}\\
\cdot & \cdot & \cdot & \cdot & \cdot\\
b_{\nu 1} & b_{\nu 2} & b_{\nu 3} & \dots & b_{\nu n}
\end{vmatrix}.
\]

Dabei hat das letzte Product, also das Product der beiden Matrices, natürlich nur symbolische Bedeutung, insofern damit die unmittelbar vorhergehende Summe von Producten je zweier entsprechender Determinanten $\nu^{\text{ten}}$ Grades der beiden Matrices ausgedrückt sein soll. So lange also $n>\nu$, liefert das Product der beiden Matrices, gebildet durch Multiplication von Zeilen mit Zeilen, eine Determinante $R$, welche im Allgemeinen nicht verschwindet.

\textit{82. Besprechung des Falles, bei welchem in beiden Matrices die Zahl der Zeilen jene der Colonnen überschreitet.} Ist nun aber $n<\nu$, d. h. die Zahl der Zeilen grösser als die der Colonnen, dann führt der durch

\weq{G}{
\begin{vmatrix}
a_{11} & a_{12} & \dots & a_{1n}\\
a_{21} & a_{22} & \dots & a_{2n}\\
\cdot & \cdot & \cdot & \cdot\\
\cdot & \cdot & \cdot & \cdot\\
a_{\nu 1} & a_{\nu 2} & \dots & a_{\nu n}
\end{vmatrix}
\cdot
\begin{vmatrix}
b_{11} & b_{12} & \dots & b_{1n}\\
b_{21} & b_{22} & \dots & b_{2n}\\
\cdot & \cdot & \cdot & \cdot\\
\cdot & \cdot & \cdot & \cdot\\
b_{\nu 1} & b_{\nu 2} & \dots & b_{\nu n}
\end{vmatrix}
}
%CONT>>


% ===== pdf idx 101, printed 87.0 =====
% !PRINTED 87
\weberpage{87}

%<<CONT
symbolisch angedeutete Multiplicationsprocess (wenn wir wieder die Elemente der Productdeterminante durch Multiplication von Zeilen mit Zeilen entstehen lassen) auf eine Determinante

\[
R=
\begin{vmatrix}
a_{11}b_{11}\cdots+a_{1n}b_{1n}, & a_{21}b_{11}\cdots+a_{2n}b_{1n}, & \cdots, & a_{\nu 1}b_{11}\cdots+a_{\nu n}b_{1n}\\
a_{11}b_{21}\cdots+a_{1n}b_{2n}, & a_{21}b_{21}\cdots+a_{2n}b_{2n}, & \cdots, & a_{\nu 1}b_{21}\cdots+a_{\nu n}b_{2n}\\
a_{11}b_{31}\cdots+a_{1n}b_{3n}, & \cdot & \cdot & \cdot\\
\cdot & \cdot & \cdot & \cdot\\
a_{11}b_{\nu 1}\cdots+a_{1n}b_{\nu n}, & a_{21}b_{\nu 1}\cdots+a_{2n}b_{\nu n}, & \cdots, & a_{\nu 1}b_{\nu 1}\cdots+a_{\nu n}b_{\nu n}
\end{vmatrix},
\]

welche aber identisch verschwindet. Denn genau dieselbe Determinante würden wir erhalten, wenn wir vorher jede der beiden Matrices in $(G)$ durch die nöthige Zahl von Colonnen, deren sämmtliche Elemente Null sind, zu Matrices erweitern, welche Determinanten repräsentiren, und nun die beiden Determinanten in der angegebenen Weise multipliciren. Da aber hier die Factoren den Werth Null besitzen, so verschwindet auch das Product, d. h. die Determinante $R$ ist identisch Null, wenn $\nu>n$.

Fassen wir also die erhaltenen Resultate zusammen, so gelangen wir zu dem Satze:

\glqq Die Determinante $R$ ist gleich dem Producte zweier $n$-gliederiger Determinanten, wenn $\nu=n$; sie zerfällt in die Summe von $\binom{n}{\nu}$ Producten je zweier entsprechender $\nu$-gliederiger Determinanten, wenn $n>\nu$; sie ist endlich identisch Null, wenn $n<\nu$.\grqq{}

\textit{83. Beispiel 1.}

\[
\begin{vmatrix}
x_1 & x_2 & x_3\\
a_1 & a_2 & a_3
\end{vmatrix}
\cdot
\begin{vmatrix}
y_1 & y_2 & y_3\\
b_1 & b_2 & b_3
\end{vmatrix}
\]
\[
=
\begin{vmatrix}
x_1y_1+x_2y_2+x_3y_3, & a_1y_1+a_2y_2+a_3y_3\\
x_1b_1+x_2b_2+x_3b_3, & a_1b_1+a_2b_2+a_3b_3
\end{vmatrix}
\]
\[
=
\begin{vmatrix}x_1 & x_2\\ b_1 & b_2\end{vmatrix}\cdot\begin{vmatrix}y_1 & y_2\\ b_1 & b_2\end{vmatrix}
+\begin{vmatrix}x_1 & x_3\\ a_1 & a_3\end{vmatrix}\cdot\begin{vmatrix}y_1 & y_3\\ b_1 & b_3\end{vmatrix}
+\begin{vmatrix}x_2 & x_3\\ a_2 & a_3\end{vmatrix}\cdot\begin{vmatrix}y_2 & y_3\\ b_2 & b_3\end{vmatrix}.
\]

\textit{Beispiel 2.}

\[
\begin{vmatrix}u_1 & v_1\\ u_2 & v_2\\ u_3 & v_3\end{vmatrix}
\cdot
\begin{vmatrix}a_1 & b_1\\ a_2 & b_2\\ a_3 & b_3\end{vmatrix}
=
\begin{vmatrix}
v_1a_1+v_1b_1, & u_2a_1+v_2b_1, & u_3a_1+v_3b_1\\ % [sic]
v_1a_2+v_1b_2, & u_2a_2+v_2b_2, & u_3a_2+v_3b_2\\
v_1a_3+v_1b_3, & u_2a_3+v_2b_3, & u_3a_3+v_3b_3
\end{vmatrix}
\]
\[
=
\begin{vmatrix}u_1 & v_1 & 0\\ u_2 & v_2 & 0\\ u_3 & v_3 & 0\end{vmatrix}
\cdot
\begin{vmatrix}a_1 & b_1 & 0\\ a_2 & b_2 & 0\\ a_3 & b_3 & 0\end{vmatrix}
=0.
\]

\textit{84. Beispiel 3.} Die Darstellung einer Determinante wie in Beispiel 1 durch eine Summe von Producten wollen wir noch verwenden,
%CONT>>


% ===== pdf idx 102, printed 88.0 =====
% !PRINTED 88
\weberpage{88}

%<<CONT
um eine Formel der analytischen Geometrie aufzustellen. Bekanntlich erfüllen die drei Richtungscosinusse einer Geraden im Raume die Relation

\[
1=\cos^2\alpha+\cos^2\beta+\cos^2\gamma,
\]

und ferner ist der Neigungswinkel dieser Geraden mit einer zweiten von den Richtungscosinussen $\cos\alpha_1$, $\cos\beta_1$, $\cos\gamma_1$ gegeben durch

\[
\cos v=\cos\alpha\cos\alpha_1+\cos\beta\cos\beta_1+\cos\gamma\cos\gamma_1.
\]

Gemäss diesen beiden Relationen kann man die Determinante

\[
\begin{vmatrix}\cos v & 1\\ 1 & \cos v\end{vmatrix}=\sin^2 v
\]

auch schreiben

\[
\begin{vmatrix}
\cos^2\alpha+\cos^2\beta+\cos^2\gamma, & \cos\alpha\cos\alpha_1+\cos\beta\cos\beta_1+\cos\gamma+\cos\gamma_1\\
\cos\alpha\cos\alpha_1+\cos\beta\cos\beta_1+\cos\gamma\cos\gamma_1, & \cos^2\alpha_1+\cos^2\beta_1+\cos^2\gamma_1
\end{vmatrix}. % [sic] top-right entry prints "cos gamma + cos gamma_1" not "cos gamma cos gamma_1"
\]

Diese zerfällt aber nach Beispiel 1 in die drei Producte je zweier Determinanten, so dass man erhält

\[
\sin^2v=
\begin{vmatrix}\cos\alpha & \cos\beta\\ \cos\alpha_1 & \cos\beta_1\end{vmatrix}^2
+\begin{vmatrix}\cos\beta & \cos\gamma\\ \cos\beta_1 & \cos\gamma_1\end{vmatrix}^2
+\begin{vmatrix}\cos\gamma & \cos\alpha\\ \cos\gamma_1 & \cos\alpha_1\end{vmatrix}^2.
\]

\begin{center}
\S.\ 7.\\
\gesperrt{Adjungirte Determinanten und correspondirende Matrices.}
\end{center}

\textit{85. Definition der adjungirten Determinante.} Eine Determinante von $n^2$ Elementen $a_{ik}$ besitzt auch $n^2$ Unterdeterminanten $A_{ik}$, aus welchen man, wie aus den Elementen $a_{ik}$ selbst, die Matrix einer Determinante zusammensetzen kann. Ist die Matrix der ursprünglichen Determinante

\[
\Delta=
\begin{vmatrix}
a_{11} & a_{12} & a_{13} & \dots & a_{1n}\\
a_{21} & a_{22} & \dots & \dots & a_{2n}\\
a_{31} & \cdot & \cdot & \cdot & \cdot\\
\cdot & \cdot & \cdot & \cdot & \cdot\\
a_{n1} & a_{n2} & \dots & \dots & a_{nn}
\end{vmatrix},
\]

so ist die aus den entsprechenden Unterdeterminanten gebildete

\[
D=
\begin{vmatrix}
A_{11} & A_{12} & A_{13} & \dots & A_{1n}\\
A_{21} & A_{22} & \dots & \dots & A_{2n}\\
A_{31} & \cdot & \cdot & \cdot & \cdot\\
\cdot & \cdot & \cdot & \cdot & \cdot\\
A_{n1} & A_{n2} & \dots & \dots & A_{nn}
\end{vmatrix}.
\]

Man nennt $D$ die \textit{adjungirte} Determinante zu $\Delta$. Bei der Untersuchung dieser adjungirten Determinante drängen sich zunächst zwei Fragen
%CONT>>


% ===== pdf idx 103, printed 89.0 =====
% !PRINTED 89
\weberpage{89}
\S.\,7. Adjungirte Determinanten und correspondirende Matrices.

auf: 1) Welches ist der Werth der Determinante selbst, 2) welches sind die Werthe der Unterdeterminanten sowohl der ersten als auch der übrigen Ordnungen?

86. \textit{Berechnung des Werthes der adjungirten Determinante.} Zur Beantwortung der ersten Frage bilden wir das Product der beiden Determinanten $D$ und $\Delta$. Wir erhalten:
\[
\begin{vmatrix}
a_{11}A_{11}+a_{12}A_{12}+\cdot a_{1n}A_{1n}, & a_{11}A_{21}+a_{12}A_{22}+\cdots+a_{1n}A_{2n}, & \cdot\cdot, & a_{11}A_{n1}\cdots+a_{1n}A_{nn} \\
a_{21}A_{11}+a_{22}A_{12}\cdots+a_{2n}A_{1n}, & a_{21}A_{21}+a_{22}A_{22}\cdots+a_{2n}A_{2n}, & \cdot\cdot, & a_{21}A_{n1}\cdots+a_{2n}A_{nn} \\
a_{31}A_{11}+\cdots+a_{3n}A_{1n}, & \cdot\;\cdot\;\cdot\;\cdot\;\cdot\;\cdot\;\cdot\;\cdot\;\cdot, & \cdot\cdot, & \cdot\;\cdot\;\cdot\;\cdot \\
\cdot\;\cdot\;\cdot\;\cdot\;\cdot\;\cdot\;\cdot\;\cdot\;\cdot\;\cdot\;\cdot\;\cdot\;\cdot\;\cdot\;\cdot\;\cdot\;\cdot & & & \\
a_{n1}A_{11}\cdots+a_{nn}A_{1n}, & a_{n1}A_{21}+a_{n2}A_{22}\cdots+a_{nn}A_{1n}, & \cdot\cdot, & a_{n1}A_{n1}+a_{nn}A_{nn}
\end{vmatrix}
\]
% ?? uncertain: row 1 above prints an abbreviated "+\cdot a_{1n}" between the second and third terms, unlike the "+\cdots+" ellipsis used in the other rows; reproduced as printed
\[
= D\cdot\Delta.
\]
In dieser Determinante verschwinden aber wegen der Relationen $(C)$ Nr. 32 alle ausserhalb der Diagonale stehenden Elemente, während die $n$ in ihr befindlichen Elemente den Werth $\Delta$ besitzen; sie reducirt sich also auf ihr Diagonalglied, so dass die Beziehung entsteht
\weq{1}{D\cdot\Delta=\Delta^n,}
oder
\[
D=\Delta^{n-1}.
\]

87. \textit{Berechnung eines Minors beliebiger Ordnung der adjungirten Determinante.} Auch die zweite Frage lässt sich mit Hilfe des Productsatzes für Determinanten erledigen. Wir erweitern zu diesem Zwecke nur zuerst die Matrix der Unterdeterminante $(n-\varrho)^{\text{ter}}$ Ordnung zu einer Matrix von $n^2$ Elementen, wozu ja der Satz von Laplace in einfacher Weise die Mittel liefert. Sei in erster Linie die gegebene Unterdeterminante von $D$ dargestellt durch
\[
D^{(n-\varrho)}=
\begin{vmatrix}
A_{11} & A_{12} & A_{13} & \ldots & A_{1\varrho} \\
A_{21} & A_{22} & \ldots & \ldots & A_{2\varrho} \\
A_{31} & \cdot & \cdot & \cdot & \cdot \\
\cdot & \cdot & \cdot & \cdot & \cdot \\
A_{\varrho 1} & A_{\varrho 2} & \ldots & \ldots & A_{\varrho\varrho}
\end{vmatrix},
\]
so bildet man die Erweiterung dieser Matrix dadurch, dass man zunächst in der Richtung der Diagonale noch $(n-\varrho)$ Elemente vom Werthe $1$ anfügt; die übrigen Stellen unterhalb der $\varrho^{\text{ten}}$ Zeile füllt man mit dem Elemente $0$ aus, während in die rechts der $\varrho^{\text{ten}}$ Colonne befindlichen Stellen diejenigen Elemente $A_{ik}$ einrücken, welche bereits in der ursprünglichen Determinante $D$ dieselben innehaben. Dadurch wird die Matrix von $D^{(n-\varrho)}$:
%CONT>>


% ===== pdf idx 104, printed 90.0 =====
% !PRINTED 90
\weberpage{90}
Erster Theil. Theorie der Determinanten.

%<<CONT
\weq{2}{D^{(n-\varrho)}=
\begin{vmatrix}
A_{11}\;A_{12}\;\ldots\;A_{1\varrho}\;A_{1,\varrho+1}\;\ldots\;A_{1n} \\
A_{21}\;A_{22}\;\ldots\;A_{2\varrho}\;A_{2,\varrho+1}\;\ldots\;A_{2n} \\
\vdots\;\vdots\;\ \ \vdots\;\vdots\;\ \ \vdots \\
A_{\varrho 1}\;A_{\varrho 2}\;\ldots\;A_{\varrho\varrho}\;A_{\varrho,\varrho+1}\;\ldots\;A_{\varrho n} \\
0\;0\;\ldots\;0\;1\;0\;\ldots\;0 \\
0\;0\;\ldots\;0\;0\;1\;\ldots\;0 \\
\cdot\;\cdot\;\cdot\;\cdot\;\cdot\;\cdot\;\cdot\;\cdot \\
0\;0\;\ldots\;0\;0\;0\;\ldots\;1
\end{vmatrix}.}
Die so erhaltene Determinante multipliciren wir nun wieder mit der Determinante $\Delta$.

Das Product ist dargestellt durch:
\[
\Delta\cdot D^{(n-\varrho)}=
\begin{vmatrix}
a_{11}\;a_{12}\;a_{13}\;\ldots\;a_{1n} \\
a_{21}\;a_{22}\;a_{23}\;\ldots\;a_{2n} \\
a_{31}\;\cdot\;\cdot\;\cdot\;\cdot\;\cdot\;\cdot\;\cdot\;\cdot\;a_{3n} \\
\cdot\;\cdot\;\cdot\;\cdot\;\cdot\;\cdot\;\cdot\;\cdot\;\cdot\;\cdot\;\cdot\;\cdot\;\cdot\;\cdot \\
a_{\varrho 1}\;a_{\varrho 2}\;a_{\varrho 3}\;\ldots\;a_{\varrho n} \\
a_{\varrho+1,1}\;\ldots\;\ldots\;\ldots\;a_{\varrho+1,n} \\
\cdot\;\cdot\;\cdot\;\cdot\;\cdot\;\cdot\;\cdot\;\cdot\;\cdot\;\cdot\;\cdot\;\cdot\;\cdot \\
a_{n1}\;a_{n2}\;\ldots\;\ldots\;a_{nn}
\end{vmatrix}
\cdot
\begin{vmatrix}
A_{11}\;A_{12}\;\ldots\;A_{1\varrho}\;A_{1,\varrho+1}\;\ldots\;A_{1n} \\
A_{21}\;A_{22}\;\ldots\;A_{2\varrho}\;A_{2,\varrho+1}\;\ldots\;A_{2n} \\
A_{31}\;\cdot\;\cdot\;\cdot\;\cdot\;\cdot\;\cdot\;\cdot\;\cdot \\
\cdot\;\cdot\;\cdot\;\cdot\;\cdot\;\cdot\;\cdot\;\cdot\;\cdot\;\cdot\;\cdot\;\cdot \\
A_{\varrho 1}\;A_{\varrho 2}\;\ldots\;A_{\varrho\varrho}\;A_{\varrho,\varrho+1}\;\ldots\;A_{\varrho n} \\
0\;0\;\ldots\;0\;1\;0\;\ldots\;0 \\
0\;0\;\ldots\;0\;0\;1\;\ldots\;0 \\
0\;0\;\ldots\;0\;0\;0\;\ldots\;1
\end{vmatrix}
\]
\[
=
\begin{vmatrix}
a_{11}A_{11}\cdots+a_{1n}A_{1n},\ \ldots,\ a_{11}A_{\varrho1}\cdots+a_{1n}A_{\varrho n},\ a_{1,\varrho+1},\ a_{1,\varrho+2},\ \cdots,\ a_{1n} \\
a_{21}A_{11}\cdots+a_{2n}A_{1n},\ \ldots,\ a_{21}A_{\varrho1}\cdots+a_{2n}A_{\varrho n},\ a_{2,\varrho+1},\ a_{2,\varrho+2},\ \cdots,\ a_{2n} \\
a_{31}A_{11}\ \cdot\ \cdot\ \cdot\ \cdot\ \cdot\ \cdot\ \cdot\ \cdot\ \cdot\ \cdot\ ,\ a_{3,\varrho+1},\ \cdots,\ \cdots,\ a_{3n} \\
\cdot\;\cdot\;\cdot\;\cdot\;\cdot\;\cdot\;\cdot\;\cdot\;\cdot\;\cdot\;\cdot\;\cdot\;\cdot \\
a_{\varrho1}A_{11}\cdots+a_{\varrho n}A_{1n},\ \ldots,\ a_{\varrho1}A_{\varrho,1}\cdots+a_{\varrho n}A_{\varrho n},\ a_{\varrho,\varrho+1},\ a_{\varrho,\varrho+2},\ \cdots,\ a_{\varrho n} \\
a_{\varrho+1,1}A_{11}\ \cdot\ \cdot\ \cdot\ \cdot\ \cdot\ \cdot\ \cdot\ \cdot\ \cdot\ \cdot\ \cdot\ \cdot\ a_{\varrho+1,\varrho+1}\ \cdot\ \cdot\ \cdot\ ,\ a_{\varrho+1,n} \\
\cdot\;\cdot\;\cdot\;\cdot\;\cdot\;\cdot\;\cdot\;\cdot\;\cdot\;\cdot\;\cdot\;\cdot\;\cdot\;\cdot \\
a_{n1}A_{11}\cdots+a_{nn}A_{1n},\ \ldots,\ a_{n1}A_{\varrho1}\cdots+a_{nn}A_{\varrho n},\ a_{n,\varrho+1},\ a_{n,\varrho+2},\ \cdots,\ a_{nn}
\end{vmatrix}
\]
% ?? uncertain: the penultimate entry of the fifth row above ("$a_{\varrho,\varrho+2}$") is faint in the scan and could alternatively read "$a_{\varrho,\varrho-2}$"; reproduced here consistent with the pattern of the surrounding rows
\[
=
\begin{vmatrix}
\Delta\;0\;0\;\ldots\;0,\ a_{1,\varrho+1},\ a_{1,\varrho+2},\ \ldots,\ a_{1n} \\
0\;\Delta\;0\;\ldots\;0,\ a_{2,\varrho+1},\ a_{2,\varrho+2},\ \ldots,\ a_{2n} \\
0\;0\;\Delta\;\ldots\;0,\ a_{3,\varrho+1},\ \cdot\ \cdot\ \cdot\ ,\ \cdot \\
\cdot\;\cdot\;\cdot\;\cdot\;\cdot\;\cdot\;\cdot\;\cdot\;\cdot\;\cdot \\
0\;0\;0\;\ldots\;\Delta,\ a_{\varrho,\varrho+1},\ \cdot\ \cdot\ \cdot\ ,\ a_{\varrho n} \\
0\;0\;0\;\ldots\;0,\ a_{\varrho+1,\varrho+1},\ \cdot\ \cdot\ \cdot\ ,\ a_{\varrho+1,n} \\
\cdot\ \cdot\ \cdot\ \cdot\ 0\ \cdot\ \cdot\ \cdot\ \cdot\ \cdot\ \cdot\ \cdot\ \cdot \\
0\;0\;0\;\ldots\;0,\ a_{n,\varrho+1},\ \cdot\ \cdot\ \cdot\ \cdot\ ,\ a_{nn}
\end{vmatrix}.
\]


% ===== pdf idx 105, printed 91.0 =====
% !PRINTED 91
\weberpage{91}
\S.\,7. Adjungirte Determinanten und correspondirende Matrices.

Die zuletzt erhaltene Determinante reducirt sich nach Laplace auf das Product $\Delta^\varrho\cdot A^{12\ldots\varrho}_{12\ldots\varrho}$, so dass wir erhalten:
\weq{3}{\Delta\cdot D^{(n-\varrho)}=\Delta\cdot
\begin{vmatrix}
A_{11}\;A_{12}\;\ldots\;A_{1\varrho} \\
A_{21}\;A_{22}\;\ldots\;A_{2\varrho} \\
\cdot\;\cdot\;\cdot\;\cdot\;\cdot \\
\cdot\;\cdot\;\cdot\;\cdot\;\cdot \\
A_{\varrho 1}\;A_{\varrho 2}\;\ldots\;A_{\varrho\varrho}
\end{vmatrix}
=\Delta^\varrho\cdot
\begin{vmatrix}
a_{\varrho+1,\varrho+1}\;a_{\varrho+1,\varrho+2}\;\ldots\;a_{\varrho+1,n} \\
a_{\varrho+2,\varrho+1}\;a_{\varrho+2,\varrho+2}\;\ldots\;a_{\varrho+2,n} \\
\cdot\;\cdot\;\cdot\;\cdot\;\cdot\;\cdot\;\cdot\;\cdot \\
a_{n,\varrho+1}\;a_{n,\varrho+2}\;\ldots\;a_{nn}
\end{vmatrix}}
oder nach Division mit $\Delta$:
\weq{4}{D^{(n-\varrho)}=A^{12\ldots\varrho}_{12\ldots\varrho}\cdot\Delta^{\varrho-1}.}
Die Determinante $A^{12\ldots\varrho}_{12\ldots\varrho}$ entsteht aber aus $\Delta$ durch Unterdrückung derjenigen Zeilen und Colonnen, welche sich in den Elementen der Matrix
\[
\begin{vmatrix}
a_{11}\;a_{12}\;\ldots\;a_{1\varrho} \\
a_{21}\;a_{22}\;\ldots\;a_{2\varrho} \\
\cdot\;\cdot\;\cdot\;\cdot\;\cdot \\
a_{\varrho 1}\;a_{\varrho 2}\;\ldots\;a_{\varrho\varrho}
\end{vmatrix}
\]
schneiden, d.\,h. in jenen Elementen von $\Delta$, deren Minoren erster Ordnung gerade die Elemente des Minors $D^{(n-\varrho)}$ bilden.

88. \textit{Besprechung des allgemeinen Falles.} Die Relation (4) gilt zunächst für den speciellen Minor $(n-\varrho)^{\text{ter}}$ Ordnung,* dessen Elemente die $\varrho^2$ ersten Plätze in $D$ einnehmen.
% ?? uncertain: footnote marker "*" after "Ordnung," on this page — no footnote text is visible within the captured top/mid/bot bands
Ist der Minor beliebig gewählt, etwa
\[
D^{(n-\varrho)}=
\begin{vmatrix}
A_{i_1k_1}\;A_{i_1k_2}\;\ldots\;A_{i_1k_\varrho} \\
A_{i_2k_1}\;A_{i_2k_2}\;\ldots\;A_{i_2k_\varrho} \\
A_{i_3k_1}\;\cdot\;\cdot\;\cdot\;\cdot \\
\cdot\;\cdot\;\cdot\;\cdot\;\cdot \\
A_{i_\varrho k_1}\;A_{i_\varrho k_2}\;\ldots\;A_{i_\varrho k_\varrho}
\end{vmatrix},
\]
dann können wir auf ganz analogem Wege seinen Werth berechnen. Wir erweitern wieder die Matrix dieses Minors in der unter. Nr. 87 % [sic]
%CONT>>


% ===== pdf idx 106, printed 92.0 =====
% !PRINTED 92
\weberpage{92}
Erster Theil. Theorie der Determinanten.

%<<CONT
geschilderten Weise zu einer Matrix von $n^2$ Elementen, indem wir zunächst jene $(n-\varrho)$ Colonnen anfügen, durch deren Unterdrückung die Colonnen des Minors entstanden sind, natürlich möglichst gut geordnet unter Beachtung der Anordnung ihrer Elemente nach Zeilen, welche jenen des Minors entsprechen. Sodann erweitern wir die Anzahl der Zeilen der so erhaltenen Matrix gerade so wie im speciellen Falle. Die Determinante $D^{(n-\varrho)}$ ist alsdann inclusive Vorzeichen durch die Matrix dargestellt:
\[
\begin{vmatrix}
A_{i_1k_1}\;A_{i_1k_2}\;A_{i_1k_3}\;\ldots\;A_{i_1k_\varrho}\;A_{i_1k_{\varrho+1}}\;A_{i_1k_{\varrho+2}}\;\ldots\;A_{i_1k_n} \\
A_{i_2k_1}\;A_{i_2k_2}\;\ldots\;\ldots\;A_{i_2k_\varrho}\;A_{i_2k_{\varrho+1}}\;\ldots\;\ldots\;\ldots\;A_{i_2k_n} \\
A_{i_3k_1}\;\cdot\;\cdot\;\cdot\;\cdot\;\cdot\;\cdot\;\cdot\;\cdot\;\cdot\;\cdot\;\cdot\;\cdot \\
\cdot\;\cdot\;\cdot\;\cdot\;\cdot\;\cdot\;\cdot\;\cdot\;\cdot\;\cdot\;\cdot\;\cdot\;\cdot\;\cdot \\
A_{i_\varrho k_1}\;A_{i_\varrho k_2}\;\ldots\;\ldots\;\ldots\;A_{i_\varrho k_\varrho}\;A_{i_\varrho k_{\varrho+1}}\;A_{i_\varrho k_{\varrho+2}}\;\ldots\;A_{i_\varrho k_n} \\
0\;0\;\ldots\;0\;1\;0\;\ldots\;0 \\
0\;0\;\ldots\;0\;0\;1\;\ldots\;0 \\
\cdot\;\cdot\;\cdot\;\cdot\;\cdot\;\cdot\;\cdot\;\cdot \\
0\;0\;\ldots\;0\;0\;0\;\ldots\;1
\end{vmatrix}.
\]
Bringen wir nun auch in der Determinante $\Delta$ die den Elementen $A_{i_\mu k_\mu}$ des Minors $D^{(n-\varrho)}$ entsprechenden Elemente $a_{i_\mu k_\mu}$ durch cyklische Vertauschung an die $\varrho^2$ ersten Plätze, dann hat gemäss den Entwicklungen in Nr. 48 die so erhaltene Matrix den Werth
\[
\begin{vmatrix}
a_{i_1k_1}\;a_{i_1k_2}\;\ldots\;a_{i_1k_\varrho}\;a_{i_1k_{\varrho+1}}\;\ldots\;a_{i_1k_n} \\
a_{i_2k_1}\;a_{i_2k_2}\;\ldots\;a_{i_2k_\varrho}\;a_{i_2k_{\varrho+1}}\;\ldots\;a_{i_2k_n} \\
a_{i_3k_1}\;\cdot\;\cdot\;\cdot\;\cdot\;\cdot\;\cdot\;\cdot\;\cdot\;\cdot\;\cdot\;\cdot\;\cdot \\
\cdot\;\cdot\;\cdot\;\cdot\;\cdot\;\cdot\;\cdot\;\cdot\;\cdot\;\cdot\;\cdot\;\cdot\;\cdot\;\cdot \\
a_{i_nk_1}\;a_{i_nk_2}\;\ldots\;a_{i_nk_\varrho}\;a_{i_nk_{\varrho+1}}\;\ldots\;a_{i_nk_n}
\end{vmatrix}
=(-1)^\nu\cdot\Delta,
\]
wo $\nu=(i_1\,i_2\,i_3\,\ldots\,i_\varrho)+(k_1\,k_2\,\ldots\,k_\varrho)+\displaystyle\sum_{\sigma=1}^{\sigma=\varrho}i_\sigma+\sum_{\sigma=1}^{\sigma=\varrho}k_\sigma.$

Bildet man jetzt das Product der beiden Determinanten $D^{(n-\varrho)}$ und $(-1)^\nu\Delta$, indem man die so umgestalteten Matrices derselben zu Grunde legt, so erhält man durch Multiplication der Elemente nach ihren Zeilen
%CONT>>


% ===== pdf idx 107, printed 93.0 =====
% !PRINTED 93
\weberpage{93}
\S.\,7. Adjungirte Determinanten und correspondirende Matrices.

\[
(-1)^{\nu}\cdot \Delta \cdot D^{(n-\varrho)} =
\begin{vmatrix}
\Delta & 0 & 0 & \cdots & 0 & a_{i_1 k_{\varrho+1}} & a_{i_1 k_{\varrho+2}} & \cdots & a_{i_1 k_n} \\
0 & \Delta & 0 & \cdots & 0 & a_{i_2 k_{\varrho+1}} & a_{i_2 k_{\varrho+2}} & \cdots & a_{i_2 k_n} \\
0 & 0 & \Delta & \cdots & 0 & a_{i_3 k_{\varrho+1}} & \cdot & \cdot & \cdot \\
\cdot & \cdot & \cdot & \cdot & \cdot & \cdot & \cdot & \cdot & \cdot \\
0 & 0 & 0 & \cdots & \Delta & a_{i_\varrho k_{\varrho+1}} & a_{i_\varrho k_{\varrho+2}} & \cdots & a_{i_\varrho k_n} \\
0 & 0 & 0 & \cdots & 0 & a_{i_{\varrho+1} k_{\varrho+1}} & a_{i_{\varrho+1} k_{\varrho+2}} & \cdots & a_{i_{\varrho+1} k_n} \\
0 & \cdot & \cdot & \cdot & 0 & a_{i_{\varrho+2} k_{\varrho+1}} & \cdot & \cdot & \cdot \\
\cdot & \cdot & \cdot & \cdot & \cdot & \cdot & \cdot & \cdot & \cdot \\
0 & 0 & 0 & \cdots & 0 & a_{i_n k_{\varrho+1}} & a_{i_n k_{\varrho+2}} & \cdots & a_{i_n k_n}
\end{vmatrix}\, .
\]

Diese Determinante spaltet sich nach dem Satze von Laplace wieder in das Product zweier anderer, so dass wir erhalten:

\begin{equation*}
(-1)^{\nu}\cdot \Delta \cdot D^{(n-\varrho)} = \Delta^{\varrho}\cdot
\begin{vmatrix}
a_{i_{\varrho+1} k_{\varrho+1}} & a_{i_{\varrho+1} k_{\varrho+2}} & \cdots & a_{i_{\varrho+1} k_n} \\
a_{i_{\varrho+2} k_{\varrho+1}} & a_{i_{\varrho+2} k_{\varrho+2}} & \cdots & a_{i_{\varrho+2} k_n} \\
a_{i_{\varrho+3} k_{\varrho+1}} & \cdot & \cdot & \cdot \\
\cdot & \cdot & \cdot & \cdot \\
a_{i_n k_{\varrho+1}} & a_{i_n k_{\varrho+2}} & \cdots & a_{i_n k_n}
\end{vmatrix}
\tag{$A$}
\end{equation*}

Die Matrix rechts stellt aber mit $(-1)^{\nu}$ multiplicirt gerade den Minor $A_{\substack{i_1 i_2 \ldots i_\varrho\\k_1 k_2 \ldots k_\varrho}}$ dar, d.\,h. die correspondirende Determinante zu jener Unterdeterminante von $\Delta$, deren Elemente den Elementen des Minors $D^{(n-\varrho)}$ entsprechen. Multipliciren wir also die Gleichung (A) mit $(-1)^{\nu}$, und dividiren wir sie mit $\Delta$, so kommt

\[
D^{(n-\varrho)} = \Delta^{\varrho-1}\cdot A_{\substack{i_1 i_2 \ldots i_\varrho\\k_1 k_2 \ldots k_\varrho}}\, .
\]

Wir können daher den Satz aufstellen:

\glqq Ein Minor $\varrho^{\mathrm{ten}}$ Grades der zu $\Delta$ adjungirten Determinante $D$ ist gleich dem Producte der $(\varrho-1)^{\mathrm{ten}}$ Potenz von $\Delta$ in die correspondirende Determinante jenes Minors von $\Delta$, dessen Elemente $a_{ik}$ den Elementen $A_{ik}$ des Minors von $D$ entsprechen.\grqq{}

89. \textit{Beispiel.} Es seien

\[
\begin{vmatrix}
a_{11} & a_{12} & a_{13} & \cdots & a_{15} \\
a_{21} & a_{22} & \cdot & \cdot & a_{25} \\
a_{31} & \cdot & \cdot & \cdot & \cdot \\
\cdot & \cdot & \cdot & \cdot & \cdot \\
a_{51} & \cdot & \cdot & \cdot & a_{55}
\end{vmatrix}
\quad\text{und}\quad
\begin{vmatrix}
A_{11} & A_{12} & \cdots & A_{15} \\
A_{21} & \cdot & \cdot & A_{25} \\
\cdot & \cdot & \cdot & \cdot \\
\cdot & \cdot & \cdot & \cdot \\
A_{51} & A_{52} & \cdots & A_{55}
\end{vmatrix}\, .
\]
%CONT>>


% ===== pdf idx 108, printed 94.0 =====
% !PRINTED 94
\weberpage{94}
Erster Theil. Theorie der Determinanten.

%<<CONT
die Determinante $\Delta$ und ihre adjungirte $D$; der Werth des Minors

\[
D^{(3)} =
\begin{vmatrix}
A_{31} & A_{32} \\
A_{41} & A_{42}
\end{vmatrix}
\]

bestimmt sich dann aus dem Producte

\begin{align*}
(-1)^{\nu}\, \Delta \cdot D^{(3)} &=
\begin{vmatrix}
a_{31} & a_{32} & a_{33} & a_{34} & a_{35} \\
a_{41} & a_{42} & a_{43} & a_{44} & a_{45} \\
a_{11} & a_{12} & a_{13} & a_{14} & a_{15} \\
a_{21} & a_{22} & a_{23} & a_{24} & a_{25} \\
a_{51} & a_{52} & a_{53} & a_{54} & a_{55}
\end{vmatrix}
\cdot
\begin{vmatrix}
A_{31} & A_{32} & A_{33} & A_{34} & A_{35} \\
A_{41} & A_{42} & A_{43} & A_{44} & A_{45} \\
0 & 0 & 1 & 0 & 0 \\
0 & 0 & 0 & 1 & 0 \\
0 & 0 & 0 & 0 & 1
\end{vmatrix} \\
&=
\begin{vmatrix}
\Delta & 0 & a_{33} & a_{34} & a_{35} \\
0 & \Delta & a_{43} & a_{44} & a_{45} \\
0 & 0 & a_{13} & a_{14} & a_{15} \\
0 & 0 & a_{23} & a_{24} & a_{25} \\
0 & 0 & a_{53} & a_{54} & a_{55}
\end{vmatrix}
=
\Delta^2
\begin{vmatrix}
a_{13} & a_{14} & a_{15} \\
a_{23} & a_{24} & a_{25} \\
a_{53} & a_{54} & a_{55}
\end{vmatrix}\, ,
\end{align*}

\[
\nu = (34) + (12) + (3+4) + (1+2) = 10. % ?? uncertain: "12" printed with a damaged first digit resembling a dagger mark
\]

Daher: $D^{(3)} = \Delta \cdot A_{\substack{34\\12}}$, oder:

\[
\begin{vmatrix}
A_{31} & A_{32} \\
A_{41} & A_{42}
\end{vmatrix}
=
\begin{vmatrix}
a_{11} & a_{12} & \cdots & a_{15} \\
a_{21} & a_{22} & \cdots & a_{25} \\
a_{31} & \cdot & \cdot & \cdot \\
\cdot & \cdot & \cdot & \cdot \\
a_{51} & \cdot & \cdot & a_{55}
\end{vmatrix}
\cdot
\begin{vmatrix}
a_{13} & a_{14} & a_{15} \\
a_{23} & a_{24} & a_{25} \\
a_{53} & a_{54} & a_{55}
\end{vmatrix}\, .
\]

90. \textit{Definition der correspondirenden Matrices.} Ehe wir zu den Anwendungen der Determinantentheorie übergehen, wollen wir noch einen Satz beweisen, der die Determinanten zweier solcher Matrices in Beziehung setzt, zwischen deren Elementen gewisse Relationen bestehen. Seien die beiden Matrices, deren eine $n$ Colonnen und $\mu$ Zeilen, deren andere $n$ Colonnen und $\nu$ Zeilen besitzen, dargestellt durch

\[
\begin{vmatrix}
a_{11} & a_{12} & a_{13} & \cdots & a_{1n} \\
a_{21} & a_{22} & \cdot & \cdot & a_{2n} \\
\cdot & \cdot & \cdot & \cdot & \cdot \\
a_{\mu 1} & a_{\mu 2} & \cdot & \cdot & a_{\mu n}
\end{vmatrix}
\quad\text{und}\quad
\begin{vmatrix}
b_{11} & b_{12} & b_{13} & \cdots & b_{1n} \\
b_{21} & b_{22} & \cdot & \cdot & b_{2n} \\
\cdot & \cdot & \cdot & \cdot & \cdot \\
b_{\nu 1} & b_{\nu 2} & \cdot & \cdot & b_{\nu n}
\end{vmatrix}\, ,
\]

wobei $\mu + \nu = n$ vorausgesetzt ist, dann können wir aus ihnen nach der Methode der Multiplication für Determinantenmatrices eine neue Matrix bilden, indem wir die Elemente jeder Zeile der ersten mit dem ent-
%CONT>>


% ===== pdf idx 109, printed 95.0 =====
% !PRINTED 95
\weberpage{95}
\S.\,8. Adjungirte Determinanten und correspondirende Matrices.

%<<CONT
entsprechenden jeder Zeile der zweiten multipliciren. Diese neue Matrix enthält in $\mu$ Colonnen und $\nu$ Zeilen $\mu\cdot\nu$ Elemente von der Form:

\weq{1}{a_{\varrho 1}\, b_{\sigma 1} + a_{\varrho 2}\, b_{\sigma 2} + a_{\varrho 3}\, b_{\sigma 3} + \cdots + a_{\varrho n}\, b_{\sigma n}}

\[
\varrho = 1, 2, \cdots \mu; \quad \sigma = 1, 2, \cdots \nu.
\]

Wir setzen nun im Folgenden fest, dass die Elemente $a_{ik}$, $b_{ik}$ solche Werthe besitzen, dass jede Summe, wie sie in (1) dargestellt ist, den Werth Null habe. Die Zahl der Relationen, die hierdurch zwischen den Elementen $a_{ik}$ und $b_{ik}$ bestehen soll, ist $\mu\cdot\nu$. Nun können wir aber die beiden gegebenen Matrices zu einer einzigen Matrix zusammenfügen, welche, da $\mu+\nu=n$, eine Determinante $\Delta$ repräsentirt. Jeder ihrer Unterdeterminanten von der Ordnung $\mu=n-\nu$, die nur Elemente $a_{ik}$ enthält, entspricht eine correspondirende Unterdeterminante $\nu^{\mathrm{ter}}$ Ordnung, die nur Elemente $b_{ik}$ besitzt. Wir werden demgemäss je eine Determinante der ersten und je eine der zweiten Matrix correspondirende (complementäre) nennen, wenn die Determinante der Elemente $b_{ik}$ durch Colonnen gebildet ist mit solchen Indices, die bei Bildung der ersten Determinante der Elemente $a_{ik}$ nicht verwendet werden. Beide Matrices werden überdies selbst als \textit{correspondirende} bezeichnet, wenn zwischen ihren Elementen die $\mu\cdot\nu$ Relationen bestehen:

\[
a_{\varrho 1}\, b_{\sigma 1} + a_{\varrho 2}\, b_{\sigma 2} + \cdots a_{\varrho n}\, b_{\sigma n} = 0.
\]

91. \textit{Lehrsatz über correspondirende Matrices.} Der oben erwähnte Satz lautet nun:

\glqq Die correspondirenden Determinanten zweier correspondirender Matrices sind einander proportional.\grqq{}

Wir geben zunächst von dem Satze einen ersten Beweis; ein zweiter Beweis wird sich später (vgl. Nr. 106) darbieten.

Die Matrix der Elemente $a_{ik}$ ergänzen wir durch Anfügung von $\nu$ Zeilen willkürlicher Elemente $a_{ik}$ zur Matrix einer Determinante $\Delta$, ebenso erweitern wir die zweite Matrix durch $\mu$ Zeilen willkürlicher Elemente $b_{ik}$ zur Matrix der Determinante $D$, so dass wir erhalten:

\[
\Delta =
\begin{vmatrix}
a_{11} & a_{12} & a_{13} & \cdots & a_{1n} \\
a_{21} & a_{22} & \cdot & \cdot & a_{2n} \\
\cdot & \cdot & \cdot & \cdot & \cdot \\
a_{\mu 1} & a_{\mu 2} & \cdot & \cdot & a_{\mu n} \\
a_{\mu+1,1} & a_{\mu+1,2} & \cdot & \cdot & a_{\mu+1,n} \\
\cdot & \cdot & \cdot & \cdot & \cdot \\
a_{n1} & \cdot & \cdot & \cdot & a_{nn}
\end{vmatrix}
\quad\text{und } D=
\begin{vmatrix}
b_{11} & b_{12} & b_{13} & \cdots & b_{1n} \\
b_{21} & b_{22} & \cdot & \cdot & b_{2n} \\
\cdot & \cdot & \cdot & \cdot & \cdot \\
b_{\nu 1} & b_{\nu 2} & \cdot & \cdot & b_{\nu n} \\
b_{\nu+1,1} & b_{\nu+1,2} & \cdot & \cdot & b_{\nu+1,n} \\
\cdot & \cdot & \cdot & \cdot & \cdot \\
b_{n1} & \cdot & \cdot & \cdot & b_{nn}
\end{vmatrix}\, .
\]


% ===== pdf idx 110, printed 96.0 =====
% !PRINTED 96
\weberpage{96}
Erster Theil. Theorie der Determinanten.

Das Product beider Determinanten — wir wollen es durch Multiplication von Zeilen mit Zeilen entstehen lassen — enthält Elemente von der Form

\[
c_{\varrho\sigma} = a_{\varrho 1}\, b_{\sigma 1} + a_{\varrho 2}\, b_{\sigma 2} \ldots + a_{\varrho n}\, b_{\sigma n}\, .
\]

Diese sind null, wenn sowohl $\varrho\leqq\mu$ als auch $\sigma\leqq\nu$ ist; die Matrix des Productes enthält somit als einen Theil die schon Eingangs erwähnte Matrix von $\mu\nu$ Elementen, welche verschwinden. Sie ist nämlich dargestellt durch:

\[
\Delta\cdot D =
\begin{vmatrix}
0 & 0 & 0 & \cdots & 0, & c_{1,\mu+1} & c_{1,\mu+2} & \cdots & c_{1n} \\
0 & 0 & \cdot & \cdots & 0, & c_{2,\mu+1} & c_{2,\mu+2} & \cdots & c_{2n} \\
0 & \cdot & \cdot & \cdot & \cdot & c_{3,\mu+1} & \cdot & \cdot & \cdot \\
\cdot & \cdot & \cdot & \cdot & \cdot & \cdot & \cdot & \cdot & \cdot \\
0 & 0 & 0 & \cdots & 0, & c_{\nu,\mu+1} & \cdot & \cdot & c_{\nu n} \\
c_{\nu+1,1} & c_{\nu+1,2} & \cdots & c_{\nu+1,\mu}, & c_{\nu+1,\mu+1} & \cdot & \cdot & \cdot & c_{\nu+1,n} \\
c_{\nu+2,1} & c_{\nu+2,2} & \cdot & \cdot & \cdot & \cdot & \cdot & \cdot & c_{\nu+2,n} \\
c_{\nu+3,1} & \cdot & \cdot & \cdot & \cdot & \cdot & \cdot & \cdot & \cdot \\
\cdot & \cdot & \cdot & \cdot & \cdot & \cdot & \cdot & \cdot & \cdot \\
c_{n1} & c_{n2} & \cdots & c_{n\mu}, & c_{n,\mu+1} & \cdot & \cdot & \cdot & c_{nn}
\end{vmatrix}\, .
\]

92. Daher zerfällt sie nach dem Satze von Laplace in das Product der beiden Determinanten

\[
B =
\begin{vmatrix}
c_{1,\mu+1}, & c_{1,\mu+2} & \cdots & c_{1n} \\
c_{2,\mu+1}, & c_{2,\mu+2} & \cdots & c_{2n} \\
c_{3,\mu+1} & \cdot & \cdot & \cdot \\
\cdot & \cdot & \cdot & \cdot \\
c_{\nu,\mu+1}, & c_{\nu,\mu+2} & \cdots & c_{\nu n}
\end{vmatrix}
\]

und

\[
A =
\begin{vmatrix}
c_{\nu+1,1}, & c_{\nu+1,2} & \cdots & c_{\nu+1,\mu} \\
c_{\nu+2,1}, & c_{\nu+2,2} & \cdots & c_{\nu+2,\mu} \\
c_{\nu+3,1}, & c_{\nu+3,2} & \cdots & c_{\nu+3,\mu} \\
\cdot & \cdot & \cdot & \cdot \\
c_{n1}, & c_{n2} & \cdots & c_{n\mu}
\end{vmatrix}\, .
\]

Die Determinante $A$ besitzt $\mu^2$ Elemente

\[
c_{\varrho\sigma} = a_{\varrho 1}\, b_{\nu+\sigma,1} + a_{\varrho 2}\, b_{\nu+\sigma,2} \cdots + a_{\varrho n}\, b_{\nu+\sigma,n}
\begin{cases}
\varrho = 1 \text{ bis } \mu \\
\nu + \sigma = \nu \text{ bis } n.
\end{cases}
\]

Die Determinante $B$ enthält $\nu^2$ Elemente

\[
c_{\varrho\sigma} = a_{\mu+\varrho,1}\, b_{\sigma 1} + a_{\mu+\varrho,2}\, b_{\sigma 2} \ldots + a_{\mu+\varrho,n}\, b_{\sigma n}
\begin{cases}
\mu + \varrho = \mu \text{ bis } n \\
\sigma = 1 \text{ bis } \nu.
\end{cases}
\]


% ===== pdf idx 111, printed 97.0 =====
% !PRINTED 97
\weberpage{97}
Wir erhalten somit als Product von $D \cdot \varDelta$ den Werth

\weq{2}{\varDelta \cdot D = (-1)^{\bar v} A \cdot B.}

Der Exponent $\bar v$ bestimmt sich dabei nach Nr. 52 aus der Relation

\[
\bar v = \sum_{\varrho=1}^{\varrho=v}(\mu+\varrho) + \sum_{1}^{v}\varrho = v\cdot\mu + v(v+1),
\]

oder auch aus

\[
\bar v = \sum_{1}^{\mu}\varrho + \sum_{\varrho=1}^{\varrho=\mu}(v+\varrho) = \mu(\mu+1) + v\mu.
\]

Da sowohl $v(v+1)$ als $\mu(\mu+1)$ eine gerade Zahl, so kann der Exponent $\bar v = \mu\cdot v$ gesetzt werden. Die Relation (2) wird also:

\weq{3}{\varDelta \cdot D = (-1)^{\mu\cdot v} A \cdot B}

oder

\weq{4}{\frac{\varDelta}{B} = (-1)^{\mu\cdot v}\,\frac{A}{D}\,.}

93. Nun enth\"alt aber die linke Seite dieser Gleichung weder im Nenner noch im Z\"ahler die willk\"urlich sp\"ater zugef\"ugten Elemente $b_{v+\sigma,i}$. Setzen wir also diesen Bruch gleich $\lambda$, so ist $\lambda$ unabh\"angig von diesen Elementen; die Gleichung (4) aber wird:

\weq{5}{\lambda = (-1)^{\mu\cdot v}\cdot\frac{A}{D}\,,\ \text{oder}\ (-1)^{\mu\cdot v} A = \lambda\cdot D,}

worin nun $A$ und $D$ Functionen der willk\"urlichen Elemente $b_{v+\sigma,i}$ sind, \"uber die wir nun passend verf\"ugen k\"onnen. W\"ahlen wir n\"amlich die Elemente $b_{v+1,i_1}$, $b_{v+2,i_2}\ \ldots b_{v+\mu,i_\mu}$ gleich 1, alle \"ubrigen aber gleich null, dann reducirt sich durch Unterdr\"uckung der Zeilen und Colonnen, welche das Element 1 enthalten, die Determinante

\[
D =
\begin{vmatrix}
b_{11} & b_{12} & b_{13} & \cdots & b_{1n} \\
b_{21} & b_{22} & \cdot & \cdot & b_{2n} \\
b_{31} & \cdot & \cdot & \cdot & \cdot \\
\cdot & \cdot & \cdot & \cdot & \cdot \\
b_{n1} & b_{n2} & b_{n3} & \cdots & b_{nn}
\end{vmatrix}
\]

auf eine Determinante von $v^2$ Elementen der zweiten correspondirenden Matrix, n\"amlich auf:

\[
D' = B_{\substack{v+1,\,v+2,\,\ldots\,v+\mu\\i_1\ i_2\ \ldots i_\mu}} = (-1)^{v_1}
\begin{vmatrix}
b_{1i_{\mu+1}}, & b_{1i_{\mu+2}} & \cdots & b_{1i_{\mu+v}} \\
b_{2i_{\mu+1}}, & b_{2i_{\mu+2}} & \cdots & b_{2i_{\mu+v}} \\
\cdot & \cdot & \cdots & \cdot \\
b_{vi_{\mu+1}}, & b_{vi_{\mu+2}} & \cdots & b_{vi_{\mu+v}}
\end{vmatrix},
\]

wobei:
%CONT>>


% ===== pdf idx 112, printed 98.0 =====
% !PRINTED 98
\weberpage{98}
%<<CONT
\[
v_1 = (v+1,\ v+2,\ \ldots v+\mu) + (i_1 i_2 \ldots i_\mu) + \sum_{\sigma=1}^{\sigma=\mu}(v+\sigma) + \sum_{\sigma=1}^{\sigma=\mu} i_\sigma\,,
\]

oder, weil die Indices $v+1$, $v+2$, $\ldots v+\mu$ a priori gut geordnet sind, und die Reihenfolge der Indices $i_1 i_2 \ldots i_\mu$ stets als gut geordnet vorausgesetzt werden kann:

\[
v_1 = \sum_{\sigma=1}^{\sigma=\mu}(v+\sigma) + \sum_{\sigma=1}^{\sigma=\mu} i_\sigma = \mu\cdot v + \frac{\mu(\mu+1)}{2} + \sum_{\sigma=1}^{\sigma=\mu} i_\sigma = \mu v + k.
\]

wobei

\[
k = \frac{\mu(\mu+1)}{2} + \sum_{1}^{\mu} i_\sigma\ \text{ist.}
\]

94. Nun ist anderntheils die Determinante $A$ dargestellt durch die Matrix:

\[
A =
\begin{vmatrix}
a_{11}b_{v+1,1}+\cdots a_{1n}b_{v+1,n}, & a_{11}b_{v+2,1}+\cdots a_{1n}b_{v+2,n}, & \cdots, & a_{11}b_{n1}+\cdots a_{1n}b_{nn} \\
a_{21}b_{v+1,1}+\cdots a_{2n}b_{v+1,n}, & a_{21}b_{v+2,1}+\cdots a_{2n}b_{v+2,n}, & \cdots, & a_{21}b_{n1}+\cdots a_{2n}b_{nn} \\
a_{31}b_{v+1,1}+\cdots a_{3n}b_{v+1,n}, & a_{32}b_{v+2,1}+\cdots a_{3n}b_{v+2,n}, & \cdots, & a_{31}b_{n1}+\cdots a_{3n}b_{nn} \\% [sic]
\cdot & \cdot & \cdots & \cdot \\
a_{\mu 1}b_{v+1,1}\cdots a_{\mu n}b_{v+1,n}, & a_{\mu 1}b_{v+2,1}\cdots a_{\mu n}b_{v+2,n}, & \cdots, & a_{\mu 1}b_{n1}+\cdots a_{\mu n}b_{nn}
\end{vmatrix}.
\]

Setzen wir in ihr alle Elemente $b_{v+\sigma,\varrho}$, ausgenommen die oben erw\"ahnten, gleich null, so reducirt sie sich auf

\[
A' =
\begin{vmatrix}
a_{1i_1}\cdot b_{v+1,i_1}, & a_{1i_2}\cdot b_{v+2,i_2}, & a_{1i_3}\cdot b_{v+3,i_3} & \cdots a_{1i_\mu}\cdot b_{v+\mu,i_\mu} \\
a_{2i_1}\cdot b_{v+1,i_1}, & a_{2i_2}\cdot b_{v+2,i_2}, & \cdot & \cdots a_{2i_\mu}\cdot b_{v+\mu,i_\mu} \\
a_{3i_1}\cdot b_{v+1,i_1}, & a_{3i_2}\cdot b_{v+2,i_2} & \cdot & \cdots a_{3i_\mu}\cdot b_{v+\mu,i_\mu} \\
\cdot & \cdot & \cdot & \cdots \\
a_{\mu i_1}\cdot\cdot b_{v+1,i_1}, & a_{\mu i_2}\cdot b_{v+2,i_2} & \cdot & \cdots a_{\mu i_\mu}\cdot b_{v+\mu,i_\mu}
\end{vmatrix},
\]

und weil die darin noch auftretenden Factoren $b$ durchgehends gleich 1 sind, so wird endlich

\[
A' =
\begin{vmatrix}
a_{1i_1}, & a_{1i_2}, & a_{1i_3} & \cdots a_{1i_\mu} \\
a_{2i_1}, & a_{2i_2} & \cdot & \cdots a_{2i_\mu} \\
\cdot & \cdot & \cdot & \cdots \\
a_{\mu i_1}, & a_{\mu i_2} & \cdot & \cdots a_{\mu i_\mu}
\end{vmatrix},
\]

d. h. $A$ reducirt sich auf eine Determinante der ersten correspondirenden Matrix, und zwar auf diejenige, deren Colonnen Indices f\"uhren, welche gerade bei Bildung der Determinante aus der zweiten Matrix nicht
%CONT>>


% ===== pdf idx 113, printed 99.0 =====
% !PRINTED 99
\weberpage{99}
%<<CONT
benutzt werden. Tragen wir also die so berechneten Werthe f\"ur $D$ und $A$ in die Gleichung (5) ein, so erhalten wir:

\weq{6}{(-1)^{\mu v}
\begin{vmatrix}
a_{1i_1} & a_{1i_2} & a_{1i_3} & \cdots & a_{1i_\mu} \\
a_{2i_1} & a_{2i_2} & a_{2i_3} & \cdots & a_{2i_\mu} \\
a_{3i_1} & \cdot & \cdot & \cdot & \cdot \\
\vdots & \cdot & \cdot & \cdot & \cdot \\
a_{\mu i_1} & a_{\mu i_2} & \cdot & \cdot & a_{\mu i_\mu}
\end{vmatrix}
= (-1)^{\mu v + k}\cdot\lambda
\begin{vmatrix}
b_{1 i_{\mu+1}} & b_{1 i_{\mu+2}} & b_{1 i_{\mu+3}} & \cdots & b_{1 i_{\mu+v}} \\
b_{2 i_{\mu+1}} & b_{2 i_{\mu+2}} & \cdot & \cdot & b_{2 i_{\mu+v}} \\
b_{3 i_{\mu+1}} & \cdot & \cdot & \cdot & \cdot \\
\vdots & \cdot & \cdot & \cdot & \cdot \\
b_{v i_{\mu+1}} & b_{v i_{\mu+2}} & \cdot & \cdot & b_{v i_{\mu+v}}
\end{vmatrix},}

oder nach Division mit dem Factor $(-1)^{\mu v}$:

\weq{7}{(a_{1i_1}\ a_{2i_2} \cdots a_{\mu i_\mu}) = (-1)^k\cdot\lambda\cdot(b_{1i_{\mu+1}}\ b_{2i_{\mu+2}} \cdots b_{vi_{\mu+v}}),}

wobei

\[
k = \frac{\mu(\mu+1)}{2} + \sum_{1}^{\mu} i_\sigma\,,
\]

das ist gleich der Summe aller Doppelindices des Diagonalgliedes in $(a_{1i_1}\ a_{2i_2}\ \ldots a_{\mu i_\mu})$. Dadurch ist der aufgestellte Satz bewiesen; der Proportionalit\"atsfactor $\lambda$ ist nat\"urlich unbestimmt.

95. \textit{Beispiel.} Seien gegeben die beiden Matrices, f\"ur welche $\mu = v = 2$ ist,

\[
\begin{vmatrix}
x_1 & x_2 & x_3 & x_4 \\
y_1 & y_2 & y_3 & y_4
\end{vmatrix}
\ \text{und}\
\begin{vmatrix}
u_1 & u_2 & u_3 & u_4 \\
v_1 & v_2 & v_3 & v_4
\end{vmatrix}.
\]

Zwischen den Elementen derselben sollen die vier Relationen

\[
u_x = 0\,,\quad u_y = 0\,,\quad v_x = 0\,,\quad v_y = 0
\]

bestehen, wobei

\[
u_x = u_1 x_1 + u_2 x_2 + u_3 x_3 + u_4 x_4
\]

zu setzen ist, und

\[
u_y,\ v_x\ \text{und}\ v_y,
\]

Symbole analoger Ausdr\"ucke sind.

Dann ist nach dem vorigen Satze:

\[
\begin{vmatrix}
x_1 & x_2 \\
y_1 & y_2
\end{vmatrix}
= +\lambda
\begin{vmatrix}
u_3 & u_4 \\
v_3 & v_4
\end{vmatrix},\quad
\begin{vmatrix}
x_1 & x_3 \\
y_1 & y_3
\end{vmatrix}
= -\lambda
\begin{vmatrix}
u_2 & u_4 \\
v_2 & v_4
\end{vmatrix},\ \text{u.\,s.\,w.,}
\]


% ===== pdf idx 114, printed 100.0 =====
% !PRINTED 100
\weberpage{100}
oder wenn wir der K\"urze halber setzen

\[
p_{ik} = (x_i y_k - y_i x_k)\,,
\]

und

\[
q_{ik} = (u_i v_k - v_i u_k)\,,
\]

so bestehen die sechs Relationen:

\[
1)\quad p_{12} = \lambda q_{34}\,,\quad p_{13} = \lambda q_{42}\,,\quad p_{14} = \lambda q_{23}\,,
\]
\[
2)\quad p_{23} = \lambda q_{14}\,,\quad p_{42} = \lambda q_{13}\,,\quad p_{34} = \lambda q_{12}\,.
\]

Man kann die Gr\"ossen $p_{ik}$, resp. $q_{ik}$ als Coordinaten der geraden Linie auffassen, welche als Schnitt der beiden Ebenen

\[
u_x = 0\,,\quad v_x = 0
\]

auftritt, und erkennt somit den Satz: Die Strahlencoordinaten $p_{ik}$ sind proportional den Axencoordinaten $q_{ik}$.


% ===== pdf idx 115, printed 100.001 =====
% !PRINTED none
\gesperrt{Zweiter Theil.}

\gesperrt{Anwendungen der Determinantentheorie.}

§ 8. \gesperrt{Systeme linearer Gleichungen und deren Auflösung.}

96. \textit{Auflösung eines Systemes von $n$ linearen homogenen Gleichungen
mit $n$ Unbekannten.} Die Aufgabe, ein lineares System von Gleichungen
mit mehreren Unbekannten aufzulösen, war es zunächst, welche Leibnitz
und Cramer auf die Determinantenrechnung führten. Betrachtet man
nämlich das folgende System linearer homogener Gleichungen

\weq{I}{\left\{\begin{array}{l}
a_{11}x_1+a_{12}x_2+a_{13}x_3\cdots+a_{1n}x_n=0\\
a_{21}x_1+a_{22}x_2+a_{23}x_3\cdots+a_{2n}x_n=0\\
a_{31}x_1+a_{32}x_2+a_{33}x_3\cdots+a_{3n}x_n=0\\
\cdot\ \cdot\ \cdot\ \cdot\ \cdot\ \cdot\ \cdot\ \cdot\ \cdot\ \cdot\ \cdot\ \cdot\\
a_{n1}x_1+a_{n2}x_2+a_{n3}x_3\cdots+a_{nn}x_n=0,
\end{array}\right.}

so erkennt man, dass in demselben gerade $n^2$ willkürliche Grössen $a_{ik}$
auftreten, die wir zur Matrix einer Determinante zusammensetzen können.
Wir wollen sie die Determinante $\Delta$ des simultanen Systemes (I)
nennen. Fragen wir uns nun nach denjenigen Werthen der Unbekannten
$x_1x_2\cdots x_n$, welche dieses System befriedigen, so sehen wir zunächst
unmittelbar, dass durch das System von Lösungen

\[
x_1=x_2=x_3= \text{u.\,s.\,w.} =x_n=0,
\]

das System der Gleichungen (I) in der That erfüllt wird. Wir nennen diese
Lösung eine identische, und sehen von ihr, da sie kein weiteres Interesse
erweckt, in den folgenden Betrachtungen ab.

97. \textit{Aufsuchung anderweitiger Lösungen.} Indem wir uns nun die
Aufgabe stellen, andere Lösungen des Systemes (I) zu finden, können wir
auf irgend einem elementaren Wege der gewöhnlichen Eliminationsmethoden
vorgehen. Am besten nach der Methode von Bézout, indem wir die
Gleichungen des Systemes (I) mit solchen Factoren $\lambda_1\lambda_2\cdots\lambda_n$
%CONT>>


% ===== pdf idx 116, printed 102.0 =====
% !PRINTED 102
\weberpage{102}
%<<CONT
multipliciren, dass nach Addition der so erhaltenen Producte die
Coefficienten aller Unbekannten, mit Ausnahme desjenigen von $x_\mu$,
verschwinden. Solche Factoren sind gerade die Unterdeterminanten

\[
A_{1\mu},\quad A_{2\mu},\quad A_{3\mu}\cdots A_{n\mu}
\]

der Elemente der $\mu^{\text{ten}}$ Colonne in der Determinante

\[
\Delta=
\left|
\begin{array}{cccccc}
a_{11} & a_{12} & \cdots & a_{1\mu} & \cdots & a_{1n}\\
a_{21} & a_{22} & \cdots & a_{2\mu} & \cdots & a_{2n}\\
a_{31} & \cdot & \cdot & \cdot & \cdot & a_{3n}\\
\cdot & \cdot & \cdot & \cdot & \cdot & \cdot\\
a_{n1} & a_{n2} & \cdots & a_{n\mu} & \cdots & a_{nn}
\end{array}
\right|.
\]

Multipliciren wir also die Gleichungen des Systemes (I) der Reihe nach
mit den erwähnten Unterdeterminanten und addiren, so erhalten wir:

\[
\begin{aligned}
&x_1(a_{11}A_{1\mu}+a_{21}A_{2\mu}\cdots a_{n1}A_{n\mu})+x_2(a_{12}A_{1\mu}+a_{22}A_{2\mu}+\cdots a_{n2}A_{n\mu})\\
&\qquad+\cdots+x_\mu(a_{1\mu}A_{1\mu}+a_{2\mu}A_{2\mu}+\cdots a_{n\mu}A_{n\mu})+\cdots\\
&\qquad+x_n(a_{1n}A_{1\mu}+a_{2n}A_{2\mu}+\cdots a_{nn}A_{n\mu})=0.
\end{aligned}
\]

Gemäss den Relationen $(C)$ in Nr.\,32 reducirt sich diese Gleichung auf

\weq{II}{x_\mu\cdot\Delta = x_\mu(a_{1\mu}A_{1\mu}+a_{2\mu}A_{2\mu}+\cdots a_{n\mu}A_{n\mu})=0.}

Ist hier die Determinante $\Delta$ von Null verschieden, so wird diese
Gleichung nur befriedigt für $x_\mu=0$, d.\,h. die Gleichungen des
Systemes (I) haben, wenn $\Delta \gtrless 0$, keine andere Lösung als
die identische. Verschwindet dagegen die Determinante $\Delta$, so kann
dem Systeme (I) ein System von Werthen der Unbekannten genügen, und
die Aufgabe ist, unter dieser Bedingung dasselbe zu ermitteln.

98. \textit{Resultante des Systemes} (I). Man nennt die Determinante
$\Delta$ die Resultante des Systemes der linearen Gleichungen und aus
dem eben Erwähnten folgt eine erste Eigenschaft derselben, dass sie
nämlich verschwinden muss, sobald die Gleichungen (I) eine
gemeinschaftliche Lösung haben. Umgekehrt werden wir sogleich auch
eine zweite Eigenschaft derselben erkennen, dass, wenn sie verschwindet,
ein gemeinschaftliches Lösungssystem wirklich vorhanden ist.

\textit{Anmerkung.} Doch ist die Resultante durch diese beiden
Eigenschaften nicht eindeutig definirt. Denn einerseits erkennt man aus
Gleichung (II), dass wenn $x_\mu \gtrless 0$, d.\,h. wenn eine Lösung
vorhanden, nicht nur $\Delta$, sondern auch $\Delta\cdot\varphi$
verschwindet, wo $\varphi$ eine beliebige ganze Function der
Coefficienten $a_{ik}$, so dass man also nicht be-
%CONT>>


% ===== pdf idx 117, printed 103.0 =====
% !PRINTED 103
\weberpage{103}
%<<CONT
haupten kann: Jene ganze Function der Coefficienten, die beim
Vorhandensein einer gemeinschaftlichen Lösung des Systemes (I)
verschwindet, ist die Resultante des Systemes. Andrerseits kann man
zwar nicht behaupten, dass mit $\Delta\cdot\varphi=0$ ein
Lösungssystem existirt, so lange $\varphi$ eine Function der
Coefficienten $a_{ik}$, die von $\Delta$ selbst verschieden ist. Wenn
aber

\[
\varphi=c, \quad \text{oder} \quad \varphi=c\cdot\Delta^{\varrho},
\]

d.\,h. also eine von Null verschiedene Constante oder eine Potenz von
$\Delta$ multiplicirt mit einer solchen Constanten, dann haben die
Gleichungen (I) auch eine gemeinschaftliche Lösung, wenn

\[
c\cdot\Delta^{\varrho}=0.
\]

Mit andern Worten: Beide Eigenschaften selbst zusammengenommen
definiren nicht bloss $\Delta$, sondern $c\cdot\Delta$, resp.
$c\cdot\Delta^{\varrho}$. Es ist aber $c\cdot\Delta^{\varrho}$
Resultante jedes Gleichungssystemes, das aus (I) hervorgeht, wenn wir
die Coefficienten dieses Systemes irgendwie linear transformiren. Eine
also definirte Resultante ist daher keine bestimmte lineare Function
der Coefficienten speciell dieser Gleichungen.

99. \textit{Die gemeinschaftlichen Lösungen des Systemes (I), wenn die
Resultante verschwindet.} Wir nehmen nun im Folgenden an, dass

\[
\Delta=
\left|
\begin{array}{cccc}
a_{11} & a_{12} & a_{13} & \cdots a_{1n}\\
a_{21} & a_{22} & \cdot & \cdots a_{2n}\\
a_{31} & \cdot & \cdot & \cdots\cdots\\
\cdot & \cdot & \cdot & \cdot\\
a_{n1} & \cdot & \cdots & a_{nn}
\end{array}
\right|=0,
\]

und überdies der Allgemeinheit halber, dass sämmtliche Minoren dieser
%CONT>>


% ===== pdf idx 118, printed 104.0 =====
% !PRINTED 104
\weberpage{104}
% [column spec widened array 7->8] a row carried more cells than the spec declared,
% which is a fatal alignment error. The spec was widened rather than a cell
% deleted, so no transcribed content is lost; whether Weber's grid really has
% that many columns is a question for a pass with the scan open.
%<<CONT
und überdies der Allgemeinheit halber, dass sämmtliche Minoren dieser
Determinante von der ersten bis zur $(\mu-1)^{\text{ten}}$ Ordnung
ebenfalls verschwinden. Dagegen soll von den Minoren der
$\mu^{\text{ten}}$ Ordnung wenigstens einer von Null verschieden sein.

Dieser nicht verschwindende Minor $\mu^{\text{ter}}$ Ordnung sei etwa

\[
D^{(\mu+1)}=
\left|
\begin{array}{cccc}
a_{\mu+1,\mu+1} & a_{\mu+1,\mu+2} & \cdots & a_{\mu+1,n}\\
a_{\mu+2,\mu+1} & a_{\mu+2,\mu+2} & \cdots & a_{\mu+2,n}\\
\cdot & \cdot & \cdot & \cdot\\
a_{n,\mu+1}, & a_{n,\mu+2} & \cdots & a_{nn}
\end{array}
\right|.
\]

Wir betrachten dann zunächst das System der $(n-\mu+1)$ letzten
Gleichungen des Systemes (I) in Nr.\,96

\weq{III}{\left\{\begin{array}{l}
a_{\mu1}x_1+a_{\mu2}x_2+\cdots a_{\mu,\mu}x_\mu+a_{\mu,\mu+1}x_{\mu+1}\cdots\cdots+a_{\mu n}x_n=0\\
a_{\mu+1,1}x_1+a_{\mu+1,2}x_2+\cdots a_{\mu+1,\mu}x_\mu+a_{\mu+1,\mu+1}x_{\mu+1}\cdots+a_{\mu+1,n}x_n=0\\
a_{\mu+2,1}x_1+a_{\mu+2,2}x_2+\cdots a_{\mu+2,\mu}x_\mu+a_{\mu+2,\mu+1}x_{\mu+1}\cdots+a_{\mu+2,n}x_n=0\\
\cdot\ \cdot\ \cdot\ \cdot\ \cdot\ \cdot\ \cdot\ \cdot\ \cdot\ \cdot\ \cdot\ \cdot\ \cdot\ \cdot\ \cdot\ \cdot\\
a_{n1}x_1+a_{n2}x_2+\cdots a_{n\mu}x_\mu+a_{n\mu+1}x_{\mu+1}\cdots+a_{nn}x_n=0
\end{array}\right.}

Die Matrix der Coefficienten dieses Systemes ist dargestellt durch

\weq{A}{
\left|
\begin{array}{cccccccc}
a_{\mu1} & a_{\mu2} & a_{\mu3} & \cdots & a_{\mu\mu} & a_{\mu\mu+1} & \cdots & a_{\mu n}\\
a_{\mu+1,1} & a_{\mu+1,2} & \cdots & \cdots & a_{\mu+1,\mu} & a_{\mu+1,\mu+1} & \cdots & a_{\mu+1,n}\\
a_{\mu+2,1} & \cdots & \cdots & \cdots & a_{\mu+2,\mu} & a_{\mu+2,\mu+1} & \cdots & a_{\mu+2,n}\\
\cdot & \cdot & \cdot & \cdot & \cdot & \cdot & \cdot & \cdot\\
a_{n1} & a_{n2} & a_{n3} & \cdots & a_{n\mu} & a_{n,\mu+1} & \cdots & a_{nn}
\end{array}
\right|
}

und die sämmtlichen Determinanten dieser Matrix sind nach Voraussetzung
identisch null, unter ihnen insbesondere die Determinante

\[
D^{(\mu)}=
\left|
\begin{array}{cccc}
a_{\mu\mu} & a_{\mu,\mu+1} & \cdots & a_{\mu n}\\
a_{\mu+1,\mu} & a_{\mu+1,\mu+1} & \cdots & a_{\mu+1,n}\\
a_{\mu+2,\mu} & a_{\mu+2,\mu+1} & \cdots & a_{\mu+2,n}\\
\cdot & \cdot & \cdot & \cdot\\
a_{n\mu} & a_{n,\mu+1} & \cdots & a_{nn}
\end{array}
\right|=0.
\]

Von ihren Unterdeterminanten ist wie oben vorausgesetzt, die zum
Elemente $a_{\mu\mu}$ gehörige

\[
A_{\mu\mu}=D^{(\mu+1)}
\]

von null verschieden, während alle übrigen, also auch die den Elementen
der ersten Colonne entsprechenden

\[
A_{\mu+1,\mu},\quad A_{\mu+2,\mu},\quad A_{\mu+3,\mu}\cdots A_{n\mu}
\]

beliebige Werthe, null nicht ausgeschlossen, besitzen können.

100. \textit{Erste Folgerung aus diesen Voraussetzungen: mindestens
eine Gleichung ist eine Folge der übrigen.} Multipliciren wir nun das
System der Gleichungen (III) der Reihe nach durch die
Unterdeterminanten der ersten Colonne von $D^{(\mu)}$

\[
A_{\mu\mu},\quad A_{\mu+1,\mu},\quad A_{\mu+2,\mu}\cdots A_{n\mu},
\]

und addiren dieselben, so erhalten wir die Gleichung

\weq{IV}{
\begin{aligned}
&(a_{\mu\mu}A_{\mu\mu}+a_{\mu+1,\mu}A_{\mu+1,\mu}\cdots+a_{n\mu}A_{n\mu})x_1\\
&\qquad+(a_{\mu,\mu+1}A_{\mu\mu}\cdots+a_{n,\mu+1}A_{n\mu})x_2+\cdots\\
&\qquad+(a_{\mu n}A_{\mu\mu}+a_{\mu+1,n}A_{\mu+1,\mu}\cdots+a_{nn}A_{n\mu})x_n=0.
\end{aligned}
}


% ===== pdf idx 119, printed 105.0 =====
% !PRINTED 105
\weberpage{105}

% running head: \S.\,8. Systeme linearer Gleichungen und deren Aufl\"osung.

Hierin ist der Co\"efficient von $x_1$ gem\"ass der Voraussetzung null, alle
\"ubrigen Co\"efficienten verschwinden identisch, gem\"ass den Relationen
($C$) Nr. 32. Gleichung (IV) ist also eine Identit\"at; wir k\"onnen sie
auch, wenn wir die Gleichungen des Systemes (III) einen Augenblick mit
\[
f_\mu, \quad f_{\mu+1} \cdots f_n
\]
bezeichnen, in der Form schreiben
\weq{V}{A_{\mu\mu} f_\mu + A_{\mu+1,\mu} f_{\mu+1} \cdots A_{n\mu} f_n = 0.}
% ?? uncertain: a faint mark appears above the plus sign between
% $A_{\mu+1,\mu}$ and $f_{\mu+1}$ in the print; it does not resolve into
% legible characters and is not transcribed.
Aus dieser Identit\"at schliesst man aber, dass eine der Gleichungen
$f_{\mu+i}$ eine Folge aller \"ubrigen sein muss, und zwar etwa gerade die
erste des Systemes (II) $(\overline{\mathrm{III}})$
% ?? uncertain: the printed parenthetical after "des Systemes" shows (II)
% overstruck/cancelled, followed by (III) set with an overline; the exact
% typographic intent of this correction mark is not fully legible.
\[
f_\mu = a_{\mu 1} x_1 + a_{\mu 2} x_2 \cdots + a_{\mu n} x_n = 0,
\]
da sie mit einem Factor multiplicirt auftritt, von dem wir voraus-
%CONT>>


% ===== pdf idx 120, printed 106.0 =====
% !PRINTED 106
\weberpage{106}

% running head: 106 \quad Zweiter Theil. Anwendungen der Determinantentheorie.

%<<CONT
gesetzt haben, dass er von Null verschieden sei.

101. \textit{Zweite Folgerung: $\mu$ Gleichungen sind eine Folge der
$n - \mu$ \"ubrigen.} Ersetzen wir in dem Systeme (III) die Gleichung
$f_\mu$ durch irgend eine der vorhergehenden Gleichungen $f_i$, dann ist
die Matrix dieses Systemes dargestellt durch
\weq{B}{
\begin{vmatrix}
a_{i1} & a_{i2} & a_{i3} & \cdots & a_{i\mu} & a_{i,\mu+1} & \cdots & a_{in} \\
a_{\mu+1,1} & a_{\mu+1,2} & \cdots & \cdots & a_{\mu+1,\mu} & a_{\mu+1,\mu+1} & \cdots & a_{\mu+1,n} \\
a_{\mu+2,1} & \cdots & \cdots & \cdots & \cdots & a_{\mu+2,\mu} & \cdots & a_{\mu+2,n} \\
\cdot & \cdot & \cdot & \cdot & \cdot & \cdot & \cdot & \cdot \\
a_{n1} & \cdots & \cdots & \cdots & \cdots & a_{n\mu} & a_{n,\mu+1} & \cdots a_{nn}
\end{vmatrix}.
}
Ihre Determinanten sind noch immer nach Voraussetzung null, w\"ahrend die
Unterdeterminanten derselben nicht durchwegs verschwinden; insbesondere
ist noch immer der Minor
\[
A_{\mu\mu}
\]
von Null verschieden, und daher schliessen wir ganz auf dieselbe Weise
wie vorhin, dass auch $f_i$ eine Folge der letzten $n - \mu$ Gleichungen
sein muss.

Wir sind daher berechtigt, ganz allgemein den Satz aufzustellen:
\glqq Sind in einem Systeme homogener linearer Gleichungen sowohl die
Determinante des Systemes selbst, als auch alle Unterdeterminanten bis
zur $(\mu - 1)^{\mathrm{ten}}$ Ordnung inclus. gleich Null, aber nicht
alle der n\"achsth\"oheren Ordnung, so sind $\mu$ Gleichungen die Folge
der \"ubrigen richtig gew\"ahlten $n - \mu$ Gleichungen.\grqq{}
Aus diesen $n - \mu$ Gleichungen lassen sich aber die Werthe von
$n - \mu$ Variabeln $x_i$ in Function der \"ubrigen $\mu$ Variabeln
berechnen, wie wir sogleich sehen werden.

102. \textit{Aufl\"osung eines Gleichungssystemes, in welchem die Zahl
$\mu$ der Gleichungen kleiner ist als die Zahl $n$ der Unbekannten. Erster
Fall.} Wenn wir uns nun damit befassen, Werthsysteme $x_\mu$ aufzufinden,
welche diese $n - \mu$ Gleichungen mit $\mu$ Unbekannten befriedigen, so
d\"urfen wir zun\"achst die Annahme machen, dass nicht s\"ammtliche
Determinanten der Matrix dieses Systemes von $n - \mu$ Gleichungen
verschwinden. Denn w\"urde dies eintreten, dann w\"are nach dem eben
Erw\"ahnten mindestens eine derselben eine Folge der \"ubrigen, und die
Zahl der selbstst\"andigen Gleichungen w\"urde sich eben nur noch weiter
reduciren.

\textit{Wir beginnen mit dem ersten Falle $\mu = n - 1$.} Das System der
$n - 1$ Gleichungen mit $n$ Unbekannten k\"onnen wir dann in der Form
aufstellen:
\weq{VI}{
\left\{
\begin{aligned}
a_{21} x_1 &+ a_{22} x_2 + a_{23} x_3 + \cdots a_{2n} x_n = 0 \\
a_{31} x_1 &+ a_{32} x_2 + a_{33} x_3 + \cdots a_{3n} x_n = 0 \\
a_{41} x_1 &\cdot \, \cdot \, \cdot \, \cdot \, \cdot \, \cdot \, \cdot \, \cdot \, \cdot \, a_{4n} x_n = 0 \\
&\cdot \; \cdot \; \cdot \; \cdot \; \cdot \; \cdot \; \cdot \; \cdot \; \cdot \; \cdot \\
a_{n1} x_1 &+ a_{n2} x_2 + a_{n3} x_3 + \cdots a_{nn} x_n = 0.
\end{aligned}
\right.
}
Vergleichen wir alsdann dieses System von Gleichungen mit dem in Nr. 34
aufgestellten System von Identit\"aten
\weq{D}{
\left\{
\begin{aligned}
a_{21} p_1 &+ a_{22} p_2 + a_{23} p_3 \cdots a_{2n} p_n = 0 \\
a_{31} p_1 &+ a_{32} p_2 \cdots \cdots \cdots a_{3n} p_n = 0 \\
a_{41} p_1 &\cdot \, \cdot \, \cdot \, \cdot \, \cdot \, \cdot \, \cdot \, \cdot \, \cdot = 0 \\
&\cdot \; \cdot \; \cdot \; \cdot \; \cdot \; \cdot \; \cdot \; \cdot \; \cdot \\
a_{n1} p_1 &+ a_{n2} p_2 + a_{n3} p_3 \cdots a_{nn} p_n = 0,
\end{aligned}
\right.
}
wo die Gr\"ossen $p_1 p_2 \ldots p_n$ nichts anderes sind als die mit
wechselndem Vorzeichen genommenen Determinanten der Matrix:
\weq{A}{
\begin{vmatrix}
a_{21} & a_{22} & a_{23} & \cdots & a_{2n} \\
a_{31} & a_{32} & a_{33} & \cdots & a_{3n} \\
a_{41} & a_{42} & \cdot & \cdots & a_{4n} \\
\cdot & \cdot & \cdot & \cdot & \cdot \\
a_{n1} & a_{n2} & a_{n3} & \cdot & a_{nn}
\end{vmatrix},
}
so erkennt man, dass das Gleichungssystem (VI) befriedigt wird durch
\[
x_i = \frac{1}{\varrho} \cdot p_i,
\]
%CONT>>


% ===== pdf idx 121, printed 107.0 =====
% !PRINTED 107
\weberpage{107}
wo $\dfrac{1}{\varrho}$ ein beliebiger Factor und
\[ p_i = (-1)^{i-1}(a_{21}a_{32}\ldots a_{i,i-1}a_{i+1,i+1}\ldots a_{nn}) \text{ ist.} \]

103. \textit{Eindeutigkeit der erhaltenen Lösung.} Dass diese Lösung die einzige ist, welche das System (VI) zulässt, kann man durch directe Berechnung derselben erweisen. Bezeichnet man nämlich die ersten Unterdeterminanten der Determinante
\[
p_n = (-1)^{n-1}
\begin{vmatrix}
a_{21} & a_{22} & a_{23} & \ldots & a_{2,n-1} \\
a_{31} & a_{32} & a_{33} & \ldots & a_{3,n-1} \\
a_{41} & \cdot & \cdot & \cdot & \cdot \\
 & \cdot & \cdot & \cdot & \cdot \\
a_{n1} & a_{n2} & \ldots & \ldots & a_{n,n-1}
\end{vmatrix}
\]
mit $A_{ik}$, wo $i>1$, und $k<n$, dann können wir zunächst die Gleichungen (VI) mit den Minoren ihrer ersten Colonne der Reihe nach multipliciren und addiren; dies liefert, da die Coefficienten der Unbekannten $x_2$ bis $x_{n-1}$ verschwinden, und da einestheils
\[ (a_{21}A_{21}+a_{31}A_{31}\cdots a_{n1}A_{n1}) = (-1)^{n-1}p_n, \]
und anderntheils
\[ (a_{2\,n}A_{21}+a_{3\,n}A_{31}\cdots+a_{n\,n}A_{n1}) = (-1)^{n-2}p_1, \]
die Gleichung
\weq{1}{p_n x_1 - p_1 x_n = 0.}
Ganz ebenso liefern die Multiplicationen mit den Minoren der zweiten, dritten $\ldots(n-1)^{\text{ten}}$ Colonne von $p_n$ und die daraufolgende Addition der Reihe nach die Gleichungen.
\weq{2}{p_n x_2 - p_2 x_n = 0,}
\weq{3}{p_n x_3 - p_3 x_n = 0,}
\[ \cdot\ \cdot\ \cdot\ \cdot\ \cdot\ \cdot\ \cdot \]
\weq{(n)}{p_n x_{n-1} - p_{n-1} x_n = 0.}
Aus diesen ($n$) Gleichungen folgt zunächst, dass $x_n$ proportional $p_n$ ist, so dass man hat:
\[ \varrho x_n = p_n. \]
Setzt man diesen Werth in die erhaltenen Gleichungen ein, so folgt der Reihe nach (wobei die geeigneten Vorzeichen bereits im Symbol $p_i$ mit inbegriffen)
\[ \varrho x_1 = p_1, \quad \varrho x_2 = p_2, \quad \varrho x_3 = p_3 \cdots \text{u.\,s.\,w.,} \]
also gerade die schon oben angeführten Werthe der Unbekannten. Aus
%CONT>>


% ===== pdf idx 122, printed 108.0 =====
% !PRINTED 108
\weberpage{108}
%<<CONT
diesem direct durch Berechnung erhaltenen Resultate folgt aber, dass die gewonnene Lösung die einzig mögliche ist. — Man kann daher den Satz aufstellen:

\glqq Ein System von $(n-1)$ homogenen linearen Gleichungen mit $n$ Unbekannten hat, wenn die Determinanten der Matrix des Systemes nicht durchwegs verschwinden, als einziges System von Lösungen die durch die continuirliche Proportion
\[ x_1:x_2:x_3\cdots:x_n = A_{11}:A_{12}:A_{13}\cdots:A_{1n} \]
gegebenen Werthe von $x$, wenn $A_{1i}$ die mit wechselndem Vorzeichen genommenen Determinanten der Matrix des Systemes sind.\grqq{}

104. \textit{Anwendungen. Beispiel} 1. Seien gegeben die beiden Gleichungen:
\weq{1}{u_1x_1+u_2x_2+u_3x_3=0,}
\weq{2}{v_1x_1+v_2x_2+v_3x_3=0,}
dann ist nach dem erhaltenen Satze das System von Lösungen gegeben durch:
\[ x_1:x_2:x_3 = (u_2v_3-v_2u_3):-(u_1v_3-v_1u_3):(u_1v_2-u_2v_1) \]
\[ = (u_2v_3):(u_3v_1):(u_1v_2), \]
was man unmittelbar durch die gewöhnlichen Methoden der Elimination verificiren kann. Deutet man die Gleichungen 1) und 2) geometrisch, so stellen sie zwei Gerade vor, und die Aufgabe ist, ihren Schnittpunkt suchen. Man erhält die Lösung: die Coordinaten des Schnittpunktes sind proportional den Determinanten der Liniencoordinaten beider Geraden, oder (vergleiche auch den Satz über correspondirende Matrices und Beispiel Nr. 95) die Determinanten der Matrix
\[
\begin{vmatrix} u_1 & u_2 & u_3 \\ v_1 & v_2 & v_3 \end{vmatrix}
\]
sind proportional den correspondirenden Determinanten der Matrix
\[ \begin{vmatrix} x_1, & x_2, & x_3 \end{vmatrix}. \]

\textit{Beispiel} 2. Seien ferner gegeben die drei Gleichungen
\weq{1}{u_1x_1+u_2x_2+u_3x_3+u_4x_4=0,}
\weq{2}{v_1x_1+v_2x_2+v_3x_3+v_4x_4=0,}
\weq{3}{w_1x_1+w_2x_2+w_3x_3+w_4x_4=0.}
Das gemeinschaftliche System der Lösungen ist dargestellt durch die Determinanten der Matrix
%CONT>>


% ===== pdf idx 123, printed 109.0 =====
% !PRINTED 109
\weberpage{109}
%<<CONT
\[
\begin{vmatrix}
u_1 & u_2 & u_3 & u_4 \\
v_1 & v_2 & v_3 & v_4 \\
w_1 & w_2 & w_3 & w_4
\end{vmatrix},
\]
die wir bezeichnen mit
\[ A_{11}=(u_2v_3w_4), \quad A_{12}=-(u_1v_3w_4), \quad A_{13}=(u_1v_2w_4), \]
\[ A_{14}=-(u_1v_2w_3); \]
daher sind die Coordinaten des Schnittpunktes der drei Ebenen 1), 2), 3) dargestellt durch
\[ x_1:x_2:x_3:x_4:=A_{11}:A_{12}:A_{13}:A_{14}, \]
wodurch abermals der Satz über correspondirende Matrices verificirt ist.

105. \textit{Zweiter Fall. Die Zahl $\mu$ der Gleichungen ist um mehr als 1 kleiner als die Zahl $n$ der Unbekannten.} Gehen wir nun zum allgemeinen Falle über, in welchem
\[ n > \mu + 1, \]
so treten uns sofort zwei Fragen entgegen; nämlich: 1) Wie findet man Lösungen überhaupt? 2) Welches ist die allgemeinste Lösung eines solchen Systemes von homogenen Gleichungen? Wir beantworten zunächst die erste Frage, und nehmen dabei an, dass das System von Gleichungen die Form habe
\[
\text{(I)}\quad
\left\{
\begin{aligned}
&a_{11}x_1+a_{12}x_2+a_{13}x_3+\cdots a_{1n}x_n=0 \\
&a_{21}x_1+a_{22}x_2+a_{23}x_3+\cdots a_{2n}x_n=0 \\
&a_{31}x_1+a_{32}x_2\ \cdot\ \cdot\ \cdot\ \cdot\ \cdot\ a_{3n}x_n=0 \\
&\cdot\ \cdot\ \cdot\ \cdot\ \cdot\ \cdot\ \cdot\ \cdot\ \cdot\ \cdot\ \cdot\ \cdot \\
&a_{\mu1}x_1+a_{\mu2}x_2+a_{\mu3}x_3+\cdots a_{\mu n}x_n=0,
\end{aligned}
\right.
\]
\[ \mu+1<n, \ \text{oder} \]
\[ n-\mu=\nu>1. \]
Man kann nun eine Reihe von Lösungen dadurch finden, dass man eine Anzahl der Variabeln von vornherein gleich null setzt, etwa so viele, dass in den $\mu$ Gleichungen noch $\mu+1$ Unbekannte auftreten. Löst man dieselben hierauf nach der in Nr. 102 und 103 gegebenen Methode auf, so erhält man in Verbindung mit den Variabeln, deren Werthe null sind, ein erstes Werthsystem, das die Gleichungen (I) befriedigt. Setzt man nun andere $n-(\mu+1)$ Unbekannte gleich null, so erhält man auf dieselbe Weise ein zweites System von Werthen, das vom ersten im Allgemeinen verschieden ist. Wir
%CONT>>


% ===== pdf idx 124, printed 110.0 =====
% !PRINTED 110
\weberpage{110}
%<<CONT
können uns so eine ganze Reihe von Werthsystemen verschaffen, und nehmen an, wir hätten $\nu$ von einander unabhängige Lösungen bereits gefunden, welche der Reihe nach durch die $\nu$ Zeilen der folgenden Matrix
\[
\text{(1)}\quad
\begin{vmatrix}
b_{11} & b_{12} & b_{13} & b_{14} & \ldots & b_{1n} \\
b_{21} & b_{22} & b_{23} & b_{24} & \ldots & b_{2n} \\
b_{31} & \cdot & \cdot & \cdot & \cdot & b_{3n} \\
 & \cdot & \cdot & \cdot & \cdot & \cdot \\
b_{\nu1} & b_{\nu2} & b_{\nu3} & \ldots & \ldots & b_{\nu n}
\end{vmatrix}
\]
dargestellt sein sollen. Tragen wir der Reihe nach diese $\nu$ Werthsysteme in (I) ein, so erhalten wir $\mu\cdot\nu$ Identitäten von der Form
\[
\text{(II)}\quad c_{ik}=a_{i1}b_{k1}+a_{i2}b_{k2}+a_{i3}b_{k3}\cdots+a_{in}b_{kn}=0,
\]
wo $i$ die Werthe von 1 bis $\mu$, und $k$ die Werthe von 1 bis $\nu$ annehmen kann. Vermöge dieser Gleichungen (II) sind die Elemente der Matrix (2) mit denen der Matrix des Systemes (I)
\[
\text{(2)}\quad
\begin{vmatrix}
a_{11} & a_{12} & a_{13} & a_{14} & \ldots & a_{1n} \\
a_{21} & a_{22} & a_{23} & a_{24} & \ldots & a_{2n} \\
a_{31} & \cdot & \cdot & \cdot & \cdot & a_{3n} \\
 & \cdot & \cdot & \cdot & \cdot & \cdot \\
a_{\mu1} & a_{\mu2} & a_{\mu3} & a_{\mu4} & \ldots & a_{\mu n}
\end{vmatrix}
\]
in eine Beziehung gesetzt, welche die beiden Matrices (1) und (2) als correspondirende erscheinen lässt. (Vergleiche Nr. 94.)

106. \textit{Zweiter Beweis des Satzes über correspondirende Matrices.} Wenn wir nun im Gleichungssystem (I) Nr. 105 setzen:
\[ x_\mu = lx_\mu, \quad x_{\mu+1}=lx_{\mu+1}, \quad x_{\mu+2}=lx_{\mu+2}\cdots x_n=lx_n, \]
wo $l$ den Werth 1 besitzt, dann erhalten wir
\[
\begin{aligned}
&a_{11}x_1+a_{12}x_2+\cdots a_{1,\mu-1}x_{\mu-1}+l(a_{1\mu}x_\mu+a_{1,\mu+1}x_{\mu+1}+\cdots a_{1n}x_n)=0 \\
&a_{21}x_1+a_{22}x_2+\cdots a_{2,\mu-1}x_{\mu-1}+l(a_{2\mu}x_\mu+\cdots\cdots a_{2n}x_n)=0 \\
&a_{31}x_1+a_{32}x_2+\cdots\cdots+l(a_{3\mu}x_\mu+\cdots\cdots a_{3n}x_n)=0 \\
&\cdot\ \cdot\ \cdot\ \cdot\ \cdot\ \cdot\ \cdot\ \cdot\ \cdot\ \cdot\ \cdot\ \cdot\ \cdot\ \cdot\ \cdot\ \cdot\ \cdot \\
&a_{\mu1}x_1+a_{\mu2}x_2+\cdots\cdots+l(a_{\mu\mu}x_\mu+\cdots\cdots a_{\mu n}x_n)=0.
\end{aligned}
\]
Diese Gleichungen, $\mu$ an der Zahl, sind homogen in den $\mu$ Variabeln
\[ x_1x_2x_3\cdots x_{\mu-1},\ l; \]
%CONT>>


% ===== pdf idx 125, printed 111.0 =====
% !PRINTED 111
\weberpage{111}
\noindent \S.\,8. Systeme linearer Gleichungen und deren Aufl\"osung.

%<<CONT
sollen sie also ein gemeinschaftliches System von L\"osungen besitzen, so muss ihre Resultante verschwinden, d. h. es muss sein:

\weq{III}{
\begin{vmatrix}
a_{11}\ a_{12}\cdots a_{1,\mu-1}, & a_{1\mu}x_\mu + a_{1,\mu+1}x_{\mu+1}+\cdots a_{1n}x_n \\
a_{21}\ a_{22}\cdots a_{2,\mu-1}, & a_{2\mu}x_\mu + a_{2,\mu+1}x_{\mu+1}+\cdots a_{2n}x_n \\
a_{31}\ a_{32}\cdots a_{3,\mu-1}, & \cdot\ \cdot\ \cdot\ \cdot\ \cdot\ \cdot\ \cdot\ \cdot\ \cdot\ \cdot\ \cdot\ \cdot\ \cdot\ \cdot \\
& \cdot\ \cdot\ \cdot\ \cdot\ \cdot\ \cdot\ \cdot\ \cdot\ \cdot\ \cdot\ \cdot\ \cdot\ \cdot\ \cdot \\
a_{\mu 1}\ a_{\mu 2}\cdots a_{\mu,\mu-1}, & a_{\mu\mu}x_\mu + a_{\mu,\mu+1}x_{\mu+1}+\cdots a_{\mu n}x_n
\end{vmatrix}=0.
}
% ?? uncertain: printed scan shows a faint caret/strike mark over the $x_n$ in row 2, likely a foxing or pencil mark on the exemplar rather than typeset text; transcribed as printed letterform $a_{2n}x_n$.

Da in dieser Resultante die letzte Colonne eine Summe von $n-\mu+1$ Gliedern ist, so k\"onnen wir dieselbe nach dem Additionsgesetze (vergl. Nr. 35 und 36) f\"ur Determinanten in die Summe zerlegen:

\[
\sum_{k=0}^{k=n-\mu}
\begin{vmatrix}
a_{11}\ a_{12}\cdots a_{1,\mu-1} & a_{1,\mu+k}\cdot x_{\mu+k} \\
a_{21}\ a_{22}\cdots a_{2,\mu-1} & a_{2,\mu+k}\cdot x_{\mu+k} \\
a_{31}\ a_{32}\ \cdot\ \cdot\ \cdot\ \cdot\ \cdot\ \cdot\ \cdot\ \cdot & \\
\cdot\ \cdot\ \cdot\ \cdot\ \cdot\ \cdot\ \cdot\ \cdot\ \cdot\ \cdot & \\
a_{\mu 1}\ a_{\mu 2}\cdots a_{\mu,\mu-1} & a_{\mu,\mu+k}\cdot x_{\mu+k}
\end{vmatrix}=0.
\]

Weil aber die letzte Colonne jeder Determinante dieser Summe mit $x_{\mu+k}$ multiplicirt ist, so k\"onnen wir diese Entwicklung der Resultante unter Einf\"uhrung der Bezeichnung

\[
A_{\mu+k}=
\begin{vmatrix}
a_{11} & a_{12} & a_{13} & \cdots & a_{1,\mu-1}, & a_{1,\mu+k} \\
a_{21} & a_{22} & \cdot & \cdots & a_{2,\mu-1}, & a_{2,\mu+k} \\
a_{31} & a_{32} & \cdot & \cdots & \cdot & \cdot \\
\cdot & \cdot & \cdot & \cdots & \cdot & \cdot \\
a_{\mu 1} & a_{\mu 2} & \cdot & \cdots & a_{\mu,\mu-1}, & a_{\mu,\mu+k}
\end{vmatrix}
\]

auch darstellen in der Form:

\weq{$\mathrm{III}_a$}{
\sum_{k=0}^{k=n-\mu} A_{\mu+k}x_{\mu+k} = A_\mu x_\mu + A_{\mu+1}x_{\mu+1} + A_{\mu+2}x_{\mu+2} + \cdots + A_n x_n = 0.
}

Diese Gleichung ($\mathrm{III_a}$) muss wie das System von Gleichungen (I) f\"ur die $\nu$ Werthsysteme bestehen, welche durch die Zeilen der Matrix (1) repr\"asentirt sind. Es folgen daher aus ($\mathrm{III_a}$) die $\nu$ Relationen:

\[
\begin{aligned}
A_\mu b_{1\mu} + A_{\mu+1} b_{1\mu+1} + \cdots A_n b_{1n} &= 0 \\
A_\mu b_{2\mu} + A_{\mu+1} b_{2\mu+1} + \cdots A_n b_{2n} &= 0 \\
A_\mu b_{3\mu} + A_{\mu+1} b_{3\mu+1} + \cdots A_n b_{3n} &= 0 \\
\cdot\ \cdot\ \cdot\ \cdot\ \cdot\ \cdot\ \cdot\ \cdot\ \cdot\ \cdot\ \cdot\ \cdot \\
A_\mu b_{\nu\mu} + A_{\mu+1} b_{\nu\mu+1} + \cdots A_n b_{\nu n} &= 0.
\end{aligned}
\]


% ===== pdf idx 126, printed 112.0 =====
% !PRINTED 112
\weberpage{112}
\noindent 112 \quad Zweiter Theil. Anwendungen der Determinantentheorie.

Betrachtet man in diesen Relationen die $\nu+1$ Gr\"ossen

\[
A_\mu,\quad A_{\mu+1},\ \cdots A_n
\]

als Unbekannte und eliminirt dieselben, so erh\"alt man nach Nr. 103:

\weq{A}{
\begin{aligned}
\varrho A_\mu &= B_\mu \\
\varrho A_{\mu+1} &= B_{\mu+1} \\
\varrho A_{\mu+2} &= B_{\mu+2} \\
\cdot\ \cdot\ \cdot\ \cdot \\
\varrho A_n &= B_n.
\end{aligned}
}

Dabei sind die Gr\"ossen $B_{\mu+k}$ (inclus. Vorzeichen) Determinanten der Matrix (1), und zwar ist $B_{\mu+k}$ dadurch entstanden, dass man die $n-1$ ersten und \"uberdies die $n+k^{\text{te}}$ Colonne in derselben unterdr\"uckt hat, w\"ahrend die Gr\"osse $A_{\mu+k}$ aus Matrix (2) gerade durch die hier unterdr\"uckten Colonnen gebildet ist. Die Colonnen entsprechender Gr\"ossen $A$ und $B$ bilden also gerade eine Permutation der Zahlen

\[
1,\ 2,\ 3\ \cdots n;
\]

ihr Vorzeichen entspricht der Anzahl der Derangements dieser Permutation, wie man aus den Beziehungen ($A$) erkennt, in denen die Determinanten $A_{\mu+i}$ stets positiv, die Determinanten $B_{\mu+i}$ aber mit wechselndem Vorzeichen der Reihe nach zu nehmen sind. Die Determinanten

\[
A_{\mu+k}\quad \text{und}\quad B_{\mu+k}
\]

sind also correspondirende Determinanten gem\"ass unserer Definition in Nr. 90.

H\"atte man statt der $\nu+1$ Unbekannten

\[
x_\mu,\ x_{\mu+1}\ \cdots x_n
\]

andere des Systems (I) mit dem Factor $l$ multiplicirt und die \"ubrigen auf die angegebene Weise eliminirt, so w\"are man geradeso auf eine zweite Reihe von Relationen ($A$) gekommen, so dass man erkennt: Alle correspondirenden Determinanten beider correspondirender Matrices sind einander proportional, wie wir auch fr\"uher bereits bewiesen haben.

107. \textit{Allgemeinste L\"osung des Systemes (I) homogener linearer Gleichungen.} Nehmen wir nun an, die allgemeinste L\"osung des Systemes (I) in Nr. 105 sei gefunden, und zwar soll das Werthsystem der $n$ Variabeln dargestellt sein durch Werthe:

\[
b_1,\ b_2,\ b_3\ \cdots b_n.
\]

Tragen wir diese Werthe in die Gleichungen (II) Nr. 105 ein, so m\"ussen dieselben befriedigt werden. Es ist also auch die Matrix
%CONT>>


% ===== pdf idx 127, printed 113.0 =====
% !PRINTED 113
\weberpage{113}
\noindent \S.\,8. Systeme linearer Gleichungen und deren Aufl\"osung.

%<<CONT
\weq{3}{
\begin{vmatrix}
b_{11} & b_{12} & b_{13} & \cdots & b_{1n} \\
b_{21} & b_{22} & b_{23} & \cdots & b_{2n} \\
b_{31} & \cdot & \cdot & \cdots & b_{3n} \\
\cdot & \cdot & \cdot & \cdot & \cdot \\
b_1 & b_2 & b_3 & \cdots & b_n
\end{vmatrix},
}

die aus der Matrix (1) hervorgeht, indem man in ihr die letzte Zeile ersetzt durch die Elemente

\[
b_1\ b_2\ \cdots b_n,
\]

eine correspondirende Matrix mit Matrix (2), und ihre Determinanten sind proportional den correspondirenden der Matrix (2). Daher sind sie auch proportional den entsprechenden, d. h. den aus gleichen Colonnen gebildeten Determinanten der Matrix (1).

108. Erweitern wir also diese Matrix (1) durch Anf\"ugung der Zeile

\[
b_1\ b_2\ b_2\ \cdots b_n
\]
% [sic] printed "$b_1\ b_2\ b_2\ \cdots b_n$"; expected third entry $b_3$.

zur Matrix

\weq{4}{
\begin{vmatrix}
b_{11} & b_{12} & b_{13} & \cdots & b_{1n} \\
b_{21} & b_{22} & b_{23} & \cdots & b_{2n} \\
b_{31} & \cdot & \cdot & \cdots & b_{3n} \\
\cdot & \cdot & \cdot & \cdot & \cdot \\
b_{\nu 1} & b_{\nu 2} & b_{\nu 3} & \cdots & b_{\nu n} \\
b_1 & b_2 & b_3 & \cdots & b_n
\end{vmatrix}
}

und entnehmen aus dieser Matrix irgend eine Determinante, etwa

\[
D=
\begin{vmatrix}
b_{11} & b_{12} & \cdots & b_{1,\nu+1} \\
b_{21} & b_{22} & \cdot & b_{2,\nu+1} \\
\cdot & \cdot & \cdot & \cdot \\
b_{\nu 1} & b_{\nu 2} & \cdot & b_{\nu,\nu+1} \\
b_1 & b_2 & \cdots & b_{\nu+1}
\end{vmatrix},
\]

so sind in dieser Determinante die Minoren der $\nu^{\text{ten}}$ Zeile proportional den entsprechenden Minoren der $\nu+1^{\text{ten}}$ Zeile. Es verh\"alt sich:

\[
B_{\nu 1}:B_{\nu 2}:B_{\nu 3}\ \cdots B_{\nu,\nu+1} = B_{\nu+1,1}:B_{\nu+1,2}:B_{\nu+1,3}\ \cdots B_{\nu+1,\nu+1}.
\]

Folglich ist:

\[
B_{\nu+1,i}B_{\nu k} - B_{\nu+1,k}B_{\nu i} = 0.
\]

Nach den Lehrs\"atzen \"uber adjungirte Determinanten und deren Minoren ist aber (vergl. Nr. 88)
%CONT>>


% ===== pdf idx 128, printed 114.0 =====
% !PRINTED 114
\weberpage{114}
\noindent 114 \quad Zweiter Theil. Anwendungen der Determinantentheorie.

%<<CONT
\[
\begin{vmatrix} B_{\nu i} & B_{\nu k} \\ B_{\nu+1,i} & B_{\nu+1,k} \end{vmatrix}
= D\cdot \underset{i}{\overset{}{B}}{}^{\nu,\nu+1}_{\phantom{\nu,\nu+1}k} = 0,\quad
\begin{cases} i=1\cdots n-1 \\ k=2\cdots n \end{cases},
\]

wo

\[
\underset{i}{\overset{}{B}}{}^{\nu,\nu+1}_{\phantom{\nu,\nu+1}k}
\]

jede Determinante sein kann, die sich aus den $\nu-1$ ersten Zeilen von $D$ bilden l\"asst. Nehmen wir nun auch an

\[
\underset{i}{\overset{}{B}}{}^{\nu,\nu+1}_{\phantom{\nu,\nu+1}k} = 0,
\]

so folgt immer auch $D=0$, da eine Determinante verschwinden muss, wenn alle ihre Minoren zweiter Ordnung verschwinden, nach welchen sie man ja h\"atte entwickeln k\"onnen. In derselben Weise aber, wie wir gefunden haben $D=0$, k\"onnen wir auch schliessen, dass \"uberhaupt jede Determinante der Matrix (4) verschwinden muss. Daraus folgt, dass es immer Werthe $\lambda_i$ giebt, so beschaffen, dass

\[
\begin{aligned}
b_1 &= \lambda_1 b_{11} + \lambda_2 b_{21} + \lambda_3 b_{31}\ \cdot\ \cdot\ \cdot\ \cdot\ \cdot \lambda_\nu b_{\nu 1} \\
b_2 &= \lambda_1 b_{12} + \lambda_2 b_{22} + \ \cdot\ \cdot\ \cdot\ \cdot\ \cdot \lambda_\nu b_{\nu 2} \\
b_3 &= \lambda_1 b_{13} + \lambda_2 b_{23} + \ \cdot\ \cdot\ \cdot\ \cdot\ \cdot \lambda_\nu b_{\nu 3} \\
&\cdot\ \cdot\ \cdot\ \cdot\ \cdot\ \cdot\ \cdot\ \cdot\ \cdot\ \cdot\ \cdot\ \cdot \\
b_n &= \lambda_1 b_{1n} + \lambda_2 b_{2n} + \lambda_3 b_{3n} + \cdots \lambda_\nu b_{\nu n}.
\end{aligned}
\]

Denn tragen wir diese Werthe f\"ur $b_i$ in irgend eine der Determinanten $D$ ein, so ist dieselbe Null, gem\"ass den S\"atzen in Nr. 35 bis 39.

Wir haben also den Satz gewonnen: \glqq Sind die Zeilen der Matrix (2) $\nu$ beliebige von einander linear unabh\"angige L\"osungen, so ist die allgemeinste L\"osung des Systemes von Gleichungen (I) eine lineare Function derselben.\grqq{}

109. \textit{Beispiel.} Es seien gegeben zwei Gleichungen mit vier Unbekannten, also

\[
n=4,\quad \mu=2,\quad \nu=2,
\]

etwa das System

\weq{I}{
\begin{aligned}
u_1 x_1 + u_2 x_2 + u_3 x_3 + u_4 x_4 &= 0 \\
v_1 x_1 + v_2 x_2 + v_3 x_3 + v_4 x_4 &= 0.
\end{aligned}
}

Wir nehmen an, dass sie befriedigt werden einmal durch das System von speciellen Werthen

\[
x_1\ x_2\ x_3\ x_4,
\]

dann aber auch durch das System von Werthen

\[
y_1\ y_2\ y_3\ y_4.
\]


% ===== pdf idx 129, printed 115.0 =====
% !PRINTED 115
\weberpage{115}

Während also die Matrix der Co\"efficienten des gegebenen Gleichungssystemes dargestellt ist durch

\weq{1}{\left| \begin{array}{cccc} u_1 & u_2 & u_3 & u_4 \\ v_1 & v_2 & v_3 & v_4 \end{array} \right|,}

ist die Matrix der beiden speciellen Werthsysteme, für welche die Gleichungen (I) bestehen sollen, gegeben durch

\weq{2}{\left| \begin{array}{cccc} x_1 & x_2 & x_3 & x_4 \\ y_1 & y_2 & y_3 & y_4 \end{array} \right|,}

sodass ausser den Relationen (I) auch noch die Relationen

\weq{II}{\begin{aligned} u_1y_1+u_2y_2+u_3y_3+u_4y_4&=0\\ v_1y_1+v_2y_2+v_3y_3+v_4y_4&=0 \end{aligned}}

vorhanden sind. Setzen wir dann im System (I)

\[ x_2=lx_2,\quad x_3=lx_3,\quad x_4=lx_4,\quad \text{wo } l=1, \]

so erhält man:

\[ \begin{aligned} u_1x_1+l(u_2x_2+u_3x_3+u_4x_4)&=0\\ v_1x_1+l(v_2x_2+v_3x_3+v_4x_4)&=0. \end{aligned} \]

Die Resultante dieser Gleichung muss verschwinden, wenn Werthe $x_1$ und $l$ vorhanden sein sollen, für welche sie bestehen; d.\,h. es muss sein:

\weq{III}{\left| \begin{array}{cc} u_1 & u_2x_2+u_3x_3+u_4x_4 \\ v_1 & v_2x_2+v_3x_3+v_4x_4 \end{array} \right| = 0.}

Bezeichnet man die Determinanten der Matrix (1) mit $q_{ik}$, so dass

\[ q_{ik} = u_iv_k - v_iu_k, \]

und setzt man ebenso die Determinanten der Matrix (2)

\[ p_{ik} = x_iy_k - y_ix_k, \]

so liefert die Gleichung (III) nach den Unbekannten geordnet:

\weq{IIIa}{x_2q_{12}+x_3q_{13}+x_4q_{14}=0,}

und ebenso folgt

\[ y_1q_{11}+y_3q_{13}+y_4q_{14}=0, \] % [sic] printed q_{11}; pattern of (IIIa) suggests q_{12}, but reproduced as printed

wenn wir mit dem System (II) ebenso verfahren, wie mit (I). Löst man diese Gleichungen nach den Determinanten $q_{ik}$ auf, so kommt:

\[ \begin{aligned} \varrho\cdot q_{12} &= p_{34}\\ \varrho\cdot q_{13} &= p_{42}\\ \varrho\cdot q_{14} &= p_{23}. \end{aligned} \]


% ===== pdf idx 130, printed 116.0 =====
% !PRINTED 116
\weberpage{116}

%<<CONT
Ganz ebenso folgt aber, indem wir etwa $x_1=lx_1,\ x_3=lx_3,\ x_4=lx_4$ setzen, und entsprechend die Unbekannten $y_1=ly_1$ u.\,s.\,w.,

\[ \varrho\cdot q_{23}=p_{14},\quad \varrho\cdot q_{42}=p_{13},\quad \varrho\cdot q_{34}=p_{14}. \]

Dadurch ist zunächst die Proportionalität correspondirender Determinanten der beiden Matrices (1) und (2) nachgewiesen, wie wir sie schon in Nr.\,95 erkannt haben.

110. Suchen wir nun die allgemeinste Lösung der beiden Gleichungen (I); sie sei gegeben durch

\[ z_1\ z_2\ z_3\ z_4. \]

Wir bilden die Matrix

\weq{3}{\left| \begin{array}{cccc} x_1 & x_2 & x_3 & x_4 \\ z_1 & z_2 & z_3 & z_4 \end{array} \right|,}

deren Determinanten dargestellt sein mögen durch

\[ \bar p_{ik} = x_iz_k - z_ix_k. \]

Dann bestehen auch die Gleichungen:

\[ \begin{aligned}
&\sigma\cdot q_{12}=\bar p_{34},\quad \sigma\cdot q_{13}=\bar p_{42},\quad \sigma\cdot q_{14}=\bar p_{23}\\
&\sigma\cdot q_{23}=\bar p_{14},\quad \sigma\cdot q_{42}=\bar p_{13},\quad \sigma\cdot q_{34}=\bar p_{14},
\end{aligned} \]

und folglich wenn wir setzen $\dfrac{\sigma}{\varrho}=\tau$:

\[ \tau\cdot p_{ik} = \bar p_{ik}. \]

Betrachten wir daher die Matrix

\weq{4}{\left| \begin{array}{cccc} x_1 & x_2 & x_3 & x_4 \\ y_1 & y_2 & y_3 & y_4 \\ z_1 & z_2 & z_3 & z_4 \end{array} \right|}

und in dieser wiederum die Determinante

\[ D = \left| \begin{array}{ccc} x_1 & x_2 & x_3 \\ y_1 & y_2 & y_3 \\ z_1 & z_2 & z_3 \end{array} \right|, \]

dann ist

\[ \tau\cdot p_{12}=\bar p_{12},\quad \tau\cdot p_{13}=\bar p_{13},\quad \tau\cdot p_{23}=\bar p_{23} \]

oder

\[ \tau\cdot \left|\begin{array}{cc} x_1&x_2\\ y_1&y_2\end{array}\right| = \left|\begin{array}{cc} x_1&x_2\\ z_1&z_2\end{array}\right|,\quad
\tau\cdot \left|\begin{array}{cc} x_1&x_3\\ y_1&y_3\end{array}\right| = \left|\begin{array}{cc} x_1&x_3\\ z_1&z_3\end{array}\right|,\quad
\tau\cdot \left|\begin{array}{cc} x_2&x_3\\ y_2&y_3\end{array}\right| = \left|\begin{array}{cc} x_2&x_3\\ z_2&z_3\end{array}\right|, \]

und folglich:

\[ x_1D=0,\quad x_2D=0,\quad x_3D=0, \]

d.\,h. $D=0$. Das Nämliche folgt für die drei andern Determinanten
%CONT>>


% ===== pdf idx 131, printed 117.0 =====
% !PRINTED 117
\weberpage{117}

%<<CONT
der Matrix (4). Folglich giebt es Werthe $\lambda$ und $\mu$, so dass die allgemeinste Lösung des Gleichungssystems (I) dargestellt ist durch:

\weq{B}{\begin{aligned} z_1&=\lambda x_1+\mu y_1,\quad z_2=\lambda x_2+\mu y_2\\ z_3&=\lambda x_3+\mu y_3,\quad z_4=\lambda x_4+\mu y_4. \end{aligned}}

Betrachtet man die Gleichungen (I) als die Gleichungen zweier Ebenen, und sind $(x_i)$ und $(y_i)$ zwei Punkte, welche in beiden gleichzeitig liegen, so liegt jeder andere Punkt der beiden Ebenen gemeinschaftlich auf der Verbindungslinie der beiden Punkte, eine Gerade, die in der That durch die vier Gleichungen ($B$) dargestellt ist.

111. \textit{Auflösung eines Systems nicht homogener linearer Gleichungen.} Die Aufgabe, ein System von Gleichungen zu lösen, die nicht homogen sind, lässt sich auf die Auflösung eines Systems von homogenen Gleichungen zurückführen. Wir behandeln daher nur den Fall, der zunächst von Interesse ist, nämlich die Auflösung von $n$ nicht homogenen Gleichungen mit $n$ Unbekannten. Sie seien:

\[ \begin{aligned}
a_{11}y_1+a_{12}y_2+a_{13}y_3+\cdots a_{1n}y_n+a_{1,n+1}&=0\\
a_{21}y_1+a_{22}y_2+a_{23}y_3+\cdots a_{2n}y_n+a_{2,n+1}&=0\\
a_{31}y_1+a_{32}y_2+a_{33}y_3+\cdots a_{3n}y_n+a_{3,n+1}&=0\\
&\cdots\cdots\cdots\cdots\cdots\cdots\cdots\cdots\cdots\\
a_{n1}y_1+a_{n2}y_2+a_{n3}y_3+\cdots a_{nn}y_n+a_{n,n+1}&=0.
\end{aligned} \]

Wir können dieselben in zweierlei Weise homogen machen, entweder, indem wir setzen

\[ y_k = \frac{x_k}{x_{n+1}}, \]

und mit $x_{n+1}$ jede Gleichung multipliciren, oder, indem wir den Co\"efficienten 1 der Absolutglieder als $n+1^{\text{te}}$ Unbekannte $x_{n+1}$ betrachten. Beide Male erhalten wir $n$ homogene Gleichungen mit $n+1$ Unbekannten. Da die Matrix des gegebenen Gleichungssystems dargestellt ist durch

\[ \left| \begin{array}{cccccc}
a_{11} & a_{12} & a_{13} & \cdots & a_{1n} & a_{1,n+1}\\
a_{21} & a_{22} & a_{23} & \cdots & a_{2n} & a_{2,n+1}\\
a_{31} & \cdot & \cdot & \cdot & \cdot & a_{3,n-1}\\ % [sic] printed a_{3,n-1}; pattern of the two rows above suggests a_{3,n+1}
\cdot & \cdot & \cdot & \cdot & \cdot & \cdot\\
a_{n1} & a_{n2} & a_{n3} & \cdots & \cdot & a_{nn+1} % [sic] printed without comma, as a_{nn+1}
\end{array} \right|,\]

so erhalten wir als einziges Werthsystem, das die Gleichungen befriedigt, gemäss den Entwicklungen in Nr.\,103

\[ p_{n+1}x_1=p_1x_{n+1},\quad p_{n+1}x_2=p_2x_{n+1},\quad \cdots p_{n+1}x_n=p_nx_{n+1}, \]
%CONT>>


% ===== pdf idx 132, printed 118.0 =====
% !PRINTED 118
\weberpage{118}

%<<CONT
oder weil

\[ \frac{x_k}{x_{n+1}} = y, \]

\[ y_1=\frac{p_1}{p_{n+1}},\quad y_2=\frac{p_2}{p_{n+1}},\quad \cdots y_n=\frac{p_n}{p_{n+1}}, \]

wobei die Grössen $p_i$ die mit wechselndem Zeichen genommenen Determinanten der Matrix sind, welche man erhält, indem man in dieser Matrix die $i^{\text{te}}$ Colonne unterdrückt.

112. \textit{Discussion der Lösung.} Ist die Determinante $p_{n+1}\gtrless 0$, so besitzen die Unbekannten $y_1\,y_2\,\ldots\,y_n$ endliche und bestimmte Werthe, für welche die Gleichungen zusammen bestehen können. Ist dagegen $p_{n+1}=0$, ohne dass sämmtliche übrige Determinanten verschwinden, so ergeben sich die Lösungen durchweg in der Form

\[ 0\cdot y_i = p_i\,; \]

da es aber keinen endlichen Werth von $y_i$ giebt, der diese Relation befriedigen kann, so müssen die Gleichungen unvereinbar sein. Sind endlich ausser $p_{n+1}$ auch die übrigen Grössen $p_i=0$, dann erscheinen die Werthe der Unbekannten in der unbestimmten Form

\[ 0\cdot y = 0. \]

Nun sahen wir aber in Nr.\,100, dass in diesem Falle, wo sämmtliche Determinanten der Matrix verschwinden, zum mindesten eine Gleichung eine Folge der $n-1$ übrigen ist.

Es giebt dann unendlich viele Werthsysteme $y_i$, welche zwar alle $n$ Gleichungen befriedigen, die aber schon durch die $n-1$ übrigen Gleichungen bestimmt sind.

113. \textit{Beispiel 1.} In den Gleichungen:

\[ \begin{aligned} 3x+7y-26&=0\\ 2x-y-6&=0, \end{aligned} \]

deren Matrix dargestellt ist durch

\[ \left| \begin{array}{ccc} 3 & 7 & -26 \\ 2 & -1 & -6 \end{array} \right|, \]

ist die Determinante

\[ p_1 = \left| \begin{array}{cc} 7 & -26\\ -1 & -6 \end{array} \right| = -68 \]

\[ p_2 = -\left| \begin{array}{cc} 3 & -26\\ 2 & -6 \end{array} \right| = -34 \]

\[ p_3 = \left| \begin{array}{cc} 3 & 7\\ 2 & -1 \end{array} \right| = -17, \]

also:

\[ x=\frac{p_1}{p_3}=4;\quad y=\frac{p_2}{p_3}=+2. \]


% ===== pdf idx 133, printed 119.0 =====
% !PRINTED 119
\weberpage{119}

2. Die drei Gleichungen:
\[
\begin{aligned}
4z-3y-2x-7&=0\\
7x-4y-1&=0\\
4z-9x-1&=0
\end{aligned}
\]
haben als Matrix
\[
\begin{vmatrix}
4 & -3 & -2 & -7\\
0 & -4 & +7 & -1\\
4 & 0 & -9 & -1
\end{vmatrix}.
\]
Deren Determinanten besitzen die Werthe:
\[
p_1=\begin{vmatrix} -3 & -2 & -7\\ -4 & 7 & -1\\ 0 & -9 & -1 \end{vmatrix},
\quad
p_2=-\begin{vmatrix} 4 & -2 & -7\\ 0 & 7 & -1\\ 4 & -9 & -1 \end{vmatrix},
\]
\[
p_3=\begin{vmatrix} 4 & -3 & -7\\ 0 & -4 & -1\\ 4 & 0 & -1 \end{vmatrix},
\quad
p_4=-\begin{vmatrix} 4 & -3 & -2\\ 0 & -4 & -7\\ 4 & 0 & -9 \end{vmatrix},
\]
also:
\[
p_1=-196,\quad p_2=-140,\quad p_3=-84,\quad p_4=-28.
\]
Daher:
\[
z=7,\quad y=5,\quad x=3.
\]

3. In den beiden folgenden Gleichungen
\[
\begin{aligned}
6x-15y-7&=0\\
-14x+35y+3&=0,
\end{aligned}
\]
deren Matrix
\[
\begin{vmatrix}
6 & -15 & -7\\
-14 & 35 & 3
\end{vmatrix}
\]
ist, wird $p_3=0$, während $p_1=200$ und $p_2=-80$ ist; die Unbekannten stellen sich also in der Form $0\cdot x=200$, $0\cdot y=80$, es sind also beide Gleichungen unvereinbar. In der That, dividirt man die erste Gleichung mit 3, die zweite mit $-7$, so kommt
\[
\begin{aligned}
2x-5y&=\frac{7}{3}\\
2x-5y&=\frac{3}{7},
\end{aligned}
\]
woraus man erkennt, dass beide Gleichungen sich widersprechen.

4. Dagegen ist in dem System von Gleichungen
\[
\begin{aligned}
3x_3-4x_1-2&=0\\
4x_2+5x_3-1&=0\\
20x_1+12x_2-7&=0,
\end{aligned}
\]
%CONT>>


% ===== pdf idx 134, printed 120.0 =====
% !PRINTED 120
\weberpage{120}
%<<CONT
deren Matrix
\[
\begin{vmatrix}
-4 & 0 & 3 & -2\\
0 & 4 & 5 & -1\\
20 & 12 & 0 & -7
\end{vmatrix}
\]
ist,
\[
p_4=-\begin{vmatrix} 20 & 12 & 0\\ 0 & 4 & 5\\ -4 & 0 & 3 \end{vmatrix}=0,
\]
ebenso
\[
p_1=p_2=p_3=0.
\]
Es ist also eine der drei Gleichungen eine Folge der zwei andern, etwa die dritte eine Folge der beiden ersten. Daher wird das Werthsystem, das alle drei Gleichungen befriedigt, bereits durch Auflösung der beiden ersten Gleichungen gefunden, indem wir die Gleichungen zuerst mit $x_4$ homogen machen, und dann nach dem bereits in Nr.\ 109 gegebenen Beispiele verfahren.

\begin{center}
\S.\,9.

\textbf{Functionaldeterminanten.}
\end{center}

114. \textit{Definition der Functionaldeterminante.} Eine zweite Reihe von Anwendungen der Lehre von den Determinanten eröffnet sich beim Studium der Eigenschaften gewisser Determinanten, die bereits von \gesperrt{Jacobi} aufgestellt und in einer Abhandlung \glqq de determinantibus functionalibus\grqq{} untersucht wurden. Bei der grossen Bedeutung derselben auf allen Gebieten der Algebra wurden sie bis in die neueste Zeit herein zum Gegenstand zahlreicher Forschungen gemacht, von denen wir im Folgenden die wichtigsten Resultate besprechen werden. Insbesondere aber werden wir uns damit beschäftigen, gewisse Aehnlichkeiten derselben mit Brüchen eingehender zu untersuchen.

Es seien gegeben $n$ Functionen
\[
y_1\;y_2\cdots y_n,
\]
deren jede von $n$ Variabeln
\[
x_1\;x_2\cdots x_n
\]
abhängig ist; etwa
\[
y_1=f_1(x_1x_2\ldots x_n),\quad y_2=f_2(x_1x_2\ldots x_n),\quad\cdots,\quad y_n=f_n(x_1x_2\ldots x_n).
\]
Man nennt alsdann die Grössen $x_i$ die unabhängigen Veränderlichen, im Gegensatz zu den abhängigen Grössen $y_i$.

Aus den sämmtlichen $n^2$ ersten partiellen Differentialquotienten dieser Functionen $y_i$ nach den $n$ unabhängigen Veränderlichen $x_i$ können wir eine Determinante
%CONT>>


% ===== pdf idx 135, printed 121.0 =====
% !PRINTED 121
\weberpage{121}
%<<CONT
\[
R=\begin{vmatrix}
\dfrac{\partial y_1}{\partial x_1} & \dfrac{\partial y_1}{\partial x_2} & \dfrac{\partial y_1}{\partial x_3} & \cdots & \dfrac{\partial y_1}{\partial x_n}\\[2ex]
\dfrac{\partial y_2}{\partial x_1} & \dfrac{\partial y_2}{\partial x_2} & \dfrac{\partial y_2}{\partial x_3} & \cdots & \dfrac{\partial y_2}{\partial x_n}\\[1.5ex]
\cdot & \cdot & \cdot & \cdot & \cdot\\[1ex]
\dfrac{\partial y_n}{\partial x_1} & \dfrac{\partial y_n}{\partial x_2} & \dfrac{\partial y_n}{\partial x_3} & \cdots & \dfrac{\partial y_n}{\partial x_n}
\end{vmatrix}
\]
bilden, und diese Determinante bezeichnet man nach \gesperrt{Jacobi} mit dem Namen: Functionaldeterminante des Systems der Functionen. Wir werden uns im Folgenden einer kürzeren Schreibweise derselben bedienen, indem wir setzen
\[
R=\begin{pmatrix} y_1 & y_2 & y_3 & \cdots & y_n\\ x_1 & x_2 & x_3 & \cdots & x_n \end{pmatrix},
\]
wobei also die untere Reihe stets jene Veränderlichen enthält, nach welchen jede Grösse der oberen Reihe als partiell differenzirt gedacht werden muss. Späterhin werden wir öfters die obere Reihe als Zähler, die untere Reihe als Nenner von $R$ bezeichnen.

115. \textit{Die Functionaldeterminante ist eine alternirende Function.} Aus der Definition der Functionaldeterminante erkennt man unmittelbar einmal, dass dieselbe ihr Zeichen ändert, sobald man $y_i$ mit $y_k$ vertauscht, dann aber auch, dass sie ihr Zeichen ändert bei Vertauschung zweier Veränderlicher $x_i$ und $x_k$. Denn dadurch gehen nur zwei Zeilen, resp. zwei Colonnen in einander über. In unserer abgekürzten Bezeichnungsweise ist also z.\,B.
\[
\begin{pmatrix} y_1 & y_2 & y_3 & y_4\\ x_1 & x_2 & x_3 & x_4 \end{pmatrix}
=-\begin{pmatrix} y_2 & y_1 & y_3 & y_4\\ x_1 & x_2 & x_3 & x_4 \end{pmatrix}
=-\begin{pmatrix} y_1 & y_2 & y_3 & y_4\\ x_1 & x_4 & x_3 & x_2 \end{pmatrix}.
\]

116. \textit{Darstellung der Minoren der Functionaldeterminante.} Wir wissen (vgl.\ Nr.\ 21 und 26), dass der Minor des Elementes $a_{ik}$ einer Determinante gegeben ist durch
\[
A_{ik}=(-1)^{i+k}(a_{11}a_{22}\ldots a_{i-1,i-1}a_{i+1,i}a_{i+2,i+1}\ldots a_{k,k-1}a_{k+1,k+1}\ldots a_{nn}).
\]
Daher ist der Minor des Elementes
\[
\frac{\partial y_i}{\partial x_k}
\]
hier symbolisch dargestellt durch
\weq{1}{A_{ik}=(-1)^{i+k}\begin{pmatrix} y_1 & y_2 & \cdots & y_{i-1}, & y_{i+1} & \cdots & y_n\\ x_1 & x_2 & \cdots & x_{k-1}, & x_{k+1} & \cdots & x_n \end{pmatrix},}
in welchem einmal alle partiellen Differentialquotienten nach $x_k$, dann
%CONT>>


% ===== pdf idx 136, printed 122.0 =====
% !PRINTED 122
\weberpage{122}
%<<CONT
aber auch jene von $y_i$ nach $x_1x_2\ldots x_n$ fehlen. Wir können indess diesen Minor auch in folgender Form schreiben:
\weq{2}{A_{ik}=(-1)^{i+k}\begin{pmatrix} y_1 & y_2 & \cdots & y_{i-1} & x_k & y_{i+1} & \cdots & y_n\\ x_1 & x_2 & \cdots & x_{k-1} & x_k & x_{k+1} & \cdots & x_n \end{pmatrix}.}
Denn in dieser Functionaldeterminante treten auch die partiellen Differentialquotienten von $x_k$ nach $x_1,x_2\ldots x_n$ genommen auf, die aber alle bis auf
\[
\frac{\partial x_k}{\partial x_k}=1
\]
verschwinden. Es sind also in derselben alle Elemente der $i^{\text{ten}}$ Colonne bis auf das $k^{\text{te}}$ Element, das den Werth 1 besitzt, null, so dass sich die in (2) dargestellte Determinante auf $A_{ik}$ reducirt.

Ganz ebenso kann man einen Minor zweiter Ordnung von $R$ etwa bezeichnen durch
\[
A_{\substack{ir\\ks}}=(-1)^v\begin{pmatrix} y_1 & y_2 & \cdots & y_{i-1} & x_k & y_{i+1} & \cdots & y_{r-1} & x_s & y_{r+1} & \cdots & y_n\\ x_1 & x_2 & \cdots & x_{k-1} & x_k & x_{k+1} & \cdots & x_{s-1} & x_s & x_{s+1} & \cdots & x_n \end{pmatrix},
\]
wo
\[
v=i+k+r+s,
\]
und analog den Minor $\varrho^{\text{ter}}$ Ordnung von $R$. Bei der Art unserer Bezeichnungsweise für eine Functionaldeterminante können wir also sagen:

\glqq Lässt man in dem Symbole
\[
R=\begin{pmatrix} y_1 & y_2 & \cdots & y_n\\ x_1 & x_2 & \cdots & x_n \end{pmatrix}
\]
zwei übereinanderstehende Grössen weg, so stellt der Rest einen Minor erster Ordnung von $R$ dar, und ferner: Sind in einer Functionaldeterminante zwei übereinanderstehende Grössen gleich, so kann man sie auch weglassen.\grqq{}

117. \textit{Productsatz für Functionaldeterminanten.} Stellen wir uns nun vor, dass auch $n$ Grössen $z_i$ gegeben seien, deren jede wiederum von allen Grössen $y_i$ abhängig ist, etwa:
\[
z_1=\varphi_1(y_1y_2\ldots y_n),\quad z_2=\varphi_2(y_1y_2\ldots y_n),\quad\cdots\quad z_n=\varphi_n(y_1y_2\ldots y_n).
\]
Diese $n$ Functionen $z_i$ besitzen eine Functionaldeterminante
\[
R'=\begin{vmatrix}
\dfrac{\partial z_1}{\partial y_1} & \dfrac{\partial z_1}{\partial y_2} & \dfrac{\partial z_1}{\partial y_3} & \cdots & \dfrac{\partial z_1}{\partial y_n}\\[2ex]
\dfrac{\partial z_2}{\partial y_1} & \dfrac{\partial z_2}{\partial y_2} & \dfrac{\partial z_2}{\partial y_3} & \cdots & \dfrac{\partial z_2}{\partial y_n}\\[1.5ex]
\cdot & \cdot & \cdot & \cdot & \cdot\\[1ex]
\dfrac{\partial z_n}{\partial y_1} & \dfrac{\partial z_n}{\partial y_2} & \dfrac{\partial z_n}{\partial y_3} & \cdots & \dfrac{\partial z_n}{\partial y_n}
\end{vmatrix}
=\begin{pmatrix} z_1 & z_2 & z_3 & \cdots & z_n\\ y_1 & y_2 & y_3 & \cdots & y_n \end{pmatrix}.
\]


% ===== pdf idx 137, printed 123.0 =====
% !PRINTED 123
\weberpage{123}

Denken wir uns nun aber in diese Functionen $z_i = \boldsymbol{\varphi}_i(y)$ die Werthe
von $y_k = f_k(x)$ eingetragen, so können wir auch noch eine zweite
Functionaldeterminante der Functionen $z_i$

\[
R'' = \left|
\begin{array}{cccc}
\dfrac{\partial z_1}{\partial x_1} & \dfrac{\partial z_1}{\partial x_2} & \dfrac{\partial z_1}{\partial x_3} \cdots & \dfrac{\partial z_1}{\partial x_n} \\[2ex]
\dfrac{\partial z_2}{\partial x_1} & \dfrac{\partial z_2}{\partial x_2} & \cdots\cdots & \dfrac{\partial z_2}{\partial x_n} \\[2ex]
\cdot & \cdot & \cdot\ \cdot\ \cdot\ \cdot\ \cdot & \cdot \\[2ex]
\dfrac{\partial z_n}{\partial x_1} & \dfrac{\partial z_n}{\partial x_2} & \cdots\cdots & \dfrac{\partial z_n}{\partial x_n}
\end{array}
\right|
=
\begin{pmatrix} z_1 & z_2 & z_3 & \ldots & z_n \\ x_1 & x_2 & x_3 & \ldots & x_n \end{pmatrix}
\]

bilden, und man kann leicht zeigen: Die Functionaldeterminante des
Systems der Grössen $z_i$ als Functionen von $x$ gedacht, ist gleich dem
Producte der Functionaldeterminante, welche zu den Systemen

\[
y_i = f_i(x)
\]

und

\[
z_i = \boldsymbol{\varphi}_i(y)
\]

gehören; oder symbolisch

\weq{1}{
\begin{pmatrix} y_1 & y_2 & y_3 & \ldots & y_n \\ x_1 & x_2 & x_3 & \ldots & x_n \end{pmatrix}
\begin{pmatrix} z_1 & z_2 & \ldots & z_n \\ y_1 & y_2 & \ldots & y_n \end{pmatrix}
=
\begin{pmatrix} z_1 & z_2 & z_3 & \ldots & z_n \\ x_1 & x_2 & x_3 & \ldots & x_n \end{pmatrix}.
}

In der That, bildet man das Product der beiden Determinanten $R$
und $R'$, so erhält man (wenn man vorher noch die Determinante $R$
transponirt)

\[
\left|
\begin{array}{cccc}
\dfrac{\partial z_1}{\partial y_1} & \dfrac{\partial z_1}{\partial y_2} & \cdots & \dfrac{\partial z_1}{\partial y_n} \\[2ex]
\dfrac{\partial z_2}{\partial y_1} & \dfrac{\partial z_2}{\partial y_2} & \cdots & \dfrac{\partial z_2}{\partial y_n} \\[2ex]
\cdot & \cdot & \cdots & \cdot \\[2ex]
\dfrac{\partial z_n}{\partial y_1} & \cdots\cdots & & \dfrac{\partial z_n}{\partial y_n}
\end{array}
\right|
\cdot
\left|
\begin{array}{cccc}
\dfrac{\partial y_1}{\partial x_1} & \dfrac{\partial y_2}{\partial x_1} & \dfrac{\partial y_3}{\partial x_1} \cdots & \dfrac{\partial y_n}{\partial x_1} \\[2ex]
\dfrac{\partial y_1}{\partial x_2} & \dfrac{\partial y_2}{\partial x_2} & \cdots\cdots & \dfrac{\partial y_n}{\partial x_2} \\[2ex]
\cdot & \cdot & \cdot\ \cdot\ \cdot\ \cdot & \cdot \\[2ex]
\dfrac{\partial y_1}{\partial x_n} & \dfrac{\partial y_2}{\partial x_n} & \cdots\cdots & \dfrac{\partial y_n}{\partial x_n}
\end{array}
\right|
\]

\[
=
\left|
\begin{array}{lll}
\dfrac{\partial z_1}{\partial y_1}\dfrac{\partial y_1}{\partial x_1}+\dfrac{\partial z_1}{\partial y_2}\dfrac{\partial y_2}{\partial x_1}+\dfrac{\partial z_1}{\partial y_3}\dfrac{\partial y_3}{\partial x_1}+\cdots\dfrac{\partial z_1}{\partial y_n}\dfrac{\partial y_n}{\partial x_1}, & \cdots, & \dfrac{\partial z_1}{\partial y_1}\dfrac{\partial y_1}{\partial x_n}+\dfrac{\partial z_1}{\partial y_2}\dfrac{\partial y_2}{\partial x_n}+\cdots\dfrac{\partial z_1}{\partial y_n}\dfrac{\partial y_n}{\partial x_n} \\[2ex]
\dfrac{\partial z_2}{\partial y_1}\dfrac{\partial y_1}{\partial x_1}+\dfrac{\partial z_2}{\partial y_2}\dfrac{\partial y_2}{\partial x_1}+\cdots+\dfrac{\partial z_2}{\partial y^n}\dfrac{\partial y_n}{\partial x_1}, & \cdots, & \dfrac{\partial z_2}{\partial y_1}\dfrac{\partial y_1}{\partial x_n}+\cdots\cdots\dfrac{\partial z_2}{\partial y_n}\dfrac{\partial y_n}{\partial x_n} \\[2ex]
\cdot\ \cdot\ \cdot\ \cdot\ \cdot\ \cdot\ \cdot\ \cdot\ \cdot\ \cdot, & & \cdot\ \cdot\ \cdot\ \cdot\ \cdot\ \cdot\ \cdot\ \cdot \\[2ex]
\dfrac{\partial z_n}{\partial y_1}\dfrac{\partial y_1}{\partial x_1}+\dfrac{\partial z_n}{\partial y_2}\dfrac{\partial y_2}{\partial x_1}+\cdots+\dfrac{\partial z_n}{\partial y_n}\dfrac{\partial y_n}{\partial x_1}, & \cdots, & \dfrac{\partial z_n}{\partial y_1}\dfrac{\partial y_1}{\partial x_n}+\dfrac{\partial z_n}{\partial y_2}\dfrac{\partial y_2}{\partial x_n}+\cdots\dfrac{\partial z_n}{\partial y_n}\dfrac{\partial y_n}{\partial x_n}
\end{array}
\right|.
\]
% ?? uncertain: printed exponent "y^n" on ∂z_2/∂y^n in the second row of the multiplied-out determinant (row 1 of expansion) reproduced as printed; may be a printing slip for ∂z_2/∂y_n

Die Elemente der Determinante rechts sind aber nichts anderes als
die partiellen Differentialquotienten der Function $z_i$, wenn man sich
%CONT>>


% ===== pdf idx 138, printed 124.0 =====
% !PRINTED 124
\weberpage{124}
%<<CONT
in denselben die Veränderlichen $y_i$ durch ihre Functionen $f_1(x)$ ersetzt
denkt; denn nach den Regeln der Differentialrechnung ist

\weq{2}{
\dfrac{\partial z_k}{\partial x_i} = \dfrac{\partial z_k}{\partial y_1}\cdot\dfrac{\partial y_1}{\partial x_i} + \dfrac{\partial z_k}{\partial y_2}\cdot\dfrac{\partial y_2}{\partial x_i} + \dfrac{\partial z_k}{\partial y_3}\cdot\dfrac{\partial y_3}{\partial x_i} + \cdots \dfrac{\partial z_k}{\partial y_n}\cdot\dfrac{\partial y_n}{\partial x_i}.
}

Damit ist aber der aufgestellte Satz bewiesen.

118. \textit{Specieller Fall $z_i = x_i$.} Sind die Functionen $z_i$ nichts
anderes als die inversen Functionen von $y_i = f_i(x_1\, x_2\, x_3 \ldots x_n)$, also
$z_i = \varphi(y_1\, y_2 \ldots y_n)$ mit $x_i$ übereinstimmend, dann ist

\[
\dfrac{\partial z_k}{\partial x_k} = \dfrac{\partial x_k}{\partial x_k} = 1, \quad \text{und} \quad \dfrac{\partial z_i}{\partial x_k} = \dfrac{\partial x_i}{\partial x_k} = 0 \quad (i \gtrless k),
\]

demnach auch die Functionaldeterminante

\[
\begin{pmatrix} z_1 & z_2 & z_3 & \ldots & z_n \\ x_1 & x_2 & x_3 & \ldots & x_n \end{pmatrix}
=
\begin{pmatrix} x_1 & x_2 & x_3 & \ldots & x_n \\ x_1 & x_2 & x_3 & \ldots & x_n \end{pmatrix}
= 1.
\]

Die Relation (1) in Nr. 117 geht alsdann über in

\[
\begin{pmatrix} y_1 & y_2 & y_3 & \ldots & y_n \\ x_1 & x_2 & x_3 & \ldots & x_n \end{pmatrix}
\begin{pmatrix} x_1 & x_2 & x_3 & \ldots & x_n \\ y_1 & y_2 & y_3 & \ldots & y_n \end{pmatrix}
= 1.
\]

119. \textit{Berechnung der Differentialquotienten $\dfrac{\partial x_i}{\partial y_k}$.} In dem zweiten
Factor der letzten Gleichung treten die Differentialquotienten $\dfrac{\partial x_i}{\partial y_k}$ auf,
deren Werthe wir uns auf zwei verschiedenen Wegen verschaffen
können. Einmal, indem wir uns in der inversen Function $x_i = \varphi_i(y)$
die Veränderlichen $y$ durch die Functionen in $x_k$ ersetzt denken, und
diese inverse Function nach $x_k$ differentiiren. Wir erhalten dann die
$n^2$ Relationen

\[
(1) \quad\cdot\quad \dfrac{\partial x_i}{\partial x_k} = \dfrac{\partial x_i}{\partial y_1}\cdot\dfrac{\partial y_1}{\partial x_k} + \dfrac{\partial x_i}{\partial y_2}\cdot\dfrac{\partial y_2}{\partial x_k} + \cdots \dfrac{\partial x_i}{\partial y_n}\cdot\dfrac{\partial y_n}{\partial x_k} = 1 \text{ oder } 0,
\]

je nachdem $i = k$, oder $i$ von $k$ verschieden ist. Aus diesen kann
man $n$ Gleichungen herausgreifen und für die $n$ Unbekannten $\dfrac{\partial x_i}{\partial y_k}$ die
Werthe aus ihnen berechnen. Rascher gelangen wir aber auf folgen-
dem Wege zum Ziele. Wir können $\dfrac{\partial x_i}{\partial y_k}$ darstellen durch die Functional-
determinante

\[
\dfrac{\partial x_i}{\partial y_k} = \begin{pmatrix} y_1 & y_2 & \ldots & y_{k-1} & x_i & y_{k+1} & \ldots & y_n \\ y_1 & y_2 & \ldots & y_{k-1} & y_k & y_{k+1} & \ldots & y_n \end{pmatrix},
\]

gemäss den vorhin in Nr. 116 gemachten Bemerkungen. Multipliciren
wir aber dieselbe nun mit der Functionaldeterminante
%CONT>>


% ===== pdf idx 139, printed 125.0 =====
% !PRINTED 125
\weberpage{125}
%<<CONT
\[
R = \begin{pmatrix} y_1 & y_2 & y_3 & \ldots & y_n \\ x_1 & x_2 & x_3 & \ldots & x_n \end{pmatrix},
\]

so erhalten wir gemäss unserm Satze über das Product zweier Functional-
determinanten Nr. 117 (1)

\weq{2}{
\begin{pmatrix} y_1 & y_2 & y_3 & \ldots & y_n \\ x_1 & x_2 & x_3 & \ldots & x_n \end{pmatrix}
\begin{pmatrix} y_1 & y_2 & \ldots & y_{k-1} & x_i & y_{k+1} & \ldots & y_n \\ y_1 & y_2 & \ldots & y_{k-1} & y_k & y_{k+1} & \ldots & y_n \end{pmatrix}
=
\begin{pmatrix} y_1 & y_2 & \ldots & y_{k-1} & x_i & y_{k+1} & \ldots & y_n \\ x_1 & x_2 & \ldots & x_{k-1} & x_k & x_{k+1} & \ldots & x_n \end{pmatrix}.
}

Die Determinante rechts ist aber nichts anderes als der Minor des
Elementes $\dfrac{\partial y_k}{\partial x_i}$ in der Determinante $R$ [vgl. Nr. 116 (2)]. Wir haben
somit für den partiellen Differentialquotienten

\weq{3}{
R \cdot \dfrac{\partial x_i}{\partial y_k} = \dfrac{\partial R}{\partial\!\left(\dfrac{\partial y_k}{\partial x_i}\right)} = A_{ik},
}

oder

\[
\dfrac{\partial x_i}{\partial y_k} = \dfrac{A_{ki}}{R}.
\]

119. \textit{Die Functionaldeterminante als ein Product partieller Differential-
quotienten.} % [sic] repeated section number "119." — the preceding heading in Nr. 118/119 on the previous page was also numbered 119.
Wenn $n$ Veränderliche $y$ gegeben sind, deren jede abhängig
ist von $n$ Veränderlichen $x$, so steht es uns frei, irgend $n$ von diesen
$2n$ Grössen als Function der übrigen zu betrachten. Nehmen wir
der Einfachheit halber zunächst drei Functionen $y_i$ abhängig von drei
Grössen $x_i$, etwa

\[
(1) \quad y_1 = f_1(x_1\, x_2\, x_3), \quad y_2 = f_2(x_1\, x_2\, x_3), \quad y_3 = f_3(x_1\, x_2\, x_3),
\]

so können wir festsetzen, $y_1$ als abhängig zu betrachten von $x_1\, x_2\, x_3$;
dagegen $y_2$ als Function von $y_1\, x_2\, x_3$, $y_3$ als Function von $y_1\, y_2$ und $x_3$,
etwa als Functionen:

\[
(2) \quad y_1 = f(x_1\, x_2\, x_3), \quad y_2 = \varphi(y_1\, x_2\, x_3), \quad y_3 = \psi(y_1\, y_2\, x_3).
\]

Bilden wir nun die Functionaldeterminante des gegebenen Systems (1),
so ist dieselbe

\[
R = \begin{pmatrix} y_1 & y_2 & y_3 \\ x_1 & x_2 & x_3 \end{pmatrix}.
\]

Bilden wir andererseits die partiellen Differentialquotienten von $y_1$
nach $x_1$, $y_2$ nach $x_2$, $y_3$ nach $x_3$, indem wir die durch das System (2)
ausgedrückte Abhängigkeit zu Grunde legen, so können wir dieselben
gemäss unserer Uebereinkunft betreffs der abgekürzten Schreibweise
(vgl. Nr. 116) darstellen durch
%CONT>>


% ===== pdf idx 140, printed 126.0 =====
% !PRINTED 126
\weberpage{126}
%<<CONT
\[
(\mathrm{I}) \quad \dfrac{\partial y_1}{\partial x_1} = \begin{pmatrix} y_1 & x_2 & x_3 \\ x_1 & x_2 & x_3 \end{pmatrix}, \quad
\dfrac{\partial y_2}{\partial x_2} = \begin{pmatrix} y_1 & y_2 & x_3 \\ y_1 & x_2 & x_3 \end{pmatrix}, \quad
\dfrac{\partial y_3}{\partial x_3} = \begin{pmatrix} y_1 & y_2 & y_3 \\ x_1 & x_2 & x_3 \end{pmatrix},
\]
% ?? uncertain: faint marks above $x_1\, x_2$ in the denominator of the third fraction in (I); read as plain $x_1\, x_2\, x_3$

wonach die Variabeln des Zählers als Functionen der Grössen im
Nenner stets nach denen des Nenners zu differenziren sind, und je zwei
übereinanderstehende gleiche unterdrückt werden können. Nun ist
aber [nach Nr. 117 (1)]

\[
(a) \quad \begin{pmatrix} y_1 & y_2 & y_3 \\ y_1 & y_2 & x_3 \end{pmatrix} \begin{pmatrix} y_1 & y_2 & x_3 \\ y_1 & x_2 & x_3 \end{pmatrix} = \begin{pmatrix} y_1 & y_2 & y_3 \\ y_1 & x_2 & x_3 \end{pmatrix} = D;
\]

und ebenso

\[
(b) \quad \begin{pmatrix} y_1 & y_2 & y_3 \\ y_1 & x_2 & x_3 \end{pmatrix} \begin{pmatrix} y_1 & x_2 & x_3 \\ x_1 & x_2 & x_3 \end{pmatrix} = \begin{pmatrix} y_1 & y_2 & y_3 \\ x_1 & x_2 & x_3 \end{pmatrix} = R.
\]

Substituiren wir also den Werth von $D$ in die Relation (b), so kommt

\weq{3}{
\begin{pmatrix} y_1 & y_2 & y_3 \\ x_1 & x_2 & x_3 \end{pmatrix}
=
\begin{pmatrix} y_1 & y_2 & y_3 \\ y_1 & y_2 & x_3 \end{pmatrix}
\begin{pmatrix} y_1 & y_2 & x_3 \\ y_1 & x_2 & x_3 \end{pmatrix}
\begin{pmatrix} y_1 & x_2 & x_3 \\ x_1 & x_2 & x_3 \end{pmatrix}
}

oder

\[
R = \dfrac{\partial y_1}{\partial x_1} \cdot \dfrac{\partial y_2}{\partial x_2} \cdot \dfrac{\partial y_3}{\partial x_3}, \tag{II}
\]

wobei also jeder der drei Differentialquotienten der rechten Seite in
anderem Sinne zu nehmen ist, nämlich in jenem, der durch die Vor-
stellung

\[
y_1 = f(x_1\, x_2\, x_3), \quad y_2 = \varphi(y_1\, x_2\, x_3), \quad y_3 = \psi(y_1\, y_2\, x_3)
\]

bestimmt ist.

Von dieser Beziehung (II) wird bei der Transformation dreifacher
Integrale Gebrauch gemacht. Um nämlich in dem dreifachen Integrale

\[
\int F(y_1\, y_2\, y_3)\,dy_1\,dy_2\,dy_3
\]

an Stelle der Veränderlichen $y_i$ drei neue Veränderliche $x_i$ einzuführen,
welche mit den alten durch die Relation (1) verbunden sind, kann
man sich zunächst, indem man etwa die Integration nach $y_3$ zuerst
auszuführen gedenkt, $y_3$ als Function $\psi(y_1\, y_2\, x_3)$ vorstellen. In diesem
Falle darf man alsdann das Differential $dy_3$ ersetzen durch $\dfrac{\partial \psi}{\partial x_3}\,dx_3$.
In dem neuen dreifachen Integral kann man die Integration mit $y_2$
beginnen; stellt man sich daher $y_2$ als Function von $\varphi(y_1\, x_2\, x_3)$ vor,
so darf man $dy_2$ ersetzen durch $\dfrac{\partial \varphi}{\partial x_2}\,dx_2$. So kann man endlich an
Stelle von $dy_1$ auch $\dfrac{\partial f_1}{\partial x_1}\,dx_1$ einführen und erhält zum Schlusse:

\[
\int F(y_1\, y_2\, y_3)\,dy_1\,dy_2\,dy_3 = \int F(f_1(x_1\, x_2\, x_3)\cdots)\dfrac{\partial f_1}{\partial x_1}\cdot\dfrac{\partial \varphi}{\partial x_2}\cdot\dfrac{\partial \psi}{\partial x_3}\,dx_1\,dx_2\,dx_3
\]
\[
= \int \Phi(x_1\, x_2\, x_3)\cdot R \cdot dx_1\,dx_2\,dx_3.
\]


% ===== pdf idx 141, printed 127.0 =====
% !PRINTED 127
\weberpage{127}

Allgemein ist, wenn $y_1 y_2 \ldots y_n$ Functionen der Grössen $x_1 x_2 \ldots x_n$ sind, die in (3) gegebene Relation dargestellt durch
\begin{gather*}
R = \begin{pmatrix} y_1 & y_2 & y_3 & \cdots & y_n \\ x_1 & x_2 & x_3 & \cdots & x_n \end{pmatrix}
= \begin{pmatrix} y_1 & y_2 & \cdots & y_n \\ y_1 & y_2 & \cdots & x_n \end{pmatrix}
\begin{pmatrix} y_1 & y_2 & \cdots & y_{n-1} & x_n \\ y_1 & y_2 & \cdots & x_{n-1} & x_n \end{pmatrix} \cdots \\
\cdots \begin{pmatrix} y_1 & y_2 & \cdots & y_{n-2} & x_{n-1} & x_n \\ y_1 & y_2 & \cdots & x_{n-2} & x_{n-1} & x_n \end{pmatrix}
\begin{pmatrix} y_1 & x_2 & x_3 & \cdots & x_{n-1} & x_n \\ x_1 & x_2 & x_3 & \cdots & \cdots & x_n \end{pmatrix}.
\end{gather*}

120. \textit{Differentialquotienten von Minoren der Functionaldeterminante.} Bezeichnen wir die Minoren erster Ordnung von $R$ mit $A_{11} A_{12} \ldots A_{nn}$, dann sind die Minoren der Elemente, welche sich in der $k^{\mathrm{ten}}$ Zeile befinden, dargestellt durch
\[
A_{k1} = \pm \dfrac{\partial R}{\partial\left(\dfrac{\partial y_k}{\partial x_1}\right)}, \quad
A_{k2} = \mp \dfrac{\partial R}{\partial\left(\dfrac{\partial y_k}{\partial x_2}\right)}, \quad \ldots, \quad
A_{kn} = \pm \dfrac{\partial R}{\partial\left(\dfrac{\partial y_k}{\partial x_n}\right)}.
\]
Bildet man nun für jeden solchen Minor $A_{ki}$ den dem Index $i$ entsprechenden Differentialquotienten nach $x_i$: $\dfrac{\partial A_{ki}}{\partial x_i}$, so existirt der Satz:

\glqq Die Summe der partiellen Differentialquotienten
\[
\frac{\partial A_{ki}}{\partial x_i}, \quad i = 1, 2 \ldots n,
\]
ist identisch Null.\grqq{} Der Beweis lässt sich einfach rechnerisch führen\footnote{In der Invariantentheorie würde dieser Satz selbstverständlich erscheinen. Denn da dort die Functionaldeterminante symbolisch dargestellt ist durch
\[
R = (abc \ldots q)\, a_x^{n-1} b_x^{m-1} \cdots q_x^{s-1},
\]
so würde die Summe der entsprechenden Differentialquotienten aller Minoren einer Zeile nur den Ausdruck ergeben
\[
\sum \left[ (abc \ldots a)\, a_x^{n-2} b_x^{n-1} \cdots \right],
\]
von dem man aus den Klammerfactoren unmittelbar erkennt, dass er verschwindet. Ist z.\,B. $R = (abc)\, a_x^6 b_x^6 c_x^6$, so ist die Summe der Differentialquotienten aller Minoren der Elemente in der zweiten Zeile (immer von Zahlenfactoren abgesehen)
\[
- \frac{\partial(a_2 c_3)\, a_x^6 c_x^6}{\partial x_1} + \frac{\partial(a_1 c_3)\, a_x^6 c_x^6}{\partial x_2} - \frac{\partial(a_1 c_2)\, a_x^6 c_x^6}{\partial x_3}
= -(a_1 a_2 c_3)\, a_x^5 c_x^6 - (c_1 a_2 c_3)\, a_x^6 c_x^5.
\]
In dem Ausdruck rechts ist aber jedes Glied Null, da die Determinanten $(a_1 a_2 c_3)$, $(c_1 a_2 c_3)$ gleiche Zeilen enthalten.}. Der Minor $A_{ik}$ ist eine Determinante von $(n-1)$ Colonnen und $(n-1)$ Zeilen, und in jedem Elemente $\dfrac{\partial y_l}{\partial x_m}$ derselben tritt im All-
%CONT>>


% ===== pdf idx 142, printed 128.0 =====
% !PRINTED 128
\weberpage{128}

%<<CONT
gemeinen noch $x_i$ auf. Wir differenziren also die Determinante $A_{ki}$ nach jedem ihrer Elemente und multipliciren jedesmal mit dem Differentialquotienten dieses Elementes nach $x_i$. Das Resultat ist dargestellt durch
\[
\frac{\partial A_{ki}}{\partial x_i} = \sum_{l,m} \frac{\partial A_{ki}}{\partial\left(\dfrac{\partial y_l}{\partial x_m}\right)} \cdot \frac{\partial^2 y_l}{\partial x_i \partial x_m} = \sum_{l,m} A_{\substack{k\,l\\ i\,m}} \cdot \frac{\partial^2 y_l}{\partial x_i \partial x_m}.
\]
Geben wir nun dem Index $i$ alle Werthe von 1 bis $n$ und addiren die so erhaltenen Relationen, so gelangen wir zur Gleichung
\[
\sum_{i=1}^{i=n} \frac{\partial A_{ki}}{\partial x_i} = \sum_{i,l,m} A_{\substack{k\,l\\ i\,m}} \frac{\partial^2 y_l}{\partial x_i \partial x_m}.
\]
Je zwei Glieder der Summe rechts sind aber bekanntlich gleich und entgegengesetzt; denn es ist $A_{\substack{k\,l\\ i\,m}} = -A_{\substack{k\,l\\ m\,i}}$ (vgl. Nr. 49); die rechte Seite verschwindet demnach; also auch die linke Seite dieser Gleichung.

121. \textit{Die Functionaldeterminante eines Systemes von Functionen $y$, zwischen denen eine Relation besteht, verschwindet.} Wenn $n$ Functionen $y_i$ von $n$ unabhängigen Variabeln $x_i$ gegeben sind, so kann man sich die Aufgabe stellen, $(n-1)$ der Grössen $x_i$ der Reihe nach zu eliminiren, etwa $x_1$ bis $x_{n-1}$, d.\,h. also eine der Functionen, etwa $y_1$, durch die übrigen $y_2 y_3 \ldots y_n$ und die Veränderliche $x_n$ darzustellen. Dabei kann der Fall eintreten, dass bei diesem Eliminationsprocess sich zugleich die $n^{\mathrm{te}}$ Veränderliche $x_n$ mit hinweghebt, d.\,h. dass sich $y_1$ darstellen lässt in Function der übrigen Grössen $y_2 y_3 \ldots y_n$ allein, etwa in der Form
\weq{1}{y_1 = F(y_i).}


% ===== pdf idx 143, printed 129.0 =====
% !PRINTED 129
\weberpage{129}

Es besteht alsdann in diesem Falle eine Relation zwischen den Functionen $y$, die wir mit
\weq{2}{\Theta(y) = y_1 - F(y_2 y_3 \cdots y_n) = 0}
bezeichnen wollen. Wir können dann den Satz beweisen: \glqq Besteht zwischen den $n$ Functionen $y$ eine Relation $\Theta(y) = 0$, so ist ihre Functionaldeterminante $R$ identisch Null.\grqq{} Es möge dabei erwähnt sein, dass die Relation $\Theta(y) = 0$ als Function der Grössen $y$ keineswegs eine Identität ist, wie schon aus der Art und Weise hervorgeht, auf welcher wir zu ihr gelangt sind.

122. \textit{Beweis des Satzes.} Um nun den Satz zu beweisen, denken wir uns die Veränderlichen $y_i$ durch ihre Functionen $f_i(x)$ ersetzt; die entstehende Relation in $x_i$ ist aber dann identisch null, und daher auch alle ihre partiellen Differentialquotienten
\[
\frac{\partial \Theta}{\partial x_1} = 0, \quad \frac{\partial \Theta}{\partial x_2} = 0, \quad \cdots \frac{\partial \Theta}{\partial x_n} = 0.
\]
Dieselben geben also, da die Grössen $x_i$ nur in Functionen $y_i$ auftreten, Veranlassung zu den $n$ Relationen:
\begin{gather*}
\frac{\partial \Theta}{\partial y_1} \cdot \frac{\partial y_1}{\partial x_1} + \frac{\partial \Theta}{\partial y_2} \cdot \frac{\partial y_2}{\partial x_1} + \cdots \frac{\partial \Theta}{\partial y_n} \cdot \frac{\partial y_n}{\partial x_1} = 0 \\
\frac{\partial \Theta}{\partial y_1} \cdot \frac{\partial y_1}{\partial x_2} + \frac{\partial \Theta}{\partial y_2} \cdot \frac{\partial y_2}{\partial x_2} + \cdots \frac{\partial \Theta}{\partial y_n} \cdot \frac{\partial y_n}{\partial x_2} = 0 \\
\cdots\cdots\cdots\cdots\cdots\cdots\cdots\cdots\cdots\cdots \\
\frac{\partial \Theta}{\partial y_1} \cdot \frac{\partial y_1}{\partial x_n} + \frac{\partial \Theta}{\partial y_2} \cdot \frac{\partial y_2}{\partial x_n} + \cdots \frac{\partial \Theta}{\partial y_n} \cdot \frac{\partial y_n}{\partial x_n} = 0.
\end{gather*}
Wir können dieselben als $n$ homogene Gleichungen mit $n$ Unbekannten $\dfrac{\partial \Theta}{\partial y_k}$ betrachten; dann ist das Werthsystem, das diese Gleichungen befriedigt (nach Nr. 97) dargestellt durch
\[
R \cdot \frac{\partial \Theta}{\partial y_k} = 0.
\]
Da nun aber die Grössen $\dfrac{\partial \Theta}{\partial y_k}$ nicht sämmtlich verschwinden können, — es ist ja beispielsweise $\dfrac{\partial \Theta}{\partial y_1} = 1$ wegen der Relation (1) — so muss der andere Factor verschwinden; d.\,h. in diesem Falle ist $R = 0$.

123. \textit{Umkehrung.} Ist nun von vornherein die Functionaldeter-
%CONT>>


% ===== pdf idx 144, printed 130.0 =====
% !PRINTED 130
\weberpage{130}

%<<CONT
minante $R$ identisch Null, so ist es immer möglich, bei geeignetem Vorgehen aus den Functionen $y_i$ alle Grössen $x$ zu eliminiren, d.\,h. in diesem Falle existirt eine Relation $\Theta(y) = 0$.

Um den Satz zu beweisen nehmen wir an, dass dessen Richtigkeit für $(n-1)$ Functionen von $(n-1)$ Variabeln bereits dargethan sei, und haben nur zu zeigen, dass er dann auch für $n$ Functionen mit $n$ Variabeln giltig ist. Denn für eine Function $y$ mit einer Variabeln $x$ ist er evident, da aus $\dfrac{\partial f}{\partial x} = 0$, $y = c$ hervorgeht.

Wir unterscheiden hierbei zwei Fälle, indem wir zunächst annehmen, dass neben $R = 0$ nicht zugleich alle Minoren erster Ordnung verschwinden sollen, dass also etwa
\[
A_{nn} = \begin{pmatrix} y_1 & y_2 & \cdots & y_{n-1} \\ x_1 & x_2 & \cdots & x_{n-1} \end{pmatrix} \gtrless 0
\]
sei. Benutzen wir alsdann den in Nr. 117 aufgestellten Productsatz, so ist, weil wir $A_{nn}$ auch darstellen können durch
\[
A_{nn} = \begin{pmatrix} y_1 & y_2 & \cdots & y_{n-1} & x_n \\ x_1 & x_2 & \cdots & x_{n-1} & x_n \end{pmatrix}
\]
die Functionaldeterminante $R$ gegeben durch
\[
\begin{pmatrix} y_1 & y_2 & y_3 & \cdots & y_n \\ x_1 & x_2 & x_3 & \cdots & x_n \end{pmatrix}
= \begin{pmatrix} y_1 & y_2 & \cdots & y_{n-1} & x_n \\ x_1 & x_2 & \cdots & x_{n-1} & x_n \end{pmatrix}
\begin{pmatrix} y_1 & y_2 & y_3 & \cdots & y_{n-1} & y_n \\ y_1 & y_2 & \cdots & \cdots & y_{n-1} & x_n \end{pmatrix}.
\]
Nun ist die linke Seite nach Voraussetzung gleich Null, der erste Factor rechts dagegen nach Voraussetzung von Null verschieden. Daher muss nothwendig der zweite Factor verschwinden. Derselbe aber ist nichts anderes als der Differentialquotient von $y_n$ nach $x_n$, wobei $y_n$ gemäss der Bedeutung unserer Schreibweise als Function von
\[
y_1 y_2 \ldots y_{n-1} x_n
\]
zu betrachten ist. Wenn aber dieser Differentialquotient
\[
\frac{\partial y_n}{\partial x_n} = 0
\]
sein soll, dann kann $y_n$ nicht mehr abhängig sein von $x_n$, d.\,h. $y_n$ ist eine Function der übrigen Grössen
\[
y_1 y_2 \ldots y_{n-1}
\]
allein und es existirt somit nothwendig eine Relation zwischen den Grössen $y_i$.

124. \textit{Zweiter Fall.} Nehmen wir nun an, es verschwinden neben $R = 0$ auch alle ihre Minoren erster Ordnung, deren einer gegeben ist durch
\[
\begin{pmatrix} y_1 & y_2 & \cdots & y_{k-1} & y_{k+1} & \cdots & y_n \\ x_1 & x_2 & \cdots & \cdot & \cdot & \cdots & x_{n-1} \end{pmatrix} = 0.
\]
Dieser Ausdruck stellt aber nichts anderes dar, als eine Functionaldeterminante der $(n-1)$ Functionen $y_i$, die nur von $(n-1)$ der Veränderlichen $x_i$ als abhängig zu betrachten sind. Da unser Satz nach Voraussetzung für diesen Fall gilt, so besteht zwischen den $n-1$ Functionen
\[
y_1 y_2 \ldots y_{k-1} y_{k+1} \ldots y_n
\]
eine Relation, die natürlich auch noch die Grösse $x_n$ enthalten kann. Solche Relationen erhalten wir $n$ an der Zahl je nach der Grösse $y$,
%CONT>>


% ===== pdf idx 145, printed 131.0 =====
% !PRINTED 131
\weberpage{131}
%<<CONT
die wir in der ersten Zeile des symbolischen Ausdruckes unterdrücken. Dieselben können nicht durchwegs identisch gleich sein, da ja jede derselben eine andere Grösse $y$ nicht enthält. Es müssen daher mindestens zwei der Relationen verschiedener Art sein; sie seien
\[ F_1 = 0 \quad \text{und} \quad F_2 = 0. \]
Eliminiren wir aus diesen beiden aber die Grösse $x_n$, so erhalten wir eine einzige Relation $\Theta(y) = 0$, wie behauptet war.

125. \textit{Beispiel.} Sind $n$ lineare Functionen $y_i$ mit $n$ Unbekannten gegeben
\[ a_{i1} x_1 + a_{i2} x_2 + a_{i3} x_3 + \cdots a_{in} x_n = y_i, \]
so ist nach Nr. 111
\[ p_{n+1} x_i = p_i = y_1 A_{i1} + y_2 A_{i2} + \cdots y_n A_{in}. \]
Nun ist $p_{n+1}$ nichts anderes als die Functionaldeterminante der gegebenen Functionen $y_i$; verschwindet sie, so besteht zwischen den gegebenen Functionen die lineare Relation
\[ y_1 A_{i1} + y_2 A_{i2} + \cdots y_n A_{in} = 0 \]
und wir wissen, dass in diesem Falle eine dieser Functionen $y_i$ eine Folge der übrigen ist.

\textbf{§ 10. Der grösste gemeinsame Factor.}

126. \textit{Euclid'sches Verfahren für zwei ganze Zahlen $a$ und $b$.} Wenn zwei Zahlen $a$ und $b$ gegeben sind, $a > b$, so kann man sich die Aufgabe stellen, den grössten gemeinsamen Factor der beiden Zahlen zu bestimmen. Das Euclid'sche Verfahren, diese Aufgabe zu lösen, besteht einfach darin, dieselbe auf leichtere Aufgaben zurückzuführen. Man dividirt nämlich bekanntlich mit $b$ in $a$ und erhält so einestheils einen Quotienten $q$ und anderntheils einen Rest $b_1$, so dass die Beziehung besteht:
\[ a = bq + b_1. \]
Daraus erkennt man, dass jeder gemeinsame Factor — also auch der grösste —, welcher in $a$ und $b$ enthalten ist, auch in $b_1$ enthalten sein muss, d. h. dass wir ihn auch ermitteln können, indem wir den grössten gemeinsamen Factor von den beiden kleineren Zahlen $b$ und $b_1$ suchen. Verfahren wir hier auf die gleiche Weise und setzen diese Methode fort, bis wir endlich zu einem Reste $b_\varrho$ gelangen, über welchen hinaus keiner mehr auftritt, also $b_{\varrho+1} = 0$ ist, so erhalten wir $\varrho$ Divisionsgleichungen
%CONT>>


% ===== pdf idx 146, printed 132.0 =====
% !PRINTED 132
\weberpage{132}
%<<CONT
\weq{I}{
\begin{aligned}
a &= bq + b_1 \\
b &= b_1 q_1 + b_2 \\
b_1 &= b_2 q_2 + b_3 \\
b_2 &= b_3 q_3 + b_4 \\
& \cdots\cdots\cdots\cdots \\
b_{\varrho-2} &= b_{\varrho-1} q_{\varrho-1} + b_\varrho.
\end{aligned}
}
In denselben nehmen die Reste $b_1 b_2 \ldots b_\varrho$ fortwährend ab und der letzte Rest $b_\varrho$ ist der gemeinschaftliche Factor aller vorhergehenden und also insbesondere auch von $a$ und $b$. Man hat dieses Verfahren auch mit dem Namen Kettenbruchsverfahren bezeichnet.

127. \textit{Darstellung der Reste $b_1 \ldots b_\varrho$ als lineare Functionen von $a$ und $b$.} Der Algorithmus (I) lehrt uns aber, dass die Grössen $b_1 b_2 b_3 \ldots b_\varrho$ sich ausdrücken lassen in der Form:
\weq{II}{ b_\mu = A_\mu a + B_\mu b. }
Der Beweis wird einfach durch Rechnung geführt. Es ist:
\[ 1)\quad b_1 = a - bq = A_1 a + B_1 b, \]
wobei:
\[ A_1 = 1, \quad B_1 = -q; \]
\[ 2)\quad b_2 = b - q_1(a - bq) = A_2 a + B_2 b, \]
wobei:
\[ A_2 = -q_1, \quad B_2 = +(1 + qq_1); \]
\[ 3)\quad b_3 = (a - bq) - q_2[(b - q_1(a - bq))] = A_3 a + B_3 b, \]
wobei:
\[ A_3 = +(1 + q_1q_2), \quad B_3 = -(q + q_2 + qq_1q_2); \]
\[ 4)\quad b_4 = [b - q_1(a - bq)] - q_3[(a - bq) - q_2(b - q_1(a - bq))] = A_4 a + B_4 b, \]
wobei:
\[ A_4 = -(q_1 + q_3 + q_1q_2q_3), \]
\[ B_4 = +(1 + qq_1 + qq_3 + q_2q_3 + qq_1q_2q_3) \text{ u. s. w.} \]
Allgemein ist also:
\[ b_\mu = A_\mu a + B_\mu b, \]
wobei für $A_\mu$ und $B_\mu$ die Recursionsformeln gelten:
\[ A_\mu = A_{\mu-2} - q_{\mu-1} A_{\mu-1} \]
\[ B_\mu = B_{\mu-2} - q_{\mu-1} B_{\mu-1}. \]


% ===== pdf idx 147, printed 133.0 =====
% !PRINTED 133
\weberpage{133}
Man folgert aus der Relation (II) direct, dass jeder gemeinsame Factor von $a$ und $b$ auch Factor des Restes $b_\mu$ ist. Denkt man sich mit demselben die Relation (II) dividirt, also $a$ und $b$ relativ prim zu einander, so kann man sie leicht direct beweisen mit Hilfe zweier einfacher Sätze der Lehre von den Kettenbrüchen. Sind nämlich $\frac{\beta_\mu}{\alpha_\mu}$ die Näherungswerthe des Bruches $\frac{b}{a}$, so ist nach bekannten Sätzen:
\weq{1}{ \frac{b}{a} = \frac{\beta_\mu b_{\mu-1} + \beta_{\mu-1} b_\mu}{\alpha_\mu b_{\mu-1} + \alpha_{\mu-1} b_\mu}; }
\weq{2}{ \alpha_\mu \beta_{\mu-1} - \alpha_{\mu-1} \beta_\mu = (-1)^\mu. }
Subtrahirt man von Gleichung (1) $\frac{\beta_\mu}{\alpha_\mu}$ auf beiden Seiten, so ergiebt sich nach leichter Umformung mit Hilfe der Relation (2)
\weq{III}{ (-1)^\mu \cdot b_\mu = -\beta_\mu a + \alpha_\mu b, }
eine Relation, die direct mit (II) übereinstimmt und woraus man erkennt, dass die Coefficienten $A_\mu$ und $B_\mu$ nichts anderes sind als der Zähler und Nenner (absolut genommen) des $\mu^{\mathrm{ten}}$ Näherungswerthes von $\frac{b}{a}$.

128. \textit{Grösster gemeinschaftlicher Theiler zweier Functionen $f(\overset{m}{x})$ und $\varphi(\overset{n}{x})$.} Ganz in derselben Weise wie vorhin der Theiler zweier Zahlen ermittelt wurde, kann man verfahren, um den grössten gemeinschaftlichen Theiler zweier Functionen $f(\overset{m}{x})$ und $\varphi(\overset{n}{x})$ zu bestimmen.
Sei
\[
\left.\begin{aligned}
f(\overset{m}{x}) &= a_0 x^m + a_1 x^{m-1} + a_2 x^{m-2} + \cdots a_m \\
\varphi(\overset{n}{x}) &= b_0 x^n + b_1 x^{n-1} + b_2 x^{n-2} + \cdots b_n
\end{aligned}\right\} \quad m > n,
\]
wobei also der Index über der Variabeln hier wie im Folgenden stets den Grad der betreffenden Function in $x$ bedeuten soll und $a_m$ und $b_n$ von Null verschieden sind. — Dividiren wir $f(x)$ durch $\varphi(x)$, mit dem entstehenden Rest $f_1$ wieder in $\varphi$, mit dem weitern Rest $f_2$ in $f_1$ u. s. w., dann entstehen wieder eine Reihe von Divisionsgleichungen wie vorhin. Dabei wollen wir jene Grössen erst als Reste definiren, in denen der Coefficient der höchsten Potenz von $x$ die Zahl 1 ist, und ferner wollen wir zunächst alle Coefficienten $a_i$ und $b_i$ als Variable ansehen; dann
%CONT>>


% ===== pdf idx 148, printed 134.0 =====
% !PRINTED 134
\weberpage{134}
%<<CONT
ist jeder folgende Rest immer um einen Grad niedriger als der vorhergehende.
Unsere Divisionsgleichungen werden demnach:
\[ f(\overset{m}{x}) = \overset{m-n}{q}(x)\, \overset{n}{\varphi}(x) + c_1 \overset{n-1}{f_1}(x) \]
\[ \overset{n}{\varphi}(x) = \overset{1}{q_1}(x)\, \overset{n-1}{f_1}(x) + c_2 \overset{n-2}{f_2}(x) \]
\[ \overset{n-1}{f_1}(x) = \overset{1}{q_2}(x)\, \overset{n-2}{f_2}(x) + c_3 \overset{n-3}{f_3}(x) \]
\[ \overset{n-2}{f_2}(x) = \overset{1}{q_3}(x)\, \overset{n-3}{f_3}(x) + c_4 \overset{n-4}{f_4}(x) \]
\[ \cdots\cdots\cdots\cdots\cdots\cdots \]
\[ \overset{3}{f_{n-3}}(x) = \overset{1}{q_{n-2}}(x)\, \overset{2}{f_{n-2}}(x) + c_{n-1} \overset{1}{f_{n-1}}(x) \]
\[ \overset{2}{f_{n-2}}(x) = \overset{1}{q_{n-1}}(x)\, \overset{1}{f_{n-1}}(x) + c_n \overset{0}{f_n}(x), \]
wobei
\[ \overset{0}{f_n}(x) = 1 \]
und die Constanten $c_1 c_2 \ldots c_n$ davon herrühren, dass der Coefficient der höchsten Potenz jedes Restes gleich 1 genommen wurde. Wir wollen nun wieder aus diesen Divisionsgleichungen die Reste $f_1$ bis $f_n$ in der Form darstellen
\weq{I}{ \overset{n-\mu}{f_\mu}(x) = \overset{\mu-1}{A_\mu}(x)\, \overset{m}{f}(x) + \overset{m-n+\mu-1}{B_\mu}(x)\, \overset{n}{\varphi}(x). }
Es ist:
\weq{1}{ \overset{n-1}{f_1}(x) = \frac{1}{c_1}\left\{ \overset{m}{f}(x) - \overset{m-n}{q}(x)\overset{n}{\varphi}(x) \right\} = \overset{0}{A_1}(x)\,\overset{m}{f}(x) + \overset{m-n}{B_1}(x)\,\overset{n}{\varphi}(x), }
wobei
\[ \overset{0}{A_1}(x) = +\frac{1}{c_1}, \quad \overset{n-m}{B_1}(x) = -\frac{\overset{m-n}{q}(x)}{c_1}; \]
\weq{2}{ \overset{n-2}{f_2}(x) = \frac{1}{c_2}\left\{ \overset{n}{\varphi}(x) - \overset{1}{q_1}(x)\left[ \overset{0}{A_1}(x)\overset{m}{f}(x) + \overset{m-n}{B_1}(x)\overset{n}{\varphi}(x) \right] \right\} = \overset{1}{A_2}(x)\,\overset{m}{f}(x) + \overset{m-n+1}{B_2}(x)\,\overset{n}{\varphi}(x), }
wobei
\[ \overset{1}{A_2}(x) = -\frac{1}{c_1c_2}\cdot \overset{1}{q_1}(x), \quad \overset{m-n+1}{B_2}(x) = +\frac{1}{c_1c_2}\left( c_1 + \overset{m-n}{q}(x)\overset{1}{q_1}(x) \right); \]
\weq{3}{ \overset{n-3}{f_3}(x) = \frac{1}{c_3}\left\{ \overset{n-1}{f_1}(x) - \overset{1}{q_2}(x)\left[ \overset{1}{A_2}(x)\overset{m}{f}(x) + \overset{m-n+1}{B_2}(x)\overset{n}{\varphi}(x) \right] \right\} }
%CONT>>


% ===== pdf idx 149, printed 135.0 =====
% !PRINTED 135
\weberpage{135}
\S.\,10. Der grösste gemeinsame Factor.

%<<CONT
\[
= \frac{1}{c_3}\left\{ \overset{0}{A_1}(x)\overset{m}{f}(x) + \overset{m-n}{B_1}(x)\varphi(x) - \overset{1}{q_2}(x)\overset{1}{A_2}(x)\overset{m}{f}(x) - \overset{1}{q_2}(x)\overset{m-n+1}{B_2}(x)\overset{n}{\varphi}(x) \right\} = \overset{2}{A_3}(x)\overset{m}{f}(x) + \overset{m-n+2}{B_3}(x)\overset{n}{\varphi}(x),
\]
wobei
\[
\overset{2}{A_3}(x) = \frac{1}{c_3}\left[ \overset{0}{A_1}(x) - \overset{1}{q_2}(x)\overset{1}{A_2}(x) \right] = +\frac{1}{c_1c_2c_3}\left[ c_2 + \overset{1}{q_1}(x)\overset{1}{q_2}(x) \right],
\]
\[
\overset{2}{B_3}(x) = \frac{1}{c_3}\left[ \overset{m-n}{B_1}(x) - \overset{1}{q_2}(x)\overset{m-n+1}{B_2}(x) \right] = -\frac{1}{c_1c_2c_3}\left[ c_2\overset{m-n}{q}(x) + c_1\overset{1}{q_2}(x) + \overset{n-m}{q}(x)\overset{1}{q_1}(x)\overset{1}{q_2}(x) \right] \text{u.\,s.\,w.}
\]

Man sieht hieraus, dass die Grössen $A_\mu(x)$ und $B_\mu(x)$ sich wieder
durch die Recursionsformeln darstellen lassen:

\[
\overset{\mu-1}{A_\mu}(x) = \frac{1}{c_\mu}\left\{ \overset{\mu-3}{A_{\mu-2}}(x) - \overset{1}{q_{\mu-1}}(x)\overset{\mu-2}{A_{\mu-1}}(x) \right\}
\]
\[
\overset{m-n+\mu-1}{B_\mu}(x) = \frac{1}{c_\mu}\left\{ \overset{m-n+\mu-3}{B_{\mu-2}}(x) - \overset{1}{q_{\mu-1}}(x)\overset{m-n+\mu-2}{B_{\mu-1}}(x) \right\}.
\]
Insbesondere erhalten wir als Ausdruck des letzten Restes

\weq{n}{\overset{0}{f_n}(x) = 1 = \overset{n-1}{A_n}(x)\overset{m}{f}(x) + \overset{m-1}{B_n}(x)\overset{n}{\varphi}(x).}

129. \textit{Exposition der hier zu lösenden Aufgaben.} Die rechte Seite
der Gleichungen (1) bis ($n$) ist stets eine ganze Function der $x$, ebenso
wie die linke; dagegen ist sie als Function der Coefficienten betrachtet
im Allgemeinen eine gebrochene Function derselben, so lange $a_i$ und $b_i$
variabel sind. Bezeichnen wir in dem Ausdrucke für $\overset{n-\mu}{f_\mu}(x)$ den Haupt\-nenner der ganzen rechten Seite mit $D_{n-\mu}$, so können wir die
Gleichung (I) in Nr. 128 auch schreiben

\weq{II}{D_{n-\mu}\overset{n-\mu}{f_\mu}(x) = \overset{\mu-1}{\overline{A}_\mu}(x)\overset{m}{f}(x) + \overset{m-n+\mu-1}{\overline{B}_\mu}(x)\overset{n}{\varphi}(x),}

wo nun $\overline{A}_\mu$ und $\overline{B}_\mu$ andere Functionen der Grössen $a$ und $b$ sind,
also vorhin $A_\mu$ und $B_\mu$. In dieser Relation (II) sind aber nun sowohl
$D_{n-\mu}\overset{n-\mu}{f}(x)$, als auch $\overline{A}_\mu$, $\overline{B}_\mu$, $f(x)$ und $\varphi(x)$ ganze Functionen der
Coefficienten wie der Variabeln. Der Nenner $D_{n-\mu}$ kann verschwinden;
in diesem Falle tritt überhaupt kein Rest vom Grade $n-\mu$ in $x$ auf,
%CONT>>


% ===== pdf idx 150, printed 136.0 =====
% !PRINTED 136
\weberpage{136}
Zweiter Theil. Anwendungen der Determinantentheorie.

%<<CONT
sondern es folgt auf den Rest vom Grade $n-\mu+1$ ein Rest vom
Grade $n-\mu-1$, oder von noch niedrigerem Grade, wie wir sogleich
auch noch ausführlicher sehen werden. Ist dagegen der Nenner
$D_{n-\mu}\,\overset{\geqq}{\leqq}\,0$, so können wir uns die Aufgabe stellen, den zugehörigen
Rest $f_\mu$, oder vielmehr das Product $D_\mu \cdot f_\mu$ zu bestimmen.

Es kann ferner eintreten, dass alle Nenner vom letzten $D_0$ an
gerechnet bis zum $\nu^{\text{ten}}$ $D_{n-\nu+1}$ verschwinden, $D_{n-\nu}$ aber nicht; dann
verschwinden nach dem eben Erwähnten auch alle Reste

\[
f_n, f_{n-1}\cdots f_{n-\nu+1}.
\]
Der erste nicht verschwindende Rest ist $f_{n-\nu}$ vom Grade $\nu$ in $x$, und
dieser ist nothwendig der grösste gemeinsame Factor von $\overset{m}{f}(x)$ und
$\overset{n}{\varphi}(x)$. Daraus erkennt man: Es sind $\nu$ Bedingungen nöthig

\[
D_0 = 0, \quad D_1 = 0, \quad D_2 = 0 \cdots, \quad D_{n-\nu+1} = 0,
\]
damit $f$ und $\varphi$ einen gemeinsamen Factor vom Grade $\nu$ in $x$ besitzen
und dieser gemeinsame Factor ist der letzte Rest, der nicht ver\-schwindet.

Wir haben also hier zwei Aufgaben zu lösen: einmal die Be\-dingungen aufzustellen, dass ein gemeinsamer Factor vom Grade $\nu$
vorhanden ist; sodann, diesen Factor selbst zu ermitteln.

130. \textit{Berechnung von $D_0$.} Um zur Lösung der ersten Aufgabe zu
gelangen, berechnen wir zuerst

\[
D_0 = D_0\overset{0}{f_n}(x),
\]
indem wir dabei die Identität benutzen

\weq{I}{1 = \overset{n-1}{A_n}(x)\overset{m}{f}(x) + \overset{m-1}{B_n}(x)\overset{n}{\varphi}(x).}

Wir folgen hierbei ganz den Entwicklungen der diesbezüglichen Arbeit
\gesperrt{Nöther}'s, veröffentlicht in \gesperrt{Faà di Bruno}'s Einleitung zur Invarianten\-theorie pag. 58 (deutsche Uebersetzung von \gesperrt{Walter}. Teubner 1881).
Die Bestimmung von $A_n$ und $B_n$, welche wir als Functionen vom
Grade $n-1$, bezw. $m-1$ darstellen können durch

\[
\begin{aligned}
A_n &= p_0 x^{n-1} + p_1 x^{n-2} + \cdots + p_{n-1}\\
B_n &= q_0 x^{m-1} + q_1 x^{m-2} + \cdots + q_{n-1}
\end{aligned}
\]
geschieht aus der Identität (I) durch Coefficientenvergleichung. Man
erhält aus ihr für die $m+n$ Unbekannten

\[
p_0, p_1, \cdots p_{n-1}, \quad q_0, q_1 \cdots q_{m-1}
\]
die $m+n$ linearen Gleichungen:
%CONT>>


% ===== pdf idx 151, printed 137.0 =====
% !PRINTED 137
\weberpage{137}
%<<CONT
\S.\,10. Der grösste gemeinsame Factor.

\[
\text{(II)}\quad
\left\{
\begin{aligned}
0 &= a_0 p_0 \cdots\cdots\cdots\cdots\cdots\cdots\cdots\cdots\cdots\cdots + b_0 q_0\\
0 &= a_1 p_0 + a_0 p_1 \cdots\cdots\cdots\cdots\cdots\cdots\cdots + b_1 q_0 + b_0 q_1\\
0 &= a_2 p_0 + a_1 p_1 + a_0 p_2 \cdots\cdots\cdots + b_2 q_0 + b_1 q_1 + b_2 q_2\\
&\cdots\cdots\cdots\cdots\cdots\cdots\cdots\cdots\cdots\cdots\cdots\cdots\cdots\cdots\\
0 &= a_n p_0 + a_{n-1} p_1 + \cdots\cdots\cdots\cdots\cdots\cdots\cdots\cdots\cdots\\
&\qquad\cdots + b_n q_0 + b_{n-1} q_1 + b_{n-2} q_2\cdots + b_{n-m+1} q_{m-1}\\
0 &= a_{n+1} p_0 + a_n p_1 + \cdots\cdots\cdots\cdots\cdots\cdots\cdots\cdots\cdots\cdots\cdots\\
&\qquad\qquad + b_n q_1 + b_{n-1} q_2 \cdots + b_{n-m+2} q_{m-1}\\
&\cdots\cdots\cdots\cdots\cdots\cdots\cdots\cdots\cdots\cdots\cdots\cdots\cdots\cdots\\
0 &= a_m p_0 + a_{m-1} p_1 + \cdots + a_{m-n+2} p_{n-2} + a_{m-n+1} p_{n-1}\\
&\qquad\qquad + b_n q_{m-n} + b_{n-1} q_{m-n+1}\cdots\\
0 &= a_m p_1 + \cdots + a_{m-n+3} p_{n-2} + a_{m-n+2} p_{n-1}\\
&\qquad\qquad + b_n q_{m-n+1} + b_{n-1} q_{m-n+2}\cdots\\
&\cdots\cdots\cdots\cdots\cdots\cdots\cdots\cdots\cdots\cdots\cdots\cdots\cdots\cdots\\
0 &= a_m p_{n-2} + a_{m-1} p_{n-1}\\
&\qquad\qquad + b_n q_{m-2} + b_{n-1} q_{m-1}\\
1 &= a_m p_{n-1}\\
&\qquad\qquad + b_n q_{m-1}
\end{aligned}
\right.
\]
Löst man diese Gleichungen nach $p_i$ und $q_i$ auf, so haben die Werthe
von $p_i$ und $q_i$ durchwegs den Nenner:

% ?? uncertain: exact column spacing/blank-cell layout of the D_0 determinant below is approximated from the scan; entries and dot-rows reproduced as printed.
\[
D_0 = \left|
\begin{array}{cccccccccc}
a_0 & 0 & 0 & \cdots & b_0 & 0 & 0 & 0 & \cdots & 0\\
a_1 & a_0 & 0 & \cdots & b_1 & b_0 & 0 & 0 & \cdots & 0\\
a_2 & a_1 & a_0 & \cdots & b_2 & b_1 & b_0 & 0 & \cdots & 0\\
\cdot & \cdot & \cdot & \cdot & \cdot & \cdot & \cdot & \cdot & \cdot & \cdot\\
 & & & & & b_n & b_{n-1} & b_{n-2} & \cdots & \\
 & & & & & 0 & b_n & b_{n-1} & \cdots & \\
a_m & a_{m-1} & a_{m-2} & \cdots & 0 & 0 & b_n & & & \\
0 & a_m & a_{m-1} & \cdots & 0 & 0 & 0 & & & \\
0 & 0 & a_m & \cdots & 0 & 0 & 0 & & & \\
\cdot & \cdot & \cdot & \cdot & \cdot & \cdot & \cdot & \cdot & \cdot & \cdot\\
 & & & & & & & & & b_{n-1}\\
\cdot & \cdot & \cdot & \cdot & \cdot & \cdot & \cdot & \cdot & \cdot & \cdot\\
0 & 0 & 0 & \cdots & a_m & 0 & 0 & \cdots & \cdots & b_n
\end{array}
\right.
\]


% ===== pdf idx 152, printed 138.0 =====
% !PRINTED 138
\weberpage{138}
Zweiter Theil. Anwendungen der Determinantentheorie.

Transponirt man diese Determinante, so erhält man

% ?? uncertain: exact column spacing/blank-cell layout of the transposed determinant below is approximated from the scan; one entry in row 8 prints a bare "b" without subscript, reproduced as printed.
\[
D_0 = \left|
\begin{array}{cccccccccc}
a_0 & a_1 & a_2 & \cdots & a_m & 0 & 0 & \cdots & & 0\\
0 & a_0 & a_1 & \cdots & \cdots & a_m & 0 & \cdots & & 0\\
0 & 0 & a_0 & \cdots & \cdots & \cdots & \cdots & \cdots & & 0\\
\cdot & \cdot & \cdot & \cdot & \cdot & \cdot & \cdot & \cdot & \cdot & \cdot\\
0 & \cdots & \cdots & \cdots & a_0 & a_1 & a_2 & a_3 & \cdots & a_m\\
b_0 & b_1 & b_2 & \cdots & \cdots & b_n & 0 & 0 & \cdots & 0\\
0 & b_0 & b_1 & \cdots & \cdots & \cdots & b_n & 0 & \cdots & 0\\
0 & 0 & b_0 & \cdots & \cdots & \cdots & \cdots & b & \cdots & 0\\
\cdot & \cdot & \cdot & \cdot & \cdot & \cdot & \cdot & \cdot & \cdot & \cdot\\
0 & 0 & \cdots & \cdots & \cdots & \cdots & \cdots & b_0 & b_1 \cdots & b_n
\end{array}
\right| = R_{f,\varphi}.
\]
131. \textit{Definition der Resultante von $f(x)$ und $\varphi(x)$.} Dieser Nenner
kann im Allgemeinen nicht verschwinden, wenn die Identität (I) Nr. 130
besteht, da ja sonst nicht die Gleichungen (II) Nr. 130 zusammen
existiren könnten (vergleiche Nr. 112), die aus dieser Identität hervor\-gehen. Wir nennen diese Determinante $D_0$ \glqq die Resultante\grqq{} der beiden
Functionen $f(x)$ und $\varphi(x)$ und bezeichnen sie als solche mit $R_{f,\varphi}$.
Sie ist eine ganze Function der Coefficienten und offenbar auch homogen
sowohl in $a_i$ als $b_i$, und zwar in $a_i$ vom Grade $n$, in $b_i$ vom Grade $m$,
wie man aus der Anzahl der Zeilen in $D_0$ erkennt, die einestheils die
Elemente $a_i$, andertheils die Elemente $b_i$ enthalten.

Ersetzt man in der Identität (I) alle Coefficienten $p$ und $q$ von
$A_n$ und $B_n$ durch jene Werthe, welche sich aus dem System (II) Nr. 130
von linearen Gleichungen ergeben und multiplicirt mit dem Nenner
aller dieser Coefficienten die Identität, so erhält man die andere Form
derselben:

\[
D_0 = \overset{n-1}{\overline{A}_n}(x)\overset{m}{f}(x) + \overset{m-1}{\overline{B}_n}(x)\varphi(x) = R_{f,\varphi}.
\]
Diese Identität kann man auch direct aus der Determinante $R_{f,\varphi}$ selbst
herleiten, wenn man in dieser Determinante die erste Colonne mit
$x^{n-m-1}$, die zweite mit $x^{n-m-2}$ u. s. w. sich multiplicirt denkt und
zur letzten addirt. Es wird dann:

\[
R_{f,\varphi} = \left|
\begin{array}{ccccc}
a_0 & a_1 & a_2 & \cdots & x^{n-1}\cdot f(x)\\
0 & a_1 & a_2 & \cdots & x^{n-2}\cdot f(x)\\
\cdot & \cdot & \cdot & \cdot & \cdot\\
0 & 0 & \cdots & a_{m-1}, & f(x)\\
b_0 & b_1 & b_2 & \cdots & x^{m-1}\varphi(x)\\
0 & b_1 & b_2 & \cdots & x^{m-2}\varphi(x)\\
\cdot & \cdot & \cdot & \cdot & \cdot\\
 & & & b_{n-1}, & \varphi(x)
\end{array}
\right.
\]


% ===== pdf idx 153, printed 139.0 =====
% !PRINTED 139
\weberpage{139}

Entwickelt man diese Determinante nach der letzten Colonne, so erhält man gerade

\[
R_{f,\varphi} = \overset{m}{f}(x)\left\{\bar p_0 x^{n-1}+\bar p_1 x^{n-2}\cdots \bar p_{n-1}\right\}
\]
\[
{}+ \overset{m}{\varphi}(x)\left\{\bar q_0 x^{m-1}+\bar q_1 x^{m-2}+\cdots \bar q_{m-1}\right\}
\]
oder
\[
R_{f,\varphi} = \overset{n-1}{\bar A_n}(x)\,\overset{m}{f}(x) + \overset{m-1}{\bar B_n}(x)\cdot \overset{n}{\varphi}(x).
\]
Sind die Co\"efficienten $a_i$ und $b_i$ so beschaffen, dass $D_0$ identisch verschwindet, dann wird diese Identität

\[
0 = \overset{n-1}{\bar A_n}(x)\,\overset{m}{f}(x) + \overset{m-1}{\bar B_n}(x)\,\overset{n}{\varphi}(x).
\]
Die Functionen $f$ und $\varphi$ haben alsdann einen gemeinschaftlichen linearen Factor $f_{n-1}$, vorausgesetzt dass -- wie wir gleich noch eingehender untersuchen werden -- $D_1$ nicht verschwindet.

132. \textit{Berechnung von $D_1$.} Unsere nächste Aufgabe ist, den Nenner $D_1$ zu berechnen. Wir benutzen dazu wiederum die Identität (I) Nr. 128, indem wir dort $n-1$ für $\mu$ substituiren; man erhält

\weq{II}{\overset{1}{f}_{n-1}(x) = \overset{n-2}{A}_{n-1}(x)\,\overset{m}{f}(x) + \overset{m-2}{B}_{n-1}(x)\,\overset{n}{\varphi}(x) = x + r.}

Die Functionen $A_{n-1}$, $B_{n-1}$ kann man in der Form annehmen:

\[
A_{n-1} = p_0 x^{n-2} + p_1 x^{n-3}\cdots + p_{n-2}
\]
\[
B_{n-1} = q_0 x^{m-2} + q_1 x^{m-3}\cdots + q_{m-2}.
\]
Vergleicht man zunächst auf beiden Seiten die Co\"efficienten von $x^{m+n-2}$, $x^{m+n-3}\ldots$ bis herab zu $x^2$, so erhält man für die $m+n-2$ Unbekannten:

\[
p_0\ p_1\ p_2 \cdots p_{n-2},\ q_0,\ q_1 \cdots q_{m-2}
\]
$(m+n-3)$ lineare und homogene Gleichungen:

\weq{I}{
\left\{
\begin{gathered}
a_0\,p_0 + b_0\,q_0 = 0\\
a_1\,p_0 + a_0\,p_1 + b_1\,q_0 + b_1\,q_1 = 0\\
a_2\,p_0 + a_1\,p_1 + a_0\,p_2 + b_2\,q_0 + b_1\,q_1 + b_0\,q_2 = 0\\
\cdot\ \cdot\ \cdot\ \cdot\ \cdot\ \cdot\ \cdot\ \cdot\ \cdot\ \cdot\ \cdot\ \cdot\ \cdot\ \cdot\\
a_m\,p_0 + a_{m-1}\,p_1 + \cdots + a_{m-n+2}\,p_{n-2} + \cdots + b_{m-n+2}\,q_{n-2} = 0\\
a_m\,p_1 + \cdots + a_{m-n+3}\,p_{n-2} + \cdots\cdots\cdots + b_{m-n+3}\,q_{n-2} = 0\\
\cdot\ \cdot\ \cdot\ \cdot\ \cdot\ \cdot\ \cdot\ \cdot\ \cdot\ \cdot\ \cdot\ \cdot\ \cdot\ \cdot\\
a_m\,p_{n-4} + a_{m-1}\,p_{n-3} + a_{m-2}\,p_{n-2} + b_n\,q_{n-4} + b_{n-1}\,q_{n-3} + b_{n-2}\,q_{n-2} = 0,
\end{gathered}
\right.
}


% ===== pdf idx 154, printed 140.0 =====
% !PRINTED 140
\weberpage{140}

Hiezu tritt noch als weitere Gleichung jene, welche man durch Comparation des Co\"efficienten der ersten Potenz von $x$ in der Identität (II) gewinnt, d. i.

\weq{2}{1 = a_m\,p_{n-3} + a_{m-1}\,p_{n-2} + b_n\,q_{m-3} + b_{n-1}\,q_{n-2}.}

Beide Systeme zusammen liefern $(m+n-2)$ Gleichungen mit $m+n-2$ Unbekannten $p_i$ und $q_i$; die letzteren ergeben sich also aus ihnen in Bruchform mit dem gemeinsamen Nenner

\[
D_1 =
\left|
\begin{array}{cccccccccc}
a_0 & a_1 & a_2 & \cdot & \cdot & \cdot & a_m & 0 & \cdot & 0\\
0 & a_0 & a_1 & \cdot & \cdot & \cdot & \cdot & a_m & \cdot & 0\\
0 & 0 & a_0 & \cdot & \cdot & \cdot & \cdot & \cdot & \cdot & 0\\
\cdot & \cdot & \cdot & \cdot & \cdot & \cdot & \cdot & \cdot & \cdot & a_m\\
0 & \cdot & \cdot & \cdot & \cdot & \cdot & \cdot & a_0 & a_1 & a_2 \cdots a_{m-1}\\
b_0 & b_1 & b_2 & \cdot & \cdot & \cdot & b_n & 0 & 0 & \cdot\ 0\\
0 & b_0 & b_1 & \cdot & \cdot & \cdot & \cdot & b_n & 0 & \cdot\ 0\\
0 & 0 & b_0 & \cdot & \cdot & \cdot & \cdot & \cdot & \cdot & b_n\ \cdot\ 0\\
\cdot & \cdot & \cdot & \cdot & \cdot & \cdot & \cdot & \cdot & \cdot & \cdot\\
\cdot & \cdot & \cdot & \cdot & \cdot & \cdot & \cdot & \cdot & \cdot & b_n\\
0 & \cdot & \cdot & \cdot & \cdot & \cdot & \cdot & b_0 & b_1 & b_2\ b_3 \cdots b_{n-1}
\end{array}
\right|.
\]

Derselbe ist nichts anderes als eine Unterdeterminante zweiter Ordnung von $D_0$; man erhält sie aus $D_0$, wenn in ihr die zwei letzten Colonnen und die letzte Zeile sowohl der Elemente $a$, als die letzte der Elemente $b$ unterdrückt werden. Die Unbekannten $p_i$ und $q_i$ sind dann, wenn man die Minoren der letzten Colonne in $D_1$ mit

\[
D_{10}\ D_{11} \ldots D_{1,n-2}\ \Delta_{10}\ \Delta_{11} \ldots \Delta_{1,m-2}
\]
bezeichnet, dargestellt durch

\[
p_i = \frac{D_{1i}}{D_1},\quad q_i = \frac{\Delta_{1i}}{D_1}.
\]
Trägt man diese Werthe in die Gleichung ein, welche durch Comparation des Co\"efficienten von $x^0$ aus der Identität (II) hervorgeht, also in

\weq{3}{r = a_m\,p_{n-2} + b_n\,q_{n-2},}

so erhält man auch den Werth dieser Grösse $r$, woraus hervorgeht, dass sie nur auf eine einzige Art bestimmt werden kann. Der Rest $x + r = \overset{1}{f}_{n-1}(x)$ ist somit eindeutig ermittelt.


% ===== pdf idx 155, printed 141.0 =====
% !PRINTED 141
\weberpage{141}

Die Determinante $D_1$ ist jene Grösse, deren Verschwinden ausdrückt, dass kein Rest vom ersten Grade in $x$ existirt. Denn tragen wir in die Identität

\[
\overset{1}{f}_{n-1}(x) = x + r = \overset{n-2}{A}_{n-1}(x)\,\overset{m}{f}(x) + \overset{m-2}{B}_{n-1}(x)\,\overset{n}{\varphi}(x)
\]
die Werthe von $p_i$ und $q_i$ ein und multipliciren mit dem gemeinsamen Nenner $D_1$, so erhalten wir

\[
D_1 \cdot \overset{1}{f}_{n-1}(x) = \overset{n-2}{\bar A}_{n-1}(x)\,\overset{m}{f}(x) + \overset{m-2}{B}_{n-1}(x)\,\overset{n}{\varphi}(x),
\]
woraus die Behauptung direct hervorgeht.

133. \textit{Berechnung von $\overset{1}{f}_{n-1}(x)$.} Verschwinden $D_0$ und $D_1$ gleichzeitig, so besitzen also die beiden Functionen $f(x)$ und $\varphi(x)$ einen gemeinsamen Factor zweiten Grades. Verschwindet nur $D_0$ allein, so existirt, wie schon in Nr. 130 erwähnt, ein gemeinsamer Factor ersten Grades, der erhalten wird, wenn wir in der Identität (II) Nr. 132 den Ausdruck auf der rechten Seite in zwei Glieder zusammenfassen, eines mit dem Factor $x^1$, das andere mit dem Factor $x^0$, also diese Identität (II) in der Form schreiben

\[
\overset{1}{f}_{(n-1)}(x) = M \cdot x^1 + N \cdot x^0
\]
und nun in $M$ und $N$ die Grössen $p_i$ und $q_i$ vermöge der Relationen

\[
p_i = \frac{D_{1i}}{D_1},\quad q_i = \frac{\Delta_{1i}}{D_1}
\]
eliminiren.

134. \textit{Berechnung von $D_2$.} Wir gehen wieder aus von der Identität

\weq{III}{\overset{.2}{f}_{n-2}(x) = \overset{n-3}{A}_{n-2}(x)\,\overset{m}{f}(x) + \overset{m-3}{B}_{n-2}(x)\,\overset{n}{\varphi}(x) = x^2 + rx + s}

und vergleichen darin zunächst die Co\"efficienten von $x^{m+n-1}$, $x^{m+n-2}\ldots$ bis $x^3$, wodurch wir $(m+n-5)$ homogene lineare Gleichungen mit den $m+n-4$ Unbekannten

\[
p_0\ p_1\ p_2 \cdots p_{n-3},\ q_0\ q_1\ q_2 \cdots q_{m-3}
\]
erhalten, nämlich die Gleichungen

\weq{1}{
\left\{
\begin{gathered}
0 = a_0\,p_0 + b_0\,q_0\\
0 = a_1\,p_0 + p_1\,a_0 + b_1\,q_0 + b_0\,q_1\\
0 = a_2\,p_0 + a_1\,p_1 + a_0\,p_2 + b_2\,q_0 + b_1\,q_1 + b_0\,q_2\\
\cdot\ \cdot\ \cdot\ \cdot\ \cdot\ \cdot\ \cdot\ \cdot\ \cdot\ \cdot\ \cdot\ \cdot\ \cdot\ \cdot\ \cdot\ \cdot\ \cdot\ \cdot\\
0 = a_m\,p_{n-5} + a_{m-1}\,p_{n-4} + a_{m-2}\,p_{n-3} + b_n\,q_{n-5} + b_{n-1}\,q_{n-4} + b_{n-2}\,q_{n-3}.
\end{gathered}
\right.
}

Hiezu tritt als nächste Gleichung jene, welche durch Comparation der Co\"efficienten von $x^2$ entsteht, nämlich

\weq{2}{1 = a_m\,p_{n-4} + a_{m-1}\,p_{n-3} + b_n\,q_{n-4} + b_{n-1}\,q_{n-3}.}


% ===== pdf idx 156, printed 142.0 =====
% !PRINTED 142
\weberpage{142}

Beide Systeme zusammen geben die nöthige Zahl von Gleichungen, um die Grössen $p_i$ und $q_i$ zu berechnen. Der gemeinsame Nenner aller Werthe $p_i$ und $q_i$ ist die Determinante des Systems (1), nämlich

\[
D_2 =
\left|
\begin{array}{cccccccc}
a_0 & a_1 & a_2 & \cdots & a_m & 0 & 0 & \cdots 0\\
 & a_0 & a_1 & \cdots & \cdots & a_m & \cdots & \cdots\\
 \cdot & \cdot & \cdot & \cdot & \cdot & \cdot & \cdot & \cdot\\
 & & & & a_0 & a_1 & \cdots & a_{m-2}\\
b_0 & b_1 & b_2 & \cdots & b_n & 0 & \cdots & 0\\
 & b_0 & b_1 & \cdots & \cdots & b_n & \cdots & \cdots\\
 \cdot & \cdot & \cdot & \cdot & \cdot & \cdot & \cdot & \cdot\\
 & & & & b_0 & b_1 & \cdots & b_{n-2}
\end{array}
\right|,
\]
also ein Minor vierter Ordnung von $D_0$, den man erhält durch Unterdrückung der letzten vier Colonnen einestheils und der letzten Zeilenpaare mit den Elementen $a$ bezw. Elementen $b$ anderntheils.

135. \textit{Berechnung von $\overset{2}{f}_{n-2}(x)$.} Tragen wir die in Function von $a_i$ und $b_i$ berechneten Werthe von $p_i$ und $q_i$ in jene Gleichungen ein, die man durch Comparation der zweiten, ersten und nullten Potenz von $x$ aus der Identität (III) gewinnt, so werden dadurch die Co\"efficienten $r$ und $s$, und also auch der ganze Rest $\overset{2}{f}_{n-2}(x)$ eindeutig bestimmt. Zu demselben Resultate gelangt man, indem man die rechte Seite der Identität (III) Nr. 134 in drei Gliedern zusammenfasst, von der Form

\[
\overset{2}{f}_{n-2}(x) = Ax^2 + Bx + C
\]
und in den Co\"efficienten $A$, $B$ und $C$ die Grössen $p_i$ und $q_i$ eliminirt. Dabei erhält die rechte Seite durchweg den Nenner $D_2$, so dass nach Multiplication mit demselben die Gleichung entsteht

\[
D_2 \cdot \overset{2}{f}_{n-2}(x) = \bar A x^2 + \bar B x + \bar C,
\]
woraus man erkennt, dass mit $D_2 = 0$ auch kein Rest vom zweiten Grade in $x$ auftritt.

Wenn also die Reste $D_0$, $D_1$ und $D_2$ verschwinden, $D_3$ aber von Null verschieden ist, dann haben $f(x)$ und $\varphi(x)$ einen Factor dritten Grades in $x$ gemein.

136. \textit{Verallgemeinerung.} Aus den bisherigen Entwicklungen ist das Verfahren, nun zur Aufstellung des Restes vom Grade $\varrho$ in $x$ zu gelangen, leicht abzusehen. Man comparirt in der Identität

\[
\overset{\varrho}{f}_{n-\varrho}(x) = \overset{n-\varrho-1}{A}_{n-\varrho}(x)\,\overset{m}{f}(x) + \overset{m-\varrho-1}{B}_{n-\varrho}(x)\,\overset{n}{\varphi}(x)
\]


% ===== pdf idx 157, printed 143.0 =====
% !PRINTED 143
\weberpage{143}

\begin{center}
\S.\,10. Der grösste gemeinsame Factor.
\end{center}

die Coefficienten der $(m+n-\varrho-1)^{\mathrm{ten}}$ bis $\varrho^{\mathrm{ten}}$ Potenz von $x$; das giebt ein erstes Gleichungssystem. Hierzu tritt die Gleichung, die durch Comparation der $\varrho^{\mathrm{ten}}$ Potenz von $x$ selbst entsteht, zusammen $m+n-2\varrho$ Gleichungen mit $m+n-2\varrho$ Unbekannten, die sich sonach daraus berechnen lassen. Ihr gemeinsamer Nenner ist $D_\varrho$, und der Rest vom Grade $\varrho$ geht aus der Identität hervor, indem man die Werthe von $p_i$ und $q_i$ in dieselbe substituirt. Verschwinden alle Determinanten $D_0, D_1 \ldots D_\varrho$, dagegen $D_{\varrho+1}$ nicht, so existirt kein Rest von niedrigerem Grade als dem $(\varrho+1)^{\mathrm{ten}}$, und die beiden Functionen $f$ und $\varphi$ haben einen gemeinsamen Factor $(\varrho+1)^{\mathrm{ten}}$ Grades in $x$. Insbesondere haben wir gesehen, ist $D_0 = 0$ die nothwendige und hinreichende Bedingung, dass beide Functionen überhaupt einen linearen Factor gemein haben.

137. \textit{Beispiel.} Sei gegeben
\[
f_x = a_0x^4+a_1x^3+a_2x^2+a_3x+a_4,
\]
\[
\varphi_x = b_0x^3+b_1x^2+b_2x+b_3.
\]

Wir bestimmen zunächst $D_0$ aus der Identität:
\[
\begin{array}{l}
1 = (p_0x^2+p_1x+p_2)(a_0x^4+a_1x^3+a_2x^2+a_3x+a_4) \\
\hphantom{1} + (q_0x^3+q_1x^2+q_2x+q_3)(b_0x^3+b_1x^2+b_2x+b_3).
\end{array}
\]

Die Coefficientenvergleichung giebt:
\weq{1}{
\begin{array}{rcl}
0 &=& a_0p_0 + b_0q_0 \\
0 &=& a_1p_0+a_0p_1 + b_1q_0+b_0q_1 \\
0 &=& a_2p_0+a_1p_1+a_0p_2 + b_2q_0+b_1q_1+b_0q_2 \\
0 &=& a_3p_0+a_2p_1+a_1p_2 + b_3q_0+b_2q_1+b_1q_2+b_0q_3 \\
0 &=& a_4p_0+a_3p_1+a_2p_2 + b_3q_1+b_2q_2+b_1q_3 \\
0 &=& a_4p_1+a_3p_2 + b_3q_2+b_2q_3 \\
1 &=& a_4p_2 + b_3q_3
\end{array}}

Die Determinante des Systems wird, wenn man sie transponirt:
\[
D_0 =
\begin{vmatrix}
a_0 & a_1 & a_2 & a_3 & a_4 & 0 & 0 \\
0 & a_0 & a_1 & a_2 & a_3 & a_4 & 0 \\
0 & 0 & a_0 & a_1 & a_2 & a_3 & a_4 \\
b_0 & b_1 & b_2 & b_3 & 0 & 0 & 0 \\
0 & b_0 & b_1 & b_2 & b_3 & 0 & 0 \\
0 & 0 & b_0 & b_1 & b_2 & b_3 & 0 \\
0 & 0 & 0 & b_0 & b_1 & b_2 & b_3
\end{vmatrix}.
\]


% ===== pdf idx 158, printed 144.0 =====
% !PRINTED 144
\weberpage{144}

\begin{center}
Zweiter Theil. Anwendungen der Determinantentheorie.
\end{center}

Ihr Verschwinden drückt aus, dass $f$ und $\varphi$ einen gemeinsamen linearen Factor besitzen.

Ermitteln wir ferner $D_1$ aus der Identität:
\[
\begin{array}{l}
x+r = (p_0x+p_1)(a_0x^4+a_1x^3+a_2x^2+a_3x+a^4) % [sic] printed a^4, apparently for a_4
\\
\hphantom{x+r} + (q_0x^2+q_1x+q_2)(b_0x^3+b_1x^2+b_2x+b_4), % [sic] printed b_4, apparently for b_3
\end{array}
\]
indem wir zuerst die Coefficienten der Potenzen $x^5, x^4, x^3, x^2, x$ vergleichen, so kommt:
\weq{1}{
\begin{array}{rcl}
0 &=& a_0p_0 + b_0q_0 \\
0 &=& a_1p_0+a_0p_1 + b_1q_0+b_0q_1 \\
0 &=& a_2p_0+a_1p_1 + b_2q_0+b_1q_1+b_0q_2 \\
0 &=& a_3p_0+a_2p_1 + b_3q_0+b_2q_1+b_1q_2 \\
1 &=& a_4p_0+a_3p_1 + b_3q_1+b_2q_2
\end{array}}

Die Determinante des Systems ist:
\[
D_1 =
\begin{vmatrix}
a_0 & a_1 & a_2 & a_3 & a_4 \\
0 & a_0 & a_1 & a_2 & a_3 \\
b_0 & b_1 & b_2 & b_3 & b_4 \\
0 & b_0 & b_1 & b_2 & b_3 \\
0 & 0 & b_0 & b_1 & b_2
\end{vmatrix}.
\]

Sie geht also in der That aus $D_0$ hervor durch Unterdrückung der beiden letzten Colonnen sowie der letzten Zeile der Elemente $a$ und der letzten Zeile der Elemente $b$.

Ebenso erhalten wir endlich $D_2$ aus der Identität
\[
\begin{array}{l}
x^2+rx+s = p_0x\cdot(a_0x^4+a_1x^3+a_2x^2+a_3x+a_4) \\
\hphantom{x^2+rx+s} + (q_0x+q_1)(b_0x^3+b_1x^2+b_2x+b_3),
\end{array}
\]
indem wir das Gleichungssystem
\weq{1}{
\begin{array}{rcl}
0 &=& a_0p_0+b_0q_0 \\
0 &=& a_1p_0+b_1q_0+b_0q_1 \\
1 &=& a_2p_0+b_2q_0+b_1q_1
\end{array}}
nach $p_i$ und $q_i$ auflösen. Der Nenner aller Werthe $p_i$ und $q_i$ ist
\[
D_2 =
\begin{vmatrix}
a_0 & a_1 & a_2 \\
b_0 & b_1 & b_2 \\
0 & b_0 & b_1
\end{vmatrix},
\]
ein Minor vierter Ordnung der Determinante $D_0$. Die Werthe von $p_i$ und $q_i$ sind:
%CONT>>


% ===== pdf idx 159, printed 145.0 =====
% !PRINTED 145
\weberpage{145}

\begin{center}
\S.\,11. Eigenschaften der Resultante $R_{f,\varphi}$.
\end{center}

%<<CONT
\[
D_2p_0 =
\begin{vmatrix}
0 & b_0 & 0 \\
0 & b_1 & 0 \\
1 & b_2 & b_1
\end{vmatrix}, \quad
D_2q_0 =
\begin{vmatrix}
a_0 & 0 & 0 \\
a_1 & 0 & b_0 \\
a_2 & 1 & b_1
\end{vmatrix}, \quad
D_2q_1 =
\begin{vmatrix}
a_0 & b_0 & 0 \\
a_1 & b_1 & 0 \\
a_2 & b_2 & 0
\end{vmatrix}.
\]
Trägt man diese Werthe in die Gleichungen
\[
r = a_3p_0+b_3q_0+b_2q_1
\]
\[
s = a_4p_0+b_3q_1
\]
ein, so erhält man die Coefficienten des quadratischen Restes ausgedrückt in den Coefficienten der gegebenen Functionen $f$ und $\varphi$.

\begin{center}
\textbf{\S.\,11. Eigenschaften der Resultante $R_{f,\varphi}$.}
\end{center}

138. \textit{Darstellung der Methode zur Ermittelung dieser Eigenschaften.} Nach den Entwicklungen des vorhergehenden Paragraphen haben wir die Resultante zweier Functionen
\[
\left.
\begin{array}{l}
f(x) = a_0x^m+a_1x^{m-1}\ldots+a_m \\
\varphi(x) = b_0x^n+b_1x^{n-1}\ldots+b_n
\end{array}
\right\} \quad m>n
\]
in Form einer Determinante vom Grade $m+n$ erhalten, nämlich:
\[
R_{f,\varphi} =
\begin{vmatrix}
a_0 & a_1 & a_2 & \cdots & a_m & 0 & \cdots & 0 \\
0 & a_0 & a_1 & \cdots & \cdots & a_m & \cdots & 0 \\
\cdot & \cdot & \cdot & \cdot & \cdot & \cdot & \cdot & \cdot \\
0 & \cdots & \cdots & a_0 & a_1 & \cdots & a_m & {} \\
b_0 & b_1 & b_2 & \cdots & b_n & 0 & \cdots & 0 \\
0 & b_0 & b_1 & \cdots & \cdots & b_n & \cdots & 0 \\
\cdot & \cdot & \cdot & \cdot & \cdot & \cdot & \cdot & \cdot \\
0 & \cdots & \cdots & b_0 & b_1 & \cdots & b_n & {}
\end{vmatrix}.
\]

Gerade aus diesem Umstande lassen sich zahlreiche Eigenschaften derselben einfach dadurch herleiten, dass wir diese Determinante geeignet umformen, ohne ihren Werth zu ändern. Die umgestaltete Determinante wiederum als Resultante anderer Functionen mit entsprechenden Coefficienten betrachtet, lässt dann direct die Eigenschaften erkennen. Diese Umformungen sind zweierlei Art; einmal lediglich Vertauschungen von Zeilen oder Reihen, dann aber auch Addition von Reihen, die mit gewissen Grössen multiplicirt gedacht werden, zu andern Reihen.
%CONT>>


% ===== pdf idx 160, printed 146.0 =====
% !PRINTED 146
\weberpage{146}

\begin{center}
Zweiter Theil. Anwendungen der Determinantentheorie.
\end{center}

%<<CONT
139. \textit{Vertauschung der letzten $m$ Zeilen mit den $n$ ersten.} Bringt man in der Determinante $R_{f,\varphi}$ die Zeilen mit dem Index $n+1, n+2, \cdots n+m$ an die ersten Stellen, und fügt die $m$ ersten Zeilen an die so erhaltene letzte Zeile an, so entsteht
\[
\begin{vmatrix}
b_0 & b_1 & b_2 & \cdots & b_n & 0 & \cdots & 0 \\
0 & b_0 & b_1 & \cdots & \cdots & b_n & \cdots & 0 \\
\cdot & \cdot & \cdot & \cdot & \cdot & \cdot & \cdot & \cdot \\
0 & \cdots & \cdots & b_0 & b_1 & \cdots & b_n & {} \\
a_0 & a_1 & a_2 & \cdots & a_m & 0 & \cdots & 0 \\
0 & a_0 & a_1 & \cdots & \cdots & a_m & \cdots & 0 \\
\cdot & \cdot & \cdot & \cdot & \cdot & \cdot & \cdot & \cdot \\
0 & \cdots & \cdots & a_0 & a_1 & \cdots & a_m & {}
\end{vmatrix}.
\]
Das Diagonalglied hat in seinen ersten Indices die Permutation
\[
(n+1, n+2 \ldots n+m, 1\,2\,3 \ldots n),
\]
jede der $m$ ersten Zahlen bildet mit den folgenden $n$ Derangements, so dass die Anzahl aller $n\cdot m$ ist. Wir haben daher die Beziehung
\[
R_{f,\varphi} = (-1)^{n\cdot m} R_{\varphi,f}.
\]

140. \textit{Transposition der Resultante.} Wendet man auf die Resultante eine zweimalige Transposition an, zunächst durch eine Umklappung um die zweite Diagonale, sodann durch eine zweite Umklappung um die erste Diagonale der erhaltenen neuen Matrix, so erhalten wir eine Matrix, in welcher die Elemente genau in umgekehrter Reihenfolge auftreten, nämlich
\[
R =
\begin{vmatrix}
b_n & b_{n-1} & b_{n-2} & \cdots & b_0 & 0 & \cdots & 0 \\
0 & b_n & b_{n-1} & \cdots & \cdots & b_0 & \cdots & 0 \\
\cdot & \cdot & \cdot & \cdot & \cdot & \cdot & \cdot & \cdot \\
0 & \cdots & \cdots & b_n & b_{n-1} & \cdots & b_0 & {} \\
a_m & a_{m-1} & \cdots & \cdots & a_0 & \cdots & \cdots & 0 \\
0 & a_m & \cdots & \cdots & \cdots & a_0 & \cdots & {} \\
\cdot & \cdot & \cdot & \cdot & \cdot & \cdot & \cdot & \cdot \\
0 & \cdots & \cdots & \cdots & \cdots & a_m & \cdots & a_0
\end{vmatrix}.
\]
Da die Transposition das Vorzeichen einer Determinante nicht ändert, so ist
\weq{I}{R = R_{f,\varphi}.}
Nun ist aber diese Determinante nichts anderes als die Resultante der Gleichungen
%CONT>>


% ===== pdf idx 161, printed 147.0 =====
% !PRINTED 147
\weberpage{147}
%<<CONT
\[ b_n x^n + b_{n-1}x^{n-1} + \cdots + b_0 = x^n \cdot \varphi\left(\frac{1}{x}\right) \]
\[ a_m x^m + a_{m-1}x^{m-1} + \cdots + a_0 = x^m \cdot f\left(\frac{1}{x}\right). \]

Die Relation (I) giebt daher den Satz:

\[ R_{f,\varphi} = R_{x^n\varphi\left(\frac{1}{x}\right),\ x^mf\left(\frac{1}{x}\right)} . \]

141. \textit{Resultante von $f(\lambda x)$ und $\varphi(\lambda x)$.} Ersetzen wir in den
Functionen $f(x)$ und $\varphi(x)$ die Ver\"anderliche $x$ durch $\lambda \cdot x$, so erhalten
wir

\weq{I}{
\begin{cases}
f(\lambda x) = a_0\lambda^m x^m + a_1\lambda^{m-1}x^{m-1} + \cdots \lambda^0 a_m \\
\varphi(\lambda x) = b_0\lambda^n x^n + b_1\lambda^{n-1}x^{n-1} + \cdots \lambda^0 b_n.
\end{cases}
}

Dasselbe Resultat ergiebt sich, wenn wir in $f(x)$ und $\varphi(x)$ $a_0$ durch
$\lambda^m a_0$, $a_1$ durch $\lambda^{m-1}a_1 \ldots$ u.\,s.\,w., $b_0$ durch
$\lambda^n b_0$, $b_1$ durch $\lambda^{n-1} b_1$ u.\,s.\,w. ersetzen. Die Resultante der
Functionen $f(\lambda x)$ und $\varphi(\lambda x)$ ergiebt sich also aus der Resultante von
$f(x)$ und $\varphi(x)$, durch Substitution dieser Werthe f\"ur $a_i$ und $b_i$, d.\,h. es
wird

\[
R_{f(\lambda x),\ \varphi(\lambda x)} =
\left|
\begin{array}{ccccccc}
\lambda^m a_0 & \lambda^{m-1}a_1 & \lambda^{m-2}a_2 & \ldots & 0 \\
0 & \lambda^m a_0 & \lambda^{m-1}a_2 & \ldots & 0 \\
\cdot & \cdot & \cdot & \cdot & \cdot & \cdot & \cdot \\
0 & 0 & \ldots & \lambda^m a_0 & \ldots & \lambda^0 a_m \\
\lambda^n b_0 & \lambda^{n-1}b_1 & \ldots & \lambda^0 b_n & 0 & \ldots & 0 \\
0 & \lambda^n b_0 & \lambda^{n-1}b_1 & \ldots & 0 \\
\cdot & \cdot & \cdot & \cdot & \cdot & \cdot & \cdot \\
0 & \ldots & \ldots & \lambda^n b_0 & \ldots & \ldots & \lambda_0 b_n
\end{array}
\right| .
\]

Genau dieselbe Resultante erh\"alt man aber, wenn wir in $R_{f,\varphi}$ die erste Colonne
mit $\lambda^{m+n-1}$, die zweite mit $\lambda^{m+n-2}$, die letzte mit $\lambda^0$, also
insgesammt mit $\lambda^{\frac{(m+n)(m+n-1)}{2}}$ multipliciren und hierauf die erste Zeile
mit $\lambda^{n-1}$, die zweite mit $\lambda^{n-2}\ldots$, die $n^{\text{te}}$ mit $\lambda^0$,
die $(n+1)^{\text{te}}$ mit $\lambda^{m-1}\ldots$, die $(n+m)^{\text{te}}$ mit $\lambda^0$,
also insgesammt mit $\lambda^{\frac{n(n-1)}{2}+\frac{m(m-1)}{2}}$ dividiren. Dadurch erh\"alt
aber $R_{f,\varphi}$ schliesslich den Factor $\lambda^{m\cdot n}$, so dass wir zu dem
Resultate gelangen:

\[ R_{f(\lambda x),\ \varphi(\lambda x)} = \lambda^{m\cdot n} R_{f,\varphi}. \]

Daraus geht zun\"achst hervor, dass jedes Glied der Resultante $R_{f,\varphi}$ so beschaffen
sein muss, dass es bei Substitution von $a_0\lambda^m$, $a_1\lambda^{m-1}\ldots$,
%CONT>>


% ===== pdf idx 162, printed 148.0 =====
% !PRINTED 148
\weberpage{148}
%<<CONT
$a_m\lambda^0$, $b_0\lambda^n$, $b_1\lambda^{n-1}\ldots$, $b_n\lambda^0$ f\"ur
$a_0a_1\ldots a_m$, $b_0b_1\ldots b_n$ den Factor $\lambda^{m\cdot n}$ erh\"alt. Aus
Gr\"unden der Symmetrie folgt, dass das N\"amliche eintritt bei den Substitutionen
$a_0\lambda^0$, $a_1\lambda^1\ldots a_m\lambda^m$, $b_0\lambda^0$, $b_1\lambda^1\ldots
b_n\lambda^n$ f\"ur $a_0a_1\ldots a_m$, $b_0b_1\ldots b_n$. Demgem\"ass haben alle Glieder der
Resultante die Eigenschaft, dass in ihnen die Summe aller Indices der darin auftretenden
Coefficienten $a_i$ und $b_i$ --- jeden Index $i$ so oftmal genommen als der bez\"ugliche
Exponent von $a_i$ oder $b_i$ angiebt --- eine constante Zahl n\"amlich gleich $m\cdot n$
ist. Man nennt diese Zahl $m\cdot n$ das Gewicht der Resultante, und sagt, die Resultante
ist eine isobare Function der Coefficienten von $f(x)$ und $\varphi(x)$.

142. \textit{Schnittpuncte zweier Curven $m^{\text{ter}}$ und $n^{\text{ter}}$ Ordnung.} Die
Eigenschaft der Resultante, eine isobare Function zu sein, k\"onnen wir benutzen, in
einfacher Weise \"uber die Zahl der Schnittpunkte zweier algebraischer Curven Aufschluss zu
erhalten. Denken wir uns n\"amlich, entsprechend den Factoren $\lambda^0$, $\lambda^1$,
$\lambda^2\ldots$ u.\,s.\,w., die Gr\"ossen $a_0$ und $b_0$ als Functionen nullten Grades in
$\xi$ und $\eta$, die Gr\"ossen $a_1$ und $b_1$ als homogene Functionen ersten, und allgemein
die Gr\"ossen $a_\varrho$ und $b_\varrho$ als homogene Functionen $\varrho^{\text{ten}}$
Grades in $\xi$ und $\eta$, so ist jedes Glied der Resultante eine homogene Function
$m\cdot n^{\text{ten}}$ Grades dieser zwei Ver\"anderlichen, wie sein Gewicht $m\cdot n$
lehrt.

Sind nun $f(x_1\overset{m}{x_2}x_3) = 0$ und $\varphi(x_1\overset{n}{x_2}x_3) = 0$ zwei
% ?? uncertain: exact placement/meaning of the small superscript degree marker over $x_2$ in $f(x_1\overset{m}{x_2}x_3)$ and $\varphi(x_1\overset{n}{x_2}x_3)$
Gleichungen, welche homogen geschrieben und nach $x_1$ geordnet die Form besitzen:

\[ f(x_1x_2x_3) = a_0x_1^m + a_1x_1^{m-1} + a_2x_1^{m-2}\ldots + a_m = 0 \]
\[ \varphi(x_1x_2x_3) = b_0x_1^n + b_1x_1^{n-1} + b_2x_2^{n-2}\ldots + b_n = 0, \] % [sic] second term prints $x_2^{n-2}$, not $x_1^{n-2}$

dann sind in ihnen die Coefficienten $a_0$, $b_0$; $a_1$, $b_1$; $a_2$, $b_2$ u.\,s.\,w.
solche homogene Functionen von $x_2$ und $x_3$ vom bezw. $0^{\text{ten}}$,
$1^{\text{ten}}$, $2^{\text{ten}}$ Grade u.\,s.\,w. Bildet man also die Resultante
$R_{f,\varphi}$ nach $x_1$, so erh\"alt man in ihr eine homogene ganze Function von $x_2$
und $x_3$, die nach dem Vorausgehenden vom Grade $m\cdot n$ ist. Deutet man daher $f$ und
$\varphi$ geometrisch als Curvengleichungen, so hat man den Satz: Zwei Curven
$m^{\text{ter}}$ bezw. $n^{\text{ter}}$ Ordnung besitzen $m\cdot n$ gemeinsame
Schnittpuncte*).

143. \textit{Die Resultante als Functionaldeterminante.} Eine weitere Reihe von
Eigenschaften der Resultante l\"asst sich durch eine Betrachtungsweise derselben erkennen,
die zuerst von Sylvester hierbei eingef\"uhrt
%CONT>>

\bigskip\noindent\rule{4cm}{0.4pt}

\smallskip\noindent
*) Dieser Satz wurde hier auf rationalem Wege abgeleitet. Eine gleiche rationale Methode, zu
zeigen dass drei Fl\"achen bezw. $m^{\text{ten}}$, $p^{\text{ten}}$ und $q^{\text{ten}}$
Grades sich in $mpq$ Punkten schneiden, besitzen wir noch nicht.


% ===== pdf idx 163, printed 149.0 =====
% !PRINTED 149
\weberpage{149}
%<<CONT
wurde. Bildet man sich aus den Functionen $f(x)$ und $\varphi(x)$ der Reihe nach die
Functionen

\[ x^0\cdot f,\quad x^1\cdot f,\quad x^2\cdot f,\quad x^3\cdot f \ldots x^{n-1}\cdot f \]
\[ x^0\cdot \varphi,\quad x^1\cdot \varphi,\quad x^2\cdot \varphi,\quad x^3\cdot \varphi
\ldots x^{m-1}\cdot \varphi, \]

so kann man sich darin die Potenzen $x^k$ durch $\xi_k$ ersetzt denken, wodurch wir $m+n$
lineare Functionen der $m+n$ Gr\"ossen

\[ \xi_0,\ \xi_1,\ \xi_2 \ldots \xi_k \ldots \xi_{n+m-1} \]

erhalten. Die partiellen Differentialquotienten nach diesen $m+n$ Gr\"ossen sind gerade die
Coefficienten der Functionen $f$ und $\varphi$, und es ist daher

\[
R_{f,\varphi} =
\left|
\begin{array}{cccc}
\dfrac{\partial(x^{n-1}f)}{\partial\xi_{n+m-1}} & \dfrac{\partial(x^{n-1}f)}{\partial\xi_{n+m-2}} & \dfrac{\partial(x^{n-1}f)}{\partial\xi_{n+m-3}} \ldots & \dfrac{\partial(x^{n-1}f)}{\partial\xi_0} \\[2ex]
\dfrac{\partial(x^{n-2}f)}{\partial\xi_{n+m-1}} & \dfrac{\partial(x^{n-2}f)}{\partial\xi_{n+m-2}} & \cdot\ \cdot\ \cdot\ \cdot\ \cdot & \dfrac{\partial(x^{n-2}f)}{\partial\xi_0} \\[2ex]
\cdot & \cdot & \cdot\ \cdot\ \cdot\ \cdot\ \cdot & \cdot \\[1ex]
\dfrac{\partial(x^{m-1}\varphi)}{\partial\xi_{n+m-1}} & \dfrac{\partial(x^{m-1}\varphi)}{\partial\xi_{n+m-2}} & \dfrac{\partial(x^{m-1}\varphi)}{\partial\xi_{n+m-3}} \ldots & \dfrac{\partial(x^{m-1}\varphi)}{\partial\xi_0} \\[2ex]
\dfrac{\partial(x^{m-2}\varphi)}{\partial\xi_{n+m-1}} & \cdot & \cdot\ \cdot\ \cdot\ \cdot\ \cdot & \cdot \\[1ex]
\cdot & \cdot & \cdot\ \cdot\ \cdot\ \cdot\ \cdot & \cdot \\[1ex]
\dfrac{\partial(x^0\varphi)}{\partial\xi_{n+m-1}} & \dfrac{\partial(x^0\varphi)}{\partial\xi_{n+m-2}} & \dfrac{\partial(x^0\varphi)}{\partial\xi_{n+m-3}} \ldots & \dfrac{\partial(x^0\varphi)}{\partial\xi_0}
\end{array}
\right| .
\]

Man kann daher in diesem Sinne auch k\"urzer schreiben:

\[
R_{f,\varphi} =
\left(
\begin{array}{cccccc}
x^{n-1}f, & x^{n-2}f \ldots x^0f, & x^{m-1}\varphi, & x^{m-2}\varphi \ldots x^0\varphi \\
x^{n+m-1}, & x^{n+m-2}\ldots\ldots x^n, & x^{n-1}\ldots x^0
\end{array}
\right).
\]

Dabei ist also in der obern Reihe dieses Ausdruckes, den wir
%CONT>>


% ===== pdf idx 164, printed 150.0 =====
% !PRINTED 150
\weberpage{150}
%<<CONT
fr\"uher mit Z\"ahler bezeichnet haben, jede Potenz von $x$ durch $\xi$ mit dem
entsprechenden Index ersetzt zu denken, und die Gr\"ossen der untern Reihe, welcher wir
seinerzeit den Namen Nenner gaben, sind jene Gr\"ossen $\xi_{n+m-1}$, $\xi_{n+m-2}\ldots
\xi_0$, von welchen wir alle Functionen im Z\"ahler als linear abh\"angig zu betrachten
haben. Wir werden im Folgenden auch noch Functionaldeterminanten betrachten, in deren Nenner
selbst lineare Functionen der $\xi_k$ sich befinden, so dass alsdann die Gr\"ossen des
Z\"ahlers nicht nur linear in $\xi_k$, sondern auch noch linear in diesen linearen
Functionen der $\xi_k$ zu denken sind. (Vergleiche auch das Beispiel in Nr. 149.)

144. \textit{Methode f\"ur die folgenden Betrachtungen.} Das wesentlichste Hilfsmittel bei
allen folgenden Untersuchungen wird der bereits in
%CONT>>

%<<CONT
Nr. 117 aufgestellte Productsatz f\"ur Functionaldeterminanten bilden, welcher dort
dargestellt wurde durch

\[
\left(
\begin{array}{cc}
y_1\ y_2\ y_3 \ldots y_n \\
x_1\ x_2\ x_3 \ldots x_n
\end{array}
\right)
\left(
\begin{array}{cc}
z_1\ z_2\ z_3 \ldots z_n \\
y_1\ y_2\ y_3 \ldots y_n
\end{array}
\right)
=
\left(
\begin{array}{cc}
z_1\ z_2\ z_3 \ldots z_n \\
x_1\ x_2\ x_3 \ldots x_n
\end{array}
\right).
\]

Dabei werden wir jene Determinanten, mit denen wir multipliciren, stets so w\"ahlen, dass
sie den Werth 1 besitzen. Man erreicht dies durch solche Wahl der Functionen und ihrer
unabh\"angig Ver\"anderlichen, dass diejenigen partiellen Differentialquotienten derselben,
welche in die Diagonale zu stehen kommen, den Werth 1, alle rechts davon befindlichen aber
oder alle links in der Matrix stehenden verschwinden. An Stelle der Multiplication tritt
bisweilen auch Division, d.\,h. anstatt mit der Functionaldeterminante

\weq{1}{
\left(
\begin{array}{cc}
z_1\ z_2\ z_3 \ldots z_n \\
y_1\ y_2\ y_3 \ldots y_n
\end{array}
\right)
}

zu multipliciren, nehmen wir die gleichsam gest\"urzte Form

\weq{2}{
\left(
\begin{array}{cc}
y_1\ y_2\ y_3 \ldots y_n \\
z_1\ z_2\ z_3 \ldots z_n
\end{array}
\right)
}

als Multiplicator, eine Determinante, die nothwendig gleichfalls den Werth 1 besitzt, wenn
ihn die urspr\"ungliche besass, wie ja aus dem Satze (vgl. Nr. 118)

\[
\left(
\begin{array}{cc}
y_1\ y_2\ y_3 \ldots y_n \\
x_1\ x_2\ x_3 \ldots x_n
\end{array}
\right)
\left(
\begin{array}{cc}
x_1\ x_2\ x_3 \ldots x_n \\
y_1\ y_2\ y_3 \ldots y_n
\end{array}
\right)
= 1
\]

direct hervorgeht.

145. \textit{Berechnung von $R_{f+\varphi\psi,\varphi}$.} Sei $\psi$ irgend eine Function
von $x$

\[ \psi = c_0x^{q-1} + c_1x^{q-2}\ldots c_q, \]

wo $q \leqq m-n$ ist. Bilden wir die Producte

\[ \varphi\cdot\psi\cdot x^{n-1},\quad \varphi\cdot\psi\cdot x^{n-2}\ldots \varphi\cdot\psi\cdot x^0, \]

ersetzen darin die Potenzen $x^k$ durch $\xi_k$, so stellt die Functionaldeterminante

\[
A =
\left(
\begin{array}{cc}
x^{n-1}(f+\varphi\cdot\psi),\ x^{n-2}(f+\varphi\cdot\psi)\ldots x^0(f+\varphi\cdot\psi),\ x^{m-1}\varphi,\ x^{m-2}\varphi\ldots x^0\varphi \\
x^{n-1}f,\ x^{n-2}f,\ \ldots\ x^0f,\ x^{m-1}\varphi,\ x^{m-2}\varphi\ldots x^0\varphi
\end{array}
\right)
\]

eine Determinante dar vom Werthe 1. Denn jede Gr\"osse im Z\"ahler ist eine lineare
Function aller Gr\"ossen im Nenner, und zwar so beschaffen, dass die partiellen
Differentialquotienten nach denselben f\"ur die Diagonalelemente die Werthe 1, f\"ur die
links gelegenen die Werthe Null liefern, wie dies ein einfaches Beispiel klar machen mag.
Ich nehme an $m=3$, $q=1$, $n=2$, dann ist die erste Function im
%CONT>>


% ===== pdf idx 165, printed 151.0 =====
% !PRINTED 151
\weberpage{151}
%<<CONT
Z\"ahler, wenn wir uns die Gr\"ossen $\xi_k$ eingetragen denken, dargestellt
durch
\[
\begin{aligned}
z_1 = (a_0\xi_4 + \alpha_1\xi_3 + a_2\xi_2 + a_3\xi_1) + c_0(b_0\xi_4 + b_1\xi_3 + b_2\xi_2)\\
{} + c_1(b_0\xi_3 + b_1\xi_2 + b_2\xi_1).
\end{aligned}
\]
% [sic] second term of z_1 printed as \alpha_1, not a_1, as elsewhere
Die Variabeln der Nenner sind:
\[
\eta_1 = a_0\xi_4 + a_1\xi_3 + a_2\xi_2 + a_3\xi_1,\quad
\eta_2 = a_0\xi_3 + a_1\xi_2 + a_2\xi_1 + a_3\xi_0,
\]
\[
\eta_3 = b_0\xi_4 + b_1\xi_2 + b_2\xi_1,\quad
\eta_4 = b_0\xi_3 + b_1\xi_2 + b_2\xi_1,\quad
\eta_5 = b_0\xi_0 + b_1\xi_1 + b_2\xi_0.
\]
Die Differentiation nach diesen Gr\"ossen liefert als Elemente der ersten
Zeile der Reihe nach
\[
1,\ 0,\ c_0,\ c_1,\ 0
\]
u. s. w. Die Determinante $A$ ist also in der That eine Determinante
von der in Nr. 144 erw\"ahnten Eigenschaft: sie hat den Werth 1.

Berechnen wir nun das Product
\[
A \cdot R_{f,\varphi} = 1 \cdot R_{f,\varphi},
\]
so entsteht gem\"ass dem Productsatze
\[
R_{f,\varphi} = \left(\begin{array}{l}
x^{n-1}(f+\varphi\cdot\psi),\ x^{n-2}(f+\varphi\cdot\psi)\ldots x^0(f+\varphi\cdot\psi),\ x^{m-1}\varphi\ldots\varphi\\[2pt]
x^{m+n-1},\quad x^{m+n-2}\ \cdots\cdots\cdots\cdots\ x^2,\ x^1,\ x^0
\end{array}\right),
\]
oder
\weq{I}{R_{f,\varphi} = R_{f+\varphi\cdot\psi,\varphi},}
d. h.: \glqq Die Resultante von $f$ und $\varphi$ ist auch gleich der
Resultante von $f+\varphi\cdot\psi$ und $\varphi$, wenn $\psi$ eine
beliebige Function vom Grade $q\leqq m-n$ ist.\grqq{}

\textit{Anmerkung.} Indem wir die Einschr\"ankung machten $q\leqq m-n$,
setzen wir, wie schon mehrmals erw\"ahnt, voraus $m>n$. Es w\"are
eine sch\"atzenswerthe Arbeit, auch den Fall zu untersuchen, in welchem
$m<n$ ist, d. h. allgemein die Frage zu behandeln: Wie h\"angen die
Resultante $R_{f+\varphi\cdot\psi,\varphi}$ und $R_{f,\varphi}$ zusammen,
wenn wir \"uber den Grad der diesbez\"uglichen Functionen keinerlei
Voraussetzung machen?

146. \textit{Beispiel.} Ein besonders h\"aufiger Fall, in welchem von
diesem Satze Gebrauch gemacht wird, ist jener, in welchem die Function
$\psi$ eine Constante $\lambda$ ist. Man hat alsdann:
\[
R_{f,\varphi} = R_{f+\lambda\varphi,\varphi}.
\]
In diesem Falle ist der Satz auch leicht durch Zerlegung der
Determinante $R_{f+\lambda\varphi,\varphi}$ zu erkennen. Sei z. B.
\[
f = a_0x^2 + a_1x + a_2
\]
\[
\varphi = b_0x^2 + b_1x + b_2,
\]
dann ist
\[
R_{f+\lambda\varphi,\varphi} = \begin{vmatrix}
a_0+\lambda b_0, & a_1+\lambda b_1, & a_2+\lambda b_2, & 0\\
0 & a_0+\lambda b_0, & a_1+\lambda b_1, & a_2+\lambda b_2\\
b_0 & b_1 & b_2 & 0\\
0 & b_0 & b_1 & b_2
\end{vmatrix}.
\]


% ===== pdf idx 166, printed 152.0 =====
% !PRINTED 152
\weberpage{152}
Zerlegt man diese Determinante nach Nr. 37 in eine Summe von
Determinanten, so sieht man leicht, dass alle bis auf
\[
R_{f,\varphi} = \begin{vmatrix}
a_0 & a_1 & a_2 & 0\\
0 & a_0 & a_1 & a_2\\
b_0 & b_1 & b_2 & 0\\
0 & b_0 & b_1 & b_2
\end{vmatrix}
\]
identisch verschwinden, dass also die Beziehung besteht
\[
R_{f+\lambda\varphi,\varphi} = R_{f,\varphi}.
\]
147. \textit{Resultante von $f(x+u)$ und $\varphi(x+u)$.} Ersetzen wir
in den Functionen $f$ und $\varphi$ die Ver\"anderliche $x$ durch $x+u$,
so erhalten wir
\[
f(x+u) = a_0(x+u)^m + a_1(x+u)^{m-1}\cdots + a^m
\]
% [sic] printed a^m (superscript), not a_m
\[
\varphi(x+u) = b_0(x+u)^n + b_1(x+u)^{n-1}\cdots + b_n.
\]
Um die Resultante der so transformirten Functionen zu gewinnen, bilden
wir zun\"achst folgende drei Functionaldeterminanten:
\[
\left(\begin{array}{l}
(x+u)^{m+n-1},\ (x+u)^{m+n-2},\ \ldots (x+u)^0\\[2pt]
x^{m+n-1},\quad x^{m+n-2},\ \ \ldots\ x^0
\end{array}\right)
\]
\[
\left(\begin{array}{l}
(x+u)^{n-1}f,\ (x+u)^{n-2}f,\ldots(x+u)^0f,\ x^{m-1}\varphi,\quad x^{m-2}\varphi\ldots x^0\varphi\\[2pt]
x^{n-1}f,\quad x^{n-2}f,\ \ldots\ x^0f,\quad x^{m-1}\varphi,\quad x^{m-2}\varphi\ldots x^0\varphi
\end{array}\right)
\]
\[
\left(\begin{array}{l}
(x+u)^{n-1}f,\ (x+u)^{n-2}f,\ldots(x+u)^0f,\ (x+u)^{m-1}\varphi,\ (x+u)^{m-2}\varphi\ldots(x+u)^0\varphi\\[2pt]
(x+u)^{n-1}f,\ (x+u)^{n-2}f,\ \ldots\ldots\ldots,\quad x^{m-1}\varphi,\quad x^{m-2}\varphi\ldots x^0\varphi
\end{array}\right)
\]
und bezeichnen die erste mit $A$, die zweite mit $B$, die dritte mit
$C$. Ersetzen wir wieder in allen darin auftretenden Gr\"ossen die
$x^k$ durch $\xi_k$ und denken uns die Differentiationen nach den
Gr\"ossen im Nenner ausgef\"uhrt, so sieht man, dass alle drei
Determinanten als Diagonalelemente die Zahl 1, und links davon nur
verschwindende Elemente besitzen. So sind z. B. die Elemente der ersten
Determinante den Zeilen nach
\[
\begin{array}{lll}
1 & \binom{m+n-1}{1}u, & \binom{m+n-1}{2}u^2\ldots u^{m+n-1}\\
0 & 1 & \binom{m+n-2}{1}u\ldots u^{m+n-2}\\
0 & 0 & 1\ \ldots\ldots u^{m+n-3}
\end{array}
\]
u. s. w. Ebenso sind die ersten Zeilen der zweiten Determinante
\[
\begin{array}{lll}
1 & \binom{n-1}{1}u, & \binom{n-1}{2}u^2\ldots u^{n-1}\ 0\ 0\ldots 0\\
0 & 1 & \binom{n-2}{1}u\ldots u^{n-2}\ 0\ 0\ldots\\
0 & 0 & 1\ \ldots u^{n-3}\ 0\ 0\ldots
\end{array}
\]
u. s. w. Multipliciren wir mit den beiden letzten Determinanten die
%CONT>>


% ===== pdf idx 167, printed 153.0 =====
% !PRINTED 153
\weberpage{153}
%<<CONT
Resultante $R_{f,\varphi}$, und dividiren wir das Product durch die
erste, so erhalten wir
\[
R_{f,\varphi}\cdot B\cdot C : A = R_{f,\varphi}
\]
\[
= \left(\begin{array}{l}
(x+u)^{n-1}f,\ (x+u)^{n-2}f\ldots(x+u)^0f,\ (x+u)^{m-1}\varphi\ldots(x+u)^0\varphi\\[2pt]
(x+u)^{n+m-1},\ (x+u)^{n+m-2}\ \ldots\ldots\ldots\ldots\ (x+u)^0
\end{array}\right).
\]
Der Ausdruck rechts ist aber nichts anderes als die Resultante der
beiden Functionen $f(x+u)$, $\varphi(x+u)$; daher das Resultat:
\weq{1}{R_{f,\varphi} = R_{f(x+u),\,\varphi(x+u)}.}
Diese Relation definirt in Verbindung mit den beiden fr\"uher
erhaltenen
\weq{2}{R_{f,\varphi} = \frac{1}{\lambda^{mn}}\cdot R_{f(\lambda x),\,\varphi(\lambda x)}}
und
\weq{3}{R_{f,\varphi} = R_{x^n f\left(\frac{1}{x}\right),\; x^m\varphi\left(\frac{1}{x}\right)}}
die Resultante als Invariante der beiden Formen $f$ und $\varphi$.

Eine Invariante $J$ ist, wie wir sp\"ater noch eingehender betrachten
werden, eine solche homogene Function der Coefficienten von $f(x)$ und
$\varphi(x)$, welche sich bei linearer Transformation h\"ochstens um
eine Potenz der Transformationsdeterminante
$\Delta = \alpha_1\beta_2 - \beta_1\alpha_2$ von jener Function
unterscheidet, die genau in derselben Weise aus den Coefficienten der
linear transformirten Functionen $f(\alpha_1\xi+\alpha_2\eta)$ und
$\varphi(\beta_1\xi+\beta_2\eta)$ gebildet ist wie $J$ selbst. Aus den
drei Substitutionen
\[
x = \xi + u,\quad x = \frac{1}{\xi},\quad x = \lambda\xi
\]
l\"asst sich aber jede lineare Transformation zusammensetzen.

148. \textit{Reduction der $(m+n)$reihigen Determinante $R_{f,\varphi}$.}
Betrachten wir die Resultante
\[
A = \left(\begin{array}{l}
x^{n-1},\quad x^{n-2}\ldots x^0,\ x^{m-1}\varphi,\ x^{m-2}\varphi\ldots x^0\varphi\\[2pt]
x^{m+n-1},\quad x^{m+n-2}\ \ldots\ldots\ldots\ldots\ x^1,\ x^0
\end{array}\right)
\]
und wenden auf sie die in Nr. 139 durchgef\"uhrte Vertauschung der
Zeilen an, so erhalten wir
\[
A = (-1)^{m\cdot n}\left(\begin{array}{l}
x^{m-1}\varphi,\ x^{m-2}\varphi\ldots x^0\varphi,\ x^{n-1},\ x^{n-2}\ldots x^0\\[2pt]
x^{m+n-1},\quad x^{m+n-2}\ \ldots\ldots\ldots\ x^1,\ x^0
\end{array}\right).
\]
Der Ausdruck rechts ist eine Determinante --- immer vorausgesetzt, dass
$x^k$ durch $\xi_k$ ersetzt ist ---, deren $m$ erste Diagonalelemente
den Werth $b_0$ und deren $n$ letzte Diagonalelemente den Werth 1
haben. Alle Elemente links der Diagonale sind null; sie reducirt sich
also auf
\[
A = (-1)^{m\cdot n}\,b_0^m.
\]


% ===== pdf idx 168, printed 154.0 =====
% !PRINTED 154
\weberpage{154}
Dividiren wir daher die Resultante $R_{f,\varphi}$ mit $A$, so erhalten
wir
\[
R_{f,\varphi}\cdot\frac{1}{(-1)^{m\cdot n}b_0^m}
\]
\[
= \left(\begin{array}{l}
x^{m+n-1},\quad x^{m+n-2}\ \ldots\ldots\ x^0\\[2pt]
x^{n-1},\quad x^{n-2}\ldots x^{m-1}\varphi\ldots x^0\varphi
\end{array}\right)
\left(\begin{array}{l}
x^{n-1}f,\quad x^{n-2}f\ldots f,\ x^{m-1}\varphi\ldots\varphi\\[2pt]
x^{m+n-1},\ x^{m+n-2}\ \ldots\ldots\ x^0
\end{array}\right)
\]
\[
R_{f,\varphi} = (-1)^{mn}\cdot b_0^m\left(\begin{array}{l}
x^{n-1}f,\ x^{n-2}f\ldots x^0f,\ x^{m-1}\varphi\ldots x^0\varphi\\[2pt]
x^{n-1},\quad x^{n-2}\ \ldots x^0,\quad x^{m-1}\varphi\ldots x^0\varphi
\end{array}\right),
\]
oder endlich (vgl. Nr. 116) nach Unterdr\"uckung gleicher
\"ubereinander stehender Gr\"ossen:
\weq{I}{R_{f,\varphi} = (-1)^{m\cdot n}\,b_0^m\left(\begin{array}{l}
x^{n-1}f,\ x^{n-2}f,\ x^{n-3}f\ldots x^0f\\[2pt]
x^{n-1},\quad x^{n-2}\ \ldots\ldots\ldots\ x^0
\end{array}\right).}
Hierbei sind die Functionen im Z\"ahler noch immer als Function sowohl
der noch vorhandenen als der unterdr\"uckten Gr\"ossen zu betrachten,
so dass also keineswegs der Ausdruck rechts, wie es auf den ersten
Blick erscheinen m\"ochte, von den Coefficienten $b$ von $\varphi$
unabh\"angig ist. Die Resultante als Determinante vom Grade $m+n$ ist
somit reducirt auf eine vom Grade $n$.

149. \textit{Beispiel.} Ich will das eben erhaltene Resultat noch
durch ein Beispiel erl\"autern, indem ich die Resultante einer
cubischen und quadratischen Form auf Grund desselben aufstelle. Es sei
\[
f(x) = a_0x^3 + a_1x^2 + a_2x + a^3
\]
% [sic] printed a^3 (superscript), not a_3
\[
\varphi(x) = b_0x^2 + b_1x + b_2.
\]
Dann ist die Resultante der beiden Functionen dargestellt durch
\[
R_{f,\varphi} = (-1)^6\cdot b_0{}^3\left(\begin{array}{l}
x^1f,\ x^0f,\ x^2\varphi,\ x^1\varphi,\ x^0\varphi\\[2pt]
x^1,\quad x^0,\quad x^2\varphi,\ x^1\varphi,\ x^0\varphi
\end{array}\right)
= b_0{}^3\left(\begin{array}{l}
x^1f,\ x^0f\\[2pt]
x^1,\quad x^0
\end{array}\right).
\]
Bezeichnen wir also die Functionen im Z\"ahler mit
$y_1\ y_2\ y_3\ y_4\ y_5$ und f\"uhren im Sylvester'schen Sinne $\xi_k$
ein, so erhalten wir
\[
y_1 = a_0\xi_4 + a_1\xi_3 + a_2\xi_2 + a_3\xi_1
\]
\[
y_2 = a_0\xi_3 + a_1\xi_2 + a_2\xi_1 + a_3\xi_0
\]
\[
y_3 = b_0\xi_4 + b_1\xi_3 + b_2\xi_2
\]
\[
y_4 = b_0\xi_3 + b_1\xi_2 + b_2\xi_1
\]
\[
y_5 = b_0\xi_2 + b_1\xi_1 + b_2\xi_0.
\]
Ich eliminire nun aus $y_1$ und $y_2$ die Gr\"ossen $\xi_1\ \xi_3\ \xi_2$
vermittelst der linearen Function $y_3\ y_4\ y_5$ in folgender Weise.
Es ist
\weq{1}{y_1 : y_3 = (a_0\xi_4 + a_1\xi_3\ldots) : (b_0\xi_4 + b_1\xi_3 + b_2\xi_2) = \frac{a_0}{b_0} + f_1,}
%CONT>>


% ===== pdf idx 169, printed 155.0 =====
% !PRINTED 155
\weberpage{155}
\textbf{\S.\,11. Eigenschaften der Resultante $R_{f,\varphi}$.}

wo mit $f_1$ der noch von $\xi_3$ und $\xi_2$ abhängige Rest
\[ f_1 = \frac{1}{b_0}(b_0a_1)\xi_3 + \frac{1}{b_0}(b_0a_2)\xi_2 + a_3\xi_1 \]
bezeichnet sei. Aus diesem Reste wird $\xi_3$ durch Division $y_4$ eliminirt. Man erhält
\weq{2}{f_1 : y_4 = \frac{(b_0a_1)}{b_0{}^2} + f_2\,,}
wobei $f_2$ der Rest ist:
\[ f_2 = \frac{b_0(b_0a_2) - b_1(b_0a_1)}{b_0{}^2}\,\xi_2 + \frac{b_0{}^2a_3 - b_2(b_0a_1)}{b_0{}^2}\,\xi_1\,. \]
Eliminirt man endlich hieraus durch Division mit $y_5$ auch $\xi_2$, so kommt:
\weq{3}{\begin{aligned}
f_2:y_5 &= \frac{b_0(b_0a_2)-b_1(b_0a_1)}{b_0{}^3} + \frac{1}{b_0{}^3}\{b_0{}^3a_3-b_0b_2(b_0a_1)-b_0b_1(b_0b_2)+b_1{}^2(b_0a_1)\}\xi_1\\
&\quad -\frac{1}{b_0{}^3}\{b_2b_0(b_0a_2)-b_2b_1(b_0a_1)\}\xi_0\,.
\end{aligned}}
Trägt man diese Reste in (2) und (1) successive ein, so kommt:
\weq{4}{\begin{aligned}
y_1 &= \frac{a_0}{b_0}y_3 + \frac{(b_0a_1)}{b_0{}^2}y_4 + \frac{b_0(b_0a_2)-b_1(b_0a_1)}{b_0{}^3}y_5\\
&\quad +\frac{1}{b_0{}^3}\Big\{[b_0{}^3a_3-b_0b_2(b_0a_1)--b_0b_1(b_0a_2)+b_1{}^2(b_0a_1)]\xi_1-[b_2b_0(b_0a_2)-b_2b_1(b_0a_1)]\xi_0\Big\}\,.
\end{aligned}} % [sic] double minus as printed before b_0b_1(b_0a_2)
Damit ist die erste Grösse $xf$ des Zählers als eine lineare Function der Grössen $\xi_k$ und anderer Grössen $y_i$, die selbst wieder linear in $\xi_k$ sind, dargestellt. Ebenso erhalten wir für die zweite Function $x^0\cdot f$
\weq{5}{y_2 = \frac{a_0}{b_0}y_4+\frac{(b_0a_1)}{b_0{}^2}y_5+\frac{1}{b_0{}^2}\Big\{[b_0(b_0a_2)-b_1(b_0a_1)]\xi_1+[a_3b_0{}^2-b_2(b_0a_1)]\xi_0\Big\}\,.}
Die Resultante der beiden Functionen ist also:
\[ R_{f,\varphi} = b_0{}^3 \begin{pmatrix} y_1 & y_2 \\ x^1 & x^0 \end{pmatrix} \]
\[ = \frac{b_0{}^3}{b_0{}^3b_0{}^2}\left|\begin{matrix}
b_0{}^3a_3-b_0b_2(b_0a_1)-(b_0a_2)b_0b_1+b_1{}^2(b_0a_1), & b_2b_0(b_0a_2)-b_2b_1(b_0a_1)\\[4pt]
b_0(b_0a_2)-b_1(b_0a_1), & a_3b_0{}^2-b_2(b_0a_1)
\end{matrix}\right|\,. \]
Ersetzt man in den Gleichungen (4) und (5) $y_3$, $y_4$, $y_5$ durch $x^2\varphi$, $x\varphi$, $\varphi$ und hierauf $\xi_k$ wieder durch $x^k$, so erhält man:
\weq{6}{x^1\cdot f = \left\{\frac{a_0}{b_0}x^2+\frac{(b_0a_1)}{b_0{}^2}x+\frac{b_0(b_0a_2)-b_1(b_0a_1)}{b_0{}^3}\right\}\cdot\varphi+g_1(x)}
\[ = p_1(x)\cdot\varphi + g_1(x)\,, \]
\weq{7}{x^0\cdot f = \left\{\frac{a_0}{b_0}x+\frac{(b_0a_1)}{b_0}\right\}\varphi+g_2(x) = p_2(x)\cdot\varphi+g_2(x)\,.}
Hier sind $g_1(x)$ und $g_2(x)$ Functionen vom ersten Grade in $x$, und zwar, wie aus (6) und (7) erkennbar ist, die Reste, welche sich bei Division von $x^1\cdot f$ bezw. $x^0\cdot f$ durch $\varphi$ ergeben.


% ===== pdf idx 170, printed 156.0 =====
% !PRINTED 156
\weberpage{156}
\textbf{156\quad Zweiter Theil. Anwendungen der Determinantentheorie.}

Gehe ich also von diesem speciellen Beispiele zum allgemeinen Falle über, der ja ganz in derselben Weise behandelt werden kann, so ergiebt sich die Bedeutung der Grössen, welche den Zähler der Resultante (I) in Nr. 148 bilden. Sie sind nichts anderes als die $n$ Reste $(n-1)^{\text{ten}}$ Grades, welche man erhält durch Division von
\[ x^{n-1}f,\ x^{n-2}f,\ x^{n-3}f,\ x^0f \]
mit der Function $\varphi$, und die Resultante zweier Functionen $f$ und $\varphi$ ist bis auf einen Zahlenfactor gleich der Functionaldeterminante dieser $n$ Reste, wenn sie als lineare Functionen von $\xi_k$ gedacht werden.

150. \textit{Reduction der Determinante $R_{f,\varphi}$, wenn die beiden Functionen $f$ und $\varphi$ vom gleichen Grade sind.} Wir sind im Vorhergehenden zu einem Resultate gelangt, das für den Fall, in welchem $f$ und $\varphi$ von gleichem Grade in $x$ bereits von Bézout erreicht wurde. Ist nämlich
\[ f = a_0x^n + a_1x^{n-1} + a_2x^{n-2} + \ldots + a_n \]
\[ \varphi = b_0x^n + b_1x^{n-1} + b_2x^{n-2} + \ldots + b_n\,, \]
so gewinnen wir mit Bézout die Resultante
\[ R_{f,\varphi} = \begin{pmatrix} x^{n-1}f, & x^{n-2}f\ldots x^0f, & x^{n-1}\varphi, & x^{n-2}\varphi\ldots x^0\varphi\\ x^{2n-1}, & x^{2n-2}\ldots x^n\,, & x^{n-1} & \ldots\ldots\ldots x^0 \end{pmatrix} \]
in folgender Weise. Zunächst multipliciren wir die Determinante $R_{f,\varphi}$ mit der Functionaldeterminante
\[ \begin{pmatrix} A_{n-1}\varphi-B_{n-1}f, & A_{n-2}\varphi-B_{n-2}f,\ldots, & A_0\varphi-B_0f, & x^{n-1}\varphi,\ldots\varphi\\ x^{n-1}f\,, & x^{n-2}f\,,\ldots, & x^0f\,, & x^{n-1}\varphi,\ldots\varphi \end{pmatrix} = B\,, \]
wobei
\[ A_{n-1} = a_0x^{n-1}\cdots+a_{n-1}\,,\quad B_{n-1}=b_0x^{n-1}\cdots+b_{n-1} \]
\[ A_{n-2} = a^0x^{n-2}\cdots+a_{n-2}\,,\quad B_{n-2}=b_0x^{n-2}\cdots+b_{n-2} \] % [sic] printed a^0, not a_0
\[ \cdots\cdots\cdots\cdots\cdots\cdots\cdots\cdots\cdots \]
\[ A_0 = a_0\,,\qquad\qquad\qquad\qquad B_0=b_0\,. \]
Um den Werth der letzteren zu ermitteln, stürzen wir in dieser Determinante die Matrix der ersten $n$ Zeilen und die Matrix der letzten Zeilen und hierauf ebenso die der ersten $n$ Colonnen und die der letzten $n$ Colonnen. Dadurch wird aus $B$:
\[ \begin{pmatrix} A_0\varphi-B_0f, & A_1\varphi-B_1f,\cdots, & A_{n-1}\varphi-B_{n-1}f, & x^0\varphi,\cdots x^{n-1}\varphi\\ x^0f\,, & x^1f,\cdots, & x^{n-1}f & ,\,x^0\varphi,\cdots x^{n-1}\varphi \end{pmatrix} = B'\,. \]
Durch diese Vertauschung von Zeilen und Colonnen wird das Vorzeichen nicht geändert, wie man aus dem Diagonalglied der Determinante $B'$ leicht erkennt; es ist sonach
\[ B = B'\,. \]


% ===== pdf idx 171, printed 157.0 =====
% !PRINTED 157
\weberpage{157}
\textbf{\S.\,11. Eigenschaften der Resultante $R_{f,\varphi}$.\quad 157}

Führen wir nun die Differentiation nach den im Nenner befindlichen Grössen aus, so erhalten wir
\[
B' = \left|\begin{array}{ccccc|ccccc}
-b_0 & 0 & \cdots & \cdots & 0 & a_0 & 0 & \cdots & \cdots & 0\\
-b_1 & -b_0 & 0 & \cdots & 0 & a_1 & a_0 & 0 & \cdots & 0\\
-b_2 & -b_1 & -b_0 & 0 & \cdots\ 0 & a_2 & a_1 & a_0 & 0 & \cdots\ 0\\
\cdot & \cdot & \cdot & \cdot & \cdot & \cdot & \cdot & \cdot & \cdot & \cdot\\
-b_{n-1} & -b_{n-2} & \cdots & \cdots & b_0 & a_{n-1} & \cdots & a_1 & a_0 & \\
0 & 0 & \cdots & \cdots & 0 & 1 & 0 & \cdots & \cdots & 0\\
0 & \cdots & \cdots & \cdots & \cdots & 0 & 0 & 1 & \cdots & \cdots\\
\cdot & \cdot & \cdot & \cdot & \cdot & \cdot & \cdot & \cdot & \cdot & \cdot\\
0 & \cdots & \cdots & \cdots & \cdots & \cdots & 0 & 0 & 0 & 1
\end{array}\right|\,.
\]
Da diese Determinante in den letzten $n$ Zeilen, von den Diagonalelementen abgesehen, die durchweg 1 sind, nur verschwindende Elemente besitzt, so zerfällt sie nach dem Laplace'schen Satze in
\[
B' = \begin{vmatrix} -b_0 & 0 & 0\ldots 0\\ -b_1 & -b_0 & 0\ldots 0\\ -b_2 & -b_1 & -b_0\ldots 0\\ \cdot&\cdot&\cdot&\cdot\\ -b_{n-1}&-b^{n-2}&\ldots& b_0 \end{vmatrix}
\cdot
\begin{vmatrix} 1&0&0\ldots 0\\ 0&1&0\ldots 0\\ 0&0&1\ldots 0\\ \cdot&\cdot&\cdot&\cdot\\ 0&0&0\ldots 1 \end{vmatrix}
= (-1)^n\cdot b_0{}^n\,.
\]
Bilden wir nun das Product $B\cdot R_{f,\varphi}$, so erhalten wir
\[ (-1)^nb_0{}^n\cdot R_{f,\varphi} \]
\[ = \begin{pmatrix} (a_0x^{n-1}\cdots+a_{n-1})\varphi-(b_0x^{n-1}\cdots+b_{n-1})f\cdots, & a_0\varphi-b_0f, & x^{n-1}\varphi,\cdots,x^0\varphi\\ x^{2n-1}\,, & x^{2n-2}\cdots, & x^n,\ x^{n-1}\,,\ldots,x^0 \end{pmatrix}\,. \]
Um diese Determinante zu reduciren, dividiren wir sie mit der Determinante
\[ A = \begin{pmatrix} x^{n-1}\,, & x^{n-2}\ldots x^0\,, & x^{n-1}\varphi\,, & x^{n-2}\varphi\ldots x^0\varphi\\ x^{2n-1}\,, & x^{2n-2}\ldots\ldots x^{n-1} & \ldots\ldots\ldots & x^0 \end{pmatrix}\,, \]
welche ebenfalls den Werth
\[ A = (-1)^{n^2}b_0 = (-1)^nb_0{}^n \]
besitzt.

Wir erhalten
\[ \left((-1)^nb_0{}^n\cdot R_{f,\varphi}\right) : (-1)^nb_0{}^n \]
\[ = \begin{pmatrix} (a_0x^{n-1}\cdots+a_{n-1})\varphi-(b_0x^{n-1}+\cdots b_{n-1})f\cdots, & a_0\varphi-b_0f, & x^{n-1}\varphi,\ldots x^0\varphi\\ x^{n-1}\,, & x^{n-2},\ldots\ldots x^0, & x^{n-1}\varphi,\ldots x^0\varphi \end{pmatrix}\,, \]
%CONT>>


% ===== pdf idx 172, printed 158.0 =====
% !PRINTED 158
\weberpage{158}
\textbf{158\quad Zweiter Theil. Anwendungen der Determinantentheorie.}

%<<CONT
oder, wenn wir die übereinanderstehenden gleichen Grössen unterdrücken und die Matrix der übrigbleibenden $n$ ersten Zeilen stürzen:
{\footnotesize\setlength{\arraycolsep}{3.5pt}
\[
R_{f,\varphi} = (-1)^{\frac{n(n-1)}{2}}
\begin{pmatrix}
a_0\varphi-b_0f, & (a_0x+a_1)\varphi-(b_0x+b_1f)\cdots, & (a_0x^{n-1}\cdots+a_{n-1})\varphi-(b_0x^{n-1}+\cdots+b_{n-1})f\\
x_0\,, & x^1,\ x^2\cdots, & x^{n-1}
\end{pmatrix}\,.
\]
}
Führen wir die Differentiationen aus, so erhalten wir
\[
R_{f,\varphi} = (-1)^{\frac{n(n-1)}{2}}\,\cdot
\]
\weq{I}{\begin{array}{llll}
(a_0b_n), & (a_0b_{n-1}), & (a_0b_{n-2}) & \cdots\cdots\cdots\cdots,\ (a_0b_1)\\
(a_1b_n), & (a_1b_{n-1})+(a_0b_n), & (a_1b_{n-2})+(a_0b_{n-1}) & \cdots\cdots,\ (a_0b_2)\\
(a_2b_n), & (a_2b_{n-1})+(a_1b_n), & (a_2b_{n-2})+(a_1b_{n-1})+(a_0b_n) & \cdots,\ (a_0b_3)\\
(a_3b_n) & \cdot\ \cdot\ \cdot\ \cdot\ \cdot\ \cdot\ \cdot\ \cdot\ \cdot\ \cdot\ \cdot\ \cdot & &\\
\cdot & \cdot & \cdot & \cdot\\
\cdot & \cdot & \cdot & \cdot\\
(a_{n-1}b_n), & (a_{n-2}b_n), & (a_{n-3}b_n) & \cdots\cdots\cdots,\ (a_0b_n)
\end{array}}
Das Gesetz, nach welchem die Elemente der Determinante gebildet werden, ist hieraus bereits erkenntlich. So ist für $n=5$
\[ R_{f,\varphi} = \]
\[
\begin{vmatrix}
(a_0b_5), & (a_0b_4), & (a_0b_3), & (a_0b_2), & (a_0b_1)\\
(a_1b_5), & (a_1b_4)+(a_0b_5), & (a_1b_3)+(a_0b_4), & (a_1b_2)+(a_0b_3), & (a_0b_2)\\
(a_2b_5), & (a_2b_4)+(a_1b_5), & (a_2b_3)+(a_1b_4)+(a_0b_5), & (a_1b_3)+(a_0b_4), & (a_0b_3)\\
(a_3b_5), & (a_3b_4)+(a_2b_5), & (a_2b_4)+(a_1b_5), & (a_1b_4)+(a_0b_5), & (a_0b_4)\\
(a_4b_5), & (a_3b_5), & (a_2b_5), & (a_1b_5), & (a_0b_5)
\end{vmatrix}\,.
\]
Man nennt diese Determinante die Bézout'sche Determinante. Sie ist rücksichtlich der zweiten Diagonale symmetrisch; denn es ist das Element
\[ a_{n-i,k} = a_{n-k,i}\,. \]
Für alle wirklichen Berechnungen mittelst Determinanten ist die Darstellung der Resultante in dieser Form die wichtigste, und wir wollen daher dieselbe noch auf eine zweite Art berechnen.

151. \textit{Cayley'sche Methode für die Berechnung der Bézout'schen Determinante.} Vertauschen wir in der Determinante $R_{f,\varphi}$ Nr. 150 (I) die letzte Colonne mit der ersten, die vorletzte mit der zweiten u. s. w., dann kommen die Elemente der zweiten Diagonale in die erste zu stehen, und die Determinante ist nun symmetrisch nach dem Gesetze
\[ a_{ik} = a_{ki}\,. \]


% ===== pdf idx 173, printed 159.0 =====
% !PRINTED 159
\weberpage{159}
\S.\,11. Eigenschaften der Resultante $R_{f,\varphi}$.

Das Vorzeichen der Determinante \"andert sich hierbei nicht (vergl.
Beispiel 3 Nr. 23). Bezeichnen wir nun die $n^2$ Elemente dieser letzten
Determinante der Reihe nach mit
\[
\begin{matrix}
c_{00} & c_{01} & c_{02} \cdots c_{0\,n-1} \\
c_{10} & c_{11} & c_{12} \cdots c_{1\,n-1} \\
c_{20} & c_{21} & c_{22} \cdots c_{2\,n-1} \\
\cdot & \cdot & \cdot\ \cdot\ \cdot\ \cdot\ \cdot\ \cdot\ \cdot \\
c_{n-1,0} & c_{n-2,0} & \cdots\cdots c_{n-1,\,n-1}\,,
\end{matrix}
\]
% [sic] second entry of the last row is printed c_{n-2,0}, not c_{n-1,1}
indem wir einstweilen von der noch nicht streng bewiesenen Symmetrie
absehen, so k\"onnen wir zur Berechnung dieser Elemente den von Cayley
angegebenen Weg einschlagen. Cayley bemerkte n\"amlich, dass der Z\"ahler
in $R_{f,\varphi}$ folgende Eigenschaft besitzt. Bildet man die
Determinante
\[
\frac{1}{y-x}
\begin{vmatrix} f(x), & \varphi(x) \\ f(y), & \varphi(y) \end{vmatrix}
=
\begin{vmatrix} f(x), & \varphi(x) \\[4pt]
\dfrac{f(y)-f(x)}{y-x}, & \dfrac{\varphi(y)-\varphi(x)}{y-x} \end{vmatrix}
\]
\[
=
\begin{vmatrix}
f(x), & a_0 y^{n-1} + (a_0 x + a_1) y^{n-2} \cdots + (a_0 x_{n-1} + \cdots a_{n-1}) y^0 \\
\varphi(x), & b_0 y^{n-1} + (b_0 x + b_1) y^{n-2} \cdots + (b_0 x^{n-1} + \cdots b_{n-1}) y^0
\end{vmatrix},
\]
% [sic] first row prints x_{n-1} (subscript), second row prints x^{n-1} (superscript) -- inconsistent as printed
und ordnet ihre Entwicklung nach Potenzen von $y$, so sind, wie man
sieht, die Coefficienten dieser Entwicklung gerade die Gr\"ossen im
Z\"ahler der Determinante $R_{f,\varphi}$. Wir k\"onnen daher schreiben:
\[
\frac{1}{y-x}\left\{f(x)\varphi(y) - \varphi(x)f(y)\right\}
\]
\[
= (c_{00}x^{n-1}+c_{01}x^{n-2}+\cdots+c_{0,n-1})y^{n-1}+(c_{10}x^{n-1}+c_{11}x^{n-2}+\cdots c_{1,n-1})x^{n-2}
\]
% [sic] the second bracketed term is printed multiplied by x^{n-2}, not y^{n-2}
\[
+ \cdots + (c_{n-1,0}x^{n-1} + c_{n-1,1}x^{n-2} + \cdots + c_{n-1,n-1}x^0)y^0
\]
\weq{I}{
\sum_{\varrho=0,\,\sigma=0}^{\varrho=n-1,\,\sigma=n-1} c_{\varrho\sigma}\, x^{n-\varrho-1}\, y^{n-\sigma-1}\,.
}
Wir k\"onnen aber die Determinante $f(x)\varphi(y) - f(y)\varphi(x)$ auch noch
auf eine andere Weise entwickeln, wodurch wir Relationen f\"ur die
Coefficienten $c_{ik}$ gewinnen. Es ist n\"amlich (vgl. Nr. 82)
\[
f(x)\varphi(y) - \varphi(x)f(y) =
\begin{vmatrix} a_0 & a_1 & a_2 & \cdots & a_n \\ b_0 & b_1 & b_2 & \cdots & b_n \end{vmatrix}
\cdot
\begin{vmatrix} x^n & x^{n-1} & x^{n-2} & \cdots & x^0 \\ y^n & y^{n-1} & y^{n-2} & \cdots & y^0 \end{vmatrix}
\]
\[
= \sum d_{ik}(x^{n-i}y^{n-k} - x^{n-k}y^{n-i}) = \sum d_{ik}\, x^{n-k}y^{n-k}(x^{k-i}-y^{k-i})\,,
\]
wobei
\[
d_{ik} = a_i b_k - b_i a_k \quad \text{und} \quad k > i\,.
\]


% ===== pdf idx 174, printed 160.0 =====
% !PRINTED 160
\weberpage{160}
Zweiter Theil. Anwendungen der Determinantentheorie.

Dividiren wir diese Relation mit $(x-y)$, so erhalten wir
\[
\frac{f(x)\varphi(y)-f(y)\varphi(x)}{x-y} = \sum_{\lambda=k-i-1}^{\lambda=0} d_{ik} \cdot x^{n-k}y^{n-k} \cdot x^{k-i-\lambda-1} \cdot y^{\lambda}
\]
\[
= \sum_{\lambda=k-i-1}^{\lambda=0} d_{ik} \cdot x^{n-i-\lambda-1} \cdot y^{n-k+\lambda}\,,
\]
wo nun $\Sigma$ eine dreifache Summe bedeutet: \"uber $i$, \"uber $k$ und \"uber
$\lambda$. Setzt man in dieser Summe
\[
n-i-\lambda-1 = n-1-\varrho
\]
und
\[
n-k+\lambda \phantom{n-i-\lambda-1} = n-1-\sigma\,,
\]
also
\[
i = \varrho - \lambda
\]
und
\[
k = \sigma+\lambda+1\,,
\]
dann wird
\weq{II}{
\frac{1}{x-y} \cdot \left\{f(x)\varphi(y) - \varphi(x)f(y)\right\} = \sum d_{\varrho-\lambda,\,\sigma+\lambda+1}\, x^{n-1-\varrho}\, y^{n-1-\sigma}\,.
}
Vergleichen wir diese Entwicklung (II) mit der Entwicklung (I), so
erhalten wir die Relation
\weq{III}{
c_{\varrho\sigma} = \sum_{\lambda=0}^{\lambda=\varrho-\sigma} d_{\varrho-\lambda,\,\sigma+\lambda+1}\,.
}
Dabei kann $\lambda$ sich \"uber positive oder negative Zahlenwerthe
erstrecken, je nachdem $\varrho \gtrless \sigma$ ist; nur ist zu beachten, dass
gem\"ass ihrer Definition
\[
d_{\varrho\varrho} = 0\,, \quad \text{und} \quad d_{\varrho\sigma} = -d_{\sigma\varrho}\,.
\]
Aus der Relation (III) geht \"uberdies hervor, dass
\[
c_{\varrho\sigma} = c_{\sigma\varrho}\,,
\]
d. h., dass in der That die B\'ezout'sche Determinante eine symmetrische
ist. Denn es ist gem\"ass dieser Formel:
\weq{IV}{
\left.
\begin{aligned}
c_{\varrho\sigma} &= d_{\varrho,\sigma+1} + d_{\varrho-1,\sigma+2} + \cdots d_{\sigma,\varrho+1} \\
c_{\sigma\varrho} &= d_{\sigma,\varrho+1} + d_{\sigma-1,\varrho+2} + \cdots d_{\varrho,\sigma+1}
\end{aligned}
\right\}\,.
}
% [printed with a single right brace spanning both lines]
Die Differenz beider Relationen aber ist Null.

\textit{Anmerkung.} Aus der B\'ezout'schen Form der Resultante geht
unmittelbar eine wichtige Eigenschaft der Resultante hervor. Sind $f$
und $\varphi$ zwei Functionen $m^{\mathrm{ten}}$ Grades und $F$ und $\Phi$
irgend zwei lineare Combinationen derselben, etwa:
\[
F = \varkappa f + \lambda \varphi
\]
\[
\Phi = \mu f + \nu \varphi\,,
\]
so ist:
\[
R_{F,\Phi} = \Delta^m \cdot R_{f,\varphi}\,,
\]
wobei
\[
\Delta = \begin{vmatrix} \varkappa, & \lambda \\ \mu, & \nu \end{vmatrix}\,.
\]


% ===== pdf idx 175, printed 161.0 =====
% !PRINTED 161
\weberpage{161}
\S.\,11. Eigenschaften der Resultante $R_{f,\varphi}$.

Denn $R_{F,\Phi}$ geht direct aus $R_{f,\varphi}$ hervor, wenn wir in der
Determinante das Element
\[
d_{ik} = a_i b_k - b_i a_k
\]
ersetzen durch
\[
D_{ik} = A_i B_k - B_i A_k\,.
\]
Es ist aber:
\[
D_{ik} = \begin{vmatrix} \varkappa a_i + \lambda b_i, & \mu a_i + \nu b_i \\ \varkappa a_k + \lambda b_k, & \mu a_k + \nu b_k \end{vmatrix}
= \begin{vmatrix} \varkappa & \lambda \\ \mu & \nu \end{vmatrix} \cdot \begin{vmatrix} a_i & b_i \\ a_k & b_k \end{vmatrix}
\]
\[
= \Delta \cdot d_{ik}\,.
\]
Jedes Element der $m$-zeiligen Determinante $R_{F,\Phi}$ enth\"alt also
den Factor $\Delta$, folglich $R_{F,\Phi}$ selbst den Factor $\Delta^m$.

152. \textit{Beispiel.} Wir ermitteln unter Anwendung der ersten Formel
(IV) Nr. 151 die Elemente der B\'ezout'schen Determinante zweier
Functionen sechsten Grades
{\footnotesize
\[
(-1)^{15}
\begin{pmatrix}
a_0\varphi-b_0f, & (a_0x+a_1)\varphi-(b_0x-b_1)f, & \cdots & (a_0x^5+a_1x^4\cdots+a_5)\varphi-(b_0x^5+b_1x^4\cdots+b_5)f \\
x^5, & x^4, & x^3, & x^2, & x, & x^0
\end{pmatrix}
\]
}
% [sic: matrix is printed with a ragged number of entries per row, reproduced as printed]
\[
= -(c_{00}\ c_{11}\ c_{22}\ c_{33}\ c_{44}\ c_{55})\,.
\]
Dabei gen\"ugt es, wegen der Symmetrie der Determinante jene Elemente
$c_{\varrho\sigma}$ zu berechnen, f\"ur welche $\varrho \geqq \sigma$ ist.
Ber\"ucksichtigt man, dass
\[
d_{ii} = 0\,, \quad d_{-i,k} = 0\,, \quad d_{ik} = -d_{ki} \quad \text{und} \quad i,k \leqq \sigma\,,
\]
so erh\"alt man aus der Formel
\[
c_{\varrho\sigma} = d_{\varrho,\sigma+1} + d_{\varrho-1,\sigma+2} + d_{\varrho-2,\sigma+3} + d_{\varrho-3,\sigma+4} \cdots + d_{\sigma,\varrho+1}
\]
\[
\begin{aligned}
c_{00} &= d_{01} & c_{33} &= d_{34} + d_{25} + d_{16} \\
c_{10} &= d_{02}\,, \quad c_{11} = d_{12} + d_{03}\,, & c_{43} &= d_{35} + d_{26} \\
c_{20} &= d_{03}\,, \quad c_{21} = d_{13} + d_{04}\,, \quad c_{22} = d_{23} + d_{14} + d_{05}\,, & c_{53} &= d_{36} \\
c_{30} &= d_{04}\,, \quad c_{31} = d_{14} + d_{05}\,, \quad c_{32} = d_{24} + d_{15} + d_{06}\,, & c_{44} &= d_{45} + d_{36} \\
c_{40} &= d_{05}\,, \quad c_{41} = d_{15} + d_{06}\,, \quad c_{42} = d_{25} + d_{16}\,, & c_{54} &= d_{46} \\
c_{50} &= d_{06}\,, \quad c_{51} = d_{16}\,, \quad c_{52} = d_{26}\,, & c_{55} &= d_{56}\,.
\end{aligned}
\]
Die Resultante wird also:
\[
R_{f,\varphi} = -
\begin{vmatrix}
d_{01} & d_{02} & d_{03} & d_{04} & d_{05} & d_{06} \\
d_{02} & d_{12}+d_{03} & d_{13}+d_{04} & d_{14}+d_{05} & d_{15}+d_{06} & d_{16} \\
d_{03} & d_{13}+d_{04} & d_{23}+d_{14}+d_{05} & d_{24}+d_{15}+d_{06} & d_{25}+d_{16} & d_{26} \\
d_{04} & d_{14}+d_{05} & d_{24}+d_{15}+d_{06} & d_{34}+d_{25}+d_{16} & d_{35}+d_{26} & d_{36} \\
d_{05} & d_{15}+d_{06} & d_{25}+d_{16} & d_{35}+d_{26} & d_{45}+d_{36} & d_{46} \\
d_{06} & d_{16} & d_{26} & d_{36} & d_{46} & d_{56}
\end{vmatrix}\,.
\]


% ===== pdf idx 176, printed 162.0 =====
% !PRINTED 162
\weberpage{162}
Zweiter Theil. Anwendungen der Determinantentheorie.

153. \textit{Resultante aus einer Function $f$ und dem Producte zweier
Functionen $\varphi$ und $\psi$.} Es sei
\[
f = a_0 x^m + a_1 x^{m-1} + \cdots + a_m
\]
\[
\varphi = b_0 x^p + b_1 x^{p-1} + \cdots + b_p
\]
\[
\psi = c_0 x^q + c_1 x^{q-1} + \cdots + c_q\,,
\]
wobei die Summe $p+q$ der Grade von $\varphi$ und $\psi$ den Grad $m$
von $f$ nicht \"ubersteigen soll. Die Resultante zu $f$ und
$\varphi \cdot \psi$ ist dann dargestellt durch
\[
R_{f,\varphi\cdot\psi} =
\begin{pmatrix}
x^{p+q-1}f, & x^{p+q-2}f, & \cdots, & x^0f, & x^{m-1}\varphi\cdot\psi, & x^{m-2}\varphi\cdot\psi & \cdots, & x^0\cdot\varphi\psi \\
x^{m+p+q-1}, & x^{m+p+q-2}, & \cdots, & x^m, & x^{m-1}, & x^{m-2} & \cdots, & x^0
\end{pmatrix}\,.
\]
Diese Resultante hat die Eigenschaft, in zwei Factoren sich zu
spalten, n\"amlich in die Resultanten $R_{f,\varphi}$ und $R_{f,\psi}$.
Um dies zu zeigen, werden wir $R_{f,\varphi\cdot\psi}$ zun\"achst so
transformiren, dass sich unter Anwendung des Laplace'schen Satzes in
erster Linie die Factorenzerf\"allung ergiebt; es er\"ubrigt dann nur
noch nachzuweisen, dass die erhaltenen Factoren in der That mit
$R_{f,\varphi}$ und $R_{f,\psi}$ \"ubereinstimmen.

Die Transformationen erreichen wir, nachdem wir in $R_{f,\varphi\psi}$
eine passende Vertauschung von Zeilen mit Zeilen vorgenommen haben,
durch Multiplication mit $A =$
\[
\begin{pmatrix}
x^{p-1}f\cdot\psi, & x^{p-2}f\cdot\psi, & \cdots x^0\cdot f\psi, & x^{m-1}\varphi\cdot\psi, & x^{m-2}\varphi\cdot\psi \cdots x^0\varphi\cdot\psi, & x^{q-1}f, & x^{q-2}f, & \cdots x^0f \\
x^{p+q-1}f, & x^{p+q-2}f, & \cdots x^q\cdot f, & x^{m-1}\varphi\cdot\psi, & x^{m-2}\varphi\cdot\psi \cdots x^0\varphi\cdot\psi, & x^{q-1}f, & x^{q-2}f, & \cdots x^0f
\end{pmatrix}
\]
und durch Division mit
\[
B =
\begin{pmatrix}
x^{m+p-1}\psi, & x^{m+p-2}\psi & \cdots x^0\psi, & x^{q-1}, & x^{q-2}, & \cdots x^0 \\
x^{m+p+q-1}, & x^{m+p+q-2}, & \cdots x^m, & x^{q-1}, & x^{q-2}, & \cdots x^0
\end{pmatrix}\,.
\]
Man \"uberzeugt sich leicht, dass die Determinante $A$ den Werth
$c_0^{\,p}$, die Determinante $B$ den Werth $c_0^{\,m+p}$ besitzt.

Vertauschen wir nun in $R_{f,\varphi\psi}$ die Zeilen so, dass die $m$
letzten Zeilen sich direct an die $p$ ersten Zeilen anschliessen,
w\"ahrend die $q$ dazwischen liegenden Zeilen die letzten werden, dann
\"andert sich $R_{f,\varphi\psi}$ nach Nr. 139 um den Factor $(-1)^{mq}$;
die Reihenfolge der Gr\"ossen im Z\"ahler von $R_{f,\varphi\psi}$
stimmt aber alsdann mit der Reihenfolge der Gr\"ossen im Nenner von
$A$ \"uberein, und wir erhalten daher durch Multiplication mit $A$
\[
c_0^{\,p} \cdot R_{f,\varphi\cdot\psi} =
\]
\[
(-1)^{mq}
\begin{pmatrix}
x^{p-1}f\psi, & x^{p-2}f\psi, & \cdots, & x^0f\psi, & x^{m-1}\varphi\psi, & \cdots, & x^0\varphi\psi, & x^{q-1}f, & \cdots x^0f \\
x^{m+p+q-1}, & x^{m+p+q-2}, & \cdots, & \cdots\cdots\cdots\cdots\cdots\cdots & & & & x^0
\end{pmatrix}\,.
\]


% ===== pdf idx 177, printed 163.0 =====
% !PRINTED 163
\weberpage{163}
\S\ 11. Eigenschaften der Resultante $R_{f,\varphi}$.

Dividiren wir diese Determinante mit $B = c_0^{\,m+p}$, so kommt:
\[
\frac{c_0^{\,p}\cdot R_{f,\varphi\cdot\psi}}{c_0^{\,m+p}} =
(-1)^{mq}
\begin{pmatrix}
x^{p-1}f\cdot\psi,\ x^{p-2}f\cdot\psi,\ \ldots x^0\cdot f\psi,\ x^{m-1}\varphi\cdot\psi \ldots x^0\cdot\varphi\psi,\ x^{q-1}f..x^0f\\
x^{m+p-1}\psi,\ x^{m+p-2}\psi,\ \ldots x^m\psi\ ,\ x^{m-1}\psi \ldots x^0\cdot\psi\ ,\ x^{q-1}\ldots x^0
\end{pmatrix}.
\]
154. \textit{Berechnung und Zerf\"allung von $R_{f,\varphi\cdot\psi}$.} F\"uhrt man hier die Differentiation nach den Gr\"ossen im Nenner aus, so erh\"alt man folgendes Resultat:
\weq{I}{
R_{f,\varphi\cdot\psi} = (-1)^{m\cdot q}\cdot c_0^{\,m}\times
\begin{vmatrix}
a_0\,a_1\,a_2\ldots\ldots a_m\ \ 0\ \ 0\ldots\ldots\ 0\ \ 0\ \ 0..0\\
0\,a_0\,a_1\ldots\ldots a_{m-1}a_m\ \ 0\ldots\ldots\ 0\ \ 0\ \ 0..0\\
0\ 0\,a_0\ldots\ldots a_{m-2}a_{m-1}a_m\ldots\ldots\ 0\ \ 0\ \ 0..0\\
\cdot\ \cdot\ \cdot\ \cdot\ \cdot\ \cdot\ \cdot\ \cdot\ \cdot\ \cdot\\
0\ 0\ 0..a_0\ \ a_1\ldots\ldots\ldots\ldots a_{m-2}a_{m-1}a_m\ 0\ \ 0\ldots\\
b_0\,b_1\,b_2..b_p\ \ 0\ \ 0\ \ 0\ \ 0\ldots\ldots\ldots\ 0\ \ 0\ \ 0\ldots\\
0\,b_0\,b_1..b_{p-1}b_p\ \ 0\ \ 0\ \ 0\ldots\ldots\ldots\ 0\ \ 0\ \ 0\ldots\\
0\ 0\,b_0\ldots\ldots b_{p-1}b_p\ \ 0\ \ 0\ldots\ldots\ldots\ 0\ \ 0\ \ 0..0\\
\cdot\ \cdot\ \cdot\ \cdot\ \cdot\ \cdot\ \cdot\ \cdot\ \cdot\\
0\ 0\ldots\ldots\ldots\ldots\ldots b_0\ldots b_{p-2}b_{p-1}b_p\ \ 0\ \ 0..0\\
0\ 0\ 0\ldots\ldots e_{11}\ \ e_{12}\ \ e_{13}\ldots\ldots\ldots\ldots e_{1m}\,e_{1,m+1}..e_{1,m+q}\\
0\ 0\ 0\ldots\ldots 0\ \ e_{22}\ \ e_{23}\ldots\ldots\ldots\ldots e_{2m}\,e_{2,m+1}..e_{2,m+q}\\
0\ 0\ldots\ldots\ldots 0\ \ 0\ \ e_{33}\ldots\ldots\ldots\ldots e_{3m}\,e_{3,m+1}..e_{3,m+q}\\
\cdot\ \cdot\ \cdot\ \cdot\ \cdot\ \cdot\ \cdot\ \cdot\ \cdot\ \cdot\\
0\ 0\ 0\ldots\ldots\ldots\ldots\ldots\ldots e_{qq}\ \ e_{q,q+1}\ldots\ldots\ldots\ldots e_{q,m+q}
\end{vmatrix}
}
Die Elemente $e_{ik}$ der letzten $q$ Zeilen dieser Determinante sind die Differentialquotienten der Gr\"ossen $x^{q-1}\cdot f,\ x^{q-2}\cdot f\ldots x^0\cdot f$ nach $x^{m-1}\cdot\psi,\ x^{m-2}\cdot\psi\ldots x^0\cdot\psi,\ x^{q-1},\ x^{q-2}\ldots x^0$ und man erh\"alt dieselben, indem man wie in Nr. 149 die Functionen $x^\nu\cdot f$ sich dargestellt denkt in der Form
\[
x^\nu\cdot f = \psi\cdot g_\nu(x) + h_\nu(x),
\]
wobei $g_\nu(x)$ vom Grade $m+\nu-q$ und $h_\nu(x)$ vom Grade $q-1$ ist. Insbesondere sind die $q^2$ Elemente, welche von den letzten $q$ Colonnen ausgeschnitten werden, die Differentialquotienten der $q$ Reste $h_\nu(x)$ nach $x^{q-1},\ x^{q-2}\ldots x^0$.

Zerf\"allt man daher nach dem Satze von Laplace die erhaltene Determinante in ihre beiden Factoren, so ist der erste, wie man unmittelbar erkennt, die Resultante
%CONT>>


% ===== pdf idx 178, printed 164.0 =====
% !PRINTED 164
\weberpage{164}
Zweiter Theil. Anwendungen der Determinantentheorie.

%<<CONT
\weq{II}{
R_{f,\varphi} =
\begin{vmatrix}
a_0\,a_1\,a_2\ldots a_m\ \ 0\ \ 0\ \ldots..0\\
0\ a_0\,a_1\ldots\ldots\ a_m\ \ 0\ \ldots..0\\
0\ 0\,a_0\ldots\ldots a_{m-1}a_m\ldots..0\\
\cdot\ \cdot\ \cdot\ \cdot\ \cdot\ \cdot\ \cdot\ \cdot\ \cdot\\
0\ 0\ 0\ldots a_0\ \ a_1\ \ a_2\ldots\ldots a_m\\
b_0\,b_1\,b_2\ldots b_p^{*}\ \ 0\ \ 0\ldots\ldots\\
0\ b_0\,b_1\ldots b_{p-1}b_p\ \ 0\ \ldots\ldots\\
0\ 0\,b_0\ldots b_{p-2}b_{p-1}b_p\ldots\ldots\\
\cdot\ \cdot\ \cdot\ \cdot\ \cdot\ \cdot\ \cdot\ \cdot\ \cdot\\
0\ 0\ 0\ldots\ldots\ldots\ldots b_0\,b_1\ldots b_p
\end{vmatrix}
}
% ?? uncertain: small mark above $b_p$ in the fifth row of matrix (II) (printed as $b_p^{*}$); could be an ink artifact or a footnote-style mark, reproduced as printed.
Der zweite Factor aber ist -- vergleiche Nr. 149 -- als Functional-determinante der $q$ Reste vom Grade $q-1$, die sich bei Division von $x^{q-1}\cdot f,\ x^{q-2}\cdot f\ldots,\ x^0\cdot f$ mit $\psi$ ergeben, nichts anderes als die Resultante von $f$ und $\psi$ bis auf einen Zahlenfactor. Es ist n\"amlich nach Nr. 148
\weq{III}{
R_{f,\psi} = (-1)^{mq}\cdot c_0^{\,m}\cdot
\begin{pmatrix}
x^{q-1}\cdot f,\ x^{q-2}\cdot f\ldots x^0f\\
x^{q-1}\ \ ,\ x^{q-2}\ldots\ldots x^0
\end{pmatrix}
}
Die Functionaldeterminante rechts ist gerade der zweite Factor der Determinante (I), so dass wir, wenn wir in (I) f\"ur die Factoren ihre Werthe aus (II) und (III) einsetzen, erhalten:
\[
R_{f,\varphi\cdot\psi} = (-1)^{mq}\cdot c_0^{\,m}\cdot R_{f,\varphi}\cdot R_{f,\psi}\cdot \frac{1}{(-1)^{mq}c_0^{\,m}}
\]
oder:
\[
R_{f,\varphi\cdot\psi} = R_{f,\varphi}\cdot R_{f,\psi}.
\]
155. \textit{Beispiel} 1. Ich nehme die Resultante der Function
\[
f = a_0x^3 + a_1x^2 + a_1x + a_3 % [sic]
\]
und des Productes $\varphi\cdot\psi$, wo
\[
\varphi = b_0x^2+b_1x+b_2
\]
\[
\psi = c_0x+c_1.
\]
Dieselbe ist zun\"achst dargestellt durch
\[
R_{f,\varphi\cdot\psi} =
\begin{pmatrix}
x^2f,\ xf,\ f\ ,\ x^2\varphi\psi,\ x\varphi\psi,\ \varphi\psi\\
x^5\ \ ,\ x^4\ \ ,\ x^3,\ x^2\ \ \ ,\ x^1\ \ \ ,\ x^0
\end{pmatrix},
\]
oder nach den geeigneten Transformationen durch
\[
R_{f,\varphi\cdot\psi} = (-1)^3\cdot c_0^{\,3}
\begin{pmatrix}
xf\psi,\ f\psi\ ,\ x^2\varphi\psi,\ x\varphi\psi,\ x^0\varphi\psi,\ f\\
x^4\psi,\ x^3\psi,\ x^2\psi\ \ ,\ x^1\psi\ ,\ x^0\psi\ \ \ ,\ x^0
\end{pmatrix}.
\]
Wir bezeichnen die Gr\"ossen des Z\"ahlers mit $y_1,\ y_2\ldots y_6$, und schreiben dieselben in der Form:
%CONT>>


% ===== pdf idx 179, printed 165.0 =====
% !PRINTED 165
\weberpage{165}
\S\ 11. Eigenschaften der Resultante $R_{f,\varphi}$.

%<<CONT
\begin{gather*}
y_1 = \psi\,(x^4a_0+x^3a_1+x^2a_2+xa_3)\\
y_2 = \psi\,(x^3a_0+x^2a_1+xa_2+a_3)\\
y_3 = \psi\,(x^4b_0+x^3b_1+x^2b_2)\\
y_4 = \psi\,(x^3b_0+x^2b_1+xb_2)\\
y_5 = \psi\,(x^2b_0+xb_1+b_2)\\
y_6 = \psi\left(x^2\frac{a_0}{c_0}+x\frac{(a_1c_0)}{c_0^{\,2}}+\frac{c_0^{\,2}a_2-(c_0a_1)c_1}{c_0^{\,3}}\right)+\frac{a^3c_0^{\,3}-a_2c_0^{\,2}c_1+c_1^{\,2}c_0a_1-c_1^{\,3}a_0}{c_0^{\,3}} % [sic]
\\
= \psi\,(x^2e_{11}+xe_{12}+e_{13})+e_{14}.
\end{gather*}
Die partiellen Differentialquotienten dieser sechs Gr\"ossen nach $x^4\psi,\ x^3\psi,\ x^2\varphi,\ x^1\psi,\ x^0\psi,\ x^0$ liefern daher die Determinante:
\[
\varDelta =
\begin{vmatrix}
a_0\ a_1\ a_2\ a_3\ \ 0\ \ 0\\
0\ \ a_0\ a_1\ a_2\ a_3\ \ 0\\
b_0\ b_1\ b_2\ \ 0\ \ 0\ \ 0\\
0\ \ b_0\ b_1\ b_2\ \ 0\ \ 0\\
0\ \ 0\ \ b_0\ b_1\ b_2\ \ 0\\
0\ \ 0\ \ e_{11}\ e_{12}\ e_{13}\ e_{14}
\end{vmatrix}
\]
oder
\[
\varDelta = e_{14}\cdot
\begin{vmatrix}
a_0\ a_1\ a_2\ a_3\ \ 0\\
0\ \ a_0\ a_1\ a_2\ a_3\\
b_0\ b_1\ b_2\ \ 0\ \ 0\\
0\ \ b_0\ b_1\ b_2\ \ 0\\
0\ \ 0\ \ b_0\ b_1\ b_2
\end{vmatrix}.
\]
Der erste Factor $e_{14}$ ist aber nichts anderes als $\dfrac{(-1)^3\cdot R_{f,\psi}}{c_0^{\,3}}$; denn
\[
R_{f,\psi} =
\begin{vmatrix}
a_0\ a_1\ a_2\ a_3\\
c_0\ c_1\ 0\ \ 0\\
0\ \ c_0\ c_1\ 0\\
0\ \ 0\ \ c_0\ c_1
\end{vmatrix}
= a_0c_1^{\,3} - c_0c_1^{\,2}a_1 + c_0^{\,2}a_2c_0 - c_0^{\,3}a_3. % [sic]
\]
Wir erhalten daher in der That:
\[
R_{f,\varphi\cdot\psi} =
\begin{vmatrix}
a_0\ a_1\ a_2\ a_3\ \ 0\\
0\ \ a_0\ a_1\ a_2\ a_3\\
b_0\ b_1\ b_2\ \ 0\ \ 0\\
0\ \ b_0\ b_1\ b_2\ \ 0\\
0\ \ 0\ \ b_0\ b_1\ b_2
\end{vmatrix}
\cdot
\begin{vmatrix}
a_0\ a_1\ a_2\ a_3\\
c_0\ c_1\ 0\ \ 0\\
0\ \ c_0\ c_1\ 0\\
0\ \ 0\ \ c_0\ c_1
\end{vmatrix}
= R_{f,\varphi}\cdot R_{f,\psi}.
\]


% ===== pdf idx 180, printed 166.0 =====
% !PRINTED 166
\weberpage{166}
Zweiter Theil. Anwendungen der Determinantentheorie.

\textit{Beispiel} 2. Ist \textit{insbesondere} $\varphi$ eine Constante etwa gleich $\lambda$, so erhalten wir:
\[
R_{f,\lambda\varphi} = R_{f,\varphi}\cdot R_{f,\lambda}.
\]
Die Resultante aus einer Function $m^{\text{ten}}$ Grades und einer Constante $\lambda$ ist aber dargestellt durch
\[
R_{f,\lambda} =
\begin{vmatrix}
\lambda\ 0\ 0\ \ldots\ldots\ 0\\
0\ \lambda\ 0\ \ldots\ldots\ 0\\
0\ 0\ \lambda\ \ldots\ldots\ 0\\
\cdot\ \cdot\ \cdot\ \cdot\ \cdot\\
\cdot\ \cdot\ \cdot\ \cdot\ \cdot\\
0\ 0\ \ldots\ldots\ 0\ \lambda
\end{vmatrix}
= \lambda^m.
\]
Denn gem\"ass der Definition der Resultante in Nr. 130 und 131 ist dieselbe eine Determinante, welche die Constanten der einen Function in so vielen Zeilen enth\"alt, als der Grad der andern Function angiebt; in unserem Falle muss sie also die Coefficienten von $f$ in Null Zeilen, die von $\lambda$ in $m$ Zeilen enthalten. Demnach ist
\[
R_{f,\lambda\varphi} = \lambda^m\cdot R_{f,\varphi}.
\]

\begin{center}
\S.\,12.\\
\textbf{Fundamentalsatz der Algebra.}
\end{center}

156. \textit{Transcendenter Theil des Beweises.} Gauss hat bekanntlich den Satz bewiesen: \glqq Ist
\[
f(x) = x^n + a_1x^{n-1} + a_2x^{n-2} + \cdots + a_n = 0
\]
eine Gleichung $n^{\text{ten}}$ Grades mit durchweg reellen Coefficienten, so giebt es immer einen reellen oder imagin\"aren Werth $x=\xi_1$, f\"ur welchen die Gleichung identisch verschwindet.\grqq{}
% ?? uncertain: a marginal pencil-style mark (squiggle and a box) appears beside the word "dritten" below; it reads as a reader's annotation rather than print, the printed word itself is transcribed as shown.
Wir folgen hier dem dritten Beweise, den Gauss f\"ur diesen Satz erbracht und den Gordan wesentlich auf Eliminations- also Determinantenprocesse reducirt hat. (Vergl. M. A. Bd. X. 252.) Man weiss, dass der Beweis nicht durch rein algebraische Betrachtungen allein geliefert werden kann. Wir k\"onnen jedoch den Beweis von dem Augenblicke an im Gebiete der Algebra weiter f\"uhren, wo wir uns \"uberzeugt haben, dass jede Gleichung ungeraden Grades jedenfalls einen Werth von $x$ besitzt, f\"ur welchen sie besteht. Um aber dies zu erkennen, brauchen wir nur f\"ur $x$ einen Werth $\xi_1 > \Sigma a_g$ in $f$ zu substituiren, wo $\Sigma a_g$ die Summe aller absolut genommenen Coefficienten sein m\"oge,
%CONT>>


% ===== pdf idx 181, printed 167.0 =====
% !PRINTED 167
\weberpage{167}
dann überragt das erste Glied $(\Sigma a_g)^n$ der Gleichung alle andern an
absolutem Werthe. Je nachdem man also $x=+\Sigma a_g$ oder $x=-\Sigma a_g$
einträgt, wird das Substitutionsresultat positiv oder negativ. Es muss
also jedenfalls zwischen $+\Sigma a_g$ und $-\Sigma a_g$ ein Werth liegen, welcher
$f(x)$ verschwinden macht, und den wir dadurch ermitteln können, dass
wir diese Grenzen successive immer enger und enger ziehen. Hierin,
also in der Angabe der beiden Grenzen und der Möglichkeit sie zu
verengern, liegt das Charakteristische für die Transcendenz des Beweises.

Von hier ab jedoch können wir das Gebiet des Transcendenten
verlassen und unter der Voraussetzung, dass jede Gleichung \emph{ungeraden}
Grades mit reellen Co\"efficienten einen linearen Factor besitzt, in rein
algebraischer Weise zeigen:

\glqq Jede Gleichung besitzt einen linearen Factor.\grqq{}

157. \emph{Gedankengang für den algebraischen Theil des Beweises.} Unsere
Aufgabe ist, die Operationen anzugeben, wie wir zum linearen Factor
gelangen. Damit ist dann auch der Beweis geliefert, dass derselbe
existirt. Wir bedienen uns dabei eines gewissen Recursionsverfahrens,
indem wir die Gleichungen ordnen in Bezug auf die Schwierigkeit sie
aufzulösen. Setzen wir dann voraus, dass jede leichter auflösbare
Gleichung einen linearen Factor besitzt, so können wir zeigen, dass
auch jede schwieriger zu lösende Gleichung einen solchen haben muss.
Zu dem Zwecke werden wir den Beweis erbringen, dass jede Glei-
chung $f(x)=0$ sich immer in zwei Factoren spalten lässt, deren
mindestens einer gleich null gesetzt, leichter lösbar ist, als die Glei-
chung $f(x)=0$ selbst, ausgenommen in einem einzigen Falle, den
wir besonders behandeln werden.

158. \emph{Classification der Gleichungen nach dem Grade ihrer Auf-
lösbarkeit.} Wir können jede ganze Zahl $n$ in der Form schreiben:
\[ n = 2^{\mu}\cdot \nu\, , \]
wo $\nu$ eine ungerade Zahl, und $\mu$ alle Werthe von 0 bis $\infty$ annehmen
kann. Für $\mu=0$ erhalten wir alle ungeraden Zahlen. Alle Gleichungen,
deren Grad
\[ n = 2^{0}\cdot \nu \]
ist, sind nach Nr. 156 als leichter auflösbar zu erachten, als alle
übrigen und zwar um so leichter, je kleiner die ungerade Zahl $\nu$ ist.

An diese schliessen sich alle Gleichungen an vom Grade
\[ n = 2^{1}\cdot \nu\, , \]
sodann jene vom Grade
\[ n = 2^{2}\cdot \nu\, , \quad \text{u. s. w.,} \]
%CONT>>


% ===== pdf idx 182, printed 168.0 =====
% !PRINTED 168
\weberpage{168}
%<<CONT
bis endlich allgemein die Gleichungen folgen vom Grade
\[ n = 2^{\mu}\cdot \nu\, . \]
Es ist also in diesem Sinne beispielsweise eine Gleichung vierten
Grades als schwerer auflösbar zu betrachten als eine Gleichung sechsten
Grades, aber leichter als eine Gleichung zwölften Grades.

Wir nehmen dann an: Jede leichter auflösbare Gleichung, also
jede Gleichung, deren Grad
\[ n = 2^{\varrho}\cdot \nu \quad \text{oder} \quad n = 2^{\mu}\cdot \overline{\nu} \]
ist, wo $\varrho<\mu$, und $\overline{\nu}<\nu$, besitze einen linearen Factor. Dann lässt sich
zeigen, dass auch $f=0$ vom Grade
\[ n = 2^{\mu}\cdot \nu \]
einen solchen besitzt. Setzen wir nämlich einstweilen voraus, $f=0$
besitze überhaupt einen reellen oder imaginären Factor, so können
wir unschwer nachweisen, dass derselbe gleich null gesetzt eine leichter
auflösbare Gleichung repräsentirt als $f(x)=0$.

159. \emph{Voraussetzung: $f(x)=0$ besitze einen reellen Factor.} Gesetzt,
es sei bereits gezeigt, dass $f=0$ überhaupt einen Factor $\varphi(x)$ vom
Grade $n_1$ besitze mit reellen Co\"efficienten, dann ist:
\[ f = \varphi(x)^{n_1}\cdot \psi(x)^{n_2}\, , \]
wobei
\[ n_1 = 2^{\mu_1}\cdot \nu_1 \]
\[ n_2 = 2^{\mu_2}\cdot \nu_2\, , \]
also
\[ n = n_1+n_2 = 2^{\mu}\cdot \nu = 2^{\mu_1}\cdot \nu_1 + 2^{\mu_2}\cdot \nu_2\, , \]
Ist hier $\mu_1\leqq \mu$, dann ist $\varphi$ leichter auflösbar als $f$, da im un-
günstigsten Falle dann $\nu_1<\nu$ sein muss; es besitzt sonach $\varphi$ einen
linearen Factor und folglich auch $f$. Ist aber $\mu_1>\mu$, also etwa $\mu_1=\mu+\varrho$,
%CONT>>


% ===== pdf idx 183, printed 169.0 =====
% !PRINTED 169
\weberpage{169}
%<<CONT
günstigsten Falle dann $\nu_1<\nu$ sein muss; es besitzt sonach $\varphi$ einen
linearen Factor und folglich auch $f$. Ist aber $\mu_1>\mu$, also etwa $\mu_1=\mu+\varrho$,
dann ist
\[ n = 2^{\mu}\cdot \nu = 2^{\mu+\varrho}\cdot \nu_1 + 2^{\mu_2}\cdot \nu_2\, , \]
oder
\[ 2^{\mu_2}\cdot \nu_2 = 2^{\mu}(\nu-2^{\varrho}\cdot \nu_1)\, . \]
Nun ist $\nu-2^{\varrho}\cdot \nu_1$ jedenfalls ungerade, da $\nu$ ungerade ist und nach
Voraussetzung $\varrho$ nicht null sein kann, also ist:
\[ \mu_2=\mu \]
und
\[ \nu_2=\nu-2^{\varrho}\cdot \nu_1 \]
oder
\[ \nu_2<\nu\, , \]
folglich ist in diesem Falle der Factor $\psi(x)$ eine Gleichung, die leichter
zu lösen ist als $f(x)$, und $f(x)=0$ hat somit einen linearen Factor,
da ihn nach Voraussetzung $\psi(x)$ besitzt.

Wir haben somit das Resultat: Sobald $f(x)=0$ einen reellen
Factor $\varphi(x)$ besitzt, besitzt es auch einen linearen Factor.

160. \emph{Voraussetzung: $f(x)=0$ besitze einen imaginären Factor.}
Nehmen wir an, $f(x)=0$ habe den Factor
\[ \varphi(x) = p(x)+i\cdot q(x) = p+iq\, , \]
dann muss es nothwendig, da nur reelle Co\"efficienten in $f=0$ auf-
treten sollen, auch den Factor
\[ \chi(x) = p(x)-iq(x) = p-iq \]
besitzen. Haben dann $\varphi(x)$ und $\chi(x)$ einen Theiler gemein, so kann
dies nur davon herrühren, dass $p(x)$ und $q(x)$ denselben als gemein-
schaftlichen Theiler haben, der in Folge dessen reell sein muss. Daraus
folgt aber, dass dann auch $f(x)$ diesen reellen Factor besitzt, und so-
mit auch einen linearen Factor.

Haben aber $\varphi(x)$ und $\chi(x)$ keinen gemeinsamen Theiler, so kann
man nach Nr. 128 $(n)$ zwei Functionen $A(x)$ und $B(x)$ so finden, dass
\weq{1}{A^{n_1-1}(x)\cdot \varphi^{n_1}(x) + B^{n_1-1}(x)\chi^{n_1}(x) = 1 = A\cdot(p+iq)+B\cdot(p-iq)}
ist. Ausserdem ist $f$ darstellbar in den Formen:
\weq{2}{\begin{cases} f = [p(x)+iq(x)]L(x) = [p+iq]\cdot L \\ f = [p(x)-iq(x)]M(x) = [p-iq]\cdot M. \end{cases}}
Multipliciren wir Gleichung (1) mit $f(x)$, so erhalten wir
\[ f(x) = A\cdot f\cdot(p+iq) + B\cdot f\cdot(p-iq) \]
oder wegen der Relationen (2)
\[ f(x) = A\cdot M\cdot(p^2+q^2) + B\cdot L\cdot(p^2+q^2) \]
%CONT>>


% ===== pdf idx 184, printed 170.0 =====
% !PRINTED 170
\weberpage{170}
%<<CONT
\[ f(x) = A\cdot M\cdot(p^2+q^2) + B\cdot L\cdot(p^2+q^2) \]
oder
\[ f(x) = (p^2+q^2)(A\cdot M + B\cdot L)\, . \]
Nun ist aber $p^2+q^2$ reell, und so lange der Grad von $p^2+q^2<n$,
also der von $\varphi<\dfrac{n}{2}$ ist, besitzt folglich $f$ auch in diesem Falle einen
reellen Factor. Wäre $\varphi(x)$ vom Grade $\varrho>\dfrac{n}{2}$ gewesen, so hätten wir
die eben gemachten Schlüsse in Bezug auf den andern Factor $\psi(x)$
ziehen können, so dass also nach Nr. 159 in jedem der beiden Fälle
$f(x)$ einen linearen Factor besitzt. Der Fall, wo der Grad von
$\varphi(x)=p(x)+iq(x)$ gleich $\dfrac{n}{2}$ ist, also $p^2+q^2$ den nämlichen Grad
in $x$ hat, wie $f(x)$ selbst, bedarf einer speciellen Untersuchung. Wir
sind also zu folgendem Resultate gelangt:

\glqq Besitzt $f(x)=0$ überhaupt einen reellen Factor, oder einen
imaginären Factor vom Grade $n_1<\dfrac{n}{2}$, so besitzt $f(x)=0$
auch immer einen linearen Factor.\grqq{}

Es bleibt also noch zu beweisen: 1) dass jede Gleichung $f(x)=0$
irgend einen Factor $\varphi(x)$ besitzen muss; 2) dass, wenn dieser Factor
von der Form $p(x)+iq(x)$ ist und vom Grade $\dfrac{n}{2}$, die Annahme, $p(x)$
und $q(x)$ hätten dann keinen gemeinschaftlichen Factor, entweder zu
einem Widerspruche führt, oder wenn nicht, dass dann ein reeller
Factor a priori vorhanden ist.

161. \emph{Beweis, dass jede Gleichung in Factoren gespalten werden kann.}
Um zu zeigen, dass jede Gleichung $f(x)=0$ einen Factor besitzen
muss, bilden wir die beiden Functionen:
\weq{I}{\begin{cases} \dfrac{f(x+u)+f(x-u)}{2} = F(x) \\ \dfrac{f(x+u)-f(x-u)}{2u} = \Phi(x). \end{cases}}
Die Resultante der beiden Functionen in Bezug auf $x$ sei $R_{F,\Phi}$. Multi-
pliciren wir die zweite Gleichung mit $u$ und addiren sie zur ersten,
so kommt:
\[ f(x+u) = F(x) + u\,\Phi(x) \]
und nach dem Satze
\[ R_{f,\varphi} = R_{f+\lambda\varphi,\varphi} \quad \text{(vgl. Nr. 145)} \]
stimmt die Resultante $R_{F,\Phi}$ der beiden Functionen (I) überein mit der
Resultante $R_{f(x+u),\Phi(x)} = P_{(u)}$ der beiden Functionen
\weq{II}{\begin{cases} f(x+u) \quad \text{und} \\ \dfrac{f(x+u)-f(x-u)}{2u} = \Phi(x). \end{cases}}
Diese Resultante $R_{f,\Phi} = P_{(u)}$ der beiden letzten Functionen kann,
gleich null gesetzt, als Gleichung für $u$ aufgefasst werden und es
fragt sich jetzt, ist sie als solche leichter lösbar als $f$ oder nicht.
Im ersten Falle hat alsdann nach Voraussetzung $P_{(u)}=0$ einen linearen
Factor, d. h. man kann einen Werth $u$ berechnen, so dass $R_{f,\varphi}$ ver-
schwindet. Trägt man diesen speciellen Werth von $u$, den wir mit
$v$ bezeichnen wollen, in die beiden Functionen (II) ein, so besitzen
dieselben einen gemeinschaftlichen Factor, da $R_{f,\varphi}$ verschwindet, folg-
lich besitzt $f(x+u)$ selbst einen Factor, also auch $f(x)$. Dies gilt
auch, wenn zufälliger Weise $v=0$ ist. Nun ist aber in der That
%CONT>>


% ===== pdf idx 185, printed 171.0 =====
% !PRINTED 171
\weberpage{171}
%% [running head, not body text -- commented out 2026-08-07] Page furniture (running
%% head, folio, signature) is excluded from this edition by standing policy.
%% \S.\,12. Fundamentalsatz der Algebra. \hfill 171

$R_{f,\varphi} = P_{(u)} = 0$ eine leichter lösbare Gleichung als $f(x) = 0$. Denn
denken wir uns $f(x+u)$ nach $x$ und $u$ geordnet, so ist
\[ f(x+u) = x^n + p_1x^{n-1} + p_2x^{n-2} + \cdots p_n, \]
wo $p_1, p_2 \ldots p_n$ Functionen von $u$ vom bezw. $1^{\text{ten}}$, $2^{\text{ten}}, \ldots n^{\text{ten}}$ Grade
sind. Ebenso wird $\Phi$, in welcher sich wegen der Subtraction
$f(x+u) - f(x-u)$ einmal die höchste Potenz von $x$ weghebt, und
anderntheils wegen des Divisors $u$ auch der Grad in $u$ um 1 sich
erniedrigt, von der Form sein:
\[ \Phi = x^{n-1} + q_1x^{n-2} + q_2x^{n-3} + \cdots q_{n-1}, \]
wo $q_1, q_2 \ldots q_{n-1}$ Functionen von $u$ bezw. $1^{\text{ten}}$, $2^{\text{ten}}, \ldots (n-1)^{\text{ten}}$ Grades
sind. Die Resultante $R_{f,\varphi}$ in $u$ ist somit (vgl. Nr. 142) vom Grade
\[ n(n-1)\cdot \]
Nun enthalten aber anderntheils die beiden Functionen (I) die Grösse $u$
nur in geraden Potenzen, wovon man sich unmittelbar überzeugt, wenn
man die einzelnen Binome $(x+u)^\varrho$ entwickelt und die Addition bezw.
Subtraction ausführt. Es kann also auch die Resultante derselben,
die mit der Resultante der beiden Functionen (II) übereinstimmt, nur
eine Function von $u^2$ sein, und folglich lässt sich durch die Substitution
$u^2 = w$ ihr Grad erniedrigen auf $\dfrac{n(n-1)}{2}$. Da nun $n = 2^\mu \cdot \nu$ ist,
so ist
\[ \frac{n(n-1)}{2} = 2^{\mu-1} \cdot \nu \cdot (2^\mu \cdot \nu - 1) \]
oder, weil der Factor $2^\mu \cdot \nu - 1$ ungerade ist, solange $\mu > 0$, was immer
vorausgesetzt werden kann,
\[ \frac{n(n-1)}{2} = 2^{\mu-1} \cdot \bar{\nu}. \]
Der Grad der Resultante besitzt also ein kleineres $\mu$ als der Grad
von $f$; folglich hat die Resultante eine Wurzel, die wir mit $v = \sqrt{w}$
bezeichnen können, da $P_{(u)}$ in $u$ quadratisch ist. Wird dieser Werth
$u = v$ in $f = f(x+u)$ und $\varphi = \dfrac{f(x+u) - f(x-u)}{2u}$ eingetragen, so
verschwindet $R_{f,\varphi}$, d.\,h. $f(x+v)$ und $\varphi$ haben einen gemeinschaft-
lichen Theiler, und sonach auch $f(x)$ allein einen Factor. Ist $v$ reell, so
ist auch dieser Factor reell. Ist dagegen $v$ complex oder rein imaginär,
so wird der Factor von $f$ im Allgemeinen gleichfalls imaginär sein.

162. \textit{Der Ausnahmefall.} Es bleibt also noch der Fall zu be-
handeln, in welchem dieser Factor von $f(x)$ complex und vom Grade $\dfrac{n}{2}$
ist, nämlich
%CONT>>


% ===== pdf idx 186, printed 172.0 =====
% !PRINTED 172
\weberpage{172}
%% [running head, not body text -- commented out 2026-08-07] Page furniture (running
%% head, folio, signature) is excluded from this edition by standing policy.
%% 172 \hfill Zweiter Theil. Anwendungen der Determinantentheorie.

%<<CONT
\[ \varphi(x) = p(x) + iq(x) \]
und $p(x)$ und $q(x)$ keinen gemeinschaftlichen Theiler haben.

Nehmen wir nämlich an, $\varphi(x)$ sei vom Grade $\dfrac{n}{2}$; der andere
conjugirt complexe Factor ist alsdann
\[ \psi(x) = p(x) - iq(x) \quad \text{und} \]
\weq{1}{f(x) = \varphi(x) \cdot \psi(x) = p^2 + q^2}
sei die einzig mögliche Zerlegung von $f(x)$. Hat aber $f(x)$ den Factor
$\varphi(x)$, so besitzt $f(x+v)$ den Factor $\varphi(x+v)$ und $f(x+v)$ ist von
der Form:
\weq{2}{f(x+v) = \varphi(x+v) \cdot \psi(x+v).}
Ebenso ist aber auch
\weq{3}{f(x-v) = \varphi(x-v) \cdot \psi(x-v).}
Nun soll nach Voraussetzung $f(x+v)$ in keiner anderen Weise zer-
legbar sein, als in der unter (1) angedeuteten, und da überdies
\[ \Phi = \frac{f(x+v) - f(x-v)}{2v} \]
mit $f(x+v)$ einen gemeinschaftlichen Factor besitzt, so besitzt
diesen ausser
\[ f(x+v) = \varphi(x+v) \cdot \psi(x+v), \]
\[ \text{auch} \quad f(x-v) = \varphi(x-v) \cdot \psi(x-v). \]
Dies ist nur möglich, wenn entweder
\[ 1) \quad \varphi(x+v) = \varphi(x-v) \quad \text{oder} \]
\[ 2) \quad \psi(x+v) = \psi(x-v) \quad \text{oder} \]
\[ 3) \quad \varphi(x+v) = \psi(x-v) \quad \text{oder} \]
\[ 4) \quad \psi(x+v) = \varphi(x-v). \]
Man kann zunächst zeigen, dass die beiden ersten Annahmen auf
Widersprüche führen.

Denn im ersten Falle müsste
\[ \varphi(x-v) - \varphi(x+v) \]
identisch null sein. Das ist nur möglich, wenn die Coefficienten der
Potenzen von $x$ einzeln verschwinden. Dieser Ausdruck ist aber vom
Grade $\dfrac{n}{2} - 1$ in $x$ und das Glied mit der Potenz $x^{\frac{n}{2}-1}$ rührt einzig
und allein von der Differenz $(x-v)^{\frac{n}{2}} - (x+v)^{\frac{n}{2}}$ her, hat also den
%CONT>>


% ===== pdf idx 187, printed 173.0 =====
% !PRINTED 173
\weberpage{173}
%% [running head, not body text -- commented out 2026-08-07] Page furniture (running
%% head, folio, signature) is excluded from this edition by standing policy.
%% \S.\,12. Fundamentalsatz der Algebra. \hfill 173

%<<CONT
Coefficienten $-n \cdot v$. Da $v$ nicht null ist, so ist diese Identität unmög-
lich. Dieselbe Schlussweise lehrt, dass auch der zweite Fall nicht
eintreten kann.

Was die dritte Annahme betrifft
\[ \varphi(x+v) = \psi(x-v), \]
so setzen wir zunächst voraus, dass $v$ eine complexe Grösse sein
möge, dass also die Identität bestehe
\weq{I}{\varphi(x+\alpha+i\beta) = \psi(x-\alpha-i\beta)}
oder ausführlicher
\[ p(x+\alpha+i\beta) + iq(x+\alpha+i\beta) = p(x-\alpha-i\beta) - iq(x-\alpha-i\beta). \]
Vertauscht man hier $i$ mit $-i$, so erhält man
\[ p(x+\alpha-i\beta) - iq(x+\alpha-i\beta) = p(x-\alpha+i\beta) + iq(x-\alpha+i\beta) \]
oder
\weq{II}{\psi(x+\alpha-i\beta) = \varphi(x-\alpha+i\beta).}
Ersetzt man in (I) $x$ durch $x-2\alpha$, so erhält man:
\weq{III}{\varphi(x-\alpha+i\beta) = \psi(x-3\alpha-i\beta);}
also durch Comparation von (II) mit (III)
\[ \psi(x-i\beta+\alpha) = \psi(x-i\beta-3\alpha). \]
Da diese Identität für jeden Werth von $x$ bestehen soll, so müssen
die Coefficienten der Potenzen von $x$ einzeln verschwinden. Es muss
also insbesondere der Coefficient der höchsten in der Identität noch
auftretenden Potenz $x^{\frac{n}{2}-1} = 0$ sein. Derselbe rührt einzig und allein
von der Differenz her:
\[ \{x + (\alpha - i\beta)\}^{\frac{n}{2}} - \{x - (3\alpha + i\beta)\}^{\frac{n}{2}}, \]
ist also:
\[ \frac{n}{2}(\alpha - i\beta) + \frac{n}{2}(3\alpha + i\beta) = 2n\alpha. \]
Daraus folgt: $\alpha$ muss selbst verschwinden, d.\,h. $v$ kann keine com-
plexe Grösse sein. Ist aber $v$ rein imaginär, also
\[ v = i\beta, \]
so geht unsere Voraussetzung (3) über in
\[ \varphi(x+i\beta) = \psi(x-i\beta) \]


% ===== pdf idx 188, printed 174.0 =====
% !PRINTED 174
\weberpage{174}
%% [running head, not body text -- commented out 2026-08-07] Page furniture (running
%% head, folio, signature) is excluded from this edition by standing policy.
%% 174 \hfill Zweiter Theil. Anwendungen der Determinantentheorie.

oder: \quad $p(x+i\beta) + iq(x+i\beta) = p(x-i\beta) - iq(x-i\beta).$

Hat also die linke Seite die Form
\[ A + iB, \]
so muss die rechte Seite, da sie durch Vertauschung von $i$ mit $-i$
aus der linken hervorgeht, dargestellt sein durch
\[ A - iB. \]
Da aber beide Ausdrücke einander gleich sein sollen, so muss
\[ B = 0 \]
sein. Dann hat aber $\varphi(x+i\beta)$ nur reelle Coefficienten, also $f(x+v)$
einen reellen Factor $A(x)$ und somit auch einen linearen. Dieselben
Resultate ergeben sich auf gleichem Wege bei Discussion der vierten
Annahme.

Wir kommen zum Schlusse: Ist der imaginäre Factor $\varphi = p(x)$
$+ iq(x)$ von $f$ vom Grade $\dfrac{n}{2}$, so besitzt, wenn $p(x)$ und $q(x)$ keinen
gemeinschaftlichen Theiler haben, $f(x+v)$ einen reellen Factor, also
auch einen linearen und demnach auch $f(x)$ einen linearen.

Es ist somit der Beweis geliefert, dass $f(x) = 0$ in jedem Falle
einen linearen Factor haben muss. Wir erhalten denselben, indem
wir die in Nr. 161 geschilderten Operationen so lange anwenden, bis
wir schliesslich auf eine Resultante kommen, für welche $\mu = 0$, die
also von ungeradem Grade ist.

\textit{Anmerkung.} Aus der Thatsache, dass jede Gleichung $n^{\text{ten}}$ Grades
mindestens eine Wurzel hat, geht unmittelbar hervor, dass sie $n$
Wurzeln besitzt. Denn dividirt man die Gleichung mit dem einen
linearen Factor, so erhält man eine Function $(n-1)^{\text{ten}}$ Grades als
Quotient, welche gleich null gesetzt, ebenso wie die ursprüngliche
Gleichung, eine Wurzel besitzen muss.

\begin{center}
\S.\,13.\\
\gesperrt{Zusammenhang der Resultante mit den Wurzeln der
Gleichungen.}
\end{center}

163. \textit{Berechnung der gemeinsamen Wurzeln zweier Gleichungen.}
Sind
\[ f(x) = x^m + a_1x^{m-1} + a_2x^{m-2} + \cdots + a_m = 0 \]
\[ \varphi(x) = x^n + b_1x^{n-1} + b_2x^{n-2} + \cdots + b_n = 0 \]
zwei Gleichungen vom $m^{\text{ten}}$ und $n^{\text{ten}}$ Grade, die beide für den gemein-
schaftlichen Werth $x = \alpha$ befriedigt werden, so kann man sich den
%CONT>>


% ===== pdf idx 189, printed 175.0 =====
% !PRINTED 175
\weberpage{175}
% running head: \S.\,13. Zusammenhang der Resultante mit den Wurzeln der Gleichungen. 175
%<<CONT
Werth dieser gemeinschaftlichen Wurzel in folgender Weise verschaffen. Man bildet sich die $m+n-2$ Gleichungen:
\weq{I}{
\left\{
\begin{array}{l}
\alpha^0 f(\alpha) = 0, \quad \alpha f(\alpha) = 0, \cdots \alpha^{n-2} f(\alpha) = 0 \\
\alpha^0 \varphi(\alpha) = 0, \quad \alpha \varphi(\alpha) = 0, \cdots \alpha^{m-2} \varphi(\alpha) = 0.
\end{array}
\right.
}
Sie enthalten linear und homogen die $m+n-1$ Unbekannten
\[
\alpha^0, \quad \alpha^1, \quad \alpha^2 \ldots \alpha^{m+n-2}
\]
und man kann daher ihre Werthe nach Nr. 102 in eindeutiger Weise bestimmen.

Ist speciell $m=n$, so kann man statt des Gleichungssystemes (I) $n-1$ Gleichungen des Systemes
\weq{II}{
p_i = c_{i1}\alpha^{n-1} + c_{i2}\alpha^{n-2} + \cdots c_{in}\alpha^0 = 0
}
benutzen. Dieselben sind (vergl. Nr. 148) die Coefficienten von $y^{n-i}$ in der Entwicklung der Function:
\[
\frac{1}{y-x}\{f(x)\varphi(y) - \varphi(x)f(y)\} = \sum_{i=1}^{i=n} p_i y^{n-i}
\]
und verschwinden offenbar für $x=\alpha$, wenn gleichzeitig sowohl $f(\alpha)$ als $\varphi(\alpha)$ verschwinden.

Das System (II) ist wieder linear und homogen in den $n$ Unbekannten
\[
\alpha^0, \quad \alpha^1, \quad \alpha^2 \ldots \alpha^{n-1}.
\]
Wenn die beiden Functionen $f$ und $\varphi$ mehr als einen linearen Factor gemeinschaftlich haben, so ist es natürlich nicht möglich, dieselben auf diesem Wege zu bestimmen, wie schon daraus erhellt, dass sowohl Gleichungssystem (I) als (II) Systeme von homogenen und linearen Gleichungen sind, die als solche nur eine einzige Lösung zulassen.

164. \textit{Die Differentialquotienten der Resultante.} Wir hatten in Nr. 131 als Ausdruck für die Resultante $R_{f,\varphi}$ erhalten:
\[
R_{f,\varphi} = \overset{n-1}{\overline{A}}(x)\,\overset{m}{f}(x) + \overset{m-1}{\overline{B}}(x)\,\overset{n}{\varphi}(x).
\]
Differenziren wir diese Gleichung nach einem Coefficienten von $f$ oder $\varphi$, also etwa nach $a_\nu$, so erhalten wir:
\[
\frac{\partial R_{f,\varphi}}{\partial a_\nu} = \frac{\partial \overline{A}}{\partial a_\nu}\cdot f + \overline{A}\cdot x^{n-\nu} + \frac{\partial \overline{B}}{\partial a_\nu}\cdot \varphi.
\]
Ist nun $\alpha$ eine gemeinschaftliche Wurzel der beiden Gleichungen $f$ und $\varphi$, so geht daraus, wegen $f(\alpha)=0$, $\varphi(\alpha)=0$, die Relation hervor,
\weq{1}{
\frac{\partial R_{f,\varphi}}{\partial a_\nu} = \overset{n-1}{\overline{A}}(\alpha)\cdot \alpha^{n-\nu}.
}


% ===== pdf idx 190, printed 176.0 =====
% !PRINTED 176
\weberpage{176}
% running head: 176  Zweiter Theil. Anwendungen der Determinantentheorie.
Da man ebenso erhält:
\weq{2}{
\frac{\partial R_{f,\varphi}}{\partial a_\mu} = \overset{n-1}{\overline{A}}(\alpha)\cdot \alpha^{n-\mu},
}
so ergiebt sich aus diesen beiden Relationen die Proportion
\[
\frac{\partial R_{f,\varphi}}{\partial a_\nu} : \frac{\partial R_{f,\varphi}}{\partial a_\mu} = \alpha^{n-\nu} : \alpha^{n-\mu}.
\]
Man erkennt also, dass die partiellen Differentialquotienten der Resultante nach den Coefficienten von $f$ oder $\varphi$ der Reihe nach proportional sind den verschiedenen Potenzen der gemeinschaftlichen Wurzel $\alpha$. Man kann aus dieser Eigenschaft unter anderm auch eine Reihe von partiellen Differentialgleichungen ableiten, welche charakteristische Gleichungen sind für die Resultante.

165. \textit{Zusammenhang zwischen den Coefficienten und Wurzeln einer Gleichung.} Die Thatsache, dass jede Gleichung $n^{\mathrm{ten}}$ Grades $n$ Wurzeln $\alpha_1 \alpha_2 \ldots \alpha_n$ besitzt, sich also in ein Product von $n$ linearen Factoren zerfällen lässt, lehrt, dass die Coefficienten der Gleichung in einem innern Zusammenhange mit den Wurzeln derselben stehen. Man findet diese Beziehungen in der That direct, wenn man in den beiden Darstellungen der Gleichung $f$:
\[
f = (x-\alpha_1)(x-\alpha_2)\cdots(x-\varrho_n) = 0 % [sic] printed \varrho_n, not \alpha_n
\]
\[
f = x^n + a_1 x^{n-1} + a_2 x^{n-2} \cdots + a_{n-1}x + a_n = 0
\]
die Coefficienten gleich hoher Potenzen von $x$ vergleicht.

Man erhält so in bekannter Weise die Grössen $a_i$ ausgedrückt in symmetrischen Functionen der Wurzeln
\[
\alpha_1 \quad \alpha_2 \quad \alpha_3 \ldots \alpha_n.
\]
Um diese symmetrischen Functionen zu ermitteln, kann man auch in folgender Weise vorgehen.

Man bildet die $n$ Gleichungen:
\[
f(\alpha_1) = 0, \quad f(\alpha_2) = 0, \quad f(\alpha_3) = 0 \ldots f(\alpha_n) = 0,
\]
die linear sind in den $n$ Coefficienten $a_1, a_2 \ldots a_n$, und also zu deren Berechnung gerade ausreichen, wenn alle Wurzeln verschieden sind.

Die Matrix dieses Systems von Gleichungen für $a_0, a_1 \ldots a_n$ ist dargestellt durch:
\weq{A}{
\begin{vmatrix}
\alpha_1^{n}, & \alpha_1^{n-1}, & \alpha_1^{n-2} & \cdots & 1 \\
\alpha_2^{n}, & \alpha_2^{n-1}, & \alpha_2^{n-2} & \cdots & 1 \\
\alpha_3^{n}, & \cdot & \cdot & \cdots & 1 \\
\cdot & \cdot & \cdot & \cdots & \cdot \\
\alpha_n^{n}, & \alpha_n^{n-1}, & \alpha_n^{n-2} & \cdots & 1
\end{vmatrix}.
}


% ===== pdf idx 191, printed 177.0 =====
% !PRINTED 177
\weberpage{177}
% running head: \S.\,13. Zusammenhang der Resultante mit den Wurzeln der Gleichungen. 177
Bezeichnen wir wie früher (vergl. Nr. 112) die mit wechselndem Vorzeichen genommenen $n+1$ Determinanten derselben, welche durch Unterdrückung der bezw. $1^{\mathrm{ten}}$, $2^{\mathrm{ten}} \ldots (n+1)^{\mathrm{ten}}$ Colonne entstehen, mit
\[
p_0, \quad p_1, \quad p_2 \cdots p_n % ?? uncertain: a small mark follows p_n in the scan (ink blot or stray accent); transcribed as plain p_n
\]
so erhalten wir:
\[
a_n = \frac{p_n}{p_0}, \quad a_{n-1} = \frac{p_{n-1}}{p_0}; \quad a_{n-2} = \frac{p_{n-2}}{p_0}, \ldots a_1 = \frac{p_1}{p_0}.
\]
Nun ist die Determinante $p_0$ dargestellt durch:
\[
p_0 =
\begin{vmatrix}
\alpha_1^{n-1}, & \alpha_1^{n-2} & \cdots & 1 \\
\alpha_2^{n-1}, & \alpha_2^{n-2} & \cdots & 1 \\
\alpha_3^{n-1}, & \cdot & \cdots & 1 \\
\cdot & \cdot & \cdots & \cdot \\
\alpha_n^{n-1}, & \alpha_n^{n-2} & \cdots & 1
\end{vmatrix}.
\]
Diese Determinante ist aber (vergl. Beispiel 2 und 3 in Nr. 40) nichts anderes als das Differenzenproduct
\[
\Delta_\alpha = (\alpha_1-\alpha_2)(\alpha_1-\alpha_3)\cdots(\alpha_1-\alpha_n)
(\alpha_2-\alpha_3)\cdots(\alpha_2-\alpha_n)
(\alpha_{n-1}-\alpha_n),
\]
und man erkennt ebenso aus dem angeführten Beispiele, dass alle Determinanten der Matrix (A) durch dieses Differenzenproduct theilbar sind, so dass wir also auf diese Weise die Coefficienten $a_i$ in der That ausgedrückt erhalten durch ganze symmetrische Functionen der Wurzeln.

166. \textit{Beispiel.} Es sei
\[
f(x) = x^4 + a_1 x^3 + a_2 x^2 + a_3 x + a_4 = 0
\]
und
\[
\alpha_1, \quad \alpha_2, \quad \alpha_3, \quad \alpha_4
\]
seien die Wurzeln dieser Gleichung. Es ist dann:
\[
\begin{aligned}
f(\alpha_1) &= \alpha_1^4 + a_1\alpha_1^3 + a_2\alpha_1^2 + a_3\alpha_1 + a_4 = 0 \\
f(\alpha_2) &= \alpha_2^4 + a_1\alpha_2^3 + a_2\alpha_2^2 + a_3\alpha_2 + a_4 = 0 \\
f(\alpha_3) &= \alpha_3^4 + a_1\alpha_3^3 + a_2\alpha_3^2 + a_3\alpha_3 + a_4 = 0 \\
f(\alpha_4) &= \alpha_4^4 + a_1\alpha_4^3 + a_2\alpha_4^2 + a_3\alpha_4 + a_4 = 0.
\end{aligned}
\]
Die Matrix dieses Systems ist:
\weq{A}{
\begin{vmatrix}
\alpha_1^{4} & \alpha_1^{3} & \alpha_1^{2} & \alpha_1 & 1 \\
\alpha_2^{4} & \alpha_2^{3} & \alpha_2^{2} & \alpha_2 & 1 \\
\alpha_3^{4} & \alpha_3^{3} & \alpha_3^{2} & \alpha_3 & 1 \\
\alpha_4^{4} & \alpha_4^{3} & \alpha_4^{2} & \alpha_4 & 1
\end{vmatrix}
}
% foot: Gordan, Determinanten.   12  [author/short-title signature line at foot of page, not body text]


% ===== pdf idx 192, printed 178.0 =====
% !PRINTED 178
\weberpage{178}
% running head: 178  Zweiter Theil. Anwendungen der Determinantentheorie.
Die Determinante
\[
p_0 =
\begin{vmatrix}
\alpha_1^{3} & \alpha_1^{2} & \alpha_1 & 1 \\
\alpha_2^{3} & \alpha_2^{2} & \alpha_2 & 1 \\
\alpha_3^{3} & \alpha_3^{2} & \alpha_3 & 1 \\
\alpha_4^{3} & \alpha_4^{2} & \alpha_4 & 1
\end{vmatrix}
= \Delta_\alpha
\]
hat den Werth:
\[
p_0 = (\alpha_1-\alpha_2)(\alpha_1-\alpha_3)(\alpha_1-\alpha_4)(\alpha_2-\alpha_3)(\alpha_2-\alpha_4)(\alpha_3-\alpha_4) = \Delta_\alpha.
\]
Ebenso findet man:
\[
\begin{aligned}
p_4 &= +\alpha_1\alpha_2\alpha_3\alpha_4 \cdot \Delta_\alpha \\
p_3 &= -(\alpha_1\alpha_2\alpha_3 + \alpha_1\alpha_2\alpha_4 + \alpha_1\alpha_3\alpha_4 + \alpha_2\alpha_3\alpha_4) \cdot \Delta_\alpha \\
p_2 &= +(\alpha_1\alpha_2 + \alpha_1\alpha_3 + \alpha_1\alpha_4 + \alpha_2\alpha_3 + \alpha_2\alpha_4 + \alpha_3\alpha_4) \cdot \Delta_\alpha \\
p_1 &= -(\alpha_1 + \alpha_2 + \alpha_3 + \alpha_4) \cdot \Delta_\alpha.
\end{aligned}
\]
Folglich erhalten wir als Relationen zwischen den Coefficienten und den Wurzeln von $f$:
\[
\left.
\begin{aligned}
a_1 &= -\sum_{1}^{4} \alpha_i \\
a_2 &= +\sum_{1}^{4} \alpha_i \alpha_k \\
a_3 &= -\sum_{1}^{4} \alpha_i \alpha_k \alpha_\varrho \\
a_4 &= +\alpha_1 \alpha_2 \alpha_3 \alpha_4
\end{aligned}
\right\}
\quad i < k < \varrho.
\]
167. \textit{Die Potenzsummen der Wurzeln.} Von den weiteren verschiedenartigen Beziehungen, welche zwischen den Coefficienten $a_i$ und Wurzeln $\alpha_i$ einer Gleichung bestehen, mögen hier nur noch die Relationen zwischen den Coefficienten und den Summen
\[
s_\varrho = \alpha_1^\varrho + \alpha_2^\varrho + \alpha_3^\varrho + \cdots \alpha_n^\varrho
\]
der Potenzen aller Wurzeln einer Gleichung erwähnt sein. Bildet man nämlich die erste Ableitung von $f(x)$ einmal in der Form:
\[
f'(x) = nx^{n-1} + (n-1)a_1 x^{n-2} + (n-2)a_2 x^{n-3} + \cdots a_{n-1},
\]
dann aber auch in der Form
\[
f'(x) = \frac{f(x)}{x-\alpha_1} + \frac{f(x)}{x-\alpha_2} + \frac{f(x)}{x-\alpha_3} + \cdots \frac{f(x)}{x-\alpha_n}
\]
und vergleicht in den beiden Ausdrücken die Coefficienten gleich hoher Potenzen von $x$, so erhält man die unter dem Namen \glqq Newton'sche Identitäten\grqq{} bekannten Formeln:
%CONT>>


% ===== pdf idx 193, printed 179.0 =====
% !PRINTED 179
\weberpage{179}
\[
(\mathrm{I}) \quad
\left\{
\begin{aligned}
s_1 &+ 1 \cdot a_1 = 0 \\
s_2 &+ a_1 s_1 + 2 \cdot a_2 = 0 \\
s_3 &+ a_1 s_2 + a_2 s_3 + 3 \cdot a_3 = 0 \\
& \cdots\cdots\cdots\cdots\cdots\cdots \\
s_{n-1} &+ a_1 s_{n-2} + a_2 s_{n-3} + \cdots (n-1) a_{n-1} = 0.
\end{aligned}
\right.
\]
Hiebei ist unter $s_1, s_2, s_3 \ldots s_{n-1}$ die Summe aller ersten, zweiten,
dritten, $\ldots$ $(n-1)^{\mathrm{ten}}$ Potenzen aller Wurzeln zu verstehen.

Zu diesen Formeln (I) erhält man noch eine zweite Reihe Formeln für die
höheren Potenzsummen auf folgende Weise. Man multiplicirt $f(x) = 0$ mit
$x^m$, so kommt:
\[
x^{n+m} + a_1 x^{n+m-1} + a_2 x^{n+m-2} + \cdots a_n x^m = 0.
\]
Da diese Gleichung immer noch für jede der $n$ Wurzeln $\alpha_i$
verschwindet, so erhält man hieraus:
\[
\alpha_i^{n+m} + a_1 \alpha_i^{n+m-1} + a_2 \alpha_i^{n+m-2} + \cdots a_n \alpha_i^m = 0.
\]
Setzt man $i = 1, 2, 3 \ldots n$ und addirt alle $n$ Relationen, so kommt:
\[
s_{n+m} + a_1 s_{n+m-1} + a_2 s_{n+m-2} + \cdots a_n s_m = 0.
\]
Für $m = 0, 1, 2 \ldots$ entstehen daraus, da $s_0 = n$, die Relationen:
\[
(\mathrm{II}) \quad
\left\{
\begin{aligned}
s_n &+ a_1 s_{n-1} + a_2 s_{n-2} + \cdots n a_n = 0 \\
s_{n+1} &+ a_1 s_n + a_2 s_{n-1} + \cdots a_n s_1 = 0 \\
s_{n+2} &+ a_1 s_{n+1} + a_2 s_n + \cdots a_n s_2 = 0 \\
& \cdots\cdots\cdots\cdots\cdots\cdots \\
s_{n+m} &+ a_1 s_{n+m-1} \cdot\cdot\cdot + a_n s_m = 0.
\end{aligned}
\right.
\]
Von diesen Formeln (I) und (II) werden wir alsbald Gebrauch machen
Sie sind linear in $s$ und in $a$ und man kann daher nach den in \S.\,8
gegebenen Methoden in einfacher Weise sowohl die Werthe von $s_\varrho$ in
Function der Coëfficienten $a_i$ als auch die Werthe von $a_i$ in Function
der Potenzsummen $s_\varrho$ darstellen.

168. \textit{Die Resultante als Product der Wurzeldifferenzen.} Es seien
\[
\overset{m}{f}(x) = 0 \quad \text{und} \quad \overset{n}{\varphi}(x) = 0
\]
zwei Gleichungen $m^{\mathrm{ten}}$, bezw. $n^{\mathrm{ten}}$ Grades, deren
erste die Wurzeln
\[
\alpha_1, \quad \alpha_2, \quad \alpha_3 \ldots \alpha_m
\]
und deren zweite die Wurzeln
\[
\beta_1, \quad \beta_2, \quad \beta_3 \ldots \beta_n
\]
besitze. Wir können uns alsdann die Frage vorlegen, unter welcher
%CONT>>


% ===== pdf idx 194, printed 180.0 =====
% !PRINTED 180
\weberpage{180}
%<<CONT
Bedingung die beiden Gleichungen eine gemeinschaftliche Wurzel besitzen.
Nun ist nach dem Gauss'schen Satze die Gleichung $f(x) = 0$ darstellbar
in der Form
\[
(1) \qquad f(x) = (x - \alpha_1)(x - \alpha_2)(x - \alpha_3) \ldots (x - \alpha_m) = 0
\]
und ebenso
\[
(2) \qquad \varphi(x) = (x - \beta_1)(x - \beta_2)(x - \beta_3) \ldots (x - \beta_n) = 0.
\]
Substituiren wir in die erste Gleichung der Reihe nach die Wurzeln
$\beta_\varrho$ der zweiten, so erhalten wir die Substitutionsresultate
\[
f(\beta_\varrho) = (\beta_\varrho - \alpha_1)(\beta_\varrho - \alpha_2)(\beta_\varrho - \alpha_3) \ldots (\beta_\varrho - \alpha_m),
\]
\[
\varrho = 1, 2, 3 \ldots n.
\]
Ist irgend eine Wurzel $\beta_\varrho$ gleich einer Wurzel $\alpha_k$, so
verschwindet $f(\beta_\varrho)$, und umgekehrt ist die nothwendige und
hinreichende Bedingung, dass eine Wurzel $\beta_\varrho$ gleich einer
Wurzel $\alpha_k$ sei, das Verschwinden des Productes:
\[
(3) \qquad f(\beta_1) \cdot f(\beta_2) \cdot f(\beta_3) \ldots f(\beta_n).
\]
Ganz ebenso drückt
\[
(4) \qquad \varphi(\alpha_1) \cdot \varphi(\alpha_2) \cdot \varphi(\alpha_3) \ldots \varphi(\alpha_m) = 0
\]
aus, dass $f$ und $\varphi$ eine gemeinschaftliche Wurzel besitzen.
Bezeichnen wir das letzte Poduct (4) mit % [sic] Poduct
\[
\begin{aligned}
U = {}& (\alpha_1 - \beta_1)(\alpha_1 - \beta_2)(\alpha_1 - \beta_3) \ldots (\alpha_1 - \beta_n) \times \\
& (\alpha_2 - \beta_1)(\alpha_2 - \beta_2) \cdots\cdots (\alpha_2 - \beta_n) \times \\
& \cdots\cdots\cdots\cdots\cdots\cdots \\
& \times (\alpha_m - \beta_1)(\alpha_m - \beta_2)(\alpha_m - \beta_3) \ldots (\alpha_m - \beta_n),
\end{aligned}
\]
welches sonach $m \cdot n$ lineare Factoren besitzt, so ist das
vorhergehende Product (3) nichts anderes als $(-1)^{m \cdot n} U$
(vergleiche auch Nr. 139). Nun hatten wir aber bereits früher (vergl.
Nr. 131) eine in den Coëfficienten von $f$ und $\varphi$ rationale und
ganze Function $R_{f,\varphi}$ aufgestellt, deren Verschwinden gleichfalls
das Vorhandensein eines gemeinschaftlichen linearen Factors ankündigte.
Unsere Aufgabe ist daher, den Zusammenhang zwischen $U$ und $R_{f,\varphi}$
zu ermitteln. Es wird sich zeigen, dass direct die Beziehung besteht:
\[
U = R_{f,\varphi}.
\]
Wir geben im Folgenden drei Beweise für diese Behauptung.

169. \textit{Erster Beweis.} Derselbe stützt sich auf den in Nr. 154
aufgestellten Lehrsatz:
\[
R_{f,\varphi \cdot \psi} = R_{f,\varphi} \cdot R_{f,\psi}.
\]


% ===== pdf idx 195, printed 181.0 =====
% !PRINTED 181
\weberpage{181}
Es ist nämlich, wegen $\varphi(x) = (x - \beta_1)(x - \beta_2) \ldots (x - \beta_n),$
\[
R_{f,\varphi} = R_{f,(x-\beta_1)(x-\kappa_2)\ldots(x-\beta_n)} = R_{f,x-\beta_1} \times R_{f,x-\beta_2} \times R_{f,x-\beta_3} \ldots \times R_{f,x-\beta_n}. % [sic] printed "κ_2" for β_2
\]
Nun ist aber allgemein:
\[
R_{f,x-\beta} =
\begin{vmatrix}
1 & a_1 & a_2 & a_3 & \cdots & a_m \\
1 & -\beta & 0 & 0 & \cdots & 0 \\
0 & 1 & -\beta & 0 & \cdots & 0 \\
\cdot & \cdot & \cdot & \cdot & \cdot & \cdot \\
0 & 0 & \cdots & 1 & -\beta & 0 \\
0 & 0 & \cdots & 0 & 1 & -\beta
\end{vmatrix}.
\]
Denken wir uns die erste Colonne mit $\beta^m$, die zweite mit
$\beta^{m-1}$, die dritte mit $\beta^{m-2}$ u.\,s.\,w., die vorletzte mit
$\beta$ multiplicirt und zur letzten addirt, so erhalten wir
\[
R_{f,x-\beta} =
\begin{vmatrix}
1 & a_1 & a_2 & a_3 & \cdots & f(\beta) \\
1 & -\beta & 0 & 0 & \cdots & 0 \\
0 & 1 & -\beta & 0 & \cdots & 0 \\
\cdot & \cdot & \cdot & \cdot & \cdot & \cdot \\
0 & 0 & \cdots & 1 & -\beta & 0 \\
0 & 0 & \cdots & 0 & 1 & 0
\end{vmatrix}
= (-1)^m \cdot f(\beta).
\]
Folglich ist
\[
\begin{aligned}
R_{f,\varphi} &= (-1)^{m \cdot n} \cdot f(\beta_1) \cdot f(\beta_2) \cdot f(\beta_3) \ldots f(\beta_n) \\
&= \varphi(\alpha_1) . \varphi(\alpha_2) \ldots \varphi(\alpha_m) = U.
\end{aligned}
\]

170. \textit{Zweiter Beweis.} Bildet man zunächst die Resultante von
$f(x)$ und $\varphi(x) - y$, so erhält man:
\[
R_{f(x),\varphi(x)-y} =
\begin{vmatrix}
1 & a_1 & a_2 & a_3 & \cdots & a_m & 0 & 0 & \cdots & 0 \\
0 & 1 & a_1 & a_2 & \cdots & \cdots & a_m & 0 & \cdots & 0 \\
0 & 0 & 1 & a_1 & \cdots & \cdots & \cdots & a_m & \cdots & 0 \\
\cdot & \cdot & \cdot & \cdot & \cdot & \cdot & \cdot & \cdot & \cdot & \cdot \\
0 & 0 & \cdots & \cdots & \cdots & \cdots & \cdots & \cdots & \cdots & a_m \\
1 & b_1 & b_2 & \cdots & b_n - y & 0 & \cdots & \cdots & \cdots & 0 \\
0 & 1 & b_1 & \cdots & \cdots & \cdots & b_n - y & \cdots & \cdots & 0 \\
\cdot & \cdot & \cdot & \cdot & \cdot & \cdot & \cdot & \cdot & \cdot & \cdot \\
0 & 0 & \cdots & \cdots & \cdots & \cdots & \cdots & \cdots & b_{n-1} & b_n - y
\end{vmatrix}
= F(y).
\]
Diese Resultante $F(y)$ ist in $y$ vom Grade $m$ und der Coëfficient der
höchsten Potenz ist $(-1)^m$. Das constante Glied dieser Function $F(y)$
%CONT>>


% ===== pdf idx 196, printed 182.0 =====
% !PRINTED 182
\weberpage{182}
%<<CONT
erhält man aus dieser Determinante für $y = 0$. Es ist also gerade jene
Function, die wir mit $R_{f,\varphi}$ früher bezeichnet hatten.
Anderntheils ist aber nach den Beziehungen, die zwischen den
Coëfficienten und Wurzeln einer Gleichung bestehen (vergl. Nr. 165 und
166), dieses constante Glied in $F(y)$ bis aufs Vorzeichen $(-1)^m$
gleich dem Producte aller Wurzeln von $F(y)$. Diese Wurzeln sind aber
nichts anderes als die Grössen
\[
\varphi(\alpha_1), \quad \varphi(\alpha_2) \ldots \varphi(\alpha_m);
\]
denn die beiden Functionen
\[
f(x) = (x - \alpha_1)(x - \alpha_2) \ldots (x - \alpha_m) = 0
\]
und
\[
\varphi(x) - \varphi(\alpha_i) = 0
\]
haben den gemeinschaftlichen Factor $x - \alpha_i$, und folglich wird
\[
R_{f(x),\varphi(x)-y}
\]
für $y = \varphi(\alpha_i)$ identisch Null. Daher ist
\[
\frac{p_m}{(-1)^m} = \frac{R_{f,\varphi}}{(-1)^m} = (-1)^m \cdot \varphi(\alpha_1) \varphi(\alpha_2) \ldots \varphi(\alpha_m)
\]
oder
\[
R_{f,\varphi} = \varphi(\alpha_1) \varphi(\alpha_2) \varphi(\alpha_3) \ldots \varphi(\alpha_m) = U.
\]

171. \textit{Dritter Beweis.} Wir bilden uns einmal aus den Potenzen der
Wurzeln die Matrices:
\[
(\mathrm{I}) \quad
\begin{vmatrix}
\alpha_1^{m+n-1}, & \alpha_1^{m+n-2}, & \cdots & \alpha_1^2 & \alpha_1 & 1 \\
\alpha_2^{m+n-1}, & \alpha_2^{m+n-2}, & \cdots & \alpha_2^2 & \alpha_2 & 1 \\
\alpha_3^{m+n-1}, & \cdots & \cdots & \cdots & \cdots & 1 \\
\cdot & \cdot & \cdot & \cdot & \cdot & \cdot \\
\alpha_m^{m+n-1}, & \alpha_m^{m+n-2}, & \cdots & \alpha_m^2 & \alpha_m & 1
\end{vmatrix}.
\]
\[
(\mathrm{II}) \quad
\begin{vmatrix}
\beta_1^{m+n-1}, & \beta_1^{m+n-2}, & \beta_1^{m+n-3} & \cdots & \beta_1^2 & \beta_1 & 1 \\
\beta_2^{m+n-1}, & \beta_2^{m+n-2}, & \beta_2^{m+n-3} & \cdots & \beta_2^2 & \beta_1 & 1 \\
\cdot & \cdot & \cdot & \cdot & \cdot & \cdot & \cdot \\
\beta_n^{m+n-1}, & \beta_n^{m+n-2}, & \beta_n^{m+n-3} & \cdots & \beta_n^2 & \beta_n & 1
\end{vmatrix}
\]
% [sic] second row of (II) prints "$\beta_1$" for the final entry instead of "$\beta_2$"
und ebenso aus den Coëfficienten von $f$ und $\varphi$ die Matrices
\[
(\mathrm{III}) \quad
\begin{vmatrix}
1 & a_1 & a_2 & a_3 & \cdots & a_m & 0 & 0 & \cdots & 0 \\
0 & 1 & a_1 & a_2 & \cdots & \cdots & a_m & 0 & \cdots & 0 \\
0 & 0 & 1 & a_1 & \cdots & \cdots & \cdots & a_m & \cdots & 0 \\
\cdot & \cdot & \cdot & \cdot & \cdot & \cdot & \cdot & \cdot & \cdot & \cdot \\
0 & 0 & 0 & \cdots & \cdots & \cdots & \cdots & \cdots & a_{m-1} & a_m
\end{vmatrix}
(n \text{ Zeilen}),
\]


% ===== pdf idx 197, printed 183.0 =====
% !PRINTED 183
\weberpage{183}
\S.\,13. Zusammenhang der Resultante mit den Wurzeln der Gleichungen.

\[
\text{(IV)}\quad
\left|\begin{array}{l}
1\ \ b_1\ \ b_2\ \ b_3\ \ \ldots\ \ b_n\ \ 0\ \ 0\ \ \ldots\ \ 0\\
0\ \ 1\ \ b_1\ \ b_2\ \ \ldots\ldots\ \ b_n\ \ 0\ \ \ldots\ \ 0\\
0\ \ 0\ \ 1\ \ b_1\ \ \ldots\ldots\ldots\ \ b_n\ \ \ldots\ \ 0\\
\ldots\ldots\ldots\ldots\ldots\ldots\ldots\ldots\ldots\ldots\ldots\ldots\ldots\\
0\ \ 0\ \ 0\ \ \ldots\ldots\ldots\ldots\ldots\ldots\ldots\ b_{n-1}\ \ b_n
\end{array}\right|
\quad (m\text{ Zeilen}).
\]
Die Matrices (I) und (III) sowohl als auch die Matrices (II) und (IV)
sind hiebei correspondirende Matrices; denn wenn wir (I) mit (III)
und (II) mit (IV) durch Multiplication von Zeile mit Zeile zu je einer
Matrix verbinden, so entstehen in beiden Fällen Matrices, deren Elemente durchwegs verschwinden.

Sind daher die Grössen $\rho$ und $\sigma$ Proportionalitätsfactoren, so bestehen (vergl. Nr. 92), wenn wir die Determinanten der Matrices
(I) und (II) mit $\Delta(\alpha)$, $\Delta(\beta)$, die der correspondirenden Matrices (III)
und (IV) mit $\Delta(a)$, $\Delta(b)$ bezeichnen, die Relationen
\[
\Delta(\alpha) = (-1)^\nu\cdot\rho\cdot\Delta(a),
\]
\[
\Delta(\beta) = (-1)^\mu\cdot\sigma\cdot\Delta(b).
\]
Dabei sind $\nu$ und $\mu$ die Summen der Doppelindices des Diagonalgliedes in $\Delta(\alpha)$ bezw. $\Delta(\beta)$.

172. Nehmen wir als Determinante $\Delta(a)$ die Determinante $\Delta_1(a)$
der $n$ ersten Colonnen in Matrix (III) heraus, so ist die correspondirende die Determinante
\[
(-1)^\nu\Delta_1(\alpha) =
\left|\begin{array}{l}
\alpha_1^{m-1},\ \ \alpha_1^{m-2},\ \ \ldots\ \ \alpha_1^{2}\ \ \alpha_1\ \ 1\\
\alpha_2^{m-1},\ \ \alpha_2^{m-2},\ \ \ldots\ \ \alpha_2^{2}\ \ \alpha_2\ \ 1\\
\ldots\ldots\ldots\ldots\ldots\ldots\ldots\ldots\ldots\ldots\ldots\ldots\\
\alpha_m^{m-1},\ \ \alpha_m^{m-2},\ \ \ldots\ \ \alpha_m^{2}\ \ \alpha_m\ \ 1
\end{array}\right|
= \Delta_\alpha.
\]
Die Determinante $\Delta_1(a)$ hat aber den Werth 1; folglich erhalten wir
für den Proportionalitätsfactor $\rho$ den Werth
\[
\Delta_1(\alpha) = \Delta_\alpha = (-1)^\nu\cdot\rho\cdot\Delta_1(a) = \rho,
\]
weil
\[
\nu = 2\sum_{k=n}^{n+m} k,
\]
also jedenfalls eine gerade Zahl ist. Ebenso ergiebt sich
\[
\sigma = \Delta_\beta.
\]
Setzt man nun einestheils die Matrices (I) und (II) zur Matrix einer
Determinante zusammen:
%CONT>>


% ===== pdf idx 198, printed 184.0 =====
% !PRINTED 184
\weberpage{184}
%<<CONT
184\quad Zweiter Theil. Anwendungen der Determinantentheorie.

\[
\Delta =
\left|\begin{array}{l}
\alpha_1^{m+n-1},\ \ \alpha_1^{m+n-2},\ \ \alpha_1^{m+n-3},\ \ \ldots\ \ \alpha_1^{2}\ \ \alpha_1\ \ 1\\
\alpha_2^{m+n-1},\ \ \alpha_2^{m+n-2},\ \ \alpha_2^{m+n-3},\ \ \ldots\ \ \alpha_2^{2}\ \ \alpha_2\ \ 1\\
\ldots\ldots\ldots\ldots\ldots\ldots\ldots\ldots\ldots\ldots\ldots\ldots\\
\alpha_m^{m+n-1},\ \ \alpha_m^{m+n-2},\ \ \alpha_m^{m+n-3},\ \ \ldots\ \ \alpha_m^{2}\ \ \alpha_m\ \ 1\\
\beta_1^{m+n-1},\ \ \beta_1^{m+n-2},\ \ \beta_1^{m+n-3},\ \ \ldots\ \ \beta_1^{2}\ \ \beta_1\ \ 1\\
\beta_2^{m+n-1},\ \ \beta_2^{m+n-2},\ \ \ldots\ldots\ldots\ldots\ldots\ldots\\
\ldots\ldots\ldots\ldots\ldots\ldots\ldots\ldots\ldots\ldots\ldots\ldots\\
\beta_n^{m+n-1},\ \ \beta_n^{m+n-2},\ \ \beta_n^{m+n-3},\ \ \ldots\ \ \beta_n^{2}\ \ \beta_n\ \ 1
\end{array}\right|
\]
und ebenso die Matrices (III) und (IV) zur Matrix der Determinante
\[
D =
\left|\begin{array}{l}
1\ \ a_1\ \ a_2\ \ a_3\ \ \ldots\ \ a_m\ \ 0\ \ \ldots\ \ 0\\
0\ \ 1\ \ a_1\ \ a_2\ \ \ldots\ldots\ldots\ \ a_m\ \ \ldots\ \ 0\\
\ldots\ldots\ldots\ldots\ldots\ldots\ldots\ldots\ldots\ldots\\
0\ \ 0\ \ \ldots\ldots\ldots\ldots\ldots\ldots\ \ a_m\\
1\ \ b_1\ \ b_2\ \ b_3\ \ \ldots\ \ b_n\ \ 0\ \ \ldots\\
0\ \ 1\ \ b_1\ \ \ldots\ldots\ \ b_{n-1}\ \ b_n\ \ \ldots\ \ 0\\
\ldots\ldots\ldots\ldots\ldots\ldots\ldots\ldots\ldots\ldots\\
0\ \ 0\ \ \ldots\ldots\ldots\ldots\ldots\ldots\ \ b_n
\end{array}\right|\ ,
\]
so besitzt die erste Determinante den Werth
\[
\Delta = \Delta_\alpha\cdot\Delta_\beta\cdot U,
\]
denn sie verschwindet 1) für $\alpha_i=\alpha_k$, 2) für $\beta_i=\beta_k$, 3) für $\alpha_i=\beta_k$.
Die zweite Determinante erkennt man aber unmittelbar als die Resultante $R_{f,\varphi}$. Entwickeln wir diese Resultante $R_{f,\varphi}$ nach den
Determinanten der ersten $n$ Zeilen gemäss dem \gesperrt{Laplace'schen}
Satze, so ist ihr Werth dargestellt in der Form
\[
R_{f,\varphi} = \sum{}'(-1)^{\bar\nu}\cdot\Delta(a)\cdot\Delta(b)\,,
\]
wo $\nu$ stets als Summe aller Doppelindices des Diagonalgliedes in
$\Delta(a)$ genommen werden kann.

Ersetzen wir hierin $\Delta(a)$ durch
\[
\dfrac{\Delta(\alpha)}{(-1)^\nu\cdot\Delta_\alpha}
\]
und $\Delta(b)$ durch
\[
\dfrac{\Delta(\beta)}{(-1)^\mu\cdot\Delta_\beta}\,,
\]
%CONT>>


% ===== pdf idx 199, printed 185.0 =====
% !PRINTED 185
\weberpage{185}
\S.\,14. Die Discriminante. 185

%<<CONT
so erhalten wir
\[
R_{f,\varphi} = \sum{}'\dfrac{(-1)^{\bar\nu}\cdot\Delta(\alpha)\cdot\Delta(\beta)}{(-1)^{\mu+\nu}\cdot\Delta_\alpha\cdot\Delta_\beta}
= \dfrac{1}{\Delta_\alpha\cdot\Delta_\beta}\cdot\sum{}'(-1)^{\bar\nu}\Delta(\alpha)\Delta(\beta);
\]
denn $(-1)^{\mu+\nu}=+1$, da $\Delta(a)$ und $\Delta(b)$ correspondirende Determinanten und somit die Summe aller Doppelindices ihrer Diagonalglieder gleich der doppelten Summe aller ersten oder zweiten Indices derselben ist. Die Summe
\[
\sum{}'(-1)^{\bar\nu}\Delta(\alpha)\Delta(\beta)
\]
ist aber nichts anderes, als die \gesperrt{Laplace'sche} Entwicklung der
Determinante $\Delta$; daher erhalten wir die Relation
\[
R_{f,\varphi} = \dfrac{1}{\Delta_\alpha\cdot\Delta_\beta}\cdot\Delta = U.
\]

\begin{center}
\S.\,14.\quad\textbf{Die Discriminante.}
\end{center}

173. \textit{Definition der Discriminante.} Denken wir uns einen Augenblick die Function $n^{\text{ten}}$ Grades
\[
f(x) = x^n + a_1x^{n-1} + a_2x^{n-2} + \cdots a_{n-1}x + a_n
\]
homogen gemacht in den Variabeln $x_1$ und $x_2$, so können wir, wenn man
\[
\frac{1}{n}\frac{\partial f}{\partial x_1} = f_1 \quad\text{und}\quad \frac{1}{n}\frac{\partial f}{\partial x_2} = f_2
\]
setzt, gemäss dem \gesperrt{Euler'schen} Satze $f$ auch in der Form
schreiben:
\[
f = x_1f_1 + x_2f_2\,.
\]
Wir definiren nun als Discriminante der Function $f$ die Resultante
von $f$, $f_1$ und schreiben
\[
R_{f_1,f} = \Delta_f\,.
\]
Gemäss zweier Eigenschaften der Resultante (vgl. Nr. 146) ist
\[
R_{f_1,f} = R_{f_1,\,x_1f_1+x_2f_2} = \rho\cdot R_{f_1,f_2} = \Delta_f\,,
\]
d. h. die Discriminante ist auch die Resultante der beiden ersten Differentialquotienten nach $x_1$ und $x_2$ bis auf eine Constante $\rho$, welche
hier $= 1$ ist.

Aus der Definition der Discriminante als Resultante folgt einmal,
dass sie eine homogene ganze Function der Co\"efficienten $a_0\,a_1\ldots a_m$
vom Grade $2(m-1)$ ist; denn $\Delta_f = R_{f_1,f_2}$; $f_1$ und $f_2$ sind aber beide
vom Grade $(m-1)$ in $x$ (vergl. Nr. 131). Betrachtet man anderntheils
die Discriminante als Resultante von $f$ und $f_1$, so folgt als
%CONT>>


% ===== pdf idx 200, printed 186.0 =====
% !PRINTED 186
\weberpage{186}
186\quad Zweiter Theil. Anwendungen der Determinantentheorie.

%<<CONT
Gewicht dieser isobaren Function der Co\"efficienten das Product der
Grade beider Functionen $f$ und $f_1$, also die Zahl $m(m-1)$ (vergl.
Nr. 141).

174. \textit{Berechnung der Discriminante.} Aus der Definition
\[
\Delta_f = R_{f_1,\,f} = (-1)^{n(n-1)}R_{f,f_1} = R_{f,f_1}
\]
erhält man unmittelbar den Werth der Discriminante ausgedrückt in
den Wurzeln der gleich Null gesetzten Function $f(x)$. Es ist nämlich
nach Nr. 168, (4)
\[
R_{f,f_1} = f_1(\alpha_1)\cdot f_1(\alpha_2)\cdot f_1(\alpha_3)\ldots f_1(\alpha_n)\,.
\]
Nun ist aber, wenn wir uns $f$ als Product seiner linearen Factoren
denken:
\[
f_1(x) = \frac{1}{n}\frac{\partial f}{\partial x_1}
= \frac{1}{n}\sum{}'(x-\alpha_1)(x-\alpha_2)\ldots(x-\alpha_{k-1})(x-\alpha_{k+1})\ldots(x-\alpha_n)\,,
\]
wobei sich die Summe über alle jene Producte erstreckt, die irgend
einen der $n$ linearen Factoren nicht enthalten. Trägt man nun aber
an Stelle von $x$ eine Wurzel $\alpha_k$ der Gleichung ein, so reducirt sich
diese Summe auf das eine Product, in welchem der lineare Factor
\[
(x - \alpha_k)
\]
von vornherein fehlte. Es ist daher
\[
f_1(\alpha_k) = \frac{1}{n}(\alpha_k-\alpha_1)(\alpha_k-\alpha_2)\ldots(\alpha_k-\alpha_{k-1})(\alpha_k-\alpha_{k+1})\ldots(\alpha_k-\alpha_n)
\]
und folglich:
\[
\begin{gathered}
R_{f,f_1} = \dfrac{1}{n^n}(\alpha_1-\alpha_2)(\alpha_1-\alpha_3)\ldots(\alpha_1-\alpha_n)\times(\alpha_2-\alpha_1)(\alpha_2-\alpha_3)\ldots(\alpha_2-\alpha_n)\times{}\\
{}\times(\alpha_3-\alpha_1)(\alpha_3-\alpha_2)(\alpha_3-\alpha_4)\ldots(\alpha_3-\alpha_n)\ldots(\alpha_n-\alpha_1)(\alpha_n-\alpha_2)\ldots(\alpha_n-\alpha_{n-1})\,.
\end{gathered}
\]
In diesem Producte tritt aber jede Differenz zweimal auf, einmal
positiv und einmal negativ, und da man aus $n$ Elementen $\dfrac{n(n-1)}{2}$
solche Differenzen bilden kann, so ist:
\weq{I}{\Delta_f = R_{f,f_1} = \dfrac{(-1)^{\frac{n(n-1)}{2}}}{n^n}\cdot\prod_{\substack{i=1\\k=2}}^{\substack{i=n-1\\k=n}}(\alpha_i-\alpha_k)^2\,,\quad i<k.}
Anderntheils erhält man den Werth der Discriminante in Function
der Co\"efficienten von $f$, wegen
\[
\Delta_f = R_{f_1,f_2}\,,
\]
ausgedrückt durch die Determinante:
%CONT>>


% ===== pdf idx 201, printed 187.0 =====
% !PRINTED 187
\weberpage{187}
%<<CONT
% [column spec widened array 9->10] a row carried more cells than the spec declared,
% which is a fatal alignment error. The spec was widened rather than a cell
% deleted, so no transcribed content is lost; whether Weber's grid really has
% that many columns is a question for a pass with the scan open.

\S.\,14. Die Discriminante.

\[
R_{f_1,f_2} =
\left|
\begin{array}{cccccccccc}
1, & \dfrac{n-1}{n}a_1, & \dfrac{n-2}{n}a_2 & \cdots & \dfrac{a_{n-1}}{n} & 0 & 0 & \cdots & 0 \\[2ex]
0 & 1, & \dfrac{n-1}{n}a_1 & \cdots & \cdots & \dfrac{a_{n-1}}{n}, & 0 & \cdots & 0 \\[2ex]
\cdot & \cdot & \cdot & \cdot & \cdot & \cdot & \cdot & \cdot & \cdot \\[1ex]
0 & 0 & \cdots & \cdots & 1, & \dfrac{n-1}{n}a_1, & \dfrac{n-2}{n}a_2 & \cdot & \dfrac{a_{n-1}}{n} \\[2ex]
\dfrac{a_1}{n} & \dfrac{2a_2}{n} & \dfrac{3a_3}{n} & \cdots & \cdots & a_n & 0 & 0 & \cdots & 0 \\[2ex]
0 & \dfrac{a_1}{n} & \dfrac{2a_2}{n} & \cdots & \dfrac{(n-1)}{n}a_{n-1}, & a_n & 0 & \cdots & 0 \\[2ex]
\cdot & \cdot & \cdot & \cdot & \cdot & \cdot & \cdot & \cdot & \cdot \\[1ex]
0 & 0 & \cdots & \cdots & \dfrac{a_1}{n} & \dfrac{2a_2}{n} & \cdots & \dfrac{(n-1)}{n}a_{n-1}, & a_n
\end{array}
\right| \tag{II}
\]

Wir haben somit die Discriminante von $f$ in zwei verschiedenen Formen ausgedrückt und wollen im Folgenden noch zwei weitere Beweise geben für die Uebereinstimmung der Ausdrücke (II) und (I).

175. \gesperrt{Zweiter Beweis.} Wir wissen (vgl. Nr.\,40 Beispiel 2 und 3), dass das Differenzenproduct
\[
P = (\alpha_1 - \alpha_2)(\alpha_1 - \alpha_3)(\alpha_1 - \alpha_4) \cdots (\alpha_1 - \alpha_n) \times
\]
\[
(\alpha_2 - \alpha_3)(\alpha_2 - \alpha_4) \cdots (\alpha_2 - \alpha_n) \times
\]
\[
\cdots \cdots \cdots \cdots
\]
\[
(\alpha_{n-2} - \alpha_{n-1})(n-2 - \alpha_n) \times % [sic] printed without alpha before "n-2" in the second factor
\]
\[
(\alpha_{n-1} - \alpha_n)
\]
sich darstellen lässt durch die Determinante:
\[
P =
\left|
\begin{array}{cccccc}
1 & 1 & 1 & 1 & \cdots & 1 \\
\alpha_1 & \alpha_2 & \alpha_3 & \alpha_4 & \cdots & \alpha_n \\
\alpha_1^2 & \alpha_2^2 & \alpha_3^2 & \alpha_4^2 & \cdots & \alpha_n^2 \\
\alpha_1^3 & \alpha_2^3 & \cdot & \cdot & \cdot & \alpha_n^3 \\
\cdot & \cdot & \cdot & \cdot & \cdot & \cdot \\
\alpha_1^{n-1} & \alpha_2^{n-1} & \alpha_3^{n-1} & \alpha_4^{n-1} & \cdot\cdot & \alpha_n^{n-1}
\end{array}
\right| .
\]
Bezeichnen wir nun die Summe der $\nu^{\mathrm{ten}}$ Potenzen aller Wurzeln mit $s_\nu$, setzen also wie früher (Nr.\,168)
\[
s_\nu = \alpha_1^{\ \nu} + \alpha_2^{\ \nu} + \alpha_3^{\ \nu} + \cdots + \alpha_n^{\ \nu},
\]
so ist das Quadrat dieser Determinante, das wir durch Multiplication von Zeile mit Zeile bilden wollen, gegeben in der neuen Determinante:
%CONT>>


% ===== pdf idx 202, printed 188.0 =====
% !PRINTED 188
\weberpage{188}
%<<CONT
% [column spec widened array 8->9] a row carried more cells than the spec declared,
% which is a fatal alignment error. The spec was widened rather than a cell
% deleted, so no transcribed content is lost; whether Weber's grid really has
% that many columns is a question for a pass with the scan open.

Zweiter Theil. Anwendungen der Determinantentheorie.

\[
P^2 =
\left|
\begin{array}{ccccc}
s_0 & s_1 & s_2 & \cdots & s_{n-1} \\
s_1 & s_2 & s_3 & \cdots & s_n \\
s_2 & s_3 & s_4 & \cdots & s_{n+1} \\
\cdot & \cdot & \cdot & \cdot & \cdot \\
s_{n-1} & s_n & s_{n+1} & \cdot\cdot & s_{2n-2}
\end{array}
\right| .
\]
Wir wollen nun dieselbe so umformen, dass sie mit Hilfe der Newton\-schen Formeln in die Determinante (II) Nr.\,174, bis auf einen Zahlenfactor übergeht.

176. Stürzen wir nun zunächst die Matrix der $(n-1)$ ersten Zeilen der Determinante $P^2$. Da hiebei alle Elemente in ihren Colonnen bleiben, also in den zweiten Indices keine Derangements entstehen, so ändert sich hiebei das Vorzeichen um $(-1)^\nu$, hiebei ist $\nu$ gleich der Summe aller Derangements der $n-1$ ersten Indices dieser Matrix, also
\[
\nu = \sum_{k=1}^{k=n-2} k = \dfrac{(n-1)(n-2)}{2} .
\]
Multipliciren wir hierauf die letzte Zeile noch mit $(-1)$, so erhalten wir
\[
(-1)^{\frac{(n-1)(n-2)}{2}+1} \cdot P^2 =
\left|
\begin{array}{ccccc}
s_{n-2} & s_{n-1} & s_n & \cdots & s_{2n-3} \\
s_{n-3} & s_{n-2} & s_{n-1} & \cdots & s_{2n-4} \\
s_{n-4} & \cdot & \cdot & \cdot & s_{2n-5} \\
\cdot & \cdot & \cdot & \cdot & \cdot \\
s_0 & s_1 & s_2 & \cdots & s_{n-1} \\
-s_{n-1} & -s_n & -s_{n+1} & \cdots & -s_{2n-2}
\end{array}
\right| .
\]
Diese Determinante ist aber der eine der beiden Factoren, in welche die Determinante
\[
D =
\left|
\begin{array}{ccccccccc}
s_0 & s_1 & s_2 & \cdots & s_{n-2} & s_{n-1} & s_n & \cdots & s_{2n-3} \\
0 & s_0 & s_1 & \cdots & s_{n-3} & s_{n-2} & \cdot & \cdot & s_{2n-4} \\
0 & 0 & s_0 & \cdots & s_{n-4} & \cdot & \cdot & \cdot & s_{2n-5} \\
\cdot & \cdot & \cdot & \cdot & \cdot & \cdot & \cdot & \cdot & \cdot \\
0 & 0 & \cdots & \cdots & s_0 & s_1 & s_2 & \cdots & s_{n-1} \\
-s_1 & -s_2 & -s_3 & \cdots & -s_{n-1} & -s_n & -s_{n+1} & \cdots & -s_{2n-2} \\
s_0 & 0 & 0 & \cdots & 0 & 0 & 0 & \cdots & 0 \\
0 & s_0 & 0 & \cdots & 0 & 0 & 0 & \cdots & 0 \\
0 & 0 & s_0 & \cdots & 0 & 0 & 0 & \cdots & 0 \\
\cdot & \cdot & \cdot & \cdot & \cdot & \cdot & \cdot & \cdot & \cdot \\
0 & 0 & 0 & \cdots & s_0 & 0 & 0 & \cdots & 0
\end{array}
\right|
\]
gemäss dem Laplace'schen Satze zerfällt.


% ===== pdf idx 203, printed 189.0 =====
% !PRINTED 189
\weberpage{189}

\S.\,14. Die Discriminante.

Der andere Factor dieser Determinante ist
\[
s_0^{\,n-2} \cdot (-1)^\mu,
\]
wobei (vergleiche Nr.\,55) $\mu$ gleich ist der Gesammtsumme aller zweiten und ersten Indices der Determinante mit dem Diagonalglied $s_0^{\,n-2}$, also
\[
\mu = (1+2+\cdots n-2) + \{(n+1)+(n+2)+\cdots(2n-2)\}
\]
\[
= (n-1)(n-2) + n(n-2).
\]
Diese Summe ist aber mit $n$ gerade oder ungerade; daher ist der zweite Factor wegen $s_0 = n$
\[
(-1)^n \cdot n^{n-2}
\]
und folglich:
\[
D = (-1)^n \cdot n^{n-2} \cdot (-1)^{\frac{(n-1)(n-2)}{2}+1} \cdot P^2 = (-1)^{\frac{n(n-1)}{2}} \cdot n^{n-2} \cdot P^2 .
\]
177. Nun entsteht aber die Determinante $D$ aus der Determinante
\[
\Delta =
\left|
\begin{array}{ccccccccc}
s_0 & s_1 & s_2 & \cdots & s_{n-2} & s_{n-1} & s_n & \cdots & s_{2n-3} \\
0 & s_0 & s_1 & \cdots & s_{n-3} & s_{n-2} & s_{n-1} & \cdots & s_{2n-4} \\
0 & 0 & s_0 & \cdots & s_{n-4} & s_{n-3} & \cdot & \cdot & s_{2n-5} \\
\cdot & \cdot & \cdot & \cdot & \cdot & \cdot & \cdot & \cdot & \cdot \\
0 & 0 & 0 & \cdots & s_0 & s_1 & s_2 & \cdots & s_{n-1} \\
-s_1 & -s_2 & -s_3 & \cdots & -s_{n-1} & -s_n & -s_{n+1} & \cdots & -s_{2n-2} \\
0 & -s_1 & -s_2 & \cdots & -s_{n-2} & -s_{n-1} & -s_n & \cdots & -s_{2n-3} \\
0 & 0 & -s_1 & \cdots & -s_{n-3} & \cdot & \cdot & \cdot & -s_{2n-4} \\
\cdot & \cdot & \cdot & \cdot & \cdot & \cdot & \cdot & \cdot & \cdot \\
0 & 0 & 0 & \cdots & -s_1 & -s_2 & -s_3 & \cdots & -s_n
\end{array}
\right| ,
\]
indem wir die erste Zeile zur $(n+1)^{\mathrm{ten}}$, die zweite zur $(n+2)^{\mathrm{ten}}$ u.\,s.\,w., endlich die $(n-1)^{\mathrm{te}}$ zur $2(n-1)^{\mathrm{ten}}$, d.\,i. zur letzten Zeile addiren. Es ist also auch
\[
\Delta = (-1)^{\frac{n(n-1)}{2}} \cdot n^{n-2} \cdot P^2 .
\]
Multipliciren wir aber die Determinante $\Delta$ mit der Determinante
\[
T =
\left|
\begin{array}{cccccccccc}
1 & 0 & 0 & \cdot & \cdot & \cdot & \cdot & \cdot & \cdot & 0 \\
a_1 & 1 & 0 & \cdot & \cdot & \cdot & \cdot & \cdot & \cdot & 0 \\
a_3 & a_1 & 1 & \cdot & \cdot & \cdot & \cdot & \cdot & \cdot & 0 \\
\cdot & \cdot & \cdot & \cdot & \cdot & \cdot & \cdot & \cdot & \cdot & \cdot \\
a_n & a_{n-1} & a_{n-2} & \cdots & 1 & 0 & \cdot & \cdot & \cdot & 0 \\
0 & a_n & a_{n-1} & \cdots & a_1 & 1 & 0 & \cdot & \cdot & 0 \\
0 & 0 & a_n & \cdots & a_2 & a_1 & 1 & \cdot & \cdot & 0 \\
\cdot & \cdot & \cdot & \cdot & \cdot & \cdot & \cdot & \cdot & \cdot & \cdot \\
0 & 0 & 0 & \cdots & a_n & a_{n-1} & \cdot & \cdot & \cdot & 1
\end{array}
\right| ,
\]
%CONT>>


% ===== pdf idx 204, printed 190.0 =====
% !PRINTED 190
\weberpage{190}
% [column spec widened array 9->10] a row carried more cells than the spec declared,
% which is a fatal alignment error. The spec was widened rather than a cell
% deleted, so no transcribed content is lost; whether Weber's grid really has
% that many columns is a question for a pass with the scan open.

Zweiter Theil. Anwendungen der Determinantentheorie.

%<<CONT
welche, da rechts der Diagonale nur verschwindende Elemente stehen, sich auf ihr Diagonalglied reducirt und also den Werth 1 besitzt, so erhalten wir:
\[
\Delta \cdot T = (-1)^{\frac{n(n-1)}{2}} \cdot n^{n-2} \cdot P^2 =
\]
\[
\left|
\begin{array}{ccccc}
s_0, & s_0a_1+s_1, & s_0a_2+s_1a_1+s_2,\cdots, & s_0a_{n-1}+s_1a_{n-2}+\cdots s_{n-2}a_1+s_{n-1}\cdots \\
0 & s_0, & s_0a_1+s_1, \cdots, & s_0a_{n-2}+s_1a_{n-3}+\cdots\cdots+s_{n-2},\cdots \\
0 & 0, & s_0,\cdots\cdots\cdots\cdots, & s_0a_{n-3}+s_1a_{n-4}+\cdots\cdots+s_{n-3},\cdots \\
\cdot & \cdot & \cdot & \cdot \\
0, & 0, & 0,\cdots\cdots\cdots s_0, s_0a_1+s_1, & s_0a_2+s_1a_1+s_2,\cdots \\
-s_1, & -(s_1a_1+s_2), & -(s_1a_2+s_2a_1+s_3),\cdots, & -(s_1a_{n-1}+s_2a_{n-2}+s_3a_{n-3}\cdots+s_n),\cdots \\
0 & -s_1, & -(s_1a_1+s_2), \cdots\cdots\cdots & -(s_1a_{n-2}+s_2a_{n-3}\cdots\cdot+s_{n-1})\cdots \\
0 & 0, & -s_1, & -(s_1a_1+s_2), \\
\cdot & \cdot & \cdot & \cdot \\
0, & 0, & 0,\cdots\cdots\cdots-s_1, & -(s_1a_1+s_2),\cdots, -(s_1a_{n-1}+\cdots+s_n)
\end{array}
\right|
\]
Tragen wir aber gemäss den Newton'schen Identitäten für die einzelnen Elemente ihre Werthe ein, so erhalten wir:
\[
(-1)^{\frac{n(n-1)}{2}} \cdot n^{n-2} \cdot P^2 =
\]
\[
\left|
\begin{array}{cccccccccc}
n, & (n-1)a_1, & (n-2)a_2 & \cdots & a_{n-1}, & 0, & 0 & \cdots & 0 \\
0 & n, & (n-1)a_1 & \cdots & 2a_{n-2}, & a_{n-1} & 0 & \cdots & 0 \\
0 & 0, & n & \cdots & \cdots & \cdots & \cdots & a_{n-1} & 0 \\
\cdot & \cdot & \cdot & \cdot & \cdot & \cdot & \cdot & \cdot & \cdot \\
0 & 0 & \cdots & \cdots & n, & (n-1)a_1, & \cdots & a_{n-2} & a_{n-1} \\
a_1 & 2a_2, & 3a_3, & \cdots & \cdots & na_n & 0, & 0 & \cdots & 0 \\
0 & a_1 & 2a_1, & \cdots & \cdots & (n-1)a_{n-1}, & na_n & 0 & \cdots & 0 \\ % [sic] printed "2a_1" — pattern elsewhere suggests 2a_2
\cdot & \cdot & \cdot & \cdot & \cdot & \cdot & \cdot & \cdot & \cdot \\
0 & 0 & 0 & \cdots & \cdots & a_1 & 2a_2, & 3a_3 & \cdots \cdots na_n
\end{array}
\right|
\]
\[
= n^{2n-2} \Delta_f,
\]
und folglich:
\[
R_{f_1,f_2} = \Delta_f = \dfrac{(-1)^{\frac{n(n-1)}{2}}}{n^n} \cdot \prod_{\substack{i=1 \\ k=2}}^{\substack{k=n \\ i=n-1}} (\alpha_i - \alpha_k)^2, \quad i<k.
\]
178. \textit{Ueber das Auftreten einer mehrfachen Wurzel von $f$ in der ersten Abgeleiteten von $f$.} Aus dem Differenzenproducte erkennt man,
%CONT>>


% ===== pdf idx 205, printed 191.0 =====
% !PRINTED 191
\weberpage{191}
\S.\,14. Die Discriminante.

dass, wenn $f$ eine Doppelwurzel besitzt, die Discriminante verschwinden
muss. Da aber die Discriminante Resultante ist von $f$ und $f_1$, so ergiebt
sich, dass auch $f_1$ den Doppelfactor von $f$ mindestens einmal
besitzen muss. Es frägt sich nun, wenn $f$ einen $\nu$-fachen linearen
Factor besitzt: in welchem Grade kommt derselbe der ersten Abgeleiteten $f_1$ zu.

Um dies zu untersuchen, schreiben wir die Function $f$ in der Form:
\[
f = (x-t)^\nu \cdot \varphi(x),
\]
wobei $(x-t)$ der $\nu$-fache Factor von $f$ ist, und $\varphi(x)=0$, die
Wurzel $x=t$ nicht mehr besitzt.

Es ist dann die erste Abgeleitete von $f$:
\[
f' = \nu (x-t)^{\nu-1} \varphi(x) + (x-t)^\nu \varphi'(x)
\]
\[
= (x-t)^{\nu-1} \{\nu \varphi(x) + (x-t) \varphi'(x)\}.
\]
Da der Ausdruck $\nu \varphi(x) + (x-t)\varphi'(x)$ für $x=t$ nicht verschwindet,
weil nach Voraussetzung $\varphi(t) \gtrless 0$, so besitzt $f'$ den Factor $(x-t)$
nur $(\nu-1)$mal. Wenn also z.\,B. $(x-t)$ Doppelfactor von $f$ ist, so
ist er einfacher Factor von $f'$.

\textit{Anmerkung.} Es besteht sodann eine Relation
\[
Af + Bf_1 = 0,
\]
in welcher $A$ und $B$ keinen gemeinsamen Factor besitzen. $B$ enthält
sämmtliche vielfache Factoren:
\[
x - t_2, \qquad x - t_2 \ldots % [sic] both subscripts printed as t_2
\]
von $f$ jeden als einfachen Factor und der grösste gemeinsame Factor $\Phi$
von $f$ und $B$ ist das Product derselben.
\[
\Phi = (x-t_1)(x-t_2) \ldots
\]
179. \textit{Berechnung der Doppelwurzel einer Function $f=0$.} Wir
haben in Nr. 167 gezeigt, dass, wenn der Coefficient der höchsten
Potenz in $f=1$ ist, stets die Beziehung besteht:
\[
R_{f_1,f} = R_{f_1,f_2}.
\]
Besitzt also $f$ eine Doppelwurzel, so ist dieselbe der gemeinschaftliche Factor der beiden Differentialquotienten $f_1$ und $f_2$ von $f$.
Ist also $\alpha$ die Doppelwurzel, so wird der \gesperrt{Cayley}'sche Ausdruck
\[
\frac{1}{y-x} \{f_1(x)f_2(y) - f_1(y)f_2(x)\} = \sum_{i=1}^{i=n} p_i y^{n-i}
\]
für $x=\alpha$ identisch null, was nur dann möglich ist, wenn in der
Summe rechts die Coefficienten von $y$ einzeln verschwinden.


% ===== pdf idx 206, printed 192.0 =====
% !PRINTED 192
\weberpage{192}
Zweiter Theil. Anwendungen der Determinantentheorie.

Man hat also zur Berechnung der Doppelwurzel $x=\alpha$ die
$n$ Gleichungen
\[
p_1 = c_{11}\alpha^{n-1} + c_{12}\alpha^{n-2} \cdots + c_{1n} = 0
\]
\[
p_2 = c_{21}\alpha^{n-1} + c_{22}\alpha^{n-2} \cdots + c_{2n} = 0
\]
\[
\cdot\ \cdot\ \cdot\ \cdot\ \cdot\ \cdot\ \cdot\ \cdot\ \cdot\ \cdot\ \cdot\ \cdot
\]
\[
p_n = c_{n1}\alpha^{n-1} + c_{n2}\alpha^{n-2} \cdots + c_{nn} = 0,
\]
welche linear und homogen sind in den $n$ Grössen:
\[
\alpha^{n-1}, \alpha^{n-2}, \ldots \alpha^2, \alpha, 1.
\]
Es ist daher:
\[
\alpha^{n-1} : \alpha^{n-2} : \ldots \alpha : 1 = A_{i1} : A_{i2} \ldots A_{in},
\]
wenn die Grössen $A_{ik}$ die Minoren einer Zeile der Determinante der
$c_{ik}$, mit abwechselndem Vorzeichen genommen, sind. Auch hier ist
wieder die Berechnung eine eindeutige, und folglich kann nur eine
Doppelwurzel und keine höhere Potenz eines linearen Factors auf
diesem Wege ermittelt werden.

180. \textit{Die ersten partiellen Differentialquotienten der Discriminante.}
Die Discriminante $\Delta_f = R_{f,f_1}$ lässt sich (nach Nr. 131) auch darstellen
in der Form:
\weq{I}{\Delta_f = R_{f,f_1} = A(x)f + B(x)f_1.}
Differentiiren wir diese Gleichung nach $a_\nu$, so kommt
\[
\frac{\partial \Delta_f}{\partial a_\nu} = \frac{\partial A}{\partial a_\nu} f(x) + x^{n-\nu} A(x) + \frac{\partial B}{\partial a_\nu} f_1 + \frac{(n-\nu)}{n} x^{n-\nu-1} \cdot B(x).
\]
Wenn demnach $x=\alpha$ wieder Doppelwurzel von $f$ ist, so erhält man
aus dieser Gleichung durch Substitution $x=\alpha$, da gemäss der Anmerkung in Nr. 178 auch $B(\alpha)=0$,
\[
\frac{\partial \Delta_f}{\partial a_\nu} = \alpha^{n-\nu} A(\alpha)
\]
und ebenso
\[
\frac{\partial \Delta_f}{\partial a_\mu} = \alpha^{n-\mu} A(\alpha),
\]
also
\[
\frac{\partial \Delta_f}{\partial a_\nu} : \frac{\partial \Delta_f}{\partial a_\mu} = \alpha^{n-\nu} : \alpha^{n-\mu}.
\]
Die partiellen Differentialquotienten von $\Delta_f$ nach den Coefficienten
$a_1$ bis $a_n$ sind also der ersten bis $n^{\text{ten}}$ Potenz des Doppelfactors von $f$
proportional. Man kann die eben erhaltenen Resultate wieder benutzen, um charakteristische Differentialgleichungen für die Discriminante
aufzustellen. Diese charakteristischen Differentialgleichungen sowohl
für die Discriminante als auch für die Resultante wurden zuerst von
\gesperrt{Brioschi} aufgestellt (Crelle's Journal Bd. 53 und Annali di Matem.
Bd. 2, 1859), später von \gesperrt{Noether} (\gesperrt{Faà di Bruno}, in's Deutsche über-
%CONT>>


% ===== pdf idx 207, printed 193.0 =====
% !PRINTED 193
\weberpage{193}
\S.\,14. Die Discriminante.

%<<CONT
setzt von \gesperrt{Walter}, Teubner 1881) eingehender auf ihren inneren Zusammenhang untersucht.

181. \textit{Dritter Beweis für die Uebereinstimmung beider Darstellungen (I) und (II) der Discriminante in} Nr. 174. Denken wir uns
die Resultante von $f(x)$ und $f(x+y)$ gebildet, etwa nach der \gesperrt{Sylvester}'schen Methode, so ist dieselbe vom Grade $n^2$ in $y$, da die Coefficienten $b_i$ von $f(x+y)$ sich auf $n$ Zeilen vertheilen, und jede Zeile
in immer anderer Colonne den Coefficienten $b_n$ enthält, der in $y$ vom
$n^{\text{ten}}$ Grade ist.

Von den $n^2$ Wurzeln $y$ dieser Resultante sind in erster Linie $n$
gleich null; denn jedesmal wenn
\[
\alpha_i + y = \alpha_i, \qquad i = 1, 2 \ldots n,
\]
ist, dann verschwindet die Resultante $R_{f(x),f(x+y)}$. Die übrigen $n(n-1)$
Wurzeln sind nichts anderes als die einmal positiv und einmal negativ
genommenen Wurzeldifferenzen
\[
\alpha_k - \alpha_i;
\]
denn die Resultante $R_{f(x),f(x+y)}$ verschwindet auch, wenn
\[
\alpha_i + y = \alpha_k, \qquad i \gtrless k,
\]
also wenn
\[
y = \alpha_k - \alpha_i,
\]
oder
\[
= \alpha_i - \alpha_k
\]
ist. Nun geht aus den Eigenschaften der Resultante (vgl. Nr. 146)
hervor, dass
\[
R_{f(x),f(x+y)} = R_{f(x),\{f(x+y)-f(x)\}} = R_{f(x), \frac{y\{f(x+y)-f(x)\}}{y}}
\]
und dass ferner (vergl. Nr. 153, 154, 155)
\[
R_{f(x), y\{f(x+y)-f(x)\}} = R_{f(x),y} \cdot R_{f(x), \frac{f(x+y)-f(x)}{y}}.
\]
Weil aber $R_{f(x),y} = y^n$, so folgt endlich
\[
R_{f(x),f(x+y)} = y^n \cdot R_{f(x), \frac{f(x+y)-f(x)}{y}} = y^n \cdot F(y).
\]
Hiebei ist $F(y)$ noch vom Grade $n(n-1)$ in $y$ und besitzt die
Wurzeln $(\alpha_k - \alpha_i)$, wo $k \gtrless i$ sein kann.

Der Coefficient $c_0$ der höchsten Potenz von $y$ in dieser Function $F(y)$
ist 1. Denn da der Coefficient von $x^n$ in $f(x)$ selbst 1 ist, so besitzen
zunächst die Diagonalglieder der $n-1$ ersten Zeilen in der Resultante
\[
R_{f, \frac{f(x+y)-f(x)}{y}}
\]
den Werth 1. Und weil die Potenzen $y^{n-1}$, welche sich nur in den
%CONT>>


% ===== pdf idx 208, printed 194.0 =====
% !PRINTED 194
\weberpage{194}
Zweiter Theil. Anwendungen der Determinantentheorie.

%<<CONT
Diagonalgliedern der $n$ letzten Zeilen befinden können, als letzte
Glieder der binomischen Entwicklungen $(x+y)^n$ gleichfalls nur den
Coefficienten 1 besitzen, so ist in der That in $F(y)$ der Coefficient
$c_0 = 1$. Daher ist das letzte Glied $c_n$ in $F(y) = 0$, --- und zwar
positiv genommen, da $F(y)$ von geradem Grade ist, --- das Product
aller Wurzeln $(\alpha_k - \alpha_i)$. Dieses letzte Glied wird aber erhalten, wenn
man in
\[
F(y) = R_{f, \frac{f(x+y)-f(x)}{y}}
\]
$y=0$ setzt. Die rechte Seite dieser Gleichung geht dadurch über in
\[
R_{f, \frac{\partial f}{\partial x}} = R_{f, nf_1} = n^n R_{f,f_1},
\]
und weil nach Nr. 139
\[
R_{f,f_1} = (-1)^{n(n-1)} R_{f,f_1} = R_{f_1,f}, % ?? uncertain: middle subscript on the printed page reads ambiguous (possibly "f_r"); transcribed as f_1 by context
\]
so erhalten wir schliesslich:
\[
n^n R_{f_1,f} = (-1)^{\frac{n(n-1)}{2}} \prod (\alpha_i - \alpha_k)^2, \qquad k>i,
\]
oder
\[
\Delta_f = \frac{(-1)^{\frac{n(n-1)}{2}}}{n^n} \cdot \prod_{\substack{i=1\\k=2}}^{\substack{k=n\\i=n-1}} (\alpha_i - \alpha_k)^2, \qquad i<k.
\]
182. \textit{Beispiel.} Die Discriminante der quadratischen Gleichung
\[
f(x) = x^2 + a_1 x + a_2 = 0
\]
ist:
\[
4\Delta_f = (-1)^{\frac{2\cdot 1}{2}} (\alpha_1-\alpha_2)^2 = -(\alpha_1-\alpha_2)^2 = \begin{vmatrix} 2 & a_1 \\ a_1 & 2a_2 \end{vmatrix} = 4a_2 - a_1^2.
\]
Sind also die Wurzeln $\alpha_1$ und $\alpha_2$ reell, so ist $\Delta_f$ negativ; sind sie
imaginär, so ist $\Delta_f$ positiv; ist $\alpha_1 = \alpha_2$, so ist
\[
\Delta_f = 0.
\]
Zur Berechnung der Doppelwurzel von $f$ hat man, da dieselbe einfache Wurzel der ersten Differentialquotienten ist, entweder die
Gleichung
\[
2x + a_1 = 0,
\]
oder
\[
a_1 x + 2a_2 = 0.
\]
183. \textit{Beispiel.} Die Discriminante der cubischen Gleichung
\[
f(x) = x^3 + px + q = 0
\]
ist:
%CONT>>


% ===== pdf idx 209, printed 195.0 =====
% !PRINTED 195
\weberpage{195}
\[
27\Delta_f = (-1)^{3\cdot\frac{2}{2}}(\alpha_1-\alpha_2)^2(\alpha_1-\alpha_3)^2(\alpha_2-\alpha_3)^2
\]
\[
= -(\alpha_1-\alpha_2)^2(\alpha_1-\alpha_3)^2(\alpha_2-\alpha_3)^2
\]
oder (vergl. Nr. 174, (II)):
\[
81\Delta_f = \begin{vmatrix} 3 & 0 & p & 0 \\ 0 & 3 & 0 & p \\ 0 & 2p & 3q & 0 \\ 0 & 0 & 2p & 3q \end{vmatrix} = 3(27q^2+4p^3),
\]
also:
\[
27\Delta_f = 27q^2+4p^3.
\]
Sind alle drei Wurzeln reell, so ist das Product
\[
(\alpha_1-\alpha_2)^2(\alpha_1-\alpha_3)^2(\alpha_2-\alpha_3)^2
\]
positiv, also die Discriminante $\Delta_f$ negativ.

Sind zwei Wurzeln imaginär, so ist das Product der quadrirten Wurzeldifferenzen negativ, daher $\Delta_f$ positiv. Besitzt $f(x)=0$ eine Doppelwurzel, so ist $\Delta_f=0$, und man kann dieselbe berechnen aus dem Differentialquotienten
\[
f_2 = 2px+3q=0.
\]
Nehmen wir die cubische Gleichung in ihrer allgemeinen Form:
\[
f = x^3+a_1x^2+a_2x+a_3=0,
\]
so ist immer noch
\[
\Delta_f = -\frac{1}{27}\{(\alpha_1-\alpha_2)(\alpha_1-\alpha_3)(\alpha_2-\alpha_3)\}^2.
\]
Den Ausdruck der Discriminante in den Coefficienten von $f$ finden wir aber am besten durch die \gesperrt{Cayley'sche} Determinante. Es ist nämlich:
\[
\begin{aligned}
\frac{1}{x-y}\cdot\{f_1(x)f_2(y)-f_2(x)f_1(y)\} &=\\
&= \begin{vmatrix} 3x^2+2a_1x+a_2, & a_1x^2+2a_2x+3a_3 \\ 3y^2+2a_1y+a_2, & a_1y^2+2a_2y+3a_3 \end{vmatrix}\cdot\frac{1}{x-y}\\
&= \begin{vmatrix} 3 & 2a_1 \\ a_1 & 2a_2 \end{vmatrix}\cdot\begin{vmatrix} x^2 & x & 1 \\ y^2 & y & 1 \end{vmatrix}\cdot\frac{1}{x-y}\\
&= \begin{vmatrix} 3 & 2a_1 \\ a_1 & 2a_2 \end{vmatrix}\cdot\begin{vmatrix} x^2 & x \\ y^2 & y \end{vmatrix}\cdot\frac{1}{x-y} + \begin{vmatrix} 3 & a_2 \\ a_1 & 3a_3 \end{vmatrix}\begin{vmatrix} x^2 & 1 \\ y^2 & 1 \end{vmatrix}\frac{1}{x-y}\\
&\quad + \begin{vmatrix} 2a_1 & a_2 \\ 2a_2 & 3a \end{vmatrix}\begin{vmatrix} x & 1 \\ y & 1 \end{vmatrix}\cdot\frac{1}{x-y} = % [sic] second entry printed "3 a" without subscript, unlike the "3a_3" elsewhere
\\
&= \begin{vmatrix} 3 & 2a_1 \\ a_1 & 2a_2 \end{vmatrix} xy + \begin{vmatrix} 3 & a_2 \\ a_1 & 3a_3 \end{vmatrix}(x+y) + \begin{vmatrix} 2a_1 & a_2 \\ 2a_2 & 3a_3 \end{vmatrix}\cdot
\end{aligned}
\]


% ===== pdf idx 210, printed 196.0 =====
% !PRINTED 196
\weberpage{196}
Ordnet man hier die rechte Seite nach $y$, so erhält man:
\[
\left\{\begin{vmatrix} 3 & 2a_1 \\ a_1 & 2a_2 \end{vmatrix} x + \begin{vmatrix} 3 & a_2 \\ a_1 & 3a_3 \end{vmatrix}\right\} y
+ \left\{\begin{vmatrix} 3 & a_2 \\ a_1 & 3a_3 \end{vmatrix} x + \begin{vmatrix} 2a_1 & a_2 \\ 2a_2 & 3a_3 \end{vmatrix}\right\} = p_1y+p_2.
\]
Die Elimination von $x$ aus den beiden Gleichungen:
\[
p_1=0, \qquad p_2=0
\]
liefert die Discriminante:
\[
-81\cdot\Delta_f=
\begin{vmatrix}
\begin{vmatrix} 3 & 2a_1 \\ a_1 & 2a_2 \end{vmatrix}, & \begin{vmatrix} 3 & a_2 \\ a_1 & 3a_3 \end{vmatrix} \\
\begin{vmatrix} 3 & a_2 \\ a_1 & 3a_3 \end{vmatrix}, & \begin{vmatrix} 2a_1 & a_2 \\ 2a_2 & 3a_3 \end{vmatrix}
\end{vmatrix}
= 4(3a_2-a_1^2)(3a_3a_1-a_2^2)-(9a_3-a_1a_2)^2
\]
oder:
\[
-27\Delta_f = 18a_1a_2a_3 - 4a_1^3a_2 - 4a_2^3 + a_1^2a_2^2 - 27a_3^2.
\]
Das negative Vorzeichen links rührt davon her (vergl. Nr. 150), dass die nach der \gesperrt{Cayley'schen} Methode berechnete Resultante $R'$ mit der durch die \gesperrt{Sylvester'sche} Methode gebildeten $R$ durch die Relation zusammenhängt
\[
(-1)^{\frac{n(n-1)}{2}} R'_{f_1,f_2} = R_{f_1,f_2}.
\]
Will man die Doppelwurzel von $f$ berechnen, so kann man sich hiezu (vergl. Nr. 179) einer der beiden Gleichungen $p_1=0$, $p_2=0$ bedienen.

§. 15. Simultanes System homogener linearer diophantischer Gleichungen.

(Anhang zu § 8.)

184. \textit{Aufsuchung ganzzahliger Lösungen überhaupt.} Wir hatten in Nr. 109 die allgemeinste Lösung eines Systemes linearer Gleichungen besprochen, in denen die Zahl der Unbekannten die Zahl der Gleichungen um mehr als 1 übertrifft. In manchen Fällen, insbesondere in gewissen Hauptaufgaben der Invariantentheorie, interessirt uns nicht die allgemeinste Lösung, sondern jene Lösungen, die nur ganzzahlige
%CONT>>


% ===== pdf idx 211, printed 197.0 =====
% !PRINTED 197
\weberpage{197}
%<<CONT
Werthe der Unbekannten, ja selbst nur ganzzahlige positive enthalten. Wenn nun auch die im Folgenden entwickelte Methode für die Aufsuchung solcher Werthsysteme die Determinantentheorie nicht zu Hilfe nimmt, so scheint doch hier der geeignete Platz zu sein, jene Punkte dieser Aufgabe zu besprechen, von denen wir im zweiten und dritten Bande Gebrauch machen werden. Die Theorie, die wir hiebei für die Auflösung solcher diophantischer Gleichungen entwickeln, wird in der Praxis im allgemeinen Falle wohl kaum verfolgt werden; für die folgenden Beweise erweist sie sich aber als äusserst zweckentsprechend.

\textit{Auflösung einer einzigen diophantischen Gleichung.} Es sei gegeben die lineare Gleichung mit $n$ Unbekannten:
\weq{I}{a_1x_1 + a_2x_2 + a_3x_3 + \cdots a_nx_n = 0.}
Alle Coefficienten sind ganzzahlig und die Aufgabe ist, ganzzahlige Werthe der Unbekannten zu suchen, welche sie befriedigen. Zu dem Zwecke nehmen wir den numerisch kleinsten Coefficienten $a_i$ der Gleichung und dividiren mit ihm alle übrigen $a_k$. Ist dann $q_k$ der sich ergebende Quotient und $r_k$ der Rest, so ist
\[
a_k = q_ka_i + r_k.
\]
Trägt man diese Werthe von $a_k$ in (I), so erhält man:
\[
a_i\left(x_i + \sum_k q_kx_k\right) + r_1x_1 + r_2x_2 + r_3x_3 + \cdots r_nx_n = 0,
\]
oder, wenn wir setzen:
\[
x_i + \sum_k q_kx_k = \lambda_i,
\]
so folgt:
\weq{II}{a_i\lambda_i + r_1x_1 + r_2x_2 + \cdots r_nx_n = 0.}
Die Coefficienten dieser neuen Gleichung sind nun kleiner als die der ersten, und die Aufgabe ist somit auf eine leichtere reducirt. Dieses Verfahren wiederholt man so lange, bis irgend eine Unbekannte als kleinsten Coefficienten die Zahl 1 besitzt, bis man also die Gleichung erhält:
\weq{III}{\lambda_i = \varrho_1\lambda_1 + \varrho_3\lambda_2 + \varrho_3\lambda_3 + \cdots \varrho_{n-1}\lambda_{n-1},} % [sic] subscripts printed 1,3,3,...,n-1
wo die Grössen $\varrho_i$ irgend welche ganze Zahlen und die $\lambda_i$ theils neu eingeführte, theils noch Unbekannte der ersten Gleichung sein werden. Durch Eliminationen erhält man sodann für jede Unbekannte einen Werth:
%CONT>>


% ===== pdf idx 212, printed 198.0 =====
% !PRINTED 198
\weberpage{198}
%<<CONT
\weq{IV}{x_i = \xi_i\lambda_1 + \eta_i\lambda_2 + \cdots \pi_i\lambda_{n-1},}
wobei $\xi_i$, $\eta_i \ldots \pi_i$ ganze Zahlen sind, deren mindestens eine von null verschieden ist. Damit ist aber die allgemeinste Lösung der diophantischen Gleichung aufgestellt. Hiebei ist die Zahl der Parameter $\lambda$ in (IV) immer um 1 kleiner als die Zahl der Unbekannten.

185. \textit{Beispiel.}
\weq{1}{5x+8y+12z-43t=0.}
Da der kleinste Coefficient die Zahl 5 ist, so setzen wir
\weq{2}{x+2y+2z-9t=\lambda_1}
und erhalten:
\[
5\lambda_1 - 2y + 2z + 2t = 0.
\]
Substituirt man hierin
\weq{3}{2\lambda_1-y+z+t=\lambda_2,}
dann erhält man die Gleichung
\weq{4}{-2\lambda_2=\lambda_1.}
Trägt man diesen Werth von $\lambda_1$ in (3) ein, so kommt:
\weq{5}{t=y-z+5\lambda_2.}
Hiezu tritt:
\weq{6}{x=7y-12z+43\lambda_2}
durch Substitution des Werthes von $\lambda_1$ und $t$ aus (4) und (5) in (2). Diese beiden Gleichungen (5) und (6) stellen in Verbindung mit
\[
y=y, \qquad z=z
\]
die allgemeine Lösung der diophantischen Gleichung (1) dar. Für
\[
\begin{aligned}
y&=1, & z&=0, & \lambda_2&=0\\
z&=1, & y&=0, & \lambda_2&=0\\
\lambda_2&=1, & y&=0, & z&=0
\end{aligned}
\]
erhalten wir die Speciallösungen:
\[
\begin{aligned}
t&=1, & -1, & \quad 5\\
x&=7, & -12, & \quad 43\\
y&=1, & 0, & \quad 0 % ?? uncertain: printed glyph reads more like "q" than "y"
\\
z&=0, & 1, & \quad 0
\end{aligned}
\]
und man erkennt, dass die Coefficienten $\xi_i$, $\eta_i \ldots \pi_i$ in (IV) nichts anderes sind, als diese einfachen Lösungen der gegebenen Gleichung.


% ===== pdf idx 213, printed 199.0 =====
% !PRINTED 199
\weberpage{199}
186. \gesperrt{Auflösung von mehreren simultanen diophantischen Gleichungen.} Die allgemeine Lösung Einer Gleichung ist also dargestellt durch
\[ x_i = \lambda_1 \xi_i + \lambda_2 \eta_i + \lambda_3 \zeta_i + \cdots \lambda_{n-1} \pi_i, \quad i = 1, 2 \ldots n, \]
in der alle Grössen rechts ganzzahlig sind. Hat man alsdann zwei diophantische Gleichungen mit $n$ Unbekannten aufzulösen, so kann man zunächst die erste Gleichung nach der vorhergehenden Methode (Nr. 184) behandeln und die sich so ergebenden Werthe der $x_i$ in die zweite Gleichung eintragen. Da die Zahl der Parameter $\lambda_i$ um 1 kleiner ist als die Zahl der Unbekannten, so erhält man dadurch eine Gleichung in $\lambda_i$ mit $n - 1$ Unbekannten. Sie wird ebenso behandelt wie die erste Gleichung für $x_i$, und man wird ein Werthsystem
\[ \lambda_i = \mu_1 \alpha_i + \mu_2 \beta_i + \mu_3 \gamma_i \cdots + \mu_{n-2} \varphi_i \]
erhalten, dessen Substitution in die Lösungen für $x_i$ die allgemeine Lösung der beiden simultanen diophantischen Gleichungen ergiebt. Ganz in derselben Weise wiederholt man die Operationen bei Auflösung von drei, vier und mehr diophantischen Gleichungen und man erhält für $\varrho$ Gleichungen mit $n$ Unbekannten, $\varrho < n$:
\weq{V}{x_i = \xi_i \nu_1 + \eta_i \nu_2 + \zeta_i \nu_3 + \cdots \pi_i \nu_{n-\varrho}.}
187. \gesperrt{Ganzzahlige positive Lösungen.} Zur Aufsuchung der ganzzahligen positiven Lösungen besitzen wir keine systematische Methode wie vorhin. Man erkennt von vornherein, dass die Gleichung
\[ a_x = a_1 x_1 + a_2 x_2 + a_3 x_3 + \cdots a_n x_n = 0 \]
überhaupt keine positive Lösung zulässt, sobald die Coefficienten $a_i$ selbst durchweg positiv sind, die identische Lösung $x_i = 0$ ausgenommen. Dasselbe gilt für ein System von diophantischen Gleichungen, sobald nur eine Combination existirt, welche lauter positive Coefficienten aufweist, da ja diese durch positive Lösungen nicht befriedigt werden kann. Man darf also annehmen, dass eine Anzahl Coefficienten negativ ist; wir bezeichnen sie mit $b_i$ und schreiben daher die Gleichung (I) in der Form
\weq{VI}{a_x = b_y.}
Von dieser Gleichung behaupten wir: erstens, sie hat immer positive Lösungen; zweitens, die Zahl der einfachen positiven Lösungen, die ihr zukommen, ist eine endliche. Dabei verstehen wir unter einfachen Lösungen diejenigen Lösungen $x_i$, welche sich nicht mehr aus anderen in der Form
\[ x_i = \xi_i + \eta_i \]
%CONT>>


% ===== pdf idx 214, printed 200.0 =====
% !PRINTED 200
\weberpage{200}
%<<CONT
zusammensetzen lassen. Solche Lösungen wie diese nennen wir zusammengesetzte.

188. \gesperrt{Specielle ganzzahlige Lösungen.} Was die erste Behauptung betrifft, so ist deren Richtigkeit leicht einzusehen. Wir setzen zunächst rechts wie links alle Unbekannten bis auf je eine gleich null und erhalten
\[ a_i x_i = b_k y_k. \]
Die Gleichung (VI) wird befriedigt für die positiven Werthe
\[ x_i = b_k, \quad y_k = a_i, \quad x_1 = x_2 = \cdots y_\varrho = y_\sigma = \cdots \text{u. s. w.} = 0. \]
Solche Lösungen giebt es $m \cdot n$ an der Zahl. Ausser diesen existiren noch sehr viele positive einfache Lösungen; aber die Gesammtzahl derselben muss eine endliche sein. Denn liegt eine Lösung vor $x_i = \xi_i$ und $y_k = \eta_k$, so dass gleichzeitig
\[ \xi_i > b_k \]
\[ \eta_k > a_i \]
ist, dann ist diese Lösung keine einfache mehr, sondern eine zusammengesetzte. Die Gleichung (VI) wird nämlich auch befriedigt durch die Differenz dieser mit der vorerwähnten Lösung. Die Werthe $\xi_i, \eta_k \ldots$ u. s. w. der einfachen Lösungen haben also obere Grenzen, und können demnach nur in endlicher Zahl vorhanden sein. [Der Einwand, dass $\xi_i$ allein über alle Grenzen wachsen könne, während $\eta_k$ die Grenze $a_i$ nicht überschreite, ist a priori hinfällig, da ja die Gleichung (VI) bestehen muss.]

189. \gesperrt{Allgemeine positive ganzzahlige Lösungen.} Ist nun $N$ die Zahl aller einfachen positiven Lösungen, etwa
\[
\begin{array}{l}
x_1 = \xi_1 \quad \eta_1 \quad \zeta_1 \ldots \pi_1 \\
x_2 = \xi_2 \quad \eta_2 \quad \zeta_2 \ldots \pi_2 \\
\hphantom{x_2 =} \cdot \; \cdot \; \cdot \; \cdot \; \cdot \; \cdot \; \cdot \; \cdot \\
y_1 = \alpha_1 \quad \beta_1 \quad \gamma_1 \ldots \psi_1 \\
y_2 = \alpha_2 \quad \beta_2 \quad \gamma_2 \ldots \psi_2, \\
\hphantom{y_2 =} \cdot \; \cdot \; \cdot \; \cdot \; \cdot \; \cdot \; \cdot \; \cdot
\end{array}
\]
so erhalten wir wie vorhin (Nr. 184) die allgemeinste positive Lösung, wenn wir die Verticalreihen dieses Systems multipliciren mit
\[ \lambda_1, \quad \lambda_2, \quad \lambda_3 \ldots \lambda_N \]
und die Producte einer Zeile addiren, nämlich
\weq{VII}{\begin{cases} x_i = \xi_i \lambda_1 + \eta_i \lambda_2 + \cdots \pi_i \lambda_N \\ y_k = \alpha_k \lambda_1 + \beta_k \lambda_2 + \cdots \psi_k \lambda_N \end{cases}}


% ===== pdf idx 215, printed 201.0 =====
% !PRINTED 201
\weberpage{201}
Tritt nun zur ersten Gleichung (VI)
\[ a_x = b_y \]
noch eine simultane zweite
\[ c_x = d_y, \]
so tragen wir wie vorhin die Werthe von $x_i$ und $y_k$ aus (VII) in diese zweite Gleichung ein. Wir erhalten so eine neue Gleichung mit $N$ Unbekannten $\lambda_i$. Die Zahl ihrer einfachen positiven Lösungen ist wie vorhin nothwendig endlich, und ihre allgemeine positive Lösung von der Form
\[ \lambda_i = \mathfrak{a}_i \mu_1 + \mathfrak{b}_i \mu_2 + \mathfrak{c}_i \mu_3 + \mathfrak{d}_i \mu_4 + \cdots + \mathfrak{r}_i \mu_M. \]
Wir tragen diese Werthe in das System (VII) ein und erhalten so die allgemeine positive Lösung der beiden simultanen diophantischen Gleichungen. Aus ihr gehen die einfachen Lösungen hervor, wenn wir eine Grösse $\mu_i = 1$, alle übrigen null setzen.

In derselben Weise erkennt man, dass ein beliebiges System von diophantischen Gleichungen nur eine endliche Zahl positiver einfacher Lösungen besitzt.

\weberfigure{201}{Typographisches Schlussornament: waagrechte Zierlinie mit zentraler Raute, flankiert von zwei kleineren Rauten, markiert das Ende von \S.\,15.}


% ===== pdf idx 216, printed 201.001 =====
% !PRINTED none
\weberfigure{none}{Leere Seite (kein Drucksatz); nur vereinzelte winzige Druck-/Scanflecken sichtbar, kein Text und keine Seitenzahl.}


% ===== pdf idx 217, printed 201.002 =====
% !PRINTED none
% Publisher's advertisement page (Teubner Verlagsverzeichnis), no running folio visible in the scan crop.

\textbf{Verlag von B. G. Teubner in Leipzig.}

\rule{6em}{0.4pt}

\textbf{Clebsch}, Dr. A., Prof. an der polytechnischen Schule zu Carlsruhe [\dag{} als Professor in G\"ottingen], Theorie der Elasticit\"at fester K\"orper. [XI u. 424 S.] gr. 8. 1862. geh. \hfill n. M. 9.--

\glqq{}Der Herr Verfasser hatte als Lehrer an der polytechnischen Schule zu Karlsruhe Gelegenheit und Beruf, sich ausf\"uhrlicher mit den Anwendungen der allgemeinen Theorie der Elastizit\"at auf die in der Technik besonders wichtigen F\"alle zu besch\"aftigen. Die Resultate dieser Studien liegen uns jetzt in einem ziemlich umfangreichen Werke vor, und man kann dem Verfasser nur Dank wissen, dafs er unsere deutsche Litteratur um eine Schrift bereichert hat, welche einerseits dem Techniker das Erlernen der strengen Theorie erm\"oglicht, ihm \"uber die Genauigkeit seiner Resultate und die Zul\"assigkeit der in der Praxis \"ublichen Voraussetzungen Aufschlufs giebt, andererseits den Mathematiker belehrt, wie man von den allgemeinsten Gleichungen der Bewegungen und des Gleichgewichts elastischer K\"orper zu speziellen F\"allen gelangen kann, und ihm die grofse M\"uhe und Zeit erspart, in den Arbeiten der Techniker den Weizen von der Spreu zu sondern. Es erg\"anzt daher dieses Handbuch das ber\"uhmte Werk des franz\"osischen Physikers Lam\'e, welches vorz\"uglich die allgemeinen Differentialgleichungen, ihre eleganten Transformationen, die Theorie der krystallinischen K\"orper und ihre optischen Eigenschaften behandelt, w\"ahrend Herr Clebsch ausschliefslich unkrystallinische K\"orper und deren Verschiebungen durch \"aufsere Kr\"afte in Betrachtung zieht.\grqq{} \hfill [Literarisches Centralblatt 1863, No. 31.]

\textemdash{} Theorie der bin\"aren algebraischen Formen. [VIII u. 467 S.] gr. 8. 1871. geh. \hfill n. M. 11.--

Gegen\"uber dem von Hrn. Fiedler \"ubersetzten Salmonschen Lehrbuche, welches bisher das Hauptmittel f\"ur das Studium der neuern Algebra war, wird man in dem vorliegenden Werke einerseits eine Beschr\"ankung, andrerseits Erweiterungen finden. Die Beschr\"ankung besteht darin, dafs nur bin\"are Formen behandelt sind. Dies wurde fast geboten durch den Umstand, dafs nur die Theorie der bin\"aren Formen bis jetzt in gewissem Grade abgeschlossen scheint. Der Verfasser geht davon aus, Invarianten und Kovarianten durch symbolische Produkte zu definieren; ein Standpunkt, welcher ihm schon vor 10 Jahren als der richtigste erschien. Diese Ansicht hat im Laufe der Zeit um so mehr sich bei ihm festigen m\"ussen, als gerade hiervon ausgehend Herr Gordan den wichtigen Satz fand, dafs die Anzahl der zu jeder Form geh\"origen Invarianten und Kovarianten eine endliche sei; und mit diesem Satze waren denn wieder Fragen gegeben, welche nur bei bin\"aren Formen aufgestellt und beantwortet werden k\"onnen, und welche eben der Theorie der bin\"aren Formen jenen vorzugsweise entwickelten Charakter geben.

Indem diese und andere, auf fr\"uhern Arbeiten des Verfassers beruhenden Untersuchungen ausf\"uhrlich behandelt sind, ist dem erw\"ahnten Werke gegen\"uber eine sehr wesentliche Fortentwickelung gegeben, welche bis zu dem neuesten Stande dieser Disziplin f\"uhrt.

Andrerseits wird in dem Werke die geometrische Interpretation ber\"ucksichtigt, welche neben den algebraischen Untersuchungen hergeht, und welche zugleich die Grundlage der synthetischen Geometrie liefert, n\"amlich die Theorie der Punktreihen und Strahlenb\"uschel.

Es ist endlich ein grofses Gewicht auf ausf\"uhrliche Behandlung einiger Probleme gelegt, welche dem Studierenden als Muster f\"ur die Behandlung algebraischer Fragen dienen k\"onnen, und welche ihn so auch praktisch in einen der wichtigsten Teile der neuern Mathematik einf\"uhren.

\rule{6em}{0.4pt} Vorlesungen \"uber Geometrie. Bearbeitet und herausgegeben von Dr. Ferdinand Lindemann. Mit einem Vorwort von Felix Klein. Erster Band. [XII u. 1050 S.] gr. 8. 1875. geh. \hfill n. M. 24.--

Auch in zwei Teilen:

I. Teil [S. 1--496]. n. M. 11.20.

II. \ ,, \ [S. I--XII u. 497--1050]. n. M. 12.80.

\textbf{Es bedarf wohl kaum einer n\"ahern Erkl\"arung, wenn nach dem nur allzu fr\"uh erfolgten Hinscheiden von A. Clebsch alsbald der Gedanke entstand, seine Vorlesungen \"uber Geometrie, welche f\"ur das Studium der Wissenschaft von so hervorragendem Einflusse waren, allgemein zug\"anglich zu machen.} Es bezieht sich dies sowohl auf die ausschliefslich geometrischen, als auch auf einzelne Abschnitte anderer Vorlesungen, soweit in letzteren geometrische Probleme gelegentlich behandelt wurden. Dem Herausgeber fiel somit die Aufgabe zu, in m\"oglichstem Anschlusse an den Vortrag, wof\"ur ihm mehrere nachgeschriebene Hefte zur Verf\"ugung standen, und teilweise unter Benutzung von Clebsch herr\"uhrender Manuskripte ein zusammenh\"angendes Bild der verschiedenen geometrischen Disziplinen zu entwerfen, besonders unter Wahrung des einheitlichen Gedankenganges, welcher den Vorlesungen von Clebsch so sehr eigent\"umlich war. W\"ahrend sonach in den ersten Abteilungen die rein geometrischen Gesichtspunkte vorwalten, dr\"angt sich schon bei den Kegelschnitten, bez. den Fl\"achen 2. Ordnung die Notwendigkeit eines eingehenden Studiums der Invariantentheorie auf; und dieser algebraische Charakter kommt bei weiterem Fortschreiten immer mehr zur Geltung. Andererseits soll das Buch auch in jene \"uberaus fruchtbaren Untersuchungen einf\"uhren, welche aus der Theorie der elliptischen und Abelschen Funktionen, sowie der
%CONT>>


% ===== pdf idx 218, printed 201.003 =====
% !PRINTED none
% Publisher's advertisement page (Teubner Verlagsverzeichnis), no running folio visible in the scan crop.
%<<CONT
eindeutigen Transformationen f\"ur die Geometrie erwachsen. Die erw\"ahnten Ziele bezeichnen zugleich die Stellung dieser Vorlesungen gegen\"uber den von Herrn Fiedler \"ubersetzten und bearbeiteten Salmonschen Lehrb\"uchern, indem der Stoff, im einzelnen mehr beschr\"ankt, im ganzen doch ein weiteres Gebiet umfafst.

In dem ersten Bande wird in dem angegebenen Sinne die Geometrie der Ebene behandelt. Den Schlufs bildet die Darstellung der Grundz\"uge einer Theorie der Konnexe, jener von Clebsch noch zuletzt in die Geometrie eingef\"uhrten Gebilde.

\textbf{Clebsch}, Dr. A., und \textbf{P. Gordan}, Professor in Erlangen, Theorie der Abel'schen Functionen. [XIII u. 333 S.] gr. 8. 1866. geh. \hfill n. M. 7.20.

Die Herren Verfasser suchen in diesem Werke die Theorie der Abelschen Funktionen auf eine ganz neue Weise zu begr\"unden, welche das Interesse der Mathematiker in hohem Grade erregt und das Verst\"andnis der Theorie wesentlich gef\"ordert hat.

\textbf{Gordan}, Dr. \textbf{Paul}, ord. Professor der Mathematik an der Universit\"at zu Erlangen, \"uber das Formensystem bin\"arer Formen. [52 S.] gr. 8. 1875. geh. \hfill n. M. 2.--

\textbf{Hesse}, Dr. \textbf{Otto}, weil. Professor am k\"onigl. Polytechnikum zu M\"unchen, Vorlesungen \"uber analytische Geometrie des Raumes, insbesondere \"uber Oberfl\"achen zweiter Ordnung. Revidirt und mit Zus\"atzen versehen von Dr. S. Gundelfinger, a. o. Professor an der Universit\"at zu T\"ubingen. Dritte Auflage. [XVI u. 546 S.] gr. 8. 1876. geh. \hfill n. M. 13.--

\glqq{}Ungeachtet des bedeutenden Aufschwunges, welchen die analytische Geometrie namentlich im Verlaufe des letzten Vierteljahrhunderts einerseits durch die Erweiterung des Koordinatenbegriffes, andererseits durch die Fortschritte der algebraischen Methoden, besonders in der Theorie der Determinanten und der homogenen Funktionen, genommen hat, fehlte es doch noch bis vor kurzem an einem Lehrbuche, welches geeignet gewesen w\"are, den Studierenden der Mathematik zur Einf\"uhrung in diese neueren Disziplinen zu dienen. F\"ur die analytische Geometrie der Ebene ist zu diesem Zwecke den deutschen J\"ungern der Wissenschaft ein wichtiges Hilfsmittel in der Fiedlerschen \"Ubertragung des Salmonschen Werkes in die Hand gegeben worden; f\"ur die des Raumes wird die erw\"ahnte L\"ucke unserer mathematischen Litteratur auf eine ausgezeichnete Weise durch das vorliegende Lehrbuch ausgef\"ullt. Dafs der Verfasser desselben vor allen berechtigt war, in diese L\"ucke einzutreten, dazu hat er sich den Anspruch durch seine r\"ustige Mitwirkung am Ausbau sowohl der analytisch-geometrischen Methoden, als der hiermit im Zusammenhange stehenden Teile der Algebra erworben; ein Blick in die letzten zwanzig Jahrg\"ange von Crelles Journal wird gen\"ugen, ihm diese Berechtigung zuzuerkennen. Gegen\"uber der Stellung des Verfassers auf dem Gebiete der Wissenschaft mufs Referent von einer kritischen Besprechung des vorliegenden Werkes absehen, umsomehr, als dieselbe nur auf Anerkennung des darin dargelegten Talentes und der Meisterschaft in der Darstellungsweise hinauslaufen k\"onnte. --- [Folgt Inhaltsangabe.] --- F\"ur Leser, welche mit den n\"otigen Vorkenntnissen ausger\"ustet, Zugang zu den neueren analytisch-geometrischen Theorieen erhalten wollen, kann das vorliegende Werk als eines der wichtigsten Hilfsmittel bezeichnet werden u. s. w.\grqq{} \hfill [O. Fort in der Zeitschrift f\"ur Mathematik u. Physik.]

\rule{6em}{0.4pt} sieben Vorlesungen aus der analytischen Geometrie der Kegelschnitte. Fortsetzung der Vorlesungen aus der analytischen Geometrie der geraden Linie, des Punktes und des Kreises. [Separatausgabe aus der Zeitschrift f\"ur Mathematik und Physik.] [52 S.] gr. 8. 1874. geh. \hfill n. M. 1.60.

Vergriffen.

\rule{6em}{0.4pt} Vorlesungen aus der analytischen Geometrie der geraden Linie, des Punktes und des Kreises in der Ebene. Dritte Auflage, revidiert von Dr. S. \textsc{Gundelfinger}, Professor am Grofsherzoglichen Polytechnikum zu Darmstadt. [VIII u. 230 S.] gr. 8. 1881. geh. \hfill n. M. 5.20.

\glqq{}Man kann ohne \"Ubertreibung behaupten, dafs es sehr wenige B\"ucher giebt, die auf dem kleinen Raume von 230 Seiten eine solche F\"ulle von Material in einer so eleganten und durchaus klaren Darstellung bieten u. s. w.\grqq{} \hfill [Schl\"omilch in d. Zeitschr. f. Mathematik u. Physik.]

\rule{6em}{0.4pt}


% ===== pdf idx 219, printed 201.004 =====
% !PRINTED none
% Publisher's advertisement page (Teubner Verlagsverzeichnis), no running folio visible in the scan crop.
% New catalogue section heading, set partly in Fraktur type in the original; transcribed in plain letters.

\textbf{Neuer Verlag von B. G. Teubner in Leipzig.}

1885.

\textbf{Bardey}, Dr. \textbf{E.}, methodisch geordnete Aufgabensammlung, mehr als 8000 Aufgaben enthaltend, \"uber alle Teile der Elementar-Arithmetik, f\"ur Gymnasien, vorzugsweise Realgymnasien u. Oberrealschulen. Zw\"olfte Auflage. [XIII u. 330 S.] gr. 8. geh. M. 2.70.

\textbf{Euclidis} opera omnia. Ediderunt I. L. \textsc{Heiberg} et H. \textsc{Menge}. Euclidis elementa. Edidit et latine interpretatus est I. L. \textsc{Heiberg}, Dr. phil. Vol. IV. Libros XI---XIII continens. [VI u. 423 S.] 8. geh. M. 4.50.

Vol. III erscheint sp\"ater.

\textbf{Fiedler}, Dr. \textbf{Wilhelm}, die darstellende Geometrie in organischer Verbindung mit der Geometrie der Lage. Dritte erweiterte Auflage. II. Theil. A. u. d. Titel: Die darstellende Geometrie der krummen Linien und Fl\"achen. F\"ur Vorlesungen und zum Selbststudium. [Mit zahlreichen Figuren im Text und 16 lithographierten Tafeln.] [XXXIII u. 560 S.] gr. 8. geh. n. M. 14.--

\textbf{Graefe}, Dr. \textbf{Fr.,} Professor, Aufgaben und Lehrs\"atze aus der analytischen Geometrie des Punktes, der geraden Linie, des Kreises und der Kegelschnitte. F\"ur Studierende an Universit\"aten und technischen Hochschulen bearbeitet. [IV u. 136 S.] gr. 8. geh. n. M. 2.40.

\textbf{Harnack}, \textbf{Axel}, Naturforschung und Naturphilosophie. Vortrag gehalten in der naturwissenschaftlichen Gesellschaft zu Dresden. [27 S.] gr. 8. geh. M. --.60.

\textbf{Hecht}, Dr. \textbf{Wilhelm}, Dozent der Mathematik an der Kgl. Forstlehranstalt zu Aschaffenburg, zur Integration der Differentialgleichung $Mdx + Ndy = 0$. [40 S.] gr. 4. geh. n. M. 1.--

\textbf{Henrici}, \textbf{Julius}, die Erforschung der Schwere durch Galilei, Huygens, Newton als Grundlage der rationellen Kinematik und Dynamik, historisch-didaktisch dargestellt. [Beilage zum Jahresbericht des Heidelberger Gymnasiums f\"ur das Schuljahr 1884---85.] [40 S.] 4. geh. M. --.60.

\textbf{Herz}, Dr. \textbf{Norbert}, Assistent f\"ur Astronomie und h\"ohere Geod\"asie an der k. k. technischen Hochschule in Wien, siebenstellige Logarithmen der trigonometrischen Funktionen f\"ur jede Zeitsecunde. Zum astronomischen Gebrauche herausgegeben. [IV u. 182 S.] Lex. 8. geh. n. M. 4.--

\textbf{Killing}, Dr. \textbf{Wilhelm}, ord. Professor am k\"onigl. Lyceum Hosianum zu Braunsberg, die nicht-Euklidischen Raumformen in analytischer Behandlung. Mit einer lithographierten Tafel. [XI u. 264 S.] gr. 8. geh. n. M. 6.80.


% ===== pdf idx 220, printed 201.005 =====
% !PRINTED none
% Publisher's advertisement page (Teubner Verlagsverzeichnis), no running folio visible in the scan crop.

\textbf{Neumann}, Dr. \textbf{F.}, Professor der Physik und Mineralogie, Vorlesungen \"uber theoretische Optik, gehalten an der Universit\"at zu K\"onigsberg. Herausgegeben von Dr. E. \textsc{Dorn}, Professor an der technischen Hochschule zu Darmstadt. Mit Figuren im Text. (Mit einem Bildnis Neumanns in Lichtdruck.) [VIII u. 310 S.] gr. 8. geh. n. M. 9.60.

\textbf{Reidt}, Dr. \textbf{Friedrich}, Professor am Gymnasium und dem Realprogymnasium in Hamm, Sammlung von Aufgaben und Beispielen aus der Trigonometrie und Stereometrie. Zweiter Teil: Stereometrie. Dritte Auflage. [VIII u. 190 S.] gr. 8. geh. M. 3.--

\textbf{Serret}, \textbf{J.-A.}, membre de l'Institut et du Bureau des longitudes, Lehrbuch der Differential- und Integralrechnung. Mit Genehmigung des Verfassers deutsch bearbeitet von \textsc{Axel Harnack}, Dr., und Professor am Polytechnikum zu Dresden. Zweiter Band: Erste H\"alfte. Integralrechnung. Mit in den Text gedruckten Figuren. [VIII u. 380 S.] gr. 8. geh. \hfill n. M. 7.20.

\textemdash{} \textemdash{} Zweiter Band: zweite H\"alfte. (Schlufs.) Differentialgleichungen. Mit in den Text gedruckten Figuren. [VI u. 388 S.] gr. 8. geh. \hfill n. M. 7.20.

Der erste Band: Differentialrechnung erschien 1884 und kostet M. 10.--

\textbf{Stolz}, Dr. \textbf{Otto}, ord. Professor an der Universit\"at zu Innsbruck, Vorlesungen \"uber allgemeine Arithmetik. Nach den neueren Ansichten bearbeitet. Erster Theil: Allgemeines und Arithmetik der reellen Zahlen. [VI u. 344 S.] gr. 8. geh. n. M. 8.--

\textbf{Thieme}, Dr. \textbf{H.}, ord. Lehrer am Realgymnasium zu Posen, Sammlung von Lehrs\"atzen u. Aufgaben aus der Stereometrie. Im Anschlufs an nachgelassene Papiere des Oberlehrers Dr. \textsc{Kretschmer} bearbeitet. [VI u. 92 S.] gr. 8. kart. n. M. 1.20.

\textbf{Weyrauch}, Dr. \textbf{Jacob}, J., Professor an der polytechnischen Schule zu Stuttgart, Aufgaben zur Theorie elastischer K\"orper. Mit 110 Figuren im Text. [X u. 350 S.] gr. 8. geh. n. M. 8.--

\textemdash{} das Princip von der Erhaltung der Energie seit Robert Mayer. Zur Orientirung. [48 S.] gr. 8. geh. M. 1.--

\textbf{W\"ullner}, Dr. \textbf{Adolph}, Professor der Physik an der K\"onigl. technischen Hochschule zu Aachen, Lehrbuch der Experimentalphysik. Dritter Band. Die Lehre von der W\"arme. Vierte vielfach umgearbeitete und verbesserte Auflage. [Mit vielen in den Text gedruckten Holzschnitten.] [VIII u. 825 S.] gr. 8. geh. n. M. 12.--

\rule{6em}{0.4pt}

\end{document}
